\documentclass{amsart}
\def\PaperFinal{1}

\usepackage[T1]{fontenc}
\usepackage{lmodern}
\usepackage{microtype}
\usepackage{amssymb,mathtools,array,longtable,xcolor,mathrsfs}
\newcommand{\Wnorm}{W_{\mathrm{norm}}}
\newcommand{\Wpt}{W}
\newcommand{\epsA}{\varepsilon_A}
\newcommand{\AepsA}{\mathcal A_{\varepsilon_A}}
\newcommand{\muII}{\mu_*}
\newcommand{\mob}{\mu_{\mathrm{Mob}}}
\newcommand{\tauS}{\tau_S}
\newcommand{\CO}{C_{\mathrm{tensor}}}
\newcommand{\Bclip}{B}
\newcommand{\Hpk}{\overline H}
\newcommand{\HLup}{H_{L,\mathrm{upper}}}
\newcommand{\varthetasp}{\vartheta_{\mathrm{sp}}}
\newcommand{\Dlower}{D_{\mathrm{lower}}}
\newcommand{\Gone}{G_1}
\newcommand{\Gbar}{\overline G}
\newcommand{\Jint}{J}
\newcommand{\Ncut}{N_{\mathrm{cut}}}
\newcommand{\NI}{N_I}
\newcommand{\calPN}{\mathcal P_N}
\newcommand{\LGone}{L_{\mathrm{G1}}}
\newcommand{\deltaGone}{\delta_{\mathrm{G1}}}
\newcommand{\Apow}{A_0}
\newcommand{\CE}{C_E}
\newcommand{\Gcap}{\overline G_{\mathrm{cap}}}
\newcommand{\Jsh}{J^{\mathrm{sh}}}
\newcommand{\Gsh}{G_1^{\mathrm{sh}}}
\newcommand{\SGone}{S_{\mathrm{G1}}}
\newcommand{\Wext}{W_{\mathrm{ext}}}
\newcommand{\Aprec}{A_{\mathrm{pr}}}
\newcommand{\cG}{c_G}
\newcommand{\cU}{c_U}
\newcommand{\cUex}{c_U^{\mathrm{ex}}}
\newcommand{\gainSw}{\Delta D_{\mathrm{sw}}}
\newcommand{\gsix}{g_{\mathrm{end}}}
\newcommand{\Mfive}{M_{\mathrm{sh}}}
\newcommand{\gS}{g_{\mathrm{S6}}}
\newcommand{\gSint}{\mathfrak g_{\mathrm{S6}}}
\newcommand{\ERA}{E_{\mathrm{S6}}}
\newcommand{\vRA}{v_{\mathrm{S6}}}
\newcommand{\DB}{D_B}
\newcommand{\Etot}{E_{\mathrm{tot}}}
\newcommand{\Eret}{E_{\mathrm{ret}}}
\newcommand{\Esav}{E^{\mathrm{sav}}}
\newcommand{\Ecv}{E^{\mathrm{cv}}}
\newcommand{\Iliftrows}{\mathcal I_{\mathrm{lift}}}
\newcommand{\Eext}{E_{\mathrm{ext}}}
\newcommand{\Next}{W_{\mathrm{norm}}/4}
\newcommand{\etaTrim}{\eta_{\mathrm{trim}}}
\newcommand{\zetaBuf}{\zeta_{\mathrm{buf}}}
\newcommand{\etaL}{\eta_L}
\newcommand{\Hlog}{H_{\log}}
\newcommand{\Rup}{R_{\mathrm{up}}}
\newcommand{\dErr}{\delta_{\mathrm{err}}}
\newcommand{\Jbarsix}{\overline J_6}
\newcommand{\Khi}{K_{\mathrm{hi}}}
\newcommand{\Kbox}{K_{\mathrm{box}}}
\newcommand{\divfn}{d}

\usepackage[textwidth=42pc,textheight=58pc,centering]{geometry}
\usepackage[hidelinks]{hyperref}
\def\UrlBreaks{\do\/\do-\do.}
\newcommand{\GoldbachTitlePhrase}{1 + 1.75}
\hypersetup{
  pdftitle={Theorem (\GoldbachTitlePhrase) on the Goldbach conjecture},
  pdfauthor={Ruiyi Zhang}
}

\newtheorem{theorem}{Theorem}
\newtheorem{lemma}{Lemma}[section]
\newtheorem{proposition}[lemma]{Proposition}
\newtheorem{corollary}[lemma]{Corollary}
\newtheorem{fact}[lemma]{Fact}
\newtheorem{hypothesis}[lemma]{Hypothesis}
\theoremstyle{remark}
\newtheorem{remark}[lemma]{Remark}
\theoremstyle{plain}

\numberwithin{equation}{section}
\numberwithin{table}{section}
\linespread{0.94}

\makeatletter
\newcommand{\captionof}[1]{\def\@captype{#1}\caption}
\makeatother
\begin{document}

\title{Theorem (\GoldbachTitlePhrase) on the Goldbach conjecture}
\author{Ruiyi Zhang}
\address{Independent Researcher, Columbia, Maryland, USA}
\email{ryzfsu@gmail.com}
\date{November 12, 2026}
\subjclass[2020]{Primary 11P32; Secondary 11N36, 11N35}
\keywords{Goldbach problem, sieve methods, primes in arithmetic progressions, automorphic forms}

\begin{abstract}
Every sufficiently large even integer $N$ has a representation
$N=p+rq$, where $p,q$ are prime, $r=1$ or prime, and
$r\le q^{a-1}=q^{3/4}$. In the terminology of Li and Liu, this proves
Proposition $(1+a)$ for $a=1.75=7/4$. A weighted sieve combines thirteen signed row estimates with
a switched upper bound. The directed comparison gives
$D_B>1.967088903489691335582703173494$ and
$c_{K2,G0}I(b)<1.966859949178424046530804965344$. The exact margin
enclosure is $D_B-c_{K2,G0}I(b)>0.000228954311267289051898208151$.
\end{abstract}

\maketitle
\raggedbottom
\section{Introduction}\label{sec:introduction}

We represent large even integers by two primes and one factor.
Write $\mathbb P$ for the primes.
Chen's theorem represents every sufficiently large even integer as a prime
plus an integer with at most two prime factors
\cite{Chen2002}. Paper~II supplies the quantitative smaller-factor starting
point and the signed sieve sums used here. We change the exterior and
switched coefficients and add retained and sixth-sum gains.
Table~\ref{tab:intro-components} lists their separate costs.
We use the terminology Proposition $(1+a)$ introduced by Li and Liu
\cite[p.~2, (1.6)]{LiLiu2026}.
Runbo Li proved lower bounds for $D_{1,2}(N)$
\cite{RunboLi2024}; we do not use that result.

Cai established a small-prime variant of Chen's theorem \cite{Cai2002}.
Wu developed Chen's double sieve for Goldbach and twin-prime problems
\cite{Wu2004}. His later paper continues this double-sieve program
\cite{Wu2008}. Lichtman modified the linear sieve and studied bounds for the
twin-prime counting function \cite{Lichtman21}. He also proved results on
large-modulus distribution and Goldbach beyond the square-root barrier
\cite{Lichtman23}. Maynard's first paper proves large-modulus estimates for
fixed residue classes \cite{MaynardI}. His second gives well-factorable
estimates in that setting \cite{MaynardII}. His third treats estimates
uniform over residue classes \cite{MaynardIII}. Pascadi studies exponents of
distribution for primes and smooth numbers \cite{Pascadi25}. Li and Liu
proved Proposition $(1+1.9)$, with a lower bound for $D_{1,1.9}(N)$
\cite[Theorem~1.1]{LiLiu2026}. The present proof does not use their
estimates.

\begin{theorem}[Representation theorem]\label{thm:main}
Let $\sigma_*$ be a rational number
satisfying \eqref{eq:bill-sigma}, and let $N_0^*(\sigma_*)$ be the
threshold \eqref{eq:not-common-threshold}. Then every even
$N\ge N_0^*(\sigma_*)$ admits
\begin{equation}\label{eq:intro-target}
 N=p+rq,\qquad p,q\in\mathbb P,\qquad
 r\in\{1\}\cup\mathbb P,\qquad r\le q^{a-1}=q^{3/4}.
\end{equation}
\end{theorem}
Section~\ref{sec:retained-side} proves the retained comparison \textup{(A1)}
of Table~\ref{tab:intro-inputs}.
Section~\ref{sec:root-a} proves the sixth-sum estimate \textup{(A2)}.
Section~\ref{sec:k2} proves the exterior comparison \textup{(A3)}.
Section~\ref{sec:l3} proves the switched comparison \textup{(A4)}. The constant bounds \textup{(A5)} follow from the
cited Paper~I and Paper~II estimates, including Paper~I,
Theorem~2.6, (2.31) \cite[Theorem~2.6, (2.31)]{PaperI}.
Hence a rational $\sigma_*$ as in the theorem exists;
$N_0^*(\sigma_*)$ may be ineffective.

The proof uses Paper~I,
\cite[Theorems~2.6, 2.8, 3.2 and~4.1]{PaperI}, and Paper~II,
\cite[(1.7)--(1.12), (6.4)--(6.4d)]{PaperII}.
We restate each input below with its precise scope.

\subsection{A positive count}

Put
\begin{equation}\label{eq:intro-exponents}
 \varepsilon=\frac1{10000},\qquad a=\frac74,\qquad
 b=\frac37-\frac1{10^6}=\frac{2999993}{7000000}.
\end{equation}
Define
\begin{equation}\label{eq:intro-count}
\begin{aligned}
 \mathcal A_N&=\{N-p:p\in\mathbb P,\ p<(1-\varepsilon)N\},\\
 \mathcal P_b&=\mathbb P\cup
      \{rq:r,q\in\mathbb P,\ r\le q,\ r<N^b\},\\
 G_b(N)&=\#(\mathcal A_N\cap\mathcal P_b).
\end{aligned}
\end{equation}
For a counted semiprime,
\begin{equation}\label{eq:intro-transfer}
\begin{gathered}
 q>\varepsilon N^{1-b},\\
 q^{3/4}>\varepsilon^{3/4}N^{b+\kappa_{\mathrm{tr}}}\ge N^b>r
 \qquad(N\ge N_{\mathrm{tr}}).
\end{gathered}
\end{equation}
The transfer parameters are defined in \eqref{eq:not-transfer-parameters},
and Lemma~\ref{lem:bill-transfer} proves \eqref{eq:intro-transfer}.
A counted prime gives $r=1$.
Hence $G_b(N)>0$ implies \eqref{eq:intro-target}.

Use Paper~II's normalization in \eqref{eq:not-normalization}
\cite[(1.2), \S6.1]{PaperII}.
The convergent product gives $C_2>0$.
Thus $W(N)>0$.

\subsection{The comparisons}

Use the sums of Paper~II \cite[(1.7)--(1.11)]{PaperII}.
Put
\begin{equation}\label{eq:intro-source}
\begin{aligned}
 R={}&3S_1-2S_3-S_4-S_5-S_7+S_8+S_9+S_{10}
              -2S_{11}-S_{12}-S_{13},\\
 \Sigma_W={}&R+S_2-S_6-S_{\mathrm{sw}},\\
 B_b(N)={}&\sum_{\substack{r,q\in\mathbb P\\
              N^b\le r\le\sqrt N\\\varepsilon N<rq<N}}
                       1_{N-rq\in\mathbb P}.
\end{aligned}
\end{equation}
Here $S_2$ is the complete Paper~II sum.
All original coprimality conditions remain imposed.
The last sum counts each eligible ordered prime pair; both orientations are counted only when both satisfy the displayed cutoff, and a square is counted once.
Here $R=R_{\mathrm{fix}}$ and $\Sigma_W=\mathcal W$
in Section~\ref{sec:source-decomposition}.

Write $\overline J_6$ for Paper~II's sixth upper endpoint
\cite[(6.1), (6.3a)]{PaperII}.
For the optional sharper arithmetic comparison, record
\begin{equation}\label{eq:intro-exterior-constant}
 \begin{gathered}
 I(b)=\log\frac{1-b}{b},\\
 c=\frac{6834282367132846348358930380956825846442683629}{10^{45}}.
 \end{gathered}
\end{equation}
Both the exterior and switched comparisons use the common coefficient
\begin{equation}\label{eq:intro-fallback-constant}
 c_{K2,G0}=\frac{6836825163875517756286846321650754372931094689}{10^{45}}.
\end{equation}
Its definition and source comparison are \eqref{eq:k2-certificate}
and \eqref{eq:k2-final-fixed-sigma}.
Put
\begin{equation}\label{eq:intro-components}
\begin{aligned}
 D_{\mathrm{II}}&=\frac{45849916338009170593}{25000000000000000000},\\
 \overline H&=\frac{389253}{5000000},\\
 M_{5,\mathrm{internal}}&=-K_{\mathrm{hi}}\frac{6073714261}{10^9}-L_{G1},\\
 g_{\mathrm{RA,lo}}&=\frac{989598353}{10^{10}},\\
 \Delta D_{L3}&=\left(2-\frac{c_{K2,G0}}4\right)\overline H.
\end{aligned}
\end{equation}
The two mass constants are
\begin{equation}\label{eq:intro-G1-constants}
\begin{aligned}
 K_{\mathrm{hi}}&=
 \frac{3989480791070578466390443009830933753748442947478110112223}{10^{58}},\\
 L_{G1}&=-\frac{243950428265194656254206921939947762868871361874409552095097}
 {100433627766186892221372630771322662657637687111424552206336}.
\end{aligned}
\end{equation}
These are Paper~II, (A.10d) and (4.84), respectively
\cite[(A.10d), (4.84)]{PaperII}.
The constant $\overline H$ bounds the switched integral.
It differs from the packing ceiling $153/256$.
The exact rational $c_{\mathrm{O1}}$ is
\eqref{eq:l3-exact-certificate-constant},
with Table~\ref{tab:l3-certificate-integers}.
It and $c$ belong to the optional certificate comparison only;
neither is used to price the main count.
The sixth-sum comparison uses the certified floor $g_{\mathrm{RA,lo}}$.
The integral gain is denoted $\mathfrak g_{\mathrm{RA}}$, with certified
enclosure $[0.0989598353353,0.0989598353355]$ in
Section~\ref{sec:root-a}. Thus $g_{\mathrm{RA,lo}}$ is rational and
$g_{\mathrm{RA,lo}}\le\mathfrak g_{\mathrm{RA}}$. It is neither the integral nor its central estimate.

The required comparisons are
\begin{equation}\label{eq:intro-comparisons}
\begin{aligned}
 4G_b(N)&\ge\Sigma_W-4B_b(N)-W(N)E_{\mathrm{exc}}(N),\\
 \frac{R+S_2}{W(N)}&\ge D_{\mathrm{II}}+\overline J_6
          +2\overline H+\frac{5719}{10^6}\\
 &\qquad+M_{5,\mathrm{internal}}-E'(N;\sigma),\\
 \frac{S_6}{W(N)}&\le\overline J_6-g_{\mathrm{RA,lo}}
                                  +E_{\mathrm{RA}}(N;\sigma),\\
 \frac{S_{\mathrm{sw}}}{W(N)}&\le
           \frac{c_{K2,G0}}4\overline H
                                  +E_{\mathrm{sw}}(N;\sigma).
\end{aligned}
\end{equation}
These errors are nonnegative by \eqref{eq:retB-goal},
\eqref{eq:ra-b11-floor} and \eqref{eq:l3-remainder}; the first includes the
inherited envelope $E_{\mathrm{van}}$.
Use the following error identifications:
\begin{equation}\label{eq:intro-error-identifications}
\begin{aligned}
 E_{\mathrm{exc}}(N)&=\frac{4E_{\mathrm{pt}}(N)}{W(N)},\\
 E_{\mathrm{sw}}(N;\sigma)&=C_{\mathrm{sw}}\sigma
                     +P_{\mathrm{sw}}\sigma^9+V_{\mathrm{sw}}(N;\sigma).
\end{aligned}
\end{equation}
Here $E_{\mathrm{pt}}$ is \eqref{eq:source-pointwise-error}.
The switched quantities are \eqref{eq:l3-price} and \eqref{eq:l3-remainder}.
Their estimates form input \textup{(A4)}.
The retained error includes the conversion errors:
\begin{equation}\label{eq:intro-retained-error}
\begin{aligned}
 E'(N;\sigma)={}&E_{\mathrm{van}}(N;\sigma)
 +\sum_{j\in\mathcal H}\bigl(|\Lambda^0_{j,N}|+|\Lambda_{j,N}|\bigr)\\
 &+|\Lambda_{G,N}|.
\end{aligned}
\end{equation}
The row set $\mathcal H$ is defined in \eqref{eq:retB-goal}.
Its two qualified labels denote $S_5$'s lower part and retained $S_{10}$.
The inherited error $E_{\mathrm{van}}$ is \eqref{eq:retB-base-errors};
the saved-array contribution $\Lambda^0_{j,N}$ is included once,
as in \eqref{eq:retB-goal}.
The differences $\Delta_j$ are \eqref{eq:retB-H6-gains}.
The conversion arrays are \eqref{eq:retB-direct-Lambda}
and \eqref{eq:retB-G-array}; the equivalent notation
$\Lambda_{j,N}=\Lambda_N(\Delta_j)$ and
$\Lambda_{G,N}=\Lambda_N(\Delta_{G_1})$ is fixed in
\eqref{eq:retB-Lambda-alias}.
The saved-array identifications and conversion bounds are proved in
Section~\ref{sec:retained-side}, as part of \textup{(A1)}.
The two finite integral bounds are
Propositions~\ref{prop:app-h6-four-profiles} and~\ref{prop:app-m5-cap}.

The exterior estimate required is
\begin{equation}\label{eq:intro-exterior}
 -\frac{4B_b(N)}{W(N)}\ge-c_{K2,G0}I(b)-E_K(N;b,\sigma),
 \qquad E_K(N;b,\sigma)\ge0.
\end{equation}
Its convolution input is Paper~II,
\cite[Lemma~5.1, (5.44)--(5.52)]{PaperII}.

Define the constant by
\begin{equation}\label{eq:intro-DB}
 \boxed{\begin{aligned}
 D_B={}&D_{\mathrm{II}}+\frac{5719}{10^6}
                  +M_{5,\mathrm{internal}}\\
       &+g_{\mathrm{RA,lo}}-\frac1{10000}+\Delta D_{L3}.
 \end{aligned}}
\end{equation}
Table~\ref{tab:intro-components} records each summand.
Its decimals are rounded; the defining formulas are exact.

\begin{table}[t]
\caption{Components of $D_B$}\label{tab:intro-components}
\small
\renewcommand{\arraystretch}{1.15}
\begin{tabular}{@{}lll@{}}
\hline
Component & Value & Definition\\
\hline
$D_{\mathrm{II}}$ & $1.833996653520367$ & \eqref{eq:intro-components}\\
$5719/10^6$ & $0.005719$ & \eqref{eq:retB-H6-finite}\\
$M_{5,\mathrm{internal}}$ & $0.005874949945026$ & \eqref{eq:intro-components}, \eqref{eq:intro-G1-constants}\\
$g_{\mathrm{RA,lo}}$ & $0.0989598353$ & \eqref{eq:intro-components}\\
$-1/10000$ & $-0.0001$ & \eqref{eq:intro-DB}\\
$\Delta D_{L3}$ & $0.022638464724298$ & \eqref{eq:intro-components}\\
\hline
$D_B$ & $1.967088903489691$ & \eqref{eq:intro-DB}\\
\multicolumn{3}{c}{$1.967088903489691335582703173494<D_B$}\\
\multicolumn{3}{c}{$D_B<1.967088903489691335582703173495$}\\
\hline
\end{tabular}
\end{table}

Suppose the optional six-endpoint estimate is available.
Then Proposition~\ref{prop:retB-comparison} would add\hfill\break
$g_6=0.0020549454896034927680\ldots$ of \eqref{eq:retB-six-gain} to the
retained comparison. The assembled retained-count lower bound below would
then increase by $g_6$.
Theorem~\ref{thm:main} does not use this.
Let
\begin{equation}\label{eq:intro-total-error}
 \mathcal E_B=E'+E_{\mathrm{RA}}+E_{\mathrm{sw}}+E_K+E_{\mathrm{exc}},
 \qquad \delta_{\mathrm{err}}=\frac1{10000}.
\end{equation}
Adding the established comparisons gives
\begin{equation}\label{eq:intro-assembled}
 \frac{4G_b(N)}{W(N)}\ge
 D_B+\delta_{\mathrm{err}}-c_{K2,G0}I(b)-\mathcal E_B(N;b,\sigma).
\end{equation}
By \textup{(A1)}--\textup{(A5)}, Proposition~\ref{prop:bill-budget}
proves
\begin{equation}\label{eq:intro-error-bound}
 \mathcal E_B(N;b,\sigma_*)\le\delta_{\mathrm{err}}
 \qquad(N\ge N_0^*(\sigma_*),\ N\in2\mathbb Z).
\end{equation}
Table~\ref{tab:intro-inputs} lists the five inputs, their exact
statements and their contributions to the final constant.
The first-line source count in \eqref{eq:intro-comparisons} is
\eqref{eq:source-first-lower}, valid for
$N\ge N_{\mathrm{pt}}=10^{12}+1$. Lemma~\ref{lem:bill-transfer} gives
$N_{\mathrm{pt}}<N_{\mathrm{tr}}\le N_0^*(\sigma_*)$.
\begingroup
\small
\renewcommand{\arraystretch}{1.08}
\begin{longtable}{@{}p{.09\linewidth}p{.57\linewidth}p{.30\linewidth}@{}}
\caption{Analytic inputs and their contributions}\label{tab:intro-inputs}\\
\hline
Input & Exact statement and use & Contribution and error\\
\hline
\endfirsthead
\hline
Input & Exact statement and use & Contribution and error\\
\hline
\endhead
A1 & Retained comparison \eqref{eq:retB-goal} with $g_6$ replaced by $0$,
Proposition~\ref{prop:retB-comparison}, at $\sigma=\sigma_*$; used in line 2
of \eqref{eq:intro-comparisons}; proved in Section~\ref{sec:retained-side}.
& $+5719/10^6+M_{5,\mathrm{internal}}$; error $E'$.\\
A2 & Proposition~\ref{hyp:ra-applications} at $\sigma=\sigma_*$, giving the
sixth-sum estimate \eqref{eq:ra-b11-floor} of
Proposition~\ref{prop:ra-b11}; used in line 3 of
\eqref{eq:intro-comparisons}. &
$+g_{\mathrm{RA,lo}}$; error $E_{\mathrm{RA}}$.\\
A3 & Proposition~\ref{prop:k2-applications} at $\sigma=\sigma_*$, giving the
exterior estimate \eqref{eq:k2-final-fixed-sigma} of
Proposition~\ref{prop:k2-comparison}; used in \eqref{eq:intro-exterior}. &
$-c_{K2,G0}I(b)$; error
$A_KC_{\mathrm{AA}}\sigma+P_K\sigma^9+r_K$.\\
A4 & The sieve data and the estimates \eqref{eq:l3-source-error},
\eqref{eq:l3-disp-total} and \eqref{eq:l3-small-estimates} of
Section~\ref{sec:l3} for the source \eqref{eq:l3-source}, at
$\sigma=\sigma_*$, giving the switched estimate \eqref{eq:l3-final} and the
limit \eqref{eq:l3-remainder-limit}; used in line 4 of
\eqref{eq:intro-comparisons}. This application uses Paper~I,
Theorem~2.8, at its stated inner spectral stage. &
$+\Delta D_{L3}$; error
$C_{\mathrm{sw}}\sigma+P_{\mathrm{sw}}\sigma^9+V_{\mathrm{sw}}$.\\
A5 & The constants $C_E$ in \eqref{eq:l3-sieve-error} and
$C_{\mathrm{AA}}$ in \eqref{eq:k2-aa-constant} are finite, positive and
independent of $\sigma$ and $N$, and $C_E\ge C_E^{\mathrm{II}}$. &
Fixes the caps in \eqref{eq:bill-sigma}.\\
\hline
\end{longtable}
\endgroup

The constants $C_E$ and $C_{\mathrm{AA}}$ are distinct. Here $C_E$ is the
common maximum of the switched and Paper~II sieve constants.
The four-row K2 error constant is $C_{\mathrm{AA}}=4C_{\mathrm{AA37}}$.

The exact constants satisfy
\begin{equation}\label{eq:intro-price}
\begin{aligned}
 0.000228954311267289051898208151
   &<D_B-c_{K2,G0}I(b),\\
 D_B-c_{K2,G0}I(b)
   &<0.000228954311267289051898208152,\\
 1.749974884137892433433207484622
   &<a_{\mathrm{buf}}^*,\\
 a_{\mathrm{buf}}^*
   &<1.749974884137892433433207484623<\frac74.
\end{aligned}
\end{equation}
The last constant is defined in \eqref{eq:bill-critical-exponent}.
The rational bounds are shown in Section~\ref{sec:final-bill}.
They need the stated analytic comparisons;
the arithmetic alone gives no positive-count estimate.
The value $a^*_{\mathrm{buf}}$ is a scalar diagnostic for this fixed $b$.
It does not prove the theorem for smaller $a$; that would require
comparisons at a different cutoff. The optional sharper scalar comparison
\eqref{eq:bill-active-critical-exponent} uses a separate constant
$D_{B,c}$; it does not enter the main count.

\subsection{Order of the parameters}

Fix $b$ and all absolute error constants first.
Choose one sufficiently small rational $\sigma_*$.
Then put $\tau_*=(\sigma_*)^9$.
Freeze the finite partitions and all coefficient bounds next.
Fix every logarithmic precision before the common threshold.
Thus the order is
\begin{equation}\label{eq:intro-order}
\begin{gathered}
 \text{absolute constants}\ \longrightarrow\
 \sigma_*\in\mathbb Q\cap(0,10^{-10}]\ \longrightarrow\
 \tau_*=(\sigma_*)^9,\\
 \longrightarrow\ \text{partitions, coefficient bounds and precisions}
 \ \longrightarrow\ N_0^*(\sigma_*)\ \longrightarrow\ N.
\end{gathered}
\end{equation}
The common threshold is defined in \eqref{eq:not-common-threshold}.
It includes every length, structural and error threshold.
Every limit is taken at fixed $\sigma_*$.
These thresholds exist by \textup{(A1)}--\textup{(A5)} and
Proposition~\ref{prop:bill-budget}.
The final section combines the comparisons and the error bounds.
For the main construction, the comparison estimates used are:
\begin{enumerate}
\item \textup{(A1)}: the actual retained-source comparison
in Section~\ref{sec:retained-side}, including all saved-array conversions.
It supplies $5719/10^6+M_{5,\mathrm{internal}}$ beyond the original
retained bound.
\item \textup{(A2)}: the sixth-sum distribution and
coefficient estimates in Section~\ref{sec:root-a}, supplying
$g_{\mathrm{RA,lo}}$.
\item \textup{(A3)}: the actual semiprime-source distribution,
principal-term and peel estimates in Section~\ref{sec:k2}, giving the
exterior cost $c_{K2,G0}I(b)$.
\item \textup{(A4)}: the switched-source comparison,
packing, distribution and peel estimates in Section~\ref{sec:l3}, giving
$\Delta D_{L3}=(2-c_{K2,G0}/4)\overline H$.
\end{enumerate}
Each requirement includes its exact nonnegative error estimates and the
fixed-$\sigma_*$ limits allocated in Section~\ref{sec:final-bill}.
The source-counting identity and pointwise bound are supplied in
Section~\ref{sec:source-decomposition}. The list only records dependencies.
The cited sections prove the estimates, with precise scopes and
error bounds; the scalar arithmetic here combines them.

The use of AI tools is described at the end.

\section{Notation and parameters}\label{sec:notation}

All logarithms are natural. Put \(L=\log N\).
The integer \(N\) is even throughout.
Table~\ref{tab:notation} fixes the common notation.

\begingroup
\small
\renewcommand{\arraystretch}{1.15}
\begin{longtable}{@{}p{.27\linewidth}p{.67\linewidth}@{}}
\caption{Common notation}\label{tab:notation}\\
\hline
Symbol & Meaning\\
\hline
\(\mathbb P\) & The set of primes.\\
\(\varphi(n),d(n)\) & Euler's totient and divisor function. The inherited notation \(\tau(n)=d(n)\) has an integer argument; the scalar \(\tau=\sigma^9\) has none.\\
\(\mu_{\mathrm{Mob}}(n),\mu(n)\) & The M\"obius function; use \(\mu_*\) for the sieve cutoff.\\
\(\Omega(n)\) & Prime factors counted with multiplicity.\\
\(P^-(n)\) & Least prime factor; \(P^-(1)=\infty\).\\
\(\mathbf1_E\) & Indicator of \(E\).\\
\(\gamma,\gamma_E\) & Euler's constant, with \(\gamma_E=\gamma\).\\
\(a,b\) & Representation exponent and smaller-factor cutoff.\\
\(\varepsilon\) & Fixed truncation parameter.\\
\(\alpha,\beta,\mu_*,\lambda,t_0,u\) & Fixed sieve exponents in \eqref{eq:not-sieve-parameters}; cutoff \(\mu=\mu_*\).\\
\(C(N),U_C(N),W(N)\) & Normalizations in \eqref{eq:not-normalization}.\\
\(\sigma,\tau,\rho\) & Partition parameters: \(\tau=\sigma^9\), \(\rho=1+\tau\).\\
\(\sigma_*,\tau_*,\rho_*\) & Their fixed choices for the final argument.\\
\(\kappa_{\mathrm{tr}},N_{\mathrm{tr}}\) & Transfer parameters in \eqref{eq:not-transfer-parameters}.\\
\(\delta_{\mathrm{err}}\) & Fixed total error bound, \(10^{-4}\).\\
\(\delta_{\mathrm{pk}}\) & The parameter in the factor-partition theorem.\\
\(h_G,h,\delta_h\) & Profile displacement \(h_G=9/5000\); \(h=h_G\) in \eqref{eq:retB-parameters}, \(\delta_h=h_G/\alpha\).\\
\(\epsilon_s,\delta_s\) & Local sieve parameters, \(\delta_s=\epsilon_s^9\); applications use \(\epsilon_s=\sigma\).\\
\(\eta_L\) & The level margin, \(10^{-4}\); separate from \(\varepsilon\) and \(\delta_{\mathrm{err}}\).\\
\(r,q\) & In \eqref{eq:intro-target}, \(q\) is prime and \(r\) is either \(1\) or prime; other uses are locally bound variables.\\
\(c,I(b)\) & Optional sharper coefficient and the logarithm in \eqref{eq:intro-exterior-constant}.\\
\(c_{K2,G0}\) & Common K2/L3 coefficient in \eqref{eq:intro-fallback-constant}.\\
\(C_E,C_{\mathrm{AA}}\) & Error constants in \eqref{eq:l3-sieve-error} and \eqref{eq:k2-aa-mass}.\\
\(K_{\mathrm{SW}}(A),C_{\mathrm{K2},\sigma},C_{\mathrm{L3},\sigma}\) & The logarithmic Siegel--Walfisz level and the K2 and L3 coefficient bounds; each is local to its stated application.\\
\(R,R_{\mathrm{fix}}\) & The signed sum in \eqref{eq:intro-source}; elsewhere \(R\) is a local level or remainder.\\
\(\Sigma_W,\mathcal W\) & Equal signed source sums; see \eqref{eq:source-weight-identity}.\\
\(W_{13},W_{rs}\) & Arithmetic source weight and spectral operator, respectively. Neither denotes \(W(N)\).\\
\(\overline H,\overline H_L\) & Equal switched-integral bounds, \(389253/5000000\).\\
\(H_{\mathrm{pk,max}},H_{\max}\) & Equal packing ceilings, \(153/256\).\\
\(D_{\mathrm{II}},\overline J_6\) & Original Paper II lower constant and sixth upper endpoint; \eqref{eq:intro-components}, \eqref{eq:ra-j-functional}.\\
\(g_{\mathrm{RA,lo}}\) & Certified floor \(0.0989598353\), used in \eqref{eq:intro-DB}.\\
\(E',E_{\mathrm{van}}\) & Complete retained error and inherited envelope; \eqref{eq:intro-retained-error}, \eqref{eq:retB-base-errors}.\\
\(\begin{gathered}\Lambda_{j,N},\\\Lambda_{G,N}\end{gathered}\) & Retained conversion differences, defined in \eqref{eq:retB-direct-Lambda} and \eqref{eq:retB-G-array}.\\
\(\Delta_G,\Lambda_{G,N}\) & Joint-profile gain and conversion difference in \eqref{eq:retB-G-array}. The spelling \(G1\) in subscripts denotes the same group.\\
\(\begin{gathered}E_{\mathrm{RA}},E_{\mathrm{sw}},\\E_K,E_{\mathrm{exc}}\end{gathered}\) & Nonnegative complete errors in \eqref{eq:intro-comparisons}--\eqref{eq:intro-error-identifications} and \eqref{eq:intro-exterior}.\\
\(\mathcal E_B\) & Their total, including \(E'\), in \eqref{eq:intro-total-error}.\\
\(D_B\) & The constant in \eqref{eq:intro-DB}; evaluated in \eqref{eq:bill-constant}.\\
\(\mathcal A_N,\mathcal P_b,G_b,B_b\) & Shifted-prime source, admissible factors, positive count and exterior count in \eqref{eq:intro-count}, \eqref{eq:intro-source}.\\
\(S_1,\ldots,S_{13},S_{\mathrm{sw}}\) & Paper II arithmetic sums retained in \eqref{eq:intro-source}; definitions in Section~\ref{sec:source-decomposition}.\\
\(\mathfrak g_{\mathrm{RA}},g_{\mathrm{RA,lo}}\) & Integral gain and its smaller rational floor; only the floor enters $D_B$.\\
\(\begin{gathered}g_6,M_{5,\mathrm{internal}},\\K_{\mathrm{hi}},L_{G1}\end{gathered}\) & Optional six-endpoint gain, shifted-profile gain and mass constants; the main proof sets \(g_6=0\), see \eqref{eq:retB-six-gain}, \eqref{eq:intro-components}, \eqref{eq:intro-G1-constants}.\\
\(\Delta D_{L3}\) & Switched gain $(2-c_{K2,G0}/4)\overline H$.\\
\(G_*,G_0,U_0,c_{\mathrm{O1}}\) & Common unshifted K2/L3 functional; coefficient cap $G_0=6088440363/10^9$; optional sharper upper certificate $U_0$; its exact coefficient in \eqref{eq:l3-exact-certificate-constant}.\\
\(\mathcal H,\Lambda^0_{j,N}\) & Retained row set and saved-array contribution in \eqref{eq:retB-goal}.\\
\(\mathcal E,E_{\mathrm{pt}}\) & Pointwise exceptional set and error in \eqref{eq:source-exceptional-set}, \eqref{eq:source-pointwise-error}.\\
\(N_0^*,N_{\mathrm{floor}},T_{\bullet}\) & Common threshold, numerical floor and named threshold requirements in \eqref{eq:not-common-threshold}.\\
\(\zeta,a_{\mathrm{buf}}^*,\mathfrak E_*\) & Buffer $10^{-6}$, scalar endpoint and its exponential in \eqref{eq:bill-critical-exponent}.\\
\(W_{\mathrm{AA},K},W_*\) & K2 source-weight bound in \eqref{eq:k2-aa-mass} and Paper II bound $10^{14}$; these are distinct.\\
\(A_K,P_K,r_K\) & K2 error multiplier, peel multiplier and vanishing remainder in \eqref{eq:bill-K-error}.\\
\(C_{\mathrm{sw}},P_{\mathrm{sw}},V_{\mathrm{sw}}\) & Switched error multipliers and positive remainder in \eqref{eq:l3-price}, \eqref{eq:l3-remainder}.\\
\(L,g_{\mathrm{AA}}\) & $\log N$ and Paper II's fixed error allowance in \eqref{eq:bill-paperII-AA}.\\
\hline
\end{longtable}
\endgroup

Cusps carry the index \(\mathfrak a\) in spectral sums.
The switched blocks carry the index \(\mathfrak b\).
Partition cells are \(c\) or \(\mathfrak c\).
Coefficient sequences are \(a_n,b_m\), with dummy indices interchanged when needed.
Spectral moduli and cusp numerators are locally bound variables.
They do not redefine the constants \(c\) and \(u\).
Packing coordinates also have local scope.
Bare \(T\) and \(w\) are local endpoints or weights; the named saved and
root-A scales are \(T_{\mathrm{save}},T_{\mathrm{int}}\), and \(w_{\mathrm{RA}}\).
The exponent \(A_0>0\) in distribution estimates is a logarithmic precision.
The normalization factor is written \(2e^{-\gamma}\).

\subsection{Fixed parameters}

Keep \(a,b,\varepsilon\) from \eqref{eq:intro-exponents}.
The sieve parameters are
\begin{equation}\label{eq:not-sieve-parameters}
\begin{aligned}
 \alpha&=\frac{100}{1149},&\beta&=\frac{50}{309},&
 \mu_*&=\frac{25}{128},\\
 \lambda&=\frac{93}{400},&t_0&=\frac{57}{224},&
 u&=\frac1{500}.
\end{aligned}
\end{equation}
These follow Paper II's notation
\cite[(1.5)]{PaperII}; $\varepsilon$ is $\varepsilon_A$ in
\cite[(1.3)]{PaperII}.
Put
\begin{equation}\label{eq:not-normalization}
\begin{aligned}
 C_2&=\prod_{\substack{\ell\in\mathbb P\\\ell>2}}
       \left(1-\frac1{(\ell-1)^2}\right),\\
 C(N)&=C_2\prod_{\substack{\ell\in\mathbb P,\ \ell\mid N\\\ell>2}}
                   \frac{\ell-1}{\ell-2},\\
 U_C(N)&=\frac{C(N)N}{L^2},\qquad
 W(N)=4(1-\varepsilon)U_C(N).
\end{aligned}
\end{equation}
We often suppress the argument \(N\).

Define the transfer parameters by
\begin{equation}\label{eq:not-transfer-parameters}
\begin{aligned}
 \kappa_{\mathrm{tr}}&=a(1-b)-1=\frac7{4000000},\\
 N_{\mathrm{tr}}&=\varepsilon^{-(a-1)/\kappa_{\mathrm{tr}}}
                 =\varepsilon^{-3000000/7}.
\end{aligned}
\end{equation}
Here $a-1=3/4$ is the exponent on the larger prime in the bound for the
smaller factor.
The transfer argument appears in the final section.

\subsection{Order of choice}

First fix all \(\sigma\)-independent constants in the estimates.
These include \(C_E\) in \eqref{eq:l3-sieve-error}
and \(C_{\mathrm{AA}}\) in \eqref{eq:k2-aa-mass}.
Then choose one rational \(\sigma_*\), with
\begin{equation}\label{eq:not-sigma-choice}
 0<\sigma_*\le10^{-10},\qquad
 \tau_*=(\sigma_*)^9,\qquad \rho_*=1+\tau_*.
\end{equation}
The final section gives the remaining upper restrictions.
Apply every partition estimate with \(\sigma=\sigma_*\).
Freeze its finite partitions and coefficient bounds next.
Freeze the resulting finite-family bounds and logarithmic precisions too.
These choices precede the common threshold.

Use the following named thresholds. Each requires its listed statement
for every even $N$ above it, at the fixed $\sigma_*$. The cited comparison
propositions establish these statements. The numerical floor supplies none
of them. The resulting threshold exists but may be ineffective.
\begin{enumerate}
\item $T_{\mathrm{src}}$: the actual-source identity and first line of
\eqref{eq:intro-comparisons}, including the pointwise hypotheses in
\eqref{eq:source-pointwise} and \eqref{eq:source-first-lower}.
Include the explicit pointwise threshold
$N_{\mathrm{pt}}=1+\lceil\max\{2^{1/\alpha},\varepsilon^{-3},
\varepsilon^{-1/(1-\lambda-\beta)}\}\rceil$ from
Section~\ref{sec:source-decomposition}.
\item $T_{\mathrm{ret}}$: the retained comparison, saved-array
identifications and nonnegativity of $E_{\mathrm{van}}$ in
Section~\ref{sec:retained-side}.
\item $T_{\mathrm{RA}}$: the sixth-sum comparison
\eqref{eq:ra-b11-floor}, with every distribution and coefficient bound
specified in Section~\ref{sec:root-a}. It also covers the finite source
threshold $N_6$, chosen so that every even $N\ge N_6$ satisfies
\[
 N^{1/6}>L^{\max(B_L,B_{\mathrm{tail}})},\quad
 N^{v_{\mathrm{RA}}-1/3}>L^B,\quad
 \varepsilon N^{1/2}>N^{\mu_*},\quad
 NL^{-333}\ge N^{\mu_*}.
\]
Thus the operative $T_{\mathrm{RA}}$ is the maximum of the sixth-sum
comparison threshold and this $N_6$.
\item $T_K$: the actual-source exterior comparison
\eqref{eq:k2-final-fixed-sigma} and its nonnegative error decomposition,
the $r_K$ formula displayed after
\eqref{eq:k2-final-fixed-sigma}.
\item $T_{\mathrm{sw}}$: the switched comparison \eqref{eq:l3-final}
and the nonnegative decomposition \eqref{eq:l3-remainder}.
\item $T_{\mathrm{len}}$: all scale-conversion, zero-mode, factorization,
packing, interval-length and logarithmic-precision conditions of the
preceding applications, with the fixed finite partitions. In particular,
require $N^{\sigma_*}\ge2$,
$N^{60\sigma_*2^{-126}}\ge12L^{334}$,
\eqref{eq:k2-med-length} and \eqref{eq:l3-lengths}, and include
$N_{\mathrm{dom}},N_{z,0},N_{\mathrm{LEN,app}}$ from
\eqref{eq:l3-structural-thresholds}. The remaining application conditions
are exactly those of the cited comparison propositions. The floor implies
none of them.
\item $T_{E'},T_V,T_{E_{\mathrm{RA}}},T_r,T_{E_{\mathrm{exc}}}$:
respectively the eventual bounds
\[
\begin{gathered}
 E'\le\frac{\delta_{\mathrm{err}}}{20},\qquad
 V_{\mathrm{sw}}\le\frac{\delta_{\mathrm{err}}}{20},\qquad
 E_{\mathrm{RA}}\le\frac{\delta_{\mathrm{err}}}{10},\\
 r_K\le\frac{\delta_{\mathrm{err}}}{10},\qquad
 E_{\mathrm{exc}}\le\frac{\delta_{\mathrm{err}}}{10}.
\end{gathered}
\]
All arguments are $(N;\sigma_*)$, or $(N;b,\sigma_*)$ for $r_K$.
\end{enumerate}
These are finitely many requirements after all finite choices.
Define the common threshold once by
\begin{equation}\label{eq:not-common-threshold}
\begin{aligned}
 \mathcal T_*={}&\{T_{\mathrm{src}},T_{\mathrm{ret}},T_{\mathrm{RA}},
 T_K,T_{\mathrm{sw}},T_{\mathrm{len}},\\
 &\hspace{12mm}T_{E'},T_V,T_{E_{\mathrm{RA}}},T_r,T_{E_{\mathrm{exc}}}\},\\
 N_0^*(\sigma_*)&=2\left\lceil\frac12
 \max\bigl(\{N_{\mathrm{tr}},N_{\mathrm{floor}}\}\cup\mathcal T_*\bigr)
 \right\rceil,\\
 N_{\mathrm{floor}}&=2\left\lceil e^{2^{150}}/2\right\rceil.
\end{aligned}
\end{equation}
The floor is a convention added to Paper~I's
selected-interval length hypothesis
\cite[\S3.2, length hypothesis preceding (3.12)]{PaperI}. It proves none of
the listed estimates. The named application estimates establish this
threshold. The resulting $N_0^*$ may be ineffective.

This implements the order in \eqref{eq:intro-order}.
Afterwards, \(N\) is even and $N\ge N_0^*(\sigma_*)$.

Let \(\epsilon_{\mathrm{err}}(N;\sigma_*)\) be a real- or complex-valued error function.
Write \(\epsilon_{\mathrm{err}}(N;\sigma_*)=o_{\sigma_*}(1)\) if and only if
\begin{equation}\label{eq:not-little-o}
\begin{gathered}
 \forall\eta>0\ \exists N_\eta(\sigma_*)\ \forall N\in2\mathbb N,\\
 N\ge N_\eta(\sigma_*)\ \Longrightarrow\ |\epsilon_{\mathrm{err}}(N;\sigma_*)|\le\eta.
\end{gathered}
\end{equation}
All preceding fixed data are suppressed here.
The parameter \(\sigma_*\) remains fixed in this limit.
Uniformity in other variables is stated with each estimate.

\section{Classical inputs}\label{sec:classical-inputs}
\raggedbottom

We use counting-measure $\ell^2$ norms. Write $n\sim Y$ for $Y<n\le2Y$.
The symbols $\mu_{\rm Mob},\Lambda,\varphi$ have their usual meanings.
Write $d(n)$ for the divisor function.
The parameters in each fact are local.
Put $P^-(1)=\infty$ and $e(t)=\exp(2\pi it)$.
Statements (P1.x) and (P2.x) restate results proved in
Papers~I and~II, respectively, with their original scope.
The structured sieve result is Paper~I, Theorem~2.6.
Its classical source is Iwaniec's Theorem~4
\cite[Theorem~4]{Iwaniec1980Remainder}. The \(S\)-form displayed below
specializes the \(W_R^\pm\) estimates of Paper~I, Theorem~2.6.
It follows from them by \(S\le W_R^+\) and \(S\ge W_R^-\).
Theorem~4 does not state it verbatim.
Each new application needs its own coefficient, support and
uniformity checks.

\begin{fact}[(F1) Prime estimates]\label{fact:prime-estimates}
For every $A>0$, one may take $B_\psi(A)=A+60$ in
\begin{equation}\label{eq:classical-bv}
 \begin{gathered}
 \sum_{q\le Q}\max_{(a,q)=1}
 \left|\sum_{\substack{n\le y\\n\equiv a\ (q)}}\Lambda(n)
                -\frac y{\varphi(q)}\right|
 \ll_A y(\log y)^{-A},\\
 y\ge3,\quad Q\le y^{1/2}(\log y)^{-B_\psi(A)}.
 \end{gathered}
\end{equation}
For fixed $A,C>0$, uniformly for $q\le(\log y)^C$,
\begin{equation}\label{eq:classical-sw}
 \pi(y;q,a)=\frac{\operatorname{Li}(y)}{\varphi(q)}
       +O_{A,C}(y(\log y)^{-A}),\quad (a,q)=1,\quad y\ge3.
\end{equation}
These are Paper~I, (2.20)--(2.21) \cite[(2.20)--(2.21), p.~9]{PaperI}.
The endpoint $y$ is fixed. The second constant is ineffective.
\end{fact}

\begin{fact}[(F1$'$) Maximal-endpoint prime estimates]\label{fact:prime-estimates-maximal}
Use the prime Siegel--Walfisz estimate \eqref{eq:classical-sw}.
For every fixed $S>0$ and $X\ge3$, put
$Q_X=X^{1/2}(\log X)^{-(S+60)}$. Then
\begin{equation}\label{eq:classical-bv-maximal}
 \begin{aligned}
 \sum_{q\le Q_X}\max_{(a,q)=1}\sup_{0\le y\le X}
 \left|\sum_{\substack{n\le y\\n\equiv a\ (q)}}\Lambda(n)
                -\frac y{\varphi(q)}\right|
 &\ll_S X(\log X)^{-S},\\
 \sum_{q\le Q_X}\max_{(a,q)=1}\sup_{0\le y\le X}
 \left|\pi(y;q,a)-\frac{\operatorname{Li}(y)}{\varphi(q)}\right|
 &\ll_S X(\log X)^{-S}.
 \end{aligned}
\end{equation}
Here $\operatorname{Li}(y)=0$ for $0\le y<2$.
The level is imposed at $X$; the endpoints may be chosen separately for
each modulus. An interval difference costs at most twice the bound.
This is the classical maximal form of Bombieri--Vinogradov
\cite[Chapter~28]{Davenport1980}. Removing prime powers and partial
summation pass from the von Mangoldt form to the prime-counting form.
The constants may be ineffective through \eqref{eq:classical-sw}.
\end{fact}

\begin{fact}[(F2) Full-support bilinear distribution]\label{fact:bv-full-support}
Let $a_m,b_n$ have supports $m\sim M,n\sim Y$.
Assume $x=MY\ge3$, set $L_x=\log x$, and fix $\epsilon,B>0$.
Assume $\epsilon\le\log Y/\log x\le1-\epsilon$.
For every $A_0>0$, assume uniformly in $h,k\ge1$ and integers
$\ell\ne0$ with $(\ell,k)=1$ that
\begin{equation}\label{eq:bvfs-sw}
 \left|\sum_{\substack{n\equiv\ell\ (k)\\(n,h)=1}}b_n
 -\frac1{\varphi(k)}\sum_{(n,hk)=1}b_n\right|
 \le C_{A_0}\|b\|_2Y^{1/2}d(h)^B(\log(2Y))^{-A_0}.
\end{equation}
Define
\[
 \Delta_{\rm fs}(q,a)=\sum_{mn\equiv a\ (q)}a_mb_n
       -\frac1{\varphi(q)}\sum_{(mn,q)=1}a_mb_n.
\]
For every $A>0$, some $B_{\rm BV}=B_{\rm BV}(A,B,\epsilon)>0$ gives
\begin{equation}\label{eq:bvfs}
 \sum_{q\le Q}\max_{(a,q)=1}|\Delta_{\rm fs}(q,a)|
 \ll \|a\|_2\|b\|_2x^{1/2}L_x^{-A},\qquad
 Q\le x^{1/2}L_x^{-B_{\rm BV}}.
\end{equation}
Both sums use the full supports. The constants depend on the family constants
$C_{A_0}$. This follows from the proof of BFI Theorem~0(b)
\cite[(A1)--(A2), p.~206; Theorem~0(b), (1.5), p.~212;
proof, pp.~213--214]{BFI86}.
\end{fact}

\begin{proof}
Character orthogonality and primitive induction give
\[
\begin{aligned}
 \text{left side of \eqref{eq:bvfs}}
 &\le\sum_{e\le Q}\frac1{\varphi(e)}
       \sum_{2\le f\le Q/e}\frac1{\varphi(f)}\\
 &\qquad\cdot\sum_{\chi\ (f)}^*|A_e(\chi)B_e(\chi)|,\\
 A_e(\chi)&=\sum_{(m,e)=1}a_m\chi(m).
\end{aligned}
\]
Define $B_e$ similarly. Split at $F=L_x^K$.
Since $M\ge x^\epsilon$, large $x$ gives $F^2\le M$.
This permits the small-conductor estimate. The hypothesis handles $f\le F$.
The multiplicative large sieve handles $f>F$. Before summing $e$, the bounds are
\[
 \ll\|a\|_2\|b\|_2x^{1/2}d(e)^BF^{3/2}L_x^{-A_0},
 \qquad
 \ll L_x^2\|a\|_2\|b\|_2
       (Q+M^{1/2}+Y^{1/2}+x^{1/2}/F).
\]
Choose $B_2$ with $\sum_{e\le Q}d(e)^B/\varphi(e)\ll L_x^{B_2}$.
In the large-conductor bound, summing $1/\varphi(e)$ costs
one more logarithm, since $\sum_{e\le Q}1/\varphi(e)\ll L_x$.
Take $K=A+B_2+4$, $A_0=A+B_2+3K/2+1$, and $B_{\rm BV}=K$.
Since $M,Y\ge x^\epsilon$, these bounds prove \eqref{eq:bvfs}.
The large sieve is BFI (1.6), p.~207, restated in Paper~I, (2.19)
\cite[(1.6), p.~207]{BFI86}\cite[(2.19), p.~9]{PaperI}.
\end{proof}

We do not use BFI86, Theorem~0(a).

Fact (F2) is the standalone full-support bound.
Section~\ref{sec:l3} checks the coefficients, relative levels and rough
factors of the first three switched rows on the actual source.

\begin{fact}[(F3) The spectral gap]\label{fact:kim-sarnak}
For every congruence subgroup $\Gamma\subset\mathrm{SL}_2(\mathbb Z)$,
\begin{equation}\label{eq:classical-spectral-constant}
 \lambda_1(\Gamma)\ge\frac{975}{4096},\qquad
 0<\lambda=\frac14-t^2<\frac14,\quad t\ge0\Longrightarrow
 0\le t\le\frac7{64},\quad \vartheta=2t\le\frac7{32}.
\end{equation}
Here $\lambda_1(\Gamma)$ is the smallest nonzero Laplace eigenvalue.
The implication applies to every eigenvalue $\lambda$ of $\Delta$
on $L^2(\Gamma\backslash\mathbb H)$ in the indicated range.
The first bound is Kim--Sarnak \cite[Appendix~2]{KimSarnak2003}.
The implications are elementary.
\end{fact}

\begin{fact}[(F4) Deligne's bound]\label{fact:F4-deligne}\label{fact:F5-deligne}
Let $f=\sum_{n\ge1}\psi_f(n)e(nz)$ be a primitive holomorphic cuspidal newform.
Assume $\psi_f(1)=1$, integer weight $k\ge2$, level $M$, and character $\chi$,
with $f((az+b)/(cz+d))=\chi(d)(cz+d)^kf(z)$ on $\Gamma_0(M)$.
This agrees with Deligne's convention $\chi(a)^{-1}$.
For $p\nmid M$, Deligne proved
\begin{equation}\label{eq:deligne-roots}
 |\psi_f(p)|\le2p^{(k-1)/2}
\end{equation}
\cite[Th\'eor\`eme~(8.2), p.~302]{Deligne1974}.
The same source states that the roots $\alpha_p,\beta_p$ of
$T^2-\psi_f(p)T+\chi(p)p^{k-1}$ both have modulus $p^{(k-1)/2}$.
Since $|\alpha_p|^2=|\beta_p|^2=p^{k-1}$, we have
$\overline{\alpha_p}=p^{k-1}/\alpha_p$ and likewise for $\beta_p$.
With $\alpha_p\beta_p=\chi(p)p^{k-1}$, this gives
\begin{equation}\label{eq:deligne-conjugation}
 \begin{aligned}
 \chi(p)\overline{\psi_f(p)}
 &=\frac{\alpha_p\beta_p}{p^{k-1}}
   \left(\frac{p^{k-1}}{\alpha_p}+\frac{p^{k-1}}{\beta_p}\right)
 =\psi_f(p).
 \end{aligned}
\end{equation}
Choose $\eta_p^2=\chi(p)$, $|\eta_p|=1$.
Then $b_p=\eta_p^{-1}p^{-(k-1)/2}\psi_f(p)$ is real and $|b_p|\le2$.
With $T=p^{(k-1)/2}\eta_p Z$, the polynomial becomes
\[
 T^2-\psi_f(p)T+\chi(p)p^{k-1}
       =p^{k-1}\chi(p)(Z^2-b_pZ+1).
\]
Write $b_p=2\cos\omega_p$, $0\le\omega_p\le\pi$.
The normalized roots are $e^{\pm i\omega_p}$,
also in the repeated case $b_p=\pm2$.
An earlier exposition of the modular-form setting is
\cite{DeligneBourbaki1971}.
\end{fact}

\begin{fact}[(F5) Averaged Kloosterman sums]\label{fact:F5-di82-thm14}\label{fact:F2-di82-thm14}
For $C,M,Y\ge1$ and $\epsilon>0$,
\begin{equation}\label{eq:F2-di82-thm14}
 \sum_{c\le C}\left|\sum_{m\le M}\sum_{n\le Y}S(m,n;c)\right|
 \ll_\epsilon(CMY)^\epsilon C(C+MY)
\end{equation}
\cite[Theorem~14, equation~(8.13), p.~275]{DI82}.
The sums run over integers $1\le c\le C$, $1\le m\le M$,
and $1\le n\le Y$. Here $S(m,n;c)=\sum_{u\ (c)}^*e((mu+n\bar u)/c)$.
\end{fact}

Use the trivial character on $\Gamma_0(q)\backslash\mathbb H$.
The sum over $f$ runs through a Petersson-orthonormal basis of
$S_k(\Gamma_0(q))$. The sum over $j$ runs through an orthonormal basis of
Maass cusp forms, including exceptional forms. The sum over $\mathfrak a$
runs through inequivalent cusps with width-one scaling matrices.
The measure is $dx\,dy/y^2$.
At infinity, the DI coefficients satisfy
\begin{equation}\label{eq:di-normalization}
 \begin{aligned}
 u_j(z)&=\sum_{n\ne0}\rho_j(n)\sqrt y
            K_{i\kappa_j}(2\pi|n|y)e(nx),\quad
 \lambda_j=\tfrac14+\kappa_j^2,\\
 f(z)&=\sum_{n\ge1}\psi_f(n)e(nz),\\
 E_{\mathfrak a}(z,s)&=\text{constant terms}
 +\sqrt y\sum_{n\ne0}\frac{2\pi^s|n|^{s-1/2}}{\Gamma(s)}
       \varphi_{\mathfrak a\infty n}(s)
       K_{s-1/2}(2\pi|n|y)e(nx).
 \end{aligned}
\end{equation}
For $0<\lambda_j<1/4$, set
$\vartheta_j=2|\Im\kappa_j|=2\sqrt{1/4-\lambda_j}$.
See Paper~I, (2.2)--(2.4) \cite[(2.2)--(2.4), p.~5]{PaperI}.

\begin{fact}[(F6) Spectral large sieves]\label{fact:di-large-sieves}
For $q,K\ge1$, $Y\ge1/2$, and $a_n$ supported on $n\sim Y$,
put
\[
 \begin{aligned}
 \Sigma_{\rm hol}(K,a)&=\sum_{\substack{2\le k\le K\\k\ {\rm even}}}
 \frac{(k-1)!}{(4\pi)^{k-1}}\sum_f
 \left|\sum_{n\sim Y}a_nn^{-(k-1)/2}\psi_f(n)\right|^2,\\
 \Sigma_{\rm Maa}(K,a)&=\sum_{|\kappa_j|\le K}\frac1{\cosh(\pi\kappa_j)}
            \left|\sum_{n\sim Y}a_n\rho_j(n)\right|^2,\\
 \Sigma_{\rm Eis}(K,a)&=\sum_{\mathfrak a}\int_{-K}^K
 \left|\sum_{n\sim Y}a_nn^{it}
       \varphi_{\mathfrak a\infty n}(1/2+it)\right|^2dt.
 \end{aligned}
\]
The large-sieve estimates are
\begin{equation}\label{eq:di-large-sieves}
 \begin{gathered}
 \Sigma_{\rm hol}(K,a),\Sigma_{\rm Maa}(K,a),\Sigma_{\rm Eis}(K,a)
       \ll_\epsilon(K^2+Y^{1+\epsilon}/q)\|a\|_2^2,\\
 \sum_{j\ {\rm exceptional}}V^{\vartheta_j}
       \left|\sum_{n\sim Y}a_n\rho_j(n)\right|^2
 \ll_\epsilon(1+\sqrt{YV/q})(1+\sqrt{Y^{1+\epsilon}/q})\|a\|_2^2,
 \quad V\ge1.
 \end{gathered}
\end{equation}
These are the specialized forms in Paper~I, (2.5)--(2.8)
\cite[Theorem~2, (1.28)--(1.30), p.~230;
Theorem~5, (1.38), p.~232]{DI82}
\cite[(2.5)--(2.8), pp.~6--7]{PaperI}.
\end{fact}

For $\phi\in C_c^\infty((0,\infty))$, define
\begin{equation}\label{eq:di-transforms}
 \begin{aligned}
 \widehat\phi(t)&=\frac\pi{\sinh\pi t}\int_0^\infty
 \frac{J_{2it}(x)-J_{-2it}(x)}{2i}\phi(x)\frac{dx}x,\\
 \check\phi(t)&=\frac4\pi\cosh\pi t\int_0^\infty
 K_{2it}(x)\phi(x)\frac{dx}x,\qquad
 \widetilde\phi(l)=\int_0^\infty J_l(x)\phi(x)\frac{dx}x.
 \end{aligned}
\end{equation}
The value at $t=0$ is the removable limit.

\begin{fact}[(F7) Corrected same-cusp trace formula]\label{fact:p1-corrected}
For $q,m,n\ge1$ and the test class above,
\begin{equation}\label{eq:p1-corrected}
 \sum_{h\ge1}\frac{S(m,n;qh)}{qh}
 \phi\left(\frac{4\pi\sqrt{mn}}{qh}\right)
       =\frac{H_\phi(m,n)}\pi-M_\phi(m,n)-\frac{E_\phi(m,n)}\pi,
\end{equation}
where
\begin{align*}
 H_\phi(m,n)&=\sum_{\substack{k\ge2\\k\ {\rm even}}}
 \frac{i^k(k-1)!\widetilde\phi(k-1)}{(4\pi\sqrt{mn})^{k-1}}
       \sum_f\overline{\psi_f(m)}\psi_f(n),\\
 M_\phi(m,n)&=\sum_j\frac{\widehat\phi(\kappa_j)}{\cosh\pi\kappa_j}
                  \overline{\rho_j(m)}\rho_j(n),\\
 E_\phi(m,n)&=\sum_{\mathfrak a}\int_{\mathbb R}(n/m)^{it}
 \overline{\varphi_{\mathfrak a\infty m}(1/2+it)}
 \varphi_{\mathfrak a\infty n}(1/2+it)\widehat\phi(t)\,dt.
\end{align*}
No extra diagonal term occurs, also for $m=n$.
This is the corrected same-cusp specialization proved in Paper~I,
(2.7.T15)--(2.7.T18) \cite[(2.7.T15)--(2.7.T18)]{PaperI}. Compare
\cite[Lemma~4.5, (4.8)--(4.17)]{Drappeau2017}, in the DI convention
\cite[(1.15)--(1.17)]{DI82}.
\end{fact}

\begin{fact}[(F8) Mixed-cusp formulas]\label{fact:mixed-cusps}
For $q=rs$, $(r,s)=1$, use the trivial character and
\[
 \sigma_{1/s}=\begin{pmatrix}\sqrt r&0\\s\sqrt r&1/\sqrt r\end{pmatrix},
 \quad\mathcal C(\infty,1/s)=\{cs\sqrt r:c\ge1,\ (c,r)=1\}.
\]
DI82 (1.6), p.~224, gives the arithmetic identity
\cite[(1.6), p.~224]{DI82}
\[
 S_{\infty,1/s}(m,n;cs\sqrt r)
       =e(n\bar s/r)S(m\bar r,n;cs),\qquad m,n\in\mathbb N.
\]
Here $s\bar s\equiv1\pmod r$ and $r\bar r\equiv1\pmod{cs}$.
The accompanying trace formulas, proved in Paper~I, are
\[
 \mathcal K_+=H_{\infty,1/s}/\pi-M_{\infty,1/s}-E_{\infty,1/s}/\pi,
 \qquad
 \mathcal K_-=(\pi/2)M^-_{\infty,1/s}+(1/2)E^-_{\infty,1/s}.
\]
Here $\mathcal K_\pm=\sum_{\gamma\in\mathcal C}\gamma^{-1}
S_{\infty,1/s}(m,\pm n;\gamma)\phi(4\pi\sqrt{mn}/\gamma)$.
The plus terms have the definitions in (F7), with their second
coefficient at $1/s$.
For $M^-$ and $E^-$, replace $\widehat\phi$ by $\check\phi$ and $n$ by $-n$
only in the second Fourier coefficient.
The cusps are singular, and the scaling matrices have width one.
The classical width of $1/s$ is $r$.
For $r>1$ they are inequivalent.
The normalization is that of Paper~I, (2.13)--(2.14). Compare
\cite[Lemmas~4.3, 4.5, (4.7)--(4.19)]{Drappeau2017}
\cite[(2.13)--(2.14), p.~8]{PaperI}.
\end{fact}

\begin{fact}[(F9) Bessel representations]\label{fact:F9-dlmf}\label{fact:F7-dlmf}
For $x>0$, $n\in\mathbb Z$, and $|\Re\nu|<1$,
\begin{align*}
 J_n(x)&=\pi^{-1}\int_0^\pi\cos(x\sin u-nu)\,du\\
       &\quad=i^{-n}\pi^{-1}\int_0^\pi e^{ix\cos u}\cos(nu)\,du,\\
 J_\nu(x)&=(2/\pi)\int_0^\infty
       \sin(x\cosh u-\pi\nu/2)\cosh(\nu u)\,du,\\
 Y_\nu(x)&=-(2/\pi)\int_0^\infty
       \cos(x\cosh u-\pi\nu/2)\cosh(\nu u)\,du,\\
 K_\nu(x)&=\sec(\pi\nu/2)\int_0^\infty
       \cos(x\sinh u)\cosh(\nu u)\,du.
\end{align*}
The infinite integrals are improper.
These are DLMF (10.9.2), (10.9.8), and (10.32.7)
\cite[(10.9.2), (10.9.8), (10.32.7)]{DLMF}.
Paper~I uses them on printed p.~20, (2.7.T31)--(2.7.T32).
\end{fact}

\begin{fact}[(F10) Kontorovich--Lebedev inversion]\label{fact:kl-inversion}
For $g$ on $(0,\infty)$, the inversion pair is
\[
 f(y)=\frac{2y\sinh\pi y}{\pi^2}\int_0^\infty g(x)K_{iy}(x)\frac{dx}x,
 \qquad g(x)=\int_0^\infty f(y)K_{iy}(x)\,dy.
\]
It holds under either of two conditions.
(a) $g/x\in C^1$, with $xg$ and $x(g/x)'$ absolutely integrable.
(b) Local piecewise continuity and bounded variation, with
\[
 \int_0^{1/2}\frac{|g(x)|}x\log(1/x)\,dx<\infty,\qquad
 \int_{1/2}^\infty\frac{|g(x)|}{\sqrt x}\,dx<\infty.
\]
Under condition (b), the inverse equality is asserted at continuity points of $g$.
In our applications, the inverse equality is used only at continuity points
of $g$.
See DLMF (10.43.30)--(10.43.32)
\cite[(10.43.30)--(10.43.32)]{DLMF}.
These hypotheses hold for $g=\phi/x$, $\phi\in C_c^\infty((0,\infty))$.
They also hold for $g=xK_{2it}(x)$ when $|\Im t|\le7/64$.
Near zero, $K_{2it}(x)\ll_t x^{-2|\Im t|}(1+|\log x|)$,
so the first integral in condition (b) converges absolutely.
At infinity, $K_{2it}(x)\ll_t x^{-1/2}e^{-x}$, so the second converges too.
Both functions are smooth on compact subintervals.
\end{fact}

\begin{fact}[(F11) Hecke operators on the full space]\label{fact:hecke-full-space}
For trivial character and $(k,q)=1$, Paper~I proves
\[
 \|T_k\|_{\mathrm{Maass},q}\le d(k)k^{7/64},\qquad
 \|T_k\|_{\mathrm{hol},q}\le d(k).
\]
Their proof uses the orthogonal cuspidal representation
decomposition of Jacquet--Langlands, Chapter~III, \S16
\cite[Chapter~III, \S16]{JacquetLanglands1970}.
It also uses the holomorphic oldform classes of Atkin--Lehner
\cite[Lemmas~6, 13--15; Theorems~3(i), 5]{AtkinLehner1970}.
Atkin--Lehner's weight $k$ means classical weight $2k$.
For an unramified normalized newform, the algebra needed is
\[
 \lambda_f(p^r)=\sum_{j=0}^r\alpha_p^{r-j}\beta_p^j,\qquad
 \lambda_f(mn)=\lambda_f(m)\lambda_f(n)\quad((m,n)=1).
\]
The Maass Satake bound is $|\alpha_p|,|\beta_p|\le p^{7/64}$
\cite[Appendix~2]{KimSarnak2003}.
The holomorphic bound is (F4). Full-space diagonalization and unitary
transport give the displayed norms, as in Paper~I, (2.7.T37)--(2.7.T40)
\cite[(2.7.T37)--(2.7.T40), pp.~21--22]{PaperI}.
\end{fact}

\begin{fact}[(F12) Eisenstein inputs]\label{fact:eisenstein-inputs}
Paper~I uses three classical Eisenstein properties.
The Eisenstein family has a meromorphic continuation, which allows compact
Fourier-coefficient extraction. It has an arithmetic Fourier normalization.
There is a complete irredundant cusp set.
For the cusp $u/w$, let $d_{\rm den}$ be the lower-row denominator variable
in $cz+d_{\rm den}$ and write the printed variable as $d=-d_{\rm den}$.
The printed denominator is therefore $cz-d$, and the Fourier phase changes
with this sign choice. Thus replacing $d$ by $-d$ changes the congruence and
the phase together, not one sign in isolation.
The normalization includes the orbit conditions
\[
 (c,d)=1,\quad(c,q)=w,\quad cd\equiv uw\pmod{(w^2,q)}.
\]
With $c=w\gamma$, the last condition is
$\gamma d\equiv u\pmod{(w,q/w)}$. The cusp set has the form
\[
 u/w,\quad w\mid q,\quad u\pmod{(w,q/w)},\quad
 (u,(w,q/w))=1,
\]
with admissible lifts chosen coprime to $q$.
The sources are Iwaniec, Chapter~6 and Theorem~3.4;
DI82 (1.17), Lemmas~2.3 and~3.6;
and Kiral--Young, Proposition~3.1 and Corollary~3.2
\cite[Chapter~6; Theorem~3.4]{IwaniecSpectralMethods2002}
\cite[(1.17), p.~227; Lemma~2.3, p.~240; Lemma~3.6, p.~247]{DI82}
\cite[Proposition~3.1, Corollary~3.2]{KiralYoung2017}.
These are the conventions used in Paper~I, (2.7.T20)--(2.7.T24)
and (2.7.T41)--(2.7.T43)
\cite[(2.7.T20)--(2.7.T24), (2.7.T41)--(2.7.T43), pp.~18--19, 22--23]{PaperI}.
\end{fact}

\begin{fact}[(F13) Further spectral estimates]\label{fact:remaining-spectral-inputs}
For exceptional eigenfunctions, the first line below follows from (F3), since
$975/4096>768/4096=3/16$ and $7/32<1/2$.
Paper~I derives the interval estimate from (F5), (F6), (F7),
(P1.1) and this first line.
The last line is its supported imaginary-transform estimate:
\begin{equation}\label{eq:remaining-spectral-inputs}
 \begin{gathered}
 \lambda_j\ge3/16,\qquad 0<\vartheta_j\le1/2,\\
 \sum_{2\le q\le Q}\sum_{j\ {\rm exceptional}}^{(q)}V^{\vartheta_j}
       \left|\sum_{n\in I}\rho_j(n)\right|^2
 \ll_\epsilon(QY)^\epsilon(Q+Y+\sqrt{YV})Y,\\
 Q,Y,V\ge1,\qquad I\subset(Y,2Y]\cap\mathbb Z\quad\text{an integer interval},\\
 |\widehat\phi(i\nu)|+|\check\phi(i\nu)|
 \ll\mathfrak B(1+|\log Z|)(1+Z^{-2\nu})/(1+Z),
 \quad0<\nu\le1/4.
 \end{gathered}
\end{equation}
The last line uses the supported test class of (P1.1).
The interval argument uses DI82's proof of Theorem~7, (8.17),
with the corrected same-cusp normalization (F7)
\cite[proof of Theorem~7, (8.17)]{DI82}
\cite[(2.9), (2.7.T32), pp.~6, 20]{PaperI}.
The weaker gap $3/16$ is also DI82, Theorem~4, p.~232,
which we do not use separately.
The reconstruction of the last estimate uses (F9).
\end{fact}

\begin{fact}[(F14) Transform identities and interchanges]\label{fact:transform-interchanges}
Paper~I proves the following identities and accompanying
interchange bounds in (2.7.T26)--(2.7.T32).
Use the transforms \eqref{eq:di-transforms} without rescaling.
For $\phi\in C_c^\infty((0,\infty))$, put
\begin{equation}\label{eq:spectral-transform-applicability}
 \begin{gathered}
 \mathcal R_\phi(v)=\int_0^\infty K_{2iv}(x)\phi(x)x^{-2}\,dx,\qquad
 k_\phi(v)=\frac4{\pi^2}v\sinh(2\pi v)\mathcal R_\phi(v),\\
 \mathcal H(t,v)=\frac{\cosh\pi t}
   {\cosh\pi(t-v)\cosh\pi(t+v)},\\
 \int_{\mathbb R}k_\phi(v)K_{2iv}(x)\,dv=\phi(x)/x,\\
 \int_{\mathbb R}v\sinh(2\pi v)\mathcal H(t,v)K_{2iv}(x)\,dv
       =\cosh(\pi t)xK_{2it}(x).
 \end{gathered}
\end{equation}
By (F10), the KL hypotheses hold for both applications,
to $\phi(x)/x$ and to $xK_{2it}(x)$.
For the latter, the product-integral condition is
$|\Re(2it)|+|\Re(2iv)|<1$, with $v\in\mathbb R$.
The permitted $t$ are real or exceptional, $|\Im t|\le7/64$.
The differentiations and contour, limit, sum, and integral interchanges
in (2.7.T26)--(2.7.T32) hold in these ranges.
This includes absolute summability after the spectral estimates.
For each fixed smooth test family, the decay bound is
\[
 \mathcal R_\phi(v)\ll_{\phi,j}e^{-\pi|v|}(1+|v|)^{1-2j}
 \qquad(j\ge0).
\]
Let $\operatorname{supp}\phi\subset[Z/21,9Z]$.
The integrations by parts and differentiated estimates retain the
five-derivative seminorm $\mathfrak B$ of (P1.1), uniformly in every
auxiliary parameter.
For the following bounds $U>0$, and the oscillatory kernels are the
normalized combinations in \eqref{eq:di-transforms}; these estimates are
not assertions about an unnormalized $J_{2it}$ tail. The hyperbolic factor
in the raw Bessel representation is cancelled only after this normalization.
For truncated normalized oscillatory representations, the tails are bounded by
\[
 O\!\left(\frac{1+|t|}{x\sinh U}\right)\quad(t\in\mathbb R),\qquad
 O\!\left(\frac{(1+\nu)e^{2\nu U}}{x\sinh U}\right)
       \quad(t=i\nu,\ 0<\nu\le1/4).
\]
These tails permit the stated limiting interchanges.
The derivative and decay constants are uniform in the families
consumed by Theorem~2.8, with the losses prescribed there.
The first inversion identity follows from (F10).
The product integral, second kernel identity, and uniform
differentiation and interchange steps are proved in the same passage
\cite[(2.7.T26)--(2.7.T32), pp.~19--20]{PaperI}.
A reconstruction from the individual Bessel representations uses (F9).
\end{fact}

\begin{fact}[(P1.1) Supported real transforms]\label{fact:real-transforms}
\leavevmode\par
For $Z>0$, let $\operatorname{supp}\phi\subset[Z/21,9Z]$ and set
$\mathfrak B=\max_{0\le j\le5}\sup_{x>0}x^j|\phi^{(j)}(x)|$.
Then
\begin{equation}\label{eq:real-transform-bound}
 \begin{aligned}
 |\widehat\phi(t)|+|\check\phi(t)|
 &\ll\frac{\mathfrak B(1+|\log Z|)}{1+Z}
       \min\{1,((1+Z)/|t|)^3\},\quad t\in\mathbb R,\\
 |\widetilde\phi(l)|
 &\ll\frac{\mathfrak B(1+|\log Z|)}{1+Z}
       \min\{1,((1+Z)/l)^3\},\quad l\in\mathbb N.
 \end{aligned}
\end{equation}
The constants are absolute. The minimum at $t=0$ is $1$.
This is the real-parameter part of Paper~I, (2.7.T31)
\cite[(2.12), (2.7.T31), pp.~8, 20]{PaperI}.
The real-transform bound uses exactly the Bessel representations
in (F9). It does not cover the imaginary-parameter or interchange steps.
\end{fact}

Write $\gamma_{\!E}$ for Euler's constant. The continuous linear-sieve
functions are normalized by
\begin{equation}\label{eq:classical-delay}
 \begin{gathered}
 sF(s)=2e^{\gamma_{\!E}}\quad(0<s\le3),\qquad
 sf(s)=0\quad(0<s\le2),\\
 (sF(s))'=f(s-1)\quad(s>3),\qquad
 (sf(s))'=F(s-1)\quad(s>2).
 \end{gathered}
\end{equation}
Thus $F(1)=2e^{\gamma_{\!E}}$ and $f(s)=0$ for $0<s\le2$.
This is Paper~I, (2.23), p.~10 \cite[(2.23), p.~10]{PaperI}.

\begin{fact}[(P1.2) Structured linear sieve]\label{fact:structured-sieve}
Let $\mathcal A$ be a finite integer sequence with nonnegative
multiplicities, let $\mathcal P$ be a set of primes, and let $\mathcal X>0$.
For a multiplicative function $\omega$, set
\begin{equation}\label{eq:classical-sieve-data}
 \begin{gathered}
 P(z)=\prod_{\substack{p<z\\p\in\mathcal P}}p,\qquad
 S(\mathcal A,\mathcal P,z)=\sum_{a\in\mathcal A}1_{(a,P(z))=1},\\
 |\mathcal A_d|=\frac{\omega(d)}d\mathcal X+r(\mathcal A,d),\qquad
 V(z)=\prod_{\substack{p<z\\p\in\mathcal P}}
                       \left(1-\frac{\omega(p)}p\right).
 \end{gathered}
\end{equation}
Assume $0<\omega(p)<p$ on $\mathcal P$, $\omega(p)=0$ off $\mathcal P$,
and $\omega(p^a)\ge0$. Suppose fixed $K\ge2$ and $L_\omega>1$ satisfy,
for every $z>w\ge2$,
\begin{equation}\label{eq:classical-density}
 \begin{gathered}
 \prod_{\substack{w\le p<z\\p\in\mathcal P}}
 \left(1-\frac{\omega(p)}p\right)^{-1}
 \le\frac{\log z}{\log w}\left(1+\frac K{\log w}\right),\\
 \sum_{\substack{w\le p<z\\p\in\mathcal P}}
       \sum_{a\ge2}\frac{\omega(p^a)}{p^a}
 \le\frac{L_\omega}{\log(3w)}.
 \end{gathered}
\end{equation}
Fix $D\ge2$, $0<\epsilon_s<1/3$ and $\delta_s=\epsilon_s^9$.
Define
\begin{equation}\label{eq:classical-grid}
 \mathcal G(D)=\{D^{\epsilon_s^2(1+\delta_s)^n}:n\ge0\}.
\end{equation}
Let $\mathcal H(D)$ consist of the finite tuples
$\mathbf D=(D_1,\ldots,D_k)$ from this grid such that
$D_k\le\cdots\le D_1<D^{1/2}$. Set
\begin{equation}\label{eq:classical-tail-tuples}
 \begin{aligned}
 \mathcal D^+(D)&=\{\mathbf D\in\mathcal H(D):
 D_1\cdots D_{2l}D_{2l+1}^{3}<D\quad
             (0\le l\le(k-1)/2)\},\\
 \mathcal D^-(D)&=\{\mathbf D\in\mathcal H(D):
 D_1\cdots D_{2l-1}D_{2l}^{3}<D\quad
             (1\le l\le k/2)\}.
 \end{aligned}
\end{equation}
Each of these sets has fewer than $\exp(8\epsilon_s^{-3})$ elements.
There are coefficients $\phi_d^\nu$ and $\lambda_d^\nu(\mathbf D)$,
$\nu\in\{+,-\}$, with $|\phi_d^\nu|,|\lambda_d^\nu(\mathbf D)|\le1$, independent of
$\mathcal A,\mathcal X,\omega$, for which
\begin{equation}\label{eq:classical-structured-sieve}
 \begin{aligned}
 S(\mathcal A,\mathcal P,z)&\le \mathcal X V(z)\{F(s)+E_s\}+R_1^++R^+
                   &&(2\le z\le D),\\
 S(\mathcal A,\mathcal P,z)&\ge \mathcal X V(z)\{f(s)-E_s\}-R_1^--R^-
                   &&(2\le z\le D^{1/2}),\\
 s&=\frac{\log D}{\log z},\qquad
 E_s\ll\epsilon_s+\epsilon_s^{-8}e^{K+L_\omega}(\log D)^{-1/3}.
 \end{aligned}
\end{equation}
The implied constant in the last line is absolute. The signed remainders are
\begin{equation}\label{eq:classical-smooth-remainder}
 R_1^\nu=\sum_{\substack{d<D^{\epsilon_s}\\
                        d\mid P(D^{\epsilon_s^2})}}
                  \phi_d^\nu r(\mathcal A,d)
\end{equation}
and
\begin{equation}\label{eq:classical-tail-remainder}
 R^\nu=\sum_{\mathbf D\in\mathcal D^\nu(D)}
 \sum_{\substack{d<D^{\epsilon_s}\\d\mid P(D^{\epsilon_s^2})}}
 \lambda_d^\nu(\mathbf D)
 \sum_{\substack{D_i\le p_i<D_i^{1+\delta_s}\\p_i\mid P(z)\\
                 p_i\ne p_j\ ((i,j)\in I(\mathbf D))}}
 r(\mathcal A,dp_1\cdots p_k).
\end{equation}
Here $I(\mathbf D)$ imposes distinctness between the equal-label
positions in Iwaniec's construction. Unequal grid intervals are disjoint.
The coefficients may depend on $\mathcal P,D,z,\epsilon_s,\mathbf D$
and these distinctness conditions.
We do not claim that coefficients at different cutoffs agree. See Paper~I, Theorem~2.6,
(2.22), (2.24), (2.27)--(2.33), pp.~10--12
\cite[Theorem~2.6, (2.22), (2.24), (2.27)--(2.33), pp.~10--12]{PaperI}.
Its classical source is Iwaniec, Theorem~4
\cite[Theorem~4]{Iwaniec1980Remainder}.
\end{fact}

Below, the prime set is $\mathcal P_N=\{p:p\nmid2N\}$.
Take the nominal density $\omega(p)=1$ on $\mathcal P_N$,
$\omega(p)=0$ elsewhere and $\omega(p^a)=0$ for $a\ge2$. Then the
second sum in \eqref{eq:classical-density} is zero. Mertens' product
estimate gives one fixed absolute $K_{\mathrm M}\ge2$ with
\begin{equation}\label{eq:classical-mertens-density}
 \prod_{w\le p<z}(1-1/p)^{-1}
 \le\frac{\log z}{\log w}\left(1+\frac{K_{\mathrm M}}{\log w}\right)
 \qquad(z>w\ge2).
\end{equation}
The first product in \eqref{eq:classical-density} is no larger than this
product. Hence $K=K_{\mathrm M}$, $L_\omega=2$ are valid for
this nominal density. The factor in
\eqref{eq:classical-structured-sieve} is then $e^{K_{\mathrm M}+2}$.
See \cite[pp.~10--11]{PaperI}. Each actual sequence still needs
its own decomposition \eqref{eq:classical-sieve-data} and error estimate.

For every relative level $L_c\ge2$ we use, exactly as in Paper~I,
(2.33a)--(2.33b), p.~12 \cite[(2.33a)--(2.33b), p.~12]{PaperI},
\begin{equation}\label{eq:classical-relative-grid}
 \begin{gathered}
 0<\sigma\le10^{-10},\qquad \tau=\sigma^9,\qquad
 \epsilon_s=\sigma,\quad\delta_s=\tau,\quad D=L_c,\\
 \mathcal G(L_c)=\{L_c^{\sigma^2(1+\tau)^n}:n\ge0\},\qquad
 d<L_c^\sigma,\quad d\mid P(L_c^{\sigma^2}).
 \end{gathered}
\end{equation}
Thus the prime intervals, the smooth-divisor support and the level all
refer to the same $L_c$. First fix $\sigma>0$. Then let $N$ tend
to infinity.


\begin{fact}[(P1.3) Finite factorization]\label{fact:hb}
Let $X\ge2$, $K\in\mathbb N$, and $z\ge1$.
Assume $z^K>X$.
For every integer $1\le n\le X$,
\begin{equation}\label{eq:classical-hb}
 \Lambda(n)=\sum_{j=1}^{K}(-1)^{j-1}\binom Kj
 \sum_{\substack{m_1\cdots m_jn_1\cdots n_j=n\\m_i\le z}}
 \mu_{\rm Mob}(m_1)\cdots\mu_{\rm Mob}(m_j)\log n_1.
\end{equation}
This is Paper~I, Lemma~2.7, (2.36), p.~13
\cite[Lemma~2.7, (2.36), p.~13]{PaperI}.
\end{fact}

\begin{proof}
For $\Re s>1$, set $M_z(s)=\sum_{m\le z}\mu_{\rm Mob}(m)m^{-s}$
and $Y_z(s)=\zeta(s)M_z(s)$. The polynomial identity
$1-(1-Y_z)^K=\sum_{j=1}^K(-1)^{j-1}\binom KjY_z^j$ gives
\[
 -\frac{\zeta'}\zeta=
 \sum_{j=1}^K(-1)^{j-1}\binom Kj(-\zeta')\zeta^{j-1}M_z^j
 +\left(-\frac{\zeta'}\zeta\right)(1-Y_z)^K.
\]
Every nonzero coefficient of $Y_z-1$ has index greater than $z$.
The last term therefore has no coefficient at $n\le X<z^K$;
coefficient extraction proves \eqref{eq:classical-hb}.
\end{proof}

For $n\le2X_m$, choose $K=3$, $z=2X_m^{1/3}$.
Then $z^3=8X_m>2X_m$.
Any interval partition is imposed after \eqref{eq:classical-hb}.
The product condition and all representation multiplicities remain in
the resulting coefficients. See \cite[(2.36)--(2.37), p.~13]{PaperI}.

For the packing theorem, put $\mu_*=25/128$ and define the packing ceiling by
\[
 H_{\rm pk,max}=\frac{153}{256}.
\]
Use this constant in the packing hypotheses below. Put
\begin{equation}\label{eq:classical-packing-functions}
 \begin{gathered}
 v_{\rm pk}(x)=\min\left\{\frac{1+x}{2},\frac{32-34x}{39},
                         \frac{57}{64}-\frac32x\right\}\quad(x\ge0),\\
 \psi_{\rm pk}(x,y)=x1_{x\ge y}\quad(x\in\mathbb R,\ y\ge0),\\
 w_{\rm pk}(x_1,x_2,x_3)=\psi_{\rm pk}\left(\min\left\{
 \frac{132-68x_3}{214},\frac{153-96x_3}{256},1-2x_2\right\},\frac{1+x_1}{2}\right).
 \end{gathered}
\end{equation}
For $x_1\ge x_2\ge x_3\ge0$, define
\begin{equation}\label{eq:classical-packing-ceiling}
 \begin{aligned}
 H(x_1,x_2,x_3)=\max\{&v_{\rm pk}(x_1),v_{\rm pk}(x_2),
 v_{\rm pk}(x_1+x_2),v_{\rm pk}(x_1+x_2+x_3),\\
 &w_{\rm pk}(x_1,x_2,x_3),w_{\rm pk}(x_1,x_3,x_2),\\
 &\psi_{\rm pk}(v_{\rm pk}(x_1+x_3),x_1+2x_2+x_3),\\
 &\psi_{\rm pk}(v_{\rm pk}(x_2+x_3),2x_1+x_2+x_3)\},\\
 \widetilde H(x_1,x_2,x_3)&=
 \begin{cases}H(x_1,x_2,x_3),&x_1\le\mu_*,\\v_{\rm pk}(x_1),&x_1>\mu_*.
 \end{cases}
 \end{aligned}
\end{equation}
These are the definitions preceding Paper~I, Theorem~4.1, p.~53
\cite[p.~53]{PaperI}.

\begin{fact}[(P1.4) Three-factor packing]\label{fact:packing}
Let $0<\delta_{\rm pk}\le3\cdot10^{-9}$,
$1/2\le h\le H_{\rm pk,max}$, $d_*=h-2\delta_{\rm pk}$, and
$\delta_{\rm pk}\le g\le1/3-\delta_{\rm pk}/2$.
Let $e_1\ge\cdots\ge e_m\ge0$ be a finite labelled list.
Equal values keep distinct labels. The empty list is allowed.
Extend the list by zeros, writing $\bar e_i=0$ for $i>m$.
Suppose
\begin{equation}\label{eq:classical-packing-hypotheses}
 \sum_{i<j}e_i+2e_j\le h-2\delta_{\rm pk}\quad(1\le j\le m),\qquad
 h\le\widetilde H(\bar e_1,\bar e_2,\bar e_3).
\end{equation}
The rule of Paper~I, Section~4.1, partitions the labels as
$J_1\sqcup J_2\sqcup J_3$. With
$\Sigma_s=\sum_{i\in J_s}e_i$ and $\theta_{\rm sp}=7/32$,
\begin{equation}\label{eq:classical-packing-loads}
 \begin{aligned}
 \Sigma_1&\le g,\\
 2g+\Sigma_2+2\Sigma_3&\le1-4\delta_{\rm pk},\\
 \theta_{\rm sp}+g+\Sigma_1+5\Sigma_2+2\Sigma_3&\le2-4\delta_{\rm pk},\\
 (1+2\theta_{\rm sp})g+\Sigma_1+
 (5-2\theta_{\rm sp})\Sigma_2+
 (2+2\theta_{\rm sp})\Sigma_3
 &\le2-4(1+\theta_{\rm sp})\delta_{\rm pk}.
 \end{aligned}
\end{equation}
The partition depends only on these data.
This is Paper~I, Theorem~4.1, (4.1)--(4.3), pp.~53--54
\cite[Theorem~4.1, (4.1)--(4.3), pp.~53--54]{PaperI}.
\end{fact}


\begin{fact}[(P1.5) Continuous Hecke comparison]\label{fact:eisenstein-hecke}
Assume the three Eisenstein statements in (F12).
Let $w\mid q$, put $w'=(w,q/w)$, and let $\chi$ be a Dirichlet
character modulo $w'$.
Then Paper~I, (2.7.T41)--(2.7.T43), gives
\begin{equation}\label{eq:eisenstein-arithmetic-coefficient}
 \begin{aligned}
 \varphi_{u/w,\infty,n}(z)
  &=(w'/(wq))^z\sum_{\substack{\gamma\ge1\\(\gamma,q/w)=1}}\gamma^{-2z}
 \sum_{\substack{d\ (\gamma w)\\(d,\gamma w)=1\\\gamma d\equiv u\ (w')}}
                      e(-nd/(\gamma w)),\\
 \mathcal V_\chi(n;z)&=\sum_{u\ (w')}^*\chi(u)
                         \varphi_{u/w,\infty,n}(z),\\
 \mathcal V_\chi(kn';z)
   &=\overline{\chi(k)}\left(\sum_{d\mid k}\chi(d)^2d^{1-2z}\right)
                           \mathcal V_\chi(n';z),
       \quad(k,q)=(k,n')=1.
 \end{aligned}
\end{equation}
The first identity initially holds for $\Re z>1$.
The last extends meromorphically.
For finitely supported $v_{n'}$ in this coprime range,
\begin{equation}\label{eq:eisenstein-hecke-loss}
 \begin{aligned}
 &\sum_{\mathfrak a}\left|\sum_{n'}v_{n'}|kn'|^{it}
       \varphi_{\mathfrak a\infty,kn'}(1/2+it)\right|^2\\
 &\qquad\le d(k)^2\sum_{\mathfrak a}\left|\sum_{n'}v_{n'}|n'|^{it}
       \varphi_{\mathfrak a\infty,n'}(1/2+it)\right|^2.
 \end{aligned}
\end{equation}
It holds almost everywhere in real $t$, for either index sign
\cite[(2.7.T41)--(2.7.T43), pp.~22--23]{PaperI}.
\end{fact}

\begin{proof}
In \eqref{eq:eisenstein-arithmetic-coefficient}, write the modulus
$c_1=w\gamma=c_1'c_1''$, where $(c_1',q)=1$ and every prime of
$c_1''$ divides $q$. CRT factors the inner character sum as
$\chi(c_1')G_\chi(-n;c_1'')c_{c_1'}(n)$. Here $G_\chi$ is the
corresponding finite Gauss sum and $c_{c_1'}$ is the Ramanujan sum.
The outer character supplies another $\chi(c_1')$. The coprime factor is
\[
 L^{(q)}(2z,\chi^2)^{-1}
 \sum_{\substack{d\mid |n|\\(d,q)=1}}\chi^2(d)d^{1-2z},
\]
where $L^{(q)}$ omits Euler factors at primes dividing $q$.
For $(k,q)=(k,n')=1$, its divisor sum factors at $n=kn'$.
Changing the Gauss-sum variable supplies $\overline{\chi(k)}$.
Absolute convergence for $\Re z>1$, followed by meromorphic continuation,
proves the last identity in \eqref{eq:eisenstein-arithmetic-coefficient}.
At $z=1/2+it$, the character multiplier is
\[
 \overline{\chi(k)}\sum_{d\mid k}\chi(d)^2(k/d^2)^{it}.
\]
Its modulus is at most $d(k)$.
Orthogonality over the complete cusp set proves the inequality.
\end{proof}

\begin{fact}[(P1.6) Exceptional parameters on the full space]\label{fact:full-space-spectral-gap}
Apply (F3) to $\Gamma_0(q)$.
Every exceptional eigenfunction in the full $\Gamma_0(q)$ space then has
\[
 \vartheta_j=2\sqrt{1/4-\lambda_j}\le7/32.
\]
The bound depends only on the eigenvalue, so it covers oldforms.
These are the parameters in Paper~I, (2.7.T32), (2.7.T52), and (2.7.T56)
\cite[(2.7.T32), (2.7.T52), (2.7.T56), pp.~20, 24--25]{PaperI}.
\end{fact}

\begin{fact}[(P1.7) The spectral estimate]\label{fact:spectral-implication}
Paper~I, Theorem~2.8, proves the estimate below. Its proof uses
the large sieves (F6) and the trace formulas (F7)--(F8). It also uses the transform
bounds (P1.1), (F13)--(F14), the full-space Hecke bounds (F11), and
the continuous and exceptional comparisons (P1.5)--(P1.6).
The normalized unitary cusp transport, proved in Paper~I,
(2.7.T33)--(2.7.T40), is as follows. Choose $rw-sy=1$ and use
\begin{equation}\label{eq:unitary-cusp-transport}
 \begin{gathered}
 W_{rs}=r^{-1/2}\begin{pmatrix}r&y\\rs&rw\end{pmatrix},\qquad
 W_{rs}\Gamma_0(rs)W_{rs}^{-1}=\Gamma_0(rs),\\
 \rho_{F,1/s}(n)=e(n\bar s/r)\rho_{F\circ W_{rs},\infty}(n),\qquad
 \psi_{F,1/s}(n)=e(n\bar s/r)\psi_{F|W_{rs},\infty}(n),\\
 \varphi_{\mathfrak a,1/s,n}(z)
   =e(n\bar s/r)\varphi_{W_{rs}^{-1}\mathfrak a,\infty,n}(z),\\
 T_kF(z)=k^{-1/2}\sum_{hd=k}\sum_{j\ (d)}F((hz+j)/d),\qquad
 c_{T_kF}(n)=\sum_{h\mid(k,n)}c_F(kn/h^2).
 \end{gathered}
\end{equation}
The last line is the Maass normalization, with $(k,rs)=1$.
For holomorphic forms use the normalized coefficients of (F11).
These formulas fix the normalization of the following theorem.

Let $C,D,Y,R,S,A_T\ge1$ and $0<\epsilon_0<1/100$.
Let $a_T\in\mathbb Z$ satisfy $0<|a_T|\le A_T$.
Let $b_{n,r,s}$ be supported on
\begin{equation}\label{eq:t28-support}
 0<n\le Y,\qquad r\sim R,\qquad s\sim S,\qquad (a_T,rs)=1.
\end{equation}
Let $g$ be smooth, supported in
$(C,2C]\times(D,2D]\times\mathbb R_{>0}^3$, with
\begin{equation}\label{eq:t28-seminorms}
 |\partial^\nu g(c,d,n,r,s)|
 \le G_\nu(c^{-\nu_1}d^{-\nu_2}n^{-\nu_3}r^{-\nu_4}s^{-\nu_5})^{1-\epsilon_0}
 \quad(\nu\in\mathbb Z_{\ge0}^5).
\end{equation}
Set $\mathcal X=2+A_TCDYRS$ and $\theta_{\rm sp}=7/32$.
For every $\epsilon>0$,
\begin{equation}\label{eq:t28-conclusion}
 \left|\sum_{(r,s)=1}\sum_{0<n\le Y}b_{n,r,s}
 \sum_{\substack{c,d\ge1\\(rd,sc)=1}}
 g(c,d,n,r,s)e\left(\frac{a_T n\overline{rd}}{sc}\right)\right|
 \ll_{\epsilon,\epsilon_0,\{G_\nu\}}
       \mathcal X^{\epsilon+18\epsilon_0}J_*(a_T),
\end{equation}
where the inverse is modulo $sc$, and
\begin{equation}\label{eq:t28-main-bound}
 \begin{aligned}
 J_*(a_T)^2={}&\sum_{\substack{n''\mid a_T^\infty\\n''\le2Y}}
 A_T^{\theta_{\rm sp}}(n'')^{\theta_{\rm sp}}\\
 &\quad\times\left[CS\left(\frac Y{n''}+RS\right)(C+DR)+|a_T|YRS\right]
                 \|\widetilde b(n'')\|_2^2\\
 &+\left[CS(CD\sqrt R)^{2\theta_{\rm sp}}(Y+RS)^{1-\theta_{\rm sp}}\right.\\
 &\qquad\left.\times(C+DR)^{1-2\theta_{\rm sp}}+D^2YR\right]\|b\|_2^2,\\
 \|b\|_2^2={}&\sum_{n,r,s}|b_{n,r,s}|^2,\qquad
 \|\widetilde b(n'')\|_2^2=
       \sum_{\substack{n',r,s\\(n',a_T)=1}}|b_{n'n'',r,s}|^2.
 \end{aligned}
\end{equation}
Here $n''\mid a_T^\infty$ means that every prime factor of $n''$ divides $a_T$.
It includes $n''=1$; when $a_T=\pm1$, only $n''=1$ occurs.
The bound is uniform in $(a_T,A_T,C,D,Y,R,S,b)$.
These are Paper~I, Theorem~2.8, (2.38)--(2.43)
\cite[Theorem~2.8, (2.38)--(2.43), pp.~14--15]{PaperI}.
\end{fact}

Section~\ref{sec:l3} checks the masked coefficients, intervals and support
of the fourth switched row.


\section{Finite source identities and decomposition}
\label{sec:source-decomposition}

Throughout this section, \(N\) is an even integer with \(N\ge2\).
Further size hypotheses appear at each display or just before it.
Use \(L=\log N\) from Section~\ref{sec:notation}.
We prove the finite identities of this section below.
They count each occurrence with its multiplicity.
The later analytic estimates are separate hypotheses.
Use \(\mathcal A_N,G_b(N),\varepsilon,a,b\) from
\eqref{eq:intro-exponents}--\eqref{eq:intro-count} and
\(\alpha,\beta,\mu_*,\lambda,t_0,u\) from
\eqref{eq:not-sieve-parameters}. Take the roughness cutoff
\(z_s=N^\alpha\) from the Root--A scale definition in
\eqref{eq:ra-parameters}.
Throughout, \(n\) is a positive source integer. When
\(n\in\mathcal A_N\), \(n=N-p\) for a prime \(p\). In a source sum, \(e\)
is the integer being counted. The letters \(p,p_i,a_i,r\) denote
prime variables, \(q\) is a positive modulus, and \(m\) is the factor or
cofactor specified by its sum. Paper-II counts use
\(n=dm\); the marked count of Section~\ref{sec:root-a} uses \(m=(N-e)/p\). These roles
are local to the displayed source.
Here \(\Omega(n)\) counts prime factors with multiplicity.
For squarefree \(d\) with \((d,N)=1\), define
\begin{equation}\label{eq:source-paperii-cofactor}
 S^{\mathrm{II}}(d;y)
 =\#\{n=dm\in\mathcal A_N:P^-(m)\ge y,\ (m,Nd)=1\}.
\qquad(2\mid N,\ N\ge2).
\end{equation}
Set this count to zero for other \(d\).
For admissible \(d\), it is Paper~II's (1.7) count in our notation
\cite[(1.7)]{PaperII}. The zero value for other \(d\) is our convention.
Put \(x_i=\log p_i/\log N\).
The tenth row uses
\begin{equation}\label{eq:source-paperii-clipping}
\begin{aligned}
 C_{10}&=\frac{57}{64}-\frac1{500}-\frac{11}{5}\alpha,\\
 B_{10}(x)&=\begin{cases}
 \beta,&\beta\le x<\mu_*,\\
 \min\{\beta,\max\{\alpha,C_{10}-5x/2\}\},
     &\mu_*\le x<\lambda.
 \end{cases}
\end{aligned}
\qquad(2\mid N,\ N\ge2).
\end{equation}
Table~\ref{tab:source-paperii-rows} lists all thirteen sums with every
divisor and cutoff of \cite[(1.6)--(1.7), row table]{PaperII}.
All \(p_i\) in the table are prime.
Each admissible divisor contributes separately.

\begin{longtable}{@{}c c p{.47\textwidth} p{.19\textwidth}@{}}
\caption{The Paper~II source sums. Each row sums \(S^{\mathrm{II}}(d;y)\).}
\label{tab:source-paperii-rows}\\
\hline
Row & Coefficient & Divisor and domain & Cutoff \(y\)\\
\hline
\endfirsthead
\hline
Row & Coefficient & Divisor and domain & Cutoff \(y\)\\
\hline
\endhead
\(S_1\)&\(3\)&\(d=1\)&\(N^\alpha\)\\
\(S_2\)&\(1\)&\(d=1\)&\(N^\beta\)\\
\(S_3\)&\(-2\)&\(d=p_1,\ \alpha\le x_1<\mu_*\)&\(N^\alpha\)\\
\(S_4\)&\(-1\)&\(d=p_1,\ \mu_*\le x_1<1/4\)&\(N^\alpha\)\\
\(S_5\)&\(-1\)&\(d=p_1,\ 1/4\le x_1<t_0\)&\(N^\alpha\)\\
\(S_6\)&\(-1\)&\(d=p_1,\ t_0\le x_1<1/3\)&\(N^\alpha\)\\
\(S_7\)&\(-1\)&\(d=p_1,\ \mu_*\le x_1<\lambda\)&\(N^\alpha\)\\
\(S_8\)&\(1\)&\(d=p_1p_2,\)\newline
 \(\alpha\le x_2<x_1<\beta\)&\(N^\alpha\)\\
\(S_9\)&\(1\)&\(d=p_1p_2,\)\newline
 \(\alpha\le x_2<\beta\le x_1<\mu_*\)&\(N^\alpha\)\\
\(S_{10}\)&\(1\)&\(d=p_1p_2,\ \mu_*\le x_1<\lambda,\)\newline
 \(\alpha\le x_2<B_{10}(x_1)\)&\(N^\alpha\)\\
\(S_{11}\)&\(-2\)&\(d=p_1p_2,\ \lambda\le x_1<1/3,\)\newline
 \(x_1<x_2<(1-x_1)/2\)&\(p_2\)\\
\(S_{12}\)&\(-1\)&\(d=p_1p_2,\)\newline
 \(\alpha\le x_1<1/3\le x_2<(1-x_1)/2\)&\(p_2\)\\
\(S_{13}\)&\(-1\)&\(d=p_1p_2,\)\newline
 \(\beta\le x_1<\lambda\le x_2<(1-x_1)/2\)&
 \(\sqrt{N/(p_1p_2)}\)\\
\hline
\end{longtable}

Let \(\mathcal D_i\) be the displayed divisor domain and \(y_i(d)\) its cutoff.
Define the row occurrences by
\begin{equation}\label{eq:source-row-occurrences}
 \vartheta_i(n)=\sum_{\substack{d\in\mathcal D_i\\d\mid n\\
                  d\ \mathrm{squarefree},\ (d,N)=1}}
  1_{P^-(n/d)\ge y_i(d)}1_{(n/d,Nd)=1}.
\qquad(2\mid N,\ N\ge2).
\end{equation}
Thus \(S_i=\sum_{d\in\mathcal D_i}S^{\mathrm{II}}(d;y_i(d))\).
Keep every cofactor and endpoint restriction. Define
\begin{equation}\label{eq:source-thirteen}
\begin{aligned}
 S_i&=\sum_{n\in\mathcal A_N}\vartheta_i(n)\qquad(1\le i\le13),\\
 W_{13}(n)&=3\vartheta_1(n)+\vartheta_2(n)-2\vartheta_3(n)
 -\vartheta_4(n)-\vartheta_5(n)-\vartheta_6(n)-\vartheta_7(n)\\
 &{}+\vartheta_8(n)+\vartheta_9(n)+\vartheta_{10}(n)
 -2\vartheta_{11}(n)-\vartheta_{12}(n)-\vartheta_{13}(n),\\
 s_{\mathrm{sw}}(n)&=1_{\Omega(n)\ge3}\max\{W_{13}(n),0\},\\
 S_{\mathrm{sw}}&=\sum_{n\in\mathcal A_N}s_{\mathrm{sw}}(n),\qquad
 w(n)=W_{13}(n)-s_{\mathrm{sw}}(n),\qquad
 \mathcal W=\sum_{n\in\mathcal A_N}w(n).
\end{aligned}
\qquad(2\mid N,\ N\ge2).
\end{equation}
This is the complete switched source of Paper~II, (1.9)
\cite[(1.9)]{PaperII}.
Every divisor prime and cofactor cutoff is at least $N^\alpha$.
For $S_{13}$ the cutoff exceeds $N^{(1-\lambda)/4}>N^\alpha$.
Thus every $\vartheta_i(n)$ vanishes when $P^-(n)<N^\alpha$,
so the extra roughness indicator of Paper~II is redundant.
Use $P^-(1)=\infty$ as in Table~\ref{tab:notation}.
Finite linearity gives
\begin{equation}\label{eq:source-weight-identity}
\begin{aligned}
 \mathcal W={}&3S_1+S_2-2S_3-S_4-S_5-S_6-S_7\\
 &+S_8+S_9+S_{10}-2S_{11}-S_{12}-S_{13}-S_{\mathrm{sw}}.
\end{aligned}
\qquad(2\mid N,\ N\ge2).
\end{equation}


\subsection{The representation count}

Put
\begin{equation}\label{eq:source-bad-count}
\begin{aligned}
 g_b(e)&=\#\{(p_1,p_2)\in\mathbb P^2:
       N^b\le p_1\le\sqrt N,\ \varepsilon N<p_1p_2=e<N\},\\
 B_b&=\sum_e g_b(e)1_{N-e\in\mathbb P},\qquad
 P_2=\#\{n\in\mathcal A_N:\Omega(n)\le2\}.
\end{aligned}
\qquad(2\mid N,\ N\ge2).
\end{equation}
When \(e\) has distinct prime factors and both choices satisfy the
first-factor cutoff, both orientations are counted. For \(e=p^2\), there
is only the ordered pair \((p,p)\).
Hence
\begin{equation}\label{eq:source-multiplicity}
 0\le g_b(e)\le2,\qquad P_2-G_b(N)\le B_b
 \qquad(\varepsilon N>1).
\qquad(2\mid N,\ N\ge2).
\end{equation}
For the second inequality, write \(n=rs\), \(r\le s\), with both
factors prime.
If \(n\) is excluded from \(G_b\), then
\begin{equation}\label{eq:source-excluded-factor}
 N^b\le r\le\sqrt n<\sqrt N,\qquad g_b(n)\ge1.
\qquad(2\mid N,\ N\ge2,\ \varepsilon N>1).
\end{equation}
Prime squares are counted with their actual multiplicity.

Define
\begin{equation}\label{eq:source-exceptional-set}
 \mathcal E=\{n\in\mathcal A_N:P^-(n)\ge z_s,
       \ n\text{ is not squarefree or }(n,N)>1\}.
\qquad(2\mid N,\ N\ge2).
\end{equation}
For \(z_s=N^\alpha\ge2\),
\begin{equation}\label{eq:source-exceptional-size}
\begin{aligned}
 \#\{n\in\mathcal E:n\text{ is not squarefree}\}
 &\le N\sum_{k\ge\lceil z_s\rceil}k^{-2}\le\frac{2N}{z_s},\\
 \#\{n\le N:P^-(n)\ge z_s,\ (n,N)>1\}
 &\le\sum_{\substack{\ell\in\mathbb P,\ \ell\mid N\\\ell\ge z_s}}\frac N\ell
 \le\frac{11N}{z_s},\\
 |\mathcal E|&\le13N^{1-\alpha}.
\end{aligned}
\qquad(2\mid N,\ N\ge2^{1/\alpha}).
\end{equation}
The prime divisors in the middle sum are distinct.
Their number is at most
\(\lfloor1/\alpha\rfloor=11\).

Set
$N_{\mathrm{pt}}=1+\left\lceil\max\{2^{1/\alpha},
\varepsilon^{-3},\varepsilon^{-1/(1-\lambda-\beta)}\}\right\rceil$.
\begin{lemma}[Pointwise source weight bound]
\label{lem:source-pointwise}
For all $N\ge N_{\mathrm{pt}}$, the weight bounds are
\begin{equation}\label{eq:source-pointwise}
\begin{gathered}
 w(n)=0\quad(P^-(n)<z_s),\qquad
 w(n)\le4\,1_{\Omega(n)\le2}\quad(n\in\mathcal A_N\setminus\mathcal E),\\
 \max\{0,w(n)-4\,1_{\Omega(n)\le2}\}\le169
 \quad(n\in\mathcal E),
 \qquad E_{\mathrm{pt}}=0.
\end{gathered}
\qquad(2\mid N,\ N\ge N_{\mathrm{pt}}).
\end{equation}
\begin{proof}
Start from the roughness observation above.
If $\Omega(n)\ge3$, then $w(n)=\min\{W_{13}(n),0\}\le0$.
If $(n,N)>1$, the coprimality conditions force every row occurrence to vanish.
For $\Omega(n)\le2$ and $n\in\mathcal E$ with $(n,N)=1$,
only $n=p^2$ occurs; only $S_1,S_2$ contribute, giving $w(n)=4$.
For $n\notin\mathcal E$ with $\Omega(n)\le2$, a prime has weight $4$.
For $n=rs$, with primes $r<s$, the bands
$[\alpha,\beta)$, $[\beta,\mu_*)$, $[\mu_*,\lambda)$,
$[\lambda,1/3)$ and $[1/3,1/2)$ for $\log_Nr$
give weights $1,2,2,3,4$, respectively, directly from the row table.
This agrees with Paper~II, (1.11) \cite[(1.10)--(1.11)]{PaperII}.
Here $N_{\mathrm{pt}}$ is chosen so that
$\varepsilon N>N^{2/3}$ and
$\varepsilon N>N^{\lambda+\beta}$.
Then the larger prime exceeds $N^{1/3}$. The two-prime divisor
rows vanish, since their products are below $N^{\lambda+\beta}$ for
$S_8$--$S_{10}$ and below $N^{2/3}$ for $S_{11}$--$S_{13}$.
These observations prove all three bounds, including the stronger
exceptional bound with right-hand side zero.
\end{proof}
\end{lemma}
Put
\begin{equation}\label{eq:source-pointwise-error}
 E_{\mathrm{pt}}=\frac14\sum_{n\in\mathcal E}
       \max\{0,w(n)-4\,1_{\Omega(n)\le2}\}.
\qquad(2\mid N,\ N\ge2).
\end{equation}
The pointwise lemma gives $E_{\mathrm{pt}}=0$ for
$N\ge N_{\mathrm{pt}}$, within its stated scope.
We still record cruder bounds. Summation gives
\begin{equation}\label{eq:source-first-lower}
\begin{gathered}
 0\le E_{\mathrm{pt}}\le\frac{169}{4}|\mathcal E|
       \le\frac{2197}{4}N^{1-\alpha},\\
 \mathcal W\le4P_2+4E_{\mathrm{pt}},\qquad
 4G_b(N)\ge\mathcal W-4B_b-4E_{\mathrm{pt}}.
\end{gathered}
\qquad(2\mid N,\ N\ge N_{\mathrm{pt}}).
\end{equation}
The transfer to \(r\le q^{3/4}\) is \eqref{eq:intro-transfer}.

\subsection{Signed arithmetic sums}

Put
\begin{equation}\label{eq:source-fixed-sum}
\begin{aligned}
 R_{\mathrm{fix}}={}&3S_1-2S_3-S_4-S_5-S_7+S_8+S_9+S_{10}\\
 &-2S_{11}-S_{12}-S_{13}.
\end{aligned}
\qquad(2\mid N,\ N\ge2).
\end{equation}
Then \eqref{eq:source-weight-identity} gives
\begin{equation}\label{eq:source-finite-identity}
 \mathcal W-4B_b
 =R_{\mathrm{fix}}+S_2-S_6-S_{\mathrm{sw}}-4B_b.
\qquad(2\mid N,\ N\ge2).
\end{equation}
Define
\begin{equation}\label{eq:source-retained-definitions}
\begin{aligned}
 G_{\mathrm{rows}}&=2S_{11}+S_{12}+S_{13},\\
 \mathcal C_{\mathrm{ret}}&=3S_1+S_2-S_4-S_5-S_7+S_{10}-G_{\mathrm{rows}},\\
 \mathcal Q&=-2S_3+S_8+S_9.
\end{aligned}
\qquad(2\mid N,\ N\ge2).
\end{equation}
Coefficient comparison gives
\begin{equation}\label{eq:source-five-sums}
\begin{gathered}
 \mathcal C_{\mathrm{ret}}+\mathcal Q=R_{\mathrm{fix}}+S_2,\\
 \mathcal W-4B_b
 =\mathcal C_{\mathrm{ret}}+\mathcal Q-S_6-S_{\mathrm{sw}}-4B_b.
\end{gathered}
\qquad(2\mid N,\ N\ge2).
\end{equation}
Substitution proves both equalities coefficient by coefficient.
The coefficients \(3S_1\) and \(+S_{10}\) are fixed.
The integer supports may overlap.
These coefficient identities do not depend on later estimates.
Take the signed normalized chains \(Q_{j,k}\), their index set
\(\mathcal I\), and the stage index \(k\in\{0,\ldots,10\}\) from
Section~\ref{sec:retained-side}. See its ``original signed chains'' paragraph
and \eqref{eq:retB-telescope}. Each \(Q_{j,k}\) already includes its
signed normalized coefficient, so no extra multiplicity factor enters.
Finite telescoping gives
\begin{equation}\label{eq:source-signed-telescope}
 \sum_{j\in\mathcal I}Q_{j,0}
 =\sum_{j\in\mathcal I}Q_{j,10}
  +\sum_{j\in\mathcal I}\sum_{k=1}^{10}(Q_{j,k-1}-Q_{j,k}).
\qquad(2\mid N,\ N\ge2).
\end{equation}
This is the normalized-chain identity cited above. The arithmetic coefficient
identities above keep their unnormalized definitions.
Paper~II uses joint endpoints for
\(S_1,S_2\text{-}T_8/4,G_{\mathrm{rows}}\)
\cite[(6.4)--(6.4d), pp.~70--72]{PaperII}.
Here \(T_8\) is Paper~II's stage-eight endpoint in the \(S_2\) chain.
The notation \(S_2\text{-}T_8/4\) is a joint-chain label;
it defines no new arithmetic sum.
Their joint endpoint is kept before the differences are split.

\subsection{A finite upper decomposition}

Each original representation occurrence uses the shifted-prime set
of \(\eqref{eq:intro-count}\):
\begin{equation}\label{eq:source-A-N}
 \mathcal A_N=\{N-p:p\in\mathbb P,\ p<(1-\varepsilon)N\}.
\qquad(2\mid N,\ N\ge2).
\end{equation}
The two arithmetic sources stay separate. Take \(A_b,E_{\mathrm{top}}\)
from \eqref{eq:k2-source} and \(S_b\) from the Section~\ref{sec:k2} definition
\eqref{eq:k2-sifted-source}; \(A\) from \eqref{eq:l3-source};
\(\mathcal P_N\) from \eqref{eq:l3-constants}; and
\(b_A,P_N,S_v,S_A,E_{\mathrm{src}}\) from
\eqref{eq:l3-sieved-source}. The ordered prime coordinates
\(a_i\) and their prefix domains are those of \eqref{eq:l3-buchstab} and
the paragraph before it. Define
\begin{equation}\label{eq:source-local-sequences}
\begin{gathered}
 h_b(n)=A_b(N-n),\qquad
 \mathscr L_b=\{p\in\mathbb P:p\nmid2N\},\qquad E_b=E_{\mathrm{top}},\\
 h_{\mathrm{sw}}(n)=A(N-n),\qquad
 \mathscr L_{\mathrm{sw}}=\mathcal P_N,\qquad
 E_{\mathrm{sw}}^{\mathrm{src}}=E_{\mathrm{src}},
\end{gathered}
\qquad(2\mid N,\ N\ge2).
\end{equation}
Here \(A_b,E_{\mathrm{top}}\) are the Section~\ref{sec:k2} source objects in
\(\eqref{eq:k2-source}\). The objects \(A,\mathcal P_N,E_{\mathrm{src}}\) are
the Section~\ref{sec:l3} source objects in \(\eqref{eq:l3-source}\),
\(\eqref{eq:l3-constants}\), and \(\eqref{eq:l3-sieved-source}\). Both sequences are nonnegative.
No extra \(1_{(e,2N)=1}\) factor enters.
The script-font symbols \(\mathscr L_b\), \(\mathscr L_{\mathrm{sw}}\),
\(\mathscr L_X\), and \(\mathscr U_X\) denote the variable sieve families
and their sifted sums. The fixed prime support is written \(\mathcal P_N\).
For \(X\in\{b,\mathrm{sw}\}\), write \(h_X,\mathscr L_X,E_X\) for the
corresponding line of \(\eqref{eq:source-local-sequences}\), and put
\begin{equation}\label{eq:source-sifted-sequence}
\begin{aligned}
 P_X(y)&=\prod_{\substack{\ell<y\\\ell\in\mathscr L_X}}\ell,\\
 \mathscr U_X(d;y)&=\sum_{d\mid n}h_X(n)
              1_{(n/d,P_X(y))=1}.
\end{aligned}
\qquad(2\mid N,\ N\ge2).
\end{equation}
Take \(I_L,I_H,J,J_{\mathrm{tail}},J_t,z_s,\eta_{\mathrm m},
v_{\mathrm{RA}}\) from the Root--A definition \eqref{eq:ra-parameters}.
Take \(c_{R,K}\) from \eqref{eq:ra-coefficient}, and
\(R_L,R_{\mathrm{tail}},B_L,B_{\mathrm{tail}},B,c_{\mathrm{lo}},\Gamma\)
from \eqref{eq:ra-factories} and its preceding choice of fixed logarithmic
exponents. Take \(D_{\mathrm m}\) from \eqref{eq:ra-precision-order}.
Write \(Q_{\mathrm{RA}}\) for the level called \(Q\) in
\eqref{eq:ra-factories}. This alias is distinct from every signed chain
\(Q_{j,k}\). We do not redefine these objects here.
Both source comparisons require the finite threshold \(N_6\)
of Section~\ref{sec:notation}.
It enters the operative \(T_{\mathrm{RA}}\), the maximum of the
sixth-sum comparison threshold and \(N_6\). We do not redefine it here.

Use the common choice \(\sigma=\sigma_*\) from
\eqref{eq:not-sigma-choice}. It lies in the Section~\ref{sec:k2} range of
\eqref{eq:k2-parameters}, and the switched comparison uses it. Take
\(\tau,\rho\) from Section~\ref{sec:k2} \eqref{eq:k2-parameters}; Section~\ref{sec:l3} uses the same
values. We make no new choice or definition of
\(\sigma,\tau,\rho\) here.
Partition the logarithmic prime ranges into half-open intervals
\([t_i,\rho t_i)\), with \(t_i=u\rho^{k_i}\), \(k_i\ge0\).
Clip the intervals at \(\mu_*\).
Each prime has one cell.
For \(X=b\), use the retained cells of the Section~\ref{sec:k2}
\(\eqref{eq:k2-retention}\) retained-cell definition and the cutoff in
\(\eqref{eq:k2-four-rows}\). Write them as
\(\mathcal R_b,z_{b,c},\chi_{b,c}\). For \(X=\mathrm{sw}\), use the
Section~\ref{sec:l3} definitions \(\eqref{eq:l3-buchstab-peel}\) and
\(\eqref{eq:l3-row3-level}\). Write them as
\(\mathcal R_{\mathrm{sw}},z_{\mathrm{sw},c},\chi_{\mathrm{sw},c}\).
These definitions come from the two source sections.
They form no common envelope in this section.
We use no equality between the retained collections of the three sections.
The closed-cell infima of Sections~\ref{sec:k2} and~\ref{sec:l3} and their endpoint rules
apply only within their own source.
For \(X=b\), the packing function is \(H\), as specified in
\eqref{eq:k2-retention}. For \(X=\mathrm{sw}\), it is
\(\widetilde H\), as specified in \eqref{eq:l3-row3-level} and its
preceding definition of \(\Theta_c\). The two functions are defined in
Paper~II, equation (2.17d), and the paragraph after it.
We have \(\widetilde H=H\) only on the stated largest-coordinate
range, and use no equality outside it.
For either source write \(z_{X,c}=N^{t_3}\), with \(t_3\) the lower
endpoint of the third coordinate of its cell \(c\). Here
\(z_{\mathrm{sw},c}\) is the object called \(z_c\) in Section~\ref{sec:l3}; we use only the
source-indexed name. The Section~\ref{sec:k2} cutoff is defined
in \eqref{eq:k2-four-rows}, with the lower-endpoint convention stated
before \eqref{eq:k2-parent-functional}.

The source index here is always \(X\in\{b,\mathrm{sw}\}\).
For the $b$-source the underlying prime support is \(\mathcal P_N\), so
\(P_b(y)=P_N(y)\). Both sieve products use primes below $y$ not dividing
$2N$.
For the \(b\)-source, the substitution \(e=N-n\) identifies the
generic sieve with \(S_b\) of Section~\ref{sec:k2} on that source:
\begin{equation}\label{eq:source-ch07-sieve-bridge}
 \mathscr U_b(d;y)
 =\sum_{d\mid n}A_b(N-n)1_{(n/d,P_b(y))=1}
 =\sum_{d\mid N-e}A_b(e)1_{((N-e)/d,P_b(y))=1}
 =S_b(d;y).
\qquad(2\mid N,\ N\ge2).
\end{equation}
For \(X=\mathrm{sw}\), identify the generic cofactor sieve with the
Section~\ref{sec:l3} sieve on precisely the following rows:
\begin{itemize}
\item \(d=1\) with cutoff \(y=N^{\mu_*}\), the source-comparison row
\(\mathscr U_{\mathrm{sw}}(1;N^{\mu_*})=S_A\) from Section~\ref{sec:l3}
\eqref{eq:l3-sieved-source};
\item \(v=a_1\cdots a_j\), \(0\le j\le2\), with cutoff \(y=N^u\);
\item \(v=a_1a_2a_3\), with dynamic cutoff \(y=a_3\), or retained-cell
cutoff \(y=z_{\mathrm{sw},c}=N^{t_3}\);
\item \(vr=a_1a_2a_3r\) on a retained cell, with
\(z_{\mathrm{sw},c}\le r<a_3\), at cutoff \(y=r\) or
\(y=z_{\mathrm{sw},c}\).
\end{itemize}
Here the empty prefix is \(v=1\), and all primes belong to
\(\mathcal P_N\) and have the ordering in \eqref{eq:l3-buchstab}
and \eqref{eq:l3-peel}. Let \(d\) be the displayed prefix \(v\) or
product \(vr\). Every prime divisor of \(d\) is at least \(y\), so
\((d,P_{\mathrm{sw}}(y))=1\). The definitions give
\(P_{\mathrm{sw}}(y)=P_N(y)\). Write \(S_d(y)\) for the \(S_v(y)\) of Section~\ref{sec:l3}
from \eqref{eq:l3-sieved-source} with \(v=d\). For \(d\mid n\), the Section~\ref{sec:l3} definition
\eqref{eq:l3-sieved-source} gives
\begin{equation}\label{eq:source-ch08-sieve-bridge}
 \mathscr U_{\mathrm{sw}}(d;y)
 =\sum_{d\mid n}h_{\mathrm{sw}}(n)1_{(n/d,P_N(y))=1}
 =\sum_{d\mid n}b_A(n)1_{(n,P_{\mathrm{sw}}(y))=1}
 =S_d(y).
\qquad(2\mid N,\ N\ge2).
\end{equation}
This identity holds on these rows only, not for arbitrary \(d,y\).
These \(E\)-terms are unnormalized source-restoration costs. The two
source upper comparisons are stated separately:
Here \(B_b'\) is the prime-occurrence part of \(B_b\), and
\(E_{\mathrm{top}}\) is the removed-top cost.
\begin{lemma}[Upper-source comparisons]
\label{lem:source-upper-sources}
\begin{equation}\label{eq:source-upper-comparison}
\begin{gathered}
 B_b\le E_b+\mathscr U_b(1;N^{\mu_*}),\qquad
 S_{\mathrm{sw}}\le E_{\mathrm{sw}}^{\mathrm{src}}+
 \mathscr U_{\mathrm{sw}}(1;N^{\mu_*})\\
 \text{for even }N\ge N_6.
\end{gathered}
\end{equation}
\begin{proof}
The first inequality follows from Section~\ref{sec:k2} \eqref{eq:k2-source-comparison}
and \(E_b=E_{\mathrm{top}}\); the second is the Section~\ref{sec:l3} source comparison
\eqref{eq:l3-source-id}.
We do not identify their hypotheses or their restoration costs.
\end{proof}
We invoke the source comparison once, with the named hypotheses displayed above.
\end{lemma}
For real cutoffs \(1\le y\le w_{\mathrm{cut}}\), the least allowed prime gives
\begin{equation}\label{eq:source-buchstab-one-step}
 \mathscr U_X(d;w_{\mathrm{cut}})=\mathscr U_X(d;y)
 -\sum_{\substack{y\le r<w_{\mathrm{cut}}\\r\in\mathscr L_X}}\mathscr U_X(dr;r).
\qquad(2\mid N,\ N\ge2).
\end{equation}
Write \(n=dm\).
If \(m\) survives at \(y\) but fails at \(w_{\mathrm{cut}}\),
choose its least allowed prime in \([y,w_{\mathrm{cut}})\).
Exactly one summand records that prime.
Repeated copies remain in the cofactor.
Thus the identity also counts prime powers correctly.

Three applications give
\begin{equation}\label{eq:source-buchstab}
\begin{aligned}
 \mathscr U_X(1;N^{\mu_*})&=T_{X,0}-T_{X,1}+T_{X,2}-T_{X,3},\\
 T_{X,j}&=\sum_{\substack{N^u\le a_j<\cdots<a_1<N^{\mu_*}\\
                         a_i\in\mathscr L_X}}
                         \mathscr U_X(a_1\cdots a_j;N^u)\quad(0\le j\le2),\\
 T_{X,3}&=\sum_{\substack{N^u\le a_3<a_2<a_1<N^{\mu_*}\\
                         a_i\in\mathscr L_X}}
                         \mathscr U_X(a_1a_2a_3;a_3).
\end{aligned}
\qquad(2\mid N,\ N\ge2).
\end{equation}
For \(j=0\), the sum is \(\mathscr U_X(1;N^u)\).
The primes are ordered; no factorial occurs.
Every multiplicity remains in \(h_X\).

For \(d=a_1a_2a_3\) in a retained cell of
source \(X\), put
\begin{equation}\label{eq:source-dynamic-cutoff}
\begin{aligned}
 -\mathscr U_X(d;a_3)
 &= -\mathscr U_X(d;z_{X,c})+
       \sum_{\substack{z_{X,c}\le r<a_3\\r\in\mathscr L_X}}\mathscr U_X(dr;r)\\
 &\le-\mathscr U_X(d;z_{X,c})+
       \sum_{\substack{z_{X,c}\le r<a_3\\r\in\mathscr L_X}}\mathscr U_X(dr;z_{X,c}).
\end{aligned}
\qquad(2\mid N,\ N\ge2).
\end{equation}
The equality is \eqref{eq:source-buchstab-one-step}.
Nonnegativity proves the inequality.
Actual cell membership gives \(z_{X,c}\le a_3\).
The closed infimum creates no additional occurrences.

Let \(\sum_{\mathrm{ret},X}\) restrict the ordered triples to
\(\mathcal R_X\), equivalently to \(\chi_{X,c}=1\). For either source define
the peel terms
\begin{equation}\label{eq:source-base-peel}
\begin{aligned}
 R_{X,3}&=\sum_{\mathrm{ret},X}\mathscr U_X(a_1a_2a_3;z_{X,c}),\\
 P_{X,3}&=\sum_{\mathrm{ret},X}
       \sum_{\substack{z_{X,c}\le r<a_3\\r\in\mathscr L_X}}
             \mathscr U_X(a_1a_2a_3r;z_{X,c}).
\end{aligned}
\qquad(2\mid N,\ N\ge2).
\end{equation}
For each source \(X\in\{b,\mathrm{sw}\}\), the non-retained depth-three
terms are nonnegative and occur with a minus sign, so we drop them.
The retained terms use the dynamic inequality above:
\begin{equation}\label{eq:source-nonretained-drop}
 -T_{X,3}\le -R_{X,3}+P_{X,3}.
\qquad(2\mid N,\ N\ge2).
\end{equation}
For even \(N\ge N_6\), define the two upper expressions and their
slacks, one per source. The next comparison proves their nonnegativity:
\begin{equation}\label{eq:source-upper-slack}
 \begin{aligned}
 U_b&:=T_{b,0}-T_{b,1}+T_{b,2}-R_{b,3}+P_{b,3},&
 Z_b&:=E_b+U_b-B_b\ge0,\\
 U_{\mathrm{sw}}&:=T_{\mathrm{sw},0}-T_{\mathrm{sw},1}
       +T_{\mathrm{sw},2}-R_{\mathrm{sw},3}+P_{\mathrm{sw},3},&
 Z_{\mathrm{sw}}&:=E_{\mathrm{sw}}^{\mathrm{src}}
       +U_{\mathrm{sw}}-S_{\mathrm{sw}}\ge0.
 \end{aligned}
 \qquad(2\mid N,\ N\ge N_6).
\end{equation}
Hence
\begin{equation}\label{eq:source-sifted-upper}
 \mathscr U_X(1;N^{\mu_*})\le
 \begin{cases}
 U_b,&X=b,\\
 U_{\mathrm{sw}},&X=\mathrm{sw}.
 \end{cases}
\qquad(2\mid N,\ N\ge N_6).
\end{equation}
These inequalities assume the two source comparisons in
\eqref{eq:source-upper-comparison}. They give no estimate for either
restoration cost.

\subsection{The conditional lower expression}

Take \(A_{\mathrm{mix}}\) from the Root--A definition
\eqref{eq:ra-mix-definition}, and \(Z_6\) with its comparison from
Section~\ref{sec:root-a} \eqref{eq:ra-positive-comparison}. We redefine neither object.
This sum uses the ordinary upper Rosser weights
\(\mu_d^+(D)\) with the odd-index cubic gates of
\(\eqref{eq:ra-rosser-weight}\). Their upper-sieve property is the
Iwaniec upper-bound lemma
\cite[Lemma~1, (3.2), (3.6)]{Iwaniec1980Rosser}
\cite[Lemmas~2--3, p.~313]{Iwaniec1980Remainder}. This section
introduces no new weight.
The Section~\ref{sec:root-a} comparison gives
\begin{equation}\label{eq:source-mixed-comparison}
 S_6\le A_{\mathrm{mix}},\qquad Z_6\ge0
 \qquad(2\mid N,\ N\ge N_6).
\end{equation}
For a mark \(p\), put \(m=(N-e)/p\) and
\(P_{\mathrm{RA}}=\prod_{\ell<z_s,\,\ell\nmid2N}\ell\).
For \(D>1\), the cited Iwaniec upper-bound lemmas give, for the upper
Rosser divisor sum,
\cite[Lemma~1, (3.2), (3.6)]{Iwaniec1980Rosser}
\cite[Lemmas~2--3, p.~313]{Iwaniec1980Remainder},
\[
 \sum_{d\mid(m,P_{\mathrm{RA}})}\mu_d^+(D)
 \ge1_{(m,P_{\mathrm{RA}})=1}\ge1_{P^-(m)\ge z_s}
 \qquad(2\mid N,\ N\ge2,\ D>1).
\]
Dropping \((m,Np)=1\) enlarges the original nonnegative occurrence count.
For even \(N\ge N_6\), each of \(R_L,R_{\mathrm{tail}},Q_{\mathrm{RA}}\) exceeds
\(N^{1/3}\). Hence take \(D=R_L/p,R_{\mathrm{tail}}/p,Q_{\mathrm{RA}}/p\) on,
respectively, \(I_L\times J\), \(I_H\times J_{\mathrm{tail}}\),
\(I_H\times J_t\). These pieces partition the original source: \(e=N/4\)
belongs to \(I_L\), and the shared mark endpoint belongs to
\(J_{\mathrm{tail}}\).
Since \(p\ge N^{t_0}>z_s\), the mark is unique in \(q=pd\).
The condition \(d\mid m\) is exactly \(e\equiv N\pmod{pd}\). The original gates give
\(c_{R,K}(pd)=\mu_d^+(R/p)\). Summation gives precisely
\(S_6\le A_{\mathrm{mix}}\) and \(Z_6\ge0\), with all mark
multiplicities. A product-cell partition only splits these exact
sums. It must keep the original \(pd<R\), the cubic gates and every clipped
endpoint. In particular \(c_{\mathrm{lo}}\) and \(\Gamma\) keep the
gates at \(Q_{\mathrm{RA}}\).
Subdivide each of the three marked pieces in
\eqref{eq:ra-three-pieces} into cells in its two displayed coordinates.
Let \(\delta=L^{-D_{\mathrm m}}\), and use half-open multiplicative
intervals of relative width at most \(1+\delta\) in each coordinate.
Insert all row cuts, the points
\(N/4\), \(1/3-\eta_{\mathrm m}\), \(1/3\), and the two endpoints of
\(I_H\) as cell boundaries. Denote the resulting finite, disjoint cell
 families by \(\Theta^{\sharp}_L,\Theta^{\sharp}_{\mathrm{tail}},\Theta^{\sharp}_t\).
Use the separate piece index \(Y\in\{L,\mathrm{tail},t\}\) and the
mapping
\begin{equation}\label{eq:source-root-a-piece-map}
 (I_Y,K_Y,R_Y)=
 \begin{cases}
 (I_L,J,R_L),&Y=L,\\
 (I_H,J_{\mathrm{tail}},R_{\mathrm{tail}}),&Y=\mathrm{tail},\\
 (I_H,J_t,Q_{\mathrm{RA}}),&Y=t.
 \end{cases}
\qquad(2\mid N,\ N\ge N_6).
\end{equation}
Thus each \(C\in\Theta^{\sharp}_Y\) satisfies
 \(C\subseteq(\mathbb P\cap I_Y)\times(\mathbb P\cap N^{K_Y})\) and
\((p,N)=1\), with \(0<(N-e)/p\) on every included occurrence. Their unions
are exactly the corresponding coprime marked source pieces
\(I_L\times J\), \(I_H\times J_{\mathrm{tail}}\), and \(I_H\times J_t\).
A cell touching a boundary is kept as a boundary cell.
For \(Y\in\{L,\mathrm{tail},t\}\) and \(C\in\Theta^{\sharp}_Y\), define
\begin{equation}\label{eq:source-cell-sums}
\begin{aligned}
 S_{6,Y,C}&:=\sum_{(e,p)\in C}\mathbf1_{p\in\mathbb P}\mathbf1_{(p,N)=1}\mathbf1_{p\mid N-e}\mathbf1_{0<(N-e)/p}\mathbf1_{e\in\mathbb P}
 \mathbf1_{P^-((N-e)/p)\ge z_s}\mathbf1_{((N-e)/p,Np)=1},\\
 A_{Y,C}&:=\sum_{(e,p)\in C}\mathbf1_{p\in\mathbb P}\mathbf1_{(p,N)=1}\mathbf1_{p\mid N-e}\mathbf1_{0<(N-e)/p}\mathbf1_{e\in\mathbb P}
 \sum_{\substack{d\mid (N-e)/p\\d\mid P_{\mathrm{RA}}}}
 c_{R_Y,K_Y}(pd)\mathbf1_{e\equiv N\ (pd)},
\end{aligned}
\qquad(2\mid N,\ N\ge2).
\end{equation}
Here \(R_Y\) is \(R_L\), \(R_{\mathrm{tail}}\), or \(Q_{\mathrm{RA}}\). The
corresponding coefficient gates stay in the inner sum. The
half-open families make every marked occurrence belong to exactly one
\(C\), including \(e=N/4\) and the shared mark endpoint. The one-term
Rosser inequality at each occurrence gives \(S_{6,Y,C}\le A_{Y,C}\), cell
by cell. Reindexing the finite sums gives
\begin{equation}\label{eq:source-cell-sum}
 S_6=\sum_Y\sum_{C\in\Theta^{\sharp}_Y}S_{6,Y,C},\qquad
 \sum_Y\sum_{C\in\Theta^{\sharp}_Y}A_{Y,C}=A_{\mathrm{mix}}.
\qquad(2\mid N,\ N\ge N_6).
\end{equation}
The finite comparison uses these cellwise Rosser inequalities.
Thus the displayed comparison is the sum of finitely many explicit cell
inequalities. No product cell, endpoint, or boundary strip is
omitted. The later complete-cell estimates and boundary-strip bound are
quantitative inputs only and do not alter this finite identity.
The three slacks below are cited or defined by the source comparisons
above, before this substitution.
Substitution in \eqref{eq:source-finite-identity} gives
\begin{equation}\label{eq:source-exact-slack}
\begin{aligned}
 \mathcal W-4B_b-4E_{\mathrm{pt}}
 ={}&R_{\mathrm{fix}}+S_2-A_{\mathrm{mix}}
       -U_{\mathrm{sw}}-4U_b\\
    &-E_{\mathrm{sw}}^{\mathrm{src}}-4E_b-4E_{\mathrm{pt}}
       +Z_6+Z_{\mathrm{sw}}+4Z_b.
\end{aligned}
\qquad(2\mid N,\ N\ge N_6).
\end{equation}
This equality follows from the source identity and the cited and
local definitions. The Section~\ref{sec:root-a} comparison
\eqref{eq:ra-positive-comparison} gives \(Z_6\ge0\). The comparisons in
\eqref{eq:source-upper-slack} give \(Z_{\mathrm{sw}},Z_b\ge0\).
For even \(N\ge\max\{N_6,N_{\mathrm{pt}}\}\),
\eqref{eq:source-first-lower} and nonnegativity of these slacks yield
\begin{equation}\label{eq:source-final-lower}
\begin{aligned}
 4G_b(N)\ge{}&R_{\mathrm{fix}}+S_2-A_{\mathrm{mix}}\\
 &-(T_{\mathrm{sw},0}-T_{\mathrm{sw},1}+T_{\mathrm{sw},2}
                         -R_{\mathrm{sw},3}+P_{\mathrm{sw},3})\\
 &-4(T_{b,0}-T_{b,1}+T_{b,2}-R_{b,3}+P_{b,3})\\
 &-E_{\mathrm{sw}}^{\mathrm{src}}-4E_b-4E_{\mathrm{pt}}.
\end{aligned}
\qquad(2\mid N,\ N\ge\max\{N_6,N_{\mathrm{pt}}\}).
\end{equation}
The two expanded signs are
\begin{equation}\label{eq:source-expanded-signs}
\begin{aligned}
 -U_{\mathrm{sw}}&=-T_{\mathrm{sw},0}+T_{\mathrm{sw},1}
                  -T_{\mathrm{sw},2}+R_{\mathrm{sw},3}-P_{\mathrm{sw},3},\\
 -4U_b&=-4T_{b,0}+4T_{b,1}-4T_{b,2}+4R_{b,3}-4P_{b,3}.
\end{aligned}
\qquad(2\mid N,\ N\ge N_6).
\end{equation}
Use one bound for the complete \(S_{\mathrm{sw}}\).
Paper~II gives its coefficient-eight estimate
\cite[(1.9), (5.29)--(5.43)]{PaperII}; the Buchstab upper bound
concerns the same occurrence.
Sections~\ref{sec:k2} and \ref{sec:l3} estimate these two five-term
upper expressions (four principal rows and their peel term) separately.
The displayed expansion fixes each peel term's sign.
The finite switched-source domination is already in the
local \(Z_{\mathrm{sw}}\) identity above. The remaining distribution
and quantitative estimates for this occurrence are separate hypotheses
of the later bill. They lie outside the finite decomposition of this section.


\subsection{\texorpdfstring{The lower expression with the original $S_2$}{The lower expression with the original S2}}

Keep the original \(S_2\) throughout.
Use $\overline H$, $M_{5,\mathrm{internal}}$ and $\Delta D_{L3}$
from \eqref{eq:intro-components}.
Here $\overline H$ is the switched integral bound.
The distinct constant $153/256$ bounds the packing envelope.
The complete normalized switched error is $E_{\mathrm{sw}}(N;\sigma)$.
The source-restoration error $E_{\mathrm{sw}}^{\mathrm{src}}$ is unnormalized.
The main route takes $g_6:=0$ and excludes the optional six-endpoint
improvement.
Fix this common parameter and an even \(N\ge\max\{N_6,N_{\mathrm{pt}}\}\)
throughout the following conditional lemma. Take
\(W,E_{\mathrm{van}},\Lambda^0_{j,N},\Lambda_N,\Delta_j,\Delta_{G1}\)
from \eqref{eq:not-normalization}, \eqref{eq:intro-retained-error}, and
Section~\ref{sec:retained-side}, and \(g_{\mathrm{RA,lo}},E_{\mathrm{RA}}\)
from \eqref{eq:ra-b11-floor}. Use the complete switched and exterior
errors from \eqref{eq:l3-final} and \eqref{eq:k2-final-fixed-sigma}.
The fixed constants are defined in \eqref{eq:intro-components}
and \eqref{eq:intro-fallback-constant}.

\begin{lemma}[Addition of the four comparisons]
\label{lem:source-trackb-combination}
Assume \eqref{eq:source-pointwise} and the following four inequalities.
Write
\begin{equation}\label{eq:source-trackb-retained}
\begin{aligned}
 \frac{R_{\mathrm{fix}}+S_2}{W(N)}
 &\ge D_{\mathrm{II}}+\overline J_6
       +2\overline H+\frac{5719}{10^6}
       +M_{5,\mathrm{internal}}-E'(N;\sigma),\\
 E'(N;\sigma)&=E_{\mathrm{van}}(N;\sigma)
    +\sum_{j\in\mathcal H}
          \bigl(|\Lambda^0_{j,N}|+|\Lambda_N(\Delta_j)|\bigr)
    +|\Lambda_N(\Delta_{G1})|.
\end{aligned}
\qquad(2\mid N,\ N\ge\max\{N_6,N_{\mathrm{pt}}\}\ \text{(assumed)}).
\end{equation}
Here \(D_{\mathrm{II}}\) denotes the original Paper~II floor.
The main route does not use the six endpoint replacements.
Here $\mathcal H=\{4,5,7,10\}$, with $5$ denoting the lower
part of $S_5$ and $10$ its retained $S_{10}$ part.
Section~\ref{sec:retained-side} gives the hypotheses for this estimate.

The other three comparisons have the forms
\begin{equation}\label{eq:source-trackb-three-estimates}
\begin{aligned}
 \frac{S_6}{W(N)}
   &\le\overline J_6-g_{\mathrm{RA,lo}}+E_{\mathrm{RA}}(N;\sigma),\\
 \frac{S_{\mathrm{sw}}}{W(N)}
   &\le2\overline H-\Delta D_{L3}+E_{\mathrm{sw}}(N;\sigma),\\
 \frac{4B_b}{W(N)}&\le c_{K2,G0}I(b)+E_K(N;b,\sigma).
\end{aligned}
\qquad(2\mid N,\ N\ge\max\{N_6,N_{\mathrm{pt}}\}\ \text{(assumed)}).
\end{equation}
All error terms are nonnegative.

Then the following normalized lower bound holds:
\begin{equation}\label{eq:source-trackb-combination}
\begin{aligned}
 \frac{4G_b(N)}{W(N)}
 &\ge D_{\mathrm{II}}+\frac{5719}{10^6}
      +M_{5,\mathrm{internal}}+g_{\mathrm{RA,lo}}+\Delta D_{L3}
      -c_{K2,G0}I(b)-\mathcal E_B(N;b,\sigma),\\
 \mathcal E_B(N;b,\sigma)
 &=E'(N;\sigma)+E_{\mathrm{RA}}(N;\sigma)+E_{\mathrm{sw}}(N;\sigma)+E_K(N;b,\sigma)
       +\frac{4E_{\mathrm{pt}}}{W(N)}.
\end{aligned}
\qquad(2\mid N,\ N\ge\max\{N_6,N_{\mathrm{pt}}\}).
\end{equation}
\end{lemma}

\begin{proof}
Substitute the four assumed estimates into
\eqref{eq:source-finite-identity} and use \eqref{eq:source-first-lower}.
The terms \(\overline J_6\) and \(2\overline H\) cancel exactly.
Each complete upper source is subtracted once.
In particular, no additional switched majorant is added.

Using the same retained constant as in Section~\ref{sec:final-bill}, put
\begin{equation}\label{eq:source-trackb-floor}
\begin{aligned}
 D_B&=D_{\mathrm{II}}+\frac{5719}{10^6}
      +M_{5,\mathrm{internal}}+g_{\mathrm{RA,lo}}
      -\frac1{10000}+\Delta D_{L3},\\
 \mathcal E_B(N;b,\sigma)\le\frac1{10000}
 &\quad\Longrightarrow\quad
 \frac{4G_b(N)}{W(N)}\ge D_B-c_{K2,G0}I(b).
\end{aligned}
\qquad(2\mid N,\ N\ge\max\{N_6,N_{\mathrm{pt}}\}).
\end{equation}
\end{proof}

\section{The retained side}\label{sec:retained-side}

Keep every original source of Paper~II, including the complete sum $S_2$.
Fix $0<\sigma\le10^{-10}$ satisfying Paper~II, (1.27). Put
\begin{equation}\label{eq:retB-parameters}
 L=\log N,\qquad \tau=\sigma^9,\qquad \rho=1+\tau,
 \qquad h=h_G=\frac9{5000},\qquad \eta_L=\frac1{10000}.
\end{equation}
The choice of $\sigma$ precedes the limit.
Use $W$ from \eqref{eq:not-normalization}.
Write $R_{\mathrm{fix}}$ for \eqref{eq:source-fixed-sum}.

The comparison required in this section is
\begin{equation}\label{eq:retB-goal}
\begin{aligned}
 \frac{R_{\mathrm{fix}}+S_2}{W}
 &\ge D_{\mathrm{II}}+g_6+\overline J_6+2\overline H
       +\frac{5719}{10^6}+M_{5,\mathrm{internal}}-E'(N;\sigma),\\
 E'(N;\sigma)
 &=E_{\mathrm{van}}(N;\sigma)
   +\sum_{j\in\mathcal H}\bigl(|\Lambda^0_{j,N}|+|\Lambda_{j,N}|\bigr)
   +|\Lambda_{G,N}|,\qquad
 \mathcal H=\{4,5,7,10\}.
\end{aligned}
\end{equation}
In $\mathcal H$, indices $5$ and $10$ denote the lower $S_5$
and retained $S_{10}$ parts.
Every displayed error is nonnegative.
The notation agrees with \eqref{eq:intro-retained-error}:
\begin{equation}\label{eq:retB-Lambda-alias}
 \Lambda_{j,N}=\Lambda_N(\Delta_j),\qquad
 \Lambda_{G,N}=\Lambda_N(\Delta_{G1}).
\end{equation}
The required limit is $E'(N;\sigma)\to0$.
Proposition~\ref{prop:retB-comparison} proves \eqref{eq:retB-goal}.
Its inputs are listed in Remark~\ref{rem:retB-inputs}. The six-endpoint
lemma is optional: without it, the
same proof gives \eqref{eq:retB-goal} with $g_6$ replaced by $0$.
We display the unchanged comparisons first.
Then we specify the five incremental and four saved-array conversion terms.

\subsection{The original signed chains}

Let $\mathcal I=\mathcal I_{\mathrm{row}}\sqcup\{\mathrm{T5}\}$ index the
signed chains of $R_{\mathrm{fix}}+S_2$. Here $\mathcal I_{\mathrm{row}}$ is the
set of Paper~II's nineteen row chains. The chain $\mathrm{T5}$ is Paper~II's outer
chain for the upper part of $S_5$
\cite[(2.28), (5.10), (6.1) and \S6.1, condition~1]{PaperII}.
For $j\in\mathcal I$, let $Q_{j,k}$ include its normalized coefficient.
Paper~II uses ten stages
\cite[(6.4)--(6.4d)]{PaperII}:
\begin{equation}\label{eq:retB-telescope}
 Q_{j,0}=Q_{j,10}
       +\sum_{k=1}^{10}(Q_{j,k-1}-Q_{j,k}).
\end{equation}
The initial expression is the actual arithmetic source.
The last expression is a directed rational endpoint.
For $j=\mathrm{T5}$ the stages are those of Lemma~\ref{lem:retB-S5top}:
$Q_{\mathrm{T5},k}$ is defined there for $k\le3$, and
$Q_{\mathrm{T5},k}=Q_{\mathrm{T5},3}$ for $3\le k\le10$.
For the three joint groups, use
\begin{equation}\label{eq:retB-joint}
\begin{aligned}
 \mathcal O&\in\{S_1,S_2\text{-}T_8/4,G_1\},\qquad
 Q_{\mathcal O,9}=\sum_{j\in\mathcal O}Q_{j,9},\\
 \sum_{j\in\mathcal O}Q_{j,0}
 &=E_{\mathcal O}+(Q_{\mathcal O,9}-E_{\mathcal O})
    +\sum_{j\in\mathcal O}\sum_{k=1}^{9}
             (Q_{j,k-1}-Q_{j,k}),\\
 Q_{\mathcal O,9}-E_{\mathcal O}&\ge0.
\end{aligned}
\end{equation}
The joint signs are kept before rounding.
The $G_1$ weights are $(2,1,1)$.

\paragraph{Signed sources and parent assignment.}
Here ``parent'' records the assignment of an arithmetic occurrence and its
source-restoration share. It does not assert that each Buchstab child
is numerically bounded by the unsplit parent. Let $T_{j,r}$ be the
nonnegative plain-source terms of Paper~II, (2.19) and (3.2). Let
$\Delta_j^{\rm src}$ be its signed source correction. Then
\[
 3S_1=3(T_{1,0}-T_{1,1}+T_{1,2}-T_{1,3})
                  +3\Delta_1^{\rm src},\qquad
 S_2=T_{2,0}-T_{2,1}+T_{2,2}-T_{2,3}+\Delta_2^{\rm src}.
\]
For every third-depth S1 occurrence, $E_3\le p_3$ gives
$-3S_{\rm pl}(b,p_3)\ge-3S_{\rm pl}(b,E_3)$. For S2 we make this
replacement only on Paper~II's selected whole-$K$ blocks $\mathcal F$,
including $PF$ and $\mathcal F\cap R1$. There
$-S_{\rm pl}(b,p_3)\ge-S_{\rm pl}(b,E_3)$. On the strict complement
$\mathcal B$ keep the actual $p_3$ and Paper~II's strict-BV comparison.
The last child of S2 keeps its negative sign on all three pieces
$PF$, $\mathcal F\cap R1$, and $\mathcal B$. The strict BV piece stays in S2.

For the fixed-cutoff rows use their literal divisor domains. Their
signed sources are $-2S_3$, $-S_4$, $-S_{5,\rm low}$, $-S_7$,
$+S_8$, and $+S_9$, respectively. On each source cell, an upper sieve
bound $S\le U$ gives $cS\ge cU$ when $c<0$. A lower sieve bound
$L\le S$ gives $cS\ge cL$ when $c>0$. These comparisons keep the source
corrections and the row's allocated errors.
Paper~II, (2.28), gives the exact assignment identity
\[
 -S_5=-S_{5,\rm low}-S_{5,\rm top},\quad
 S_{5,\rm low}=\sum_{1/4\le\log_Np<t_*}S(p,N^\alpha),\quad
 S_{5,\rm top}=\sum_{t_*\le\log_Np<\tau_S}S(p,N^\alpha).
\]
The endpoint $t_*$ belongs to the upper/BV part. The split
precedes the structured sieve or Heath--Brown expansion and keeps
parent labels. S5-low and T5 both have parent S5. S5-low gets no
restoration address of its own.
For S10 the source partition
$S_{10}=S_{10,\rm Ret}+S_{10,\rm Drop}$ is exact, with
$S_{10,\rm Ret},S_{10,\rm Drop}\ge0$. Hence $+S_{10}\ge+S_{10,\rm Ret}$
before its lower sieve bound.

The joint source keeps the representation weights $(2,1,1)$.
Paper~II, (4.14), gives, with its stated exceptional-source and
small-prime allowance $E_G\ge0$,
\[
 G_1\le2S(\mathcal B_{11},N^{1/2})+
          S(\mathcal B_{12},N^{1/2})+
          S(\mathcal B_{13},N^{1/2})+E_G.
\]
Make its cut-measure replacement with the original signed correction
before applying (4.18). The four signed descendants of $-G_1$ then
have coefficients $(-1,+1,-1,+1)$ on the single weighted source.
Their last positive term is bounded from below by its retained cells.
The change from actual $a_3$ to label base $E$ keeps the signed peel
correction of Paper~II, (4.45)--(4.57). The four rows share one
joint source parent and one restoration allocation.
Thus S1, S2, S4, S5, S7, S10 and G1 are the seven parents of the
sixteen listed retained rows. S3, S8 and S9 keep their original
single-comparison parents, and T5 remains the upper part of S5.
These are sign-correct source reductions for the complete parent
chains. They do not discharge the distinct S6, bad-source or
switched-source comparisons in Section~\ref{sec:source-decomposition}.

The nineteen original row chains of $\mathcal I_{\mathrm{row}}$
are the sixteen rows of Appendix~\ref{subsec:app-retB-rows} and
Paper~II's original $S_3,S_8,S_9$ comparisons. The strict BV piece of $S_2$ lies inside
the joint $S_2$ group. The source $-S_5$ of $R_{\mathrm{fix}}$ is split at
$N^{t_*}$, $t_*=7109/28000$, before any sieve step. Its lower part is the row
$S_5$-low of $\mathcal I_{\mathrm{row}}$. Its upper part is the chain
$\mathrm{T5}$. The chains of $S_6$ and $S_{\mathrm{sw}}$ are excluded.
Thus the endpoint sum over $\mathcal I$ is
\begin{equation}\label{eq:retB-base-endpoints}
 \sum_{j\in\mathcal I} Q_{j,10}
 =B_{\mathrm{print}}-\delta_{G1}+\overline J_6+2\overline H,
 \qquad \delta_{G1}=4.98\cdot10^{-54}.
\end{equation}
This follows from Paper~II's printed thirteen-term sum
\cite[\S6.4, immediately before (6.6)]{PaperII}.
Its joint endpoint is $Q_{G1,10}=L_{G1}-\delta_{G1}$;
the printed sum uses $L_{G1}$. The term $-\overline J_5$ of
\eqref{eq:retB-Bprint} is $Q_{\mathrm{T5},10}$, by
Lemma~\ref{lem:retB-S5top}. Apart from $-\overline J_6-2\overline H$,
every other term of \eqref{eq:retB-Bprint} is an endpoint. It ends a
chain of $\mathcal I_{\mathrm{row}}$ or a joint group.
This is a finite coefficient identity.
The printed sum is
\begin{equation}\label{eq:retB-Bprint}
\begin{aligned}
 B_{\mathrm{print}}={}&\underline D_1+\underline d_2-\frac{U_{3,A}}2
 -\frac{U_4+U_{5,\mathrm{low}}+U_7}{4}\\
 &+\frac{L_{8,A}+L_{9,A}+L_{10}}4+L_{G1}
 -\overline J_5-\overline J_6-2\overline H\\
 >&1.83403742109641331345957914542194044932304169727869681073064.
\end{aligned}
\end{equation}
The letters denote the rational endpoints printed in Paper~II,
Appendix~A.2--A.6 and (A.7), (A.10), (A.13)--(A.14).
The constant $\overline J_6$ keeps its rational definition:
\begin{equation}\label{eq:retB-J6-endpoint}
 \overline J_6=1.3148581959157321533513247276782285310758770468886237619585.
\end{equation}

The upper part of $S_5$ has the following outer chain. Put
$t_*=7109/28000$, $\tau_S=57/224$, and let
\begin{equation}\label{eq:retB-J5-endpoint}
 \overline J_5=\frac14\bigl(c_5+8.84\cdot10^{-57}\bigr),\qquad
 c_5=0.03658166423843258896399704402764682632822246022922592990,
\end{equation}
which is Paper~II's directed upper endpoint for
$J_5(0)=\int_{t_*}^{\tau_S}\frac{dt}{t(1/2-t)}$
\cite[(5.16), (5.20), (A.14)--(A.14a)]{PaperII}.

\begin{lemma}[The outer chain of the upper part of $S_5$]\label{lem:retB-S5top}
Define
\[
\begin{aligned}
 Q_{\mathrm{T5},0}&=-\frac1W\sum_{t_*\le\log_Np_1<\tau_S}S(p_1,N^\alpha),&
 Q_{\mathrm{T5},1}&=-\frac{T_5}W,\\
 Q_{\mathrm{T5},2}&=-J_5(0),&
 Q_{\mathrm{T5},3}&=-\overline J_5,
\end{aligned}
\]
where the first sum is the second term of Paper~II, (2.28), and $T_5$ is
Paper~II, (5.10). Then
\[
 Q_{\mathrm{T5},0}-Q_{\mathrm{T5},1}\ge-E_{\mathrm{T5},1}(N),\qquad
 Q_{\mathrm{T5},1}-Q_{\mathrm{T5},2}\ge-E_{\mathrm{T5},2}(N),\qquad
 Q_{\mathrm{T5},2}-Q_{\mathrm{T5},3}\ge0,
\]
with $E_{\mathrm{T5},1},E_{\mathrm{T5},2}\ge0$ tending to $0$ at fixed
$\sigma$. Both are entries of $E_{\mathrm{van}}$. In particular
$Q_{\mathrm{T5},0}\ge-\overline J_5-E_{\mathrm{T5},1}(N)-E_{\mathrm{T5},2}(N)$.
\end{lemma}

\begin{proof}
Paper~II splits $-S_5$ at $N^{t_*}$ before the structured sieve or HB.
The endpoint $t_*$ belongs to the upper part \cite[(2.28) and the three
lines after it]{PaperII}. The coefficient is $-1$.

Stage one. For a prefix $p_1\in\mathcal P_N$, dropping the condition
$(m,Nb)=1$ gives $S(p_1,N^\alpha)\le S_{p_1}(N^\alpha)$ in the notation of
Paper~II, (5.9) \cite[the sentence after (5.9)]{PaperII}. Prefixes
$p_1\mid N$ are removed at their first failure. This is the error-table
entry ``Exceptional source factors'', of size
$O_\sigma(L^3/N+N^{-\alpha/2+o(1)})$ \cite[\S6.3]{PaperII}. A prime equal
to $N^{t_*}$ or $N^{\tau_S}$ divides $N$ \cite[after (5.20)]{PaperII}. With
the coefficient $-1$ this gives the first inequality.

Stage two. Paper~II applies the upper sieve (5.4) to each $\mathcal A_b$,
$b=p_1$. Its parameters are $X_{\mathcal A}=\operatorname{Li}(X)/\varphi(b)$, level
$Q=N^{1/2}L^{-B}$ ($\kappa=0$), $D_b=Q/b$ and $z_b=N^\alpha$
\cite[(5.12)]{PaperII}. For $t=\log_Np_1\in[t_*,\tau_S)$ one has
$a_0(t)=1/2-t\ge55/224$, $t_*-\alpha=5368241/32172000>0$, and
\[
 \frac{a_0}{\alpha}\in\Bigl[\frac{12639}{4480},\frac{7917759}{2800000}\Bigr]
 \subset(2,3)
\]
\cite[table after (5.12), row T5]{PaperII}. Hence eventually
$s_b=(a_0-\delta_N)/\alpha\in(2,3)$, where $F(s)=2e^\gamma/s$
\cite[between (5.4) and (5.5)]{PaperII}. By (5.14) the remainders are $o(U_C)$. The main
terms are (5.15), and (5.16) gives
$T_5\le4(1-\varepsilon)U_CJ_5(0)+o(U_C)$ with the vanishing logarithmic
loss (5.21) \cite[(5.14)--(5.16), (5.21) and the sentence ``S5 keeps its
vanishing logarithmic loss'']{PaperII}. Since $W=4(1-\varepsilon)U_C$, this
is the second inequality. Its error is the error-table entry ``S5 BV
remainder'' plus the prime-measure and Euler conversion of that row
\cite[\S6.3]{PaperII}.

Stage three. By Paper~II, (A.14)--(A.14a), $4J_5(0)$ lies in the ball of
centre $c_5$ and radius $8.84\cdot10^{-57}$, so $J_5(0)\le\overline J_5$.
Paper~II also gives the independent rational enclosure
$P_5<4J_5(0)<P_5+R_5$ inside the same ball \cite[after (A.14a)]{PaperII}.

The chain uses $\kappa=0$, so the increment $J_5(\kappa)-J_5(0)$ does not
occur. Paper~II's allowance $E_{5,\kappa}=4\cdot10^{-8}$ is a summand of
$g_{\mathrm{BV}}=0.00000575$ in \eqref{eq:retB-fixed-cost} and stays there
once. Paper~II, Appendix~A, records the hypothetical increment and its
once-booked decomposition. They are in the paragraph opening
\emph{For comparison, if S5-top used $\kappa=10^{-6}$} and in the
BV-allowance line right after it
\cite[Appendix~A, comparison paragraph and following decomposition]{PaperII}.
Equation (A.14a) gives only the nominal value at $\kappa=0$.
\end{proof}

The chain $\mathrm{T5}$ does not meet the changes of this section. Its
three stages use no sieve level $\ell_c$ and no saved profile. They use no
Class II, packing or Class I term. Paper~II's prepared Class II, packing and Class I
contributions of this row are zero \cite[\S2.4, paragraph ``S5-top and S6
use BV separately'']{PaperII}. The lift below changes the kernels on $I_4$,
$I_5=[1/4,t_*)$, $I_7$ and $\mathcal D_{10}$ only. The interval $I_5$ is disjoint
from $[t_*,\tau_S)$. The joint shift changes the $G_1$ cone only. Hence
Lemma~\ref{lem:retB-S5top} holds unchanged after both changes.

Keep the whole fixed deduction of Paper~II:
\begin{equation}\label{eq:retB-fixed-cost}
\begin{aligned}
 F_{\mathrm{net}}={}&0.000006484294227956901+0.00000575
                   +0.00002853328\\
                  &+4.98\cdot10^{-54}+2\cdot10^{-24}+10^{-390}\\
 &<0.000040767574227956903.
\end{aligned}
\end{equation}
These are Paper~II's six deductions \cite[(6.6)]{PaperII}.
The first remains paid although its source is excluded.
The joint $G_1$ reconciliation also remains paid.
Every adverse row difference has an absolute bound.
Keep their whole vanishing envelope, $E_{\mathrm{van}}$.
The row estimates required here are
\begin{equation}\label{eq:retB-base-errors}
 \sum_{j\in\mathcal I}\sum_{k=1}^{10}
       (Q_{j,k-1}-Q_{j,k})
 \ge-\bigl(F_{\mathrm{net}}-\delta_{G1}\bigr)
       -E_{\mathrm{van}}(N;\sigma),
 \qquad E_{\mathrm{van}}(N;\sigma)\longrightarrow0.
\end{equation}
The $\delta_{G1}$ part is already in the endpoint sum, so the remaining
chain deduction is $F_{\mathrm{net}}-\delta_{G1}$. Their total is exactly
$F_{\mathrm{net}}$.
The fixed AA37 share is paid under the admissible condition
$0<\sigma\le4g_{\mathrm{AA}}/(C_EW_*)$ of Paper~II, (1.27)
\cite{PaperII}. See Paper~II, \S6.5. This condition is part of the
fixed-parameter choice made before the limit.
The analytic inputs are Paper~II, (6.4b)--(6.5),
and its direct and $G_1$ estimates
\cite[(4.66)--(4.69), (6.4b)--(6.5)]{PaperII}.
Their logarithmic parameters are
\begin{equation}\label{eq:retB-precision}
 K_0=8042,\qquad A'=8047,\qquad A_{\mathrm{work}}=8051,
 \qquad K_0+2-A'=-3.
\end{equation}
These are Paper~II, (4.66), for the original coefficient families.
This section leaves them unchanged.

\begin{lemma}[Inherited chain errors]\label{lem:retB-base-errors}
The bound \eqref{eq:retB-base-errors} holds, and $E_{\mathrm{van}}(N;\sigma)\to0$
as $N\to\infty$ at fixed $\sigma$. Each row or separately listed stage
difference used below has an allocated nonnegative share in Paper~II's
error table. Denote it by $\mathcal E_j$. Paper~II, \S6, charges these shares once.
An entry may be grouped across rows, so the lemma claims nothing for
arbitrary subfamilies and does not claim pairwise disjoint entries. The only
omission bounds used are those for the whole chains or stage differences
named in the proof below.
\end{lemma}

\begin{proof}
The signed-assembly theorem of Paper~II proves
$D_{\mathrm{src}}\ge B_{\mathrm{main}}-F_{\mathrm{pers}}-E_{\mathrm{van}}$
with $E_{\mathrm{van}}\to0$ for the source (6.1)
\cite[\S6.1, (6.1)--(6.2), (6.4)--(6.4d)]{PaperII}. Its proof bounds every
adverse stage difference by an allocated share of Paper~II's error table,
\S6.3. Each entry is a nonnegative bound or a proved zero; some entries are
grouped across rows. Paper~II assigns each tuple to its first failed
condition and charges each difference or part once below
\cite[\S6.3]{PaperII}. The fixed entries are the six deductions of
\eqref{eq:retB-fixed-cost}, that is $F_{\mathrm{pers}}=F_{\mathrm{net}}$
\cite[(6.6)]{PaperII}. The others tend to zero at fixed $\sigma$, and their
sum is $E_{\mathrm{van}}$. The rows of $\mathcal I$ are rows of (6.1).
The $S_6$ level increment is still deducted when that chain is omitted.
The $\delta_{G1}$ part is already
in \eqref{eq:retB-base-endpoints}. This gives \eqref{eq:retB-base-errors}.
For $\mathrm{T5}$ the differences and entries are those of
Lemma~\ref{lem:retB-S5top}.
Separately listed entries cover the omissions used in this chapter.
The whole $S_6$ and $S_{\mathrm{sw}}$ chains use the S6 BV remainder,
S6 level increment and $E_{S_{\mathrm{sw}}}$ entries. The whole $G_1$
chain uses the \S4.6 table and (4.85) of Paper~II
\cite[\S4.6, (4.85)]{PaperII}. The stage-nine strip for the retained
part of $S_{10}$ uses the S10 selector entry. These bounds apply
to the named whole chains or stage difference after allocating their shares.
No arbitrary omitted subfamily is asserted. For the direct rows, the stage-nine entry
``Profile comparison'' is zero adverse. The stage-ten entry ``Directed
enclosure'' is zero additional \cite[\S6.3]{PaperII}. For the retained part
of $S_{10}$, the strip between the retained cells and $\mathcal D_{10}$ is the
entry ``S10 selector'', $7.160\tau$. The term $g_\tau=2\cdot10^{-24}$ replaces it
\cite[\S6.3, rows ``S10 selector'' and $g_\tau$]{PaperII}.

The entries rest on the following inputs, all used by Paper~II on these
same sources.
The linear-sieve main terms and remainders are Paper~I, Theorem~2.6,
(2.30)--(2.33b). This is the structured form of Iwaniec's Theorem~4
\cite[Theorem~4, p.~311]{Iwaniec1980Remainder}, applied with $D=L_c$,
$\epsilon_s=\sigma$ and $\delta_s=\tau$ \cite[\S6.2]{PaperII}.
Before the exterior sign, applied once, each block remainder splits as
\[
 \frac{C_c}{M_c}\langle\gamma_c,\Delta_P\rangle+NC_{\mathrm{sf},c}
 -\langle\Xi_{U,c},\Delta_P\rangle+\langle\Xi_{A,c},\Delta_0-\Delta_P\rangle,
\]
with $\Xi_{A,c}=O_c\mathbf 1_{q\ \mathrm{not\ squarefree}}$
\cite[(6.4a)--(6.4b), (4.23)]{PaperII}.
The term $NC_{\mathrm{sf}}$ vanishes on the direct rows, where
$\Delta_0=\Delta_P$. On $G_1$ it is $O_\sigma(N^{-\alpha}L^5)$ by Paper~II,
(4.64)--(4.65). The terms $\Xi_U$ and $\Xi_A$ are bounded jointly, once, by
$232\Lambda_\sigma C_{\mathrm{ret},\sigma} N^{-\ell_0\sigma^2}(1+L)^4$, with
$\ell_0=264897/2060000$ \cite[(B.1)--(B.7)]{PaperII}.
Here $C_{\mathrm{ret},\sigma}$ is the retained-side combinatorial constant;
it depends on the fixed $\sigma$ and not on $N$.
The persistent part $C_E\sigma W_{\mathrm{AA}}/4$ of Paper~II, (B.8), lies in
$g_{\mathrm{AA}}$; only its secondary part belongs to $E_{\mathrm{van}}$.
The Class II estimate (6.5) of Paper~II applies Paper~I, Theorem~3.2, (3.13),
with the precisions \eqref{eq:retB-precision}. The spectral input of that
theorem is Paper~I, Theorem~2.8; see Fact~\ref{fact:spectral-implication}.
In the large-norm case of Paper~I, (3.48), its mean-square input is
Lemma~\ref{hyp:ra-ms-transfer}(ii), in place of \cite[Theorem~0(a)]{BFI86}.
Lemma~\ref{hyp:ra-ms-transfer}(ii) applies under the Paper~I
hypotheses and selected parameter box stated there.
It uses neither the optional six-endpoint lemma nor any broader BFI86 scope.
Its premises are checked for these sources as follows. In Paper~I, conditions (3.1), (3.2), (3.6), (3.8) with
the length condition, and (3.11) display its hypotheses (H1)--(H5). Paper~II verifies them for these same sources, coefficients,
boxes and scales \cite[\S6.2, hypothesis table, row ``C1--C5, LEN'']{PaperII}.
Paper~I proves condition (3.25) for every selected factor of this form.
The proof uses the prime Siegel--Walfisz theorem (2.21) \cite[(3.25) and its
proof]{PaperI}. Conditions (3.26)--(3.27) are the definition of $\lambda_r$
and the Cauchy--Schwarz step of that proof. They hold under (H1)--(H5), with
$|\lambda_r|\le\tau(r)$ from the roughness in (H2) \cite[(3.26)--(3.27)]{PaperI}.
Stages one to eight are unchanged, so these are Paper~II's own checks. The prime-measure and Euler
conversions are the prime number theorem, Mertens' theorem and Paper~II,
(4.76)--(4.78).
The profile lift and the shifted joint profile below act at stages nine
and ten only. They change no source, cell, level, mask, box, norm or
precision at stages one to eight. Hence $E_{\mathrm{van}}$ is Paper~II's own
envelope for these rows.
\end{proof}

Paper~II's rational floor satisfies
\begin{equation}\label{eq:retB-paperII-floor}
 B_{\mathrm{print}}-F_{\mathrm{net}}
 >D_{\mathrm{II}}
 =\frac{45849916338009170593}{25000000000000000000};
\end{equation}
see \cite[(6.3a), (6.6)]{PaperII}.
The selected S2 endpoint used in this floor is $d_{T8}$.
Paper~II's tail reserve is
\begin{equation}\label{eq:retB-T8-reserve}
 \delta_{T8}=\frac1{20000},\qquad
 d_{T8}=\frac{P_--(2/e^\gamma)V_{\mathrm{full}}
                 -(2/e^\gamma)J_++C_-}{4}-\delta_{T8}.
\end{equation}
Thus $\delta_{T8}$ is deducted once inside $d_{T8}$, as in
Paper~II, (A.7)--(A.8). It is not a second deduction from
$D_{\mathrm{II}}$. The separate fixed term $275009680\tau$ remains
outside $d_{T8}$ in the S2 estimate \cite[(3.22)--(3.23), (A.7)]{PaperII}.
Combining the preceding estimates gives
\begin{equation}\label{eq:retB-base}
 \frac{R_{\mathrm{fix}}+S_2}{W}
 \ge D_{\mathrm{II}}+\overline J_6+2\overline H
                         -E_{\mathrm{van}}(N;\sigma).
\end{equation}
This deduction uses the individual signed chains.
Its analytic input is \eqref{eq:retB-base-errors}.

\subsection{Six improved endpoints}

Let
\begin{equation}\label{eq:retB-six-index}
 \mathcal J=\{1,2,4,5,7,10\}.
\end{equation}
For $j\in\mathcal J$, let $P_j$ be its Paper~II profile integral.
The indices $5$ and $10$ keep the preceding conventions.
Let $O_j$ and $N_j$ be directed lower endpoints.
The required numerical comparisons are
\begin{equation}\label{eq:retB-six-comparison}
 P_j\ge N_j\ge O_j,\qquad
 P_j-O_j=(N_j-O_j)+(P_j-N_j).
\end{equation}
The old and improved directed endpoints are the following exact terminating
decimals. They are given in the normalized D
contributions used by the signed sum. For S4, S5-low and S7 the sign $-1/4$ is included; for S10
the sign $+1/4$ is included.
\begin{center}
\small
\renewcommand{\arraystretch}{1.08}
\begin{tabular}{@{}lp{0.68\linewidth}@{}}
\hline
Endpoint source & Directed contribution in D units\\
\hline
S1, \cite[(A.4), (A.5a)]{PaperII} & $O_1=9.67653874894124412908238749318694090949698021569432709634$\\
 & $N_1=9.676557035420731097796376157212899872020518126053458552042$\\
S2, \cite[(A.8), (A.8a)]{PaperII} & $O_2=1.6195813569890835255913340433210572922007671408165$\\
 & $N_2=1.62063469556829$\\
S4, \cite[Appendix~A.2--A.6]{PaperII} & $O_4=-0.8342438370724665$\\
 & $N_4=-0.8341817006889064605174981431060625$\\
S5-low, \cite[Appendix~A.2--A.6]{PaperII} & $O_5=-0.05972210829776175$\\
 & $N_5=-0.05971885506770766576732140865035175$\\
S7, \cite[Appendix~A.2--A.6]{PaperII} & $O_7=-0.574850630119163$\\
 & $N_7=-0.574815106890481420929612425345365$\\
S10-Ret, \cite[Appendix~A.2--A.6]{PaperII} & $O_{10}=0.214711863944485$\\
 & $N_{10}=0.21559427153309934685978873940373255$\\
\hline
\end{tabular}
\end{center}
The rational endpoint differences sum to
\begin{equation}\label{eq:retB-six-gain}
\begin{aligned}
 g_6&=\sum_{j\in\mathcal J}(N_j-O_j)\\
 &=0.002054945489603492768011383006854970322770769542631455702.
\end{aligned}
\end{equation}
Exact rational addition of the six entries above gives the displayed
sum, and each entry satisfies $N_j\ge O_j$.
The values $O_j$ are Paper~II's directed endpoints, and $P_j\ge O_j$ is
proved there \cite[(A.4), (A.5a), (A.8), (A.8a), Appendix~A.2--A.6]{PaperII}.
The values $N_j$ are lower bounds read from Paper~II's printed independent
enclosures. The following optional lemma states them.
\begin{lemma}[Six directed endpoints from Paper II; optional]
\label{lem:retB-six-endpoints}
For each $j\in\mathcal J=\{1,2,4,5,7,10\}$, the Paper~II profile integral
$P_j$, with its signed normalized coefficient, satisfies $P_j\ge N_j$,
with $N_j$ as tabulated above.
\end{lemma}
\begin{proof}
For $j=1$, Paper~II, (A.5a), gives the printed $3/4$ ball; its lower endpoint
is $N_1$. For $j=2$, Paper~II, (A.8a), gives the printed interval whose
lower endpoint is $N_2$. Paper~II states that this interval covers the continuous functional and the infinite
tail. The $d_{T8}$ deduction in (6.6a) is released once in $N_2-O_2$.
The term $F_{\mathrm{net}}$ in (6.6b) contains no $\delta_{T8}$
\cite[(6.6a)--(6.6b)]{PaperII}. For $j=4,5,7$, the Appendix~A.3 balls
include the adverse $\tau=10^{-90}$ shifts. Since the profile function is
nonincreasing, the shifted raw integrals are at least the $\tau$-free ones.
Their upper endpoints with the negative coefficient therefore give
$P_j\ge N_j$. For $j=10$, the Appendix~A.3 $S_{10}$-Ret ball encloses the
$\tau$-free actual-coordinate profile on $\mathcal D_{10}$. Its lower
endpoint gives $P_{10}\ge N_{10}$. The separate $D_*\tau$ strip is paid in
the $S_{10}$ selector estimate. These are the printed enclosures in Paper~II,
(A.5a), (A.8a), and Appendix~A.3. Its numerical-certificate section records
the adverse-shift and actual-coordinate conventions
\cite[(A.5a), (A.8a), Appendix~A.3]{PaperII}. The remaining directed
endpoints $O_j$ are unchanged.
\end{proof}
With $O_j$ in place of $N_j$, \eqref{eq:retB-six-comparison} is
Paper~II's own endpoint comparison and $g_6$ is replaced by $0$.
Only the endpoint changes in \eqref{eq:retB-telescope}.
All earlier signed differences remain unchanged.
Thus \eqref{eq:retB-six-comparison} adds $g_6$ once.
The $S_5$ upper part remains separate.
The joint $G_1$ endpoint is absent from $\mathcal J$.

\subsection{Four stronger profile comparisons}

Use the sieve functions of Paper~I
\cite[(2.23), immediately before Theorem~2.6]{PaperI}. Put
\begin{equation}\label{eq:retB-hat-functions}
 \widehat F=\frac{F}{2e^\gamma},\qquad
 \widehat f=\frac{f}{2e^\gamma},\qquad
 K=\frac{57}{64}-u.
\end{equation}
Define
\begin{equation}\label{eq:retB-H6-domains}
\begin{aligned}
 s(x)&=\frac{K-5x/2}{\alpha},\qquad
 q(x,y)=\frac{K-5x/2-y}{\alpha},\\
 B_0(x)&=\min\{\beta,\max\{\alpha,K-11\alpha/5-5x/2\}\}.
\end{aligned}
\end{equation}
The domains are
\begin{equation}\label{eq:retB-H6-domain-sets}
\begin{gathered}
 I_4=[\mu_*,1/4),\quad I_5=[1/4,7109/28000),\quad
 I_7=[\mu_*,\lambda),\\
 \mathcal D_{10}=\{(x,y):\mu_*\le x<\lambda,\ \alpha\le y<B_0(x)\}.
\end{gathered}
\end{equation}
These are Paper~II's domains
\cite[(2.27), (2.29), (2.31), (2.38)--(2.40)]{PaperII}.
Set
\begin{equation}\label{eq:retB-H6-gains}
\begin{aligned}
 \Delta_j&=\int_{I_j}
  \frac{\widehat F(s(x))-\widehat F(s(x)+h/\alpha)}{\alpha x}\,dx
                      &&(j=4,5,7),\\
 \Delta_{10}&=\iint_{\mathcal D_{10}}
  \frac{\widehat f(q(x,y)+h/\alpha)-\widehat f(q(x,y))}
       {\alpha xy}\,dx\,dy.
\end{aligned}
\end{equation}
The three upper rows carry a negative sign.
The lower $S_{10}$ row carries a positive sign.
Their original coefficients are included in these gains,
so there is no further factor $1/4$.

Paper~II's actual levels give
\begin{equation}\label{eq:retB-H6-margins}
\begin{aligned}
 \ell_c&\ge\alpha s(x)+h+\frac1{10000}-3\tau
                   &&(j=4,5,7),\\
 \ell_c&\ge\alpha q(x,y)+h+\frac1{10000}
                         -\left(\frac{5\lambda}{2}+\beta\right)\tau
                   &&(j=10).
\end{aligned}
\end{equation}
See \cite[\S2.4, proof of (2.26)--(2.31);
\S2.5, proof of (2.40)]{PaperII}.
The exact split behind the displayed reserve is
\begin{equation}\label{eq:retB-level-buffer}
 u-\eta_L=\frac{19}{10000}
 =\frac9{5000}+\frac1{10000}=h_G+\frac1{10000}.
\end{equation}
After the $h_G$ lift, Paper~II's S4, S5-low (including its B branch)
and S7 comparisons leave $1/10000-3\tau$. Its S10-Ret comparison leaves
$1/10000-\tau(5\lambda/2+\beta)$.
Both are positive for $\tau\le10^{-90}$.
The G1-3 shift leaves
$g_s=1/10000-(225/256)\tau>2\mu_*\tau$ and the $\mu_*$ crossing
leaves $1/10000-(75/256)\tau>0$, by \cite[(4.72)]{PaperII}.
These are pointwise profile-numerator inequalities. They do not change
the actual level $\ell_c$, the level $D_c$, or the masks, boxes,
Class-II and AA37 inputs. The other level conditions of Paper~II
remain pointwise and unmodified. They include the G1 guard, S1--S3,
S2/S3/S8/S9, and the whole-K margin. See Paper~II, (2.21), (2.40),
(4.72), and (4.19).
Decreasing $F$ gives the first three comparisons.
Increasing $f$ gives the fourth comparison.
These operations leave the actual levels unchanged.
The H6 lift and the M5 shift act only at stages 9--10, after Class II
at stage 7. They change no Class-II, AA37, box, mask or precision input:
$\ell_c=\Theta_c-\eta_L-\sum_i t_i$ contains no saved profile, and
\eqref{eq:retB-precision} remains unchanged.

Proposition~\ref{prop:app-h6-four-profiles} gives finite integral
enclosures. With the directed calculation of the computational appendix,
they prove the following. Its exact rational interval is
\par\noindent
$[0.005720047094327988714222991186379229149381887590042773120871489672965022499
\mathbin{\pm}7.09\cdot10^{-76}]$.
This is a ball for the sum of the four rowwise lower bounds, not for
$\sum_j\Delta_j$.
We use the terminating credit $572004/10^8$. It lies below
that exact lower endpoint and leaves the existing $10^{-20}$ allowance intact.
This allowance is conservative; it is not an additional gain in the final constant.
\begin{equation}\label{eq:retB-H6-finite}
 \frac{572004}{10^8}
 <\Delta_4+\Delta_5+\Delta_7+\Delta_{10}
 <\frac{606486}{10^8},\qquad
 \frac{572004}{10^8}-10^{-20}>\frac{5719}{10^6}.
\end{equation}
The unchanged $S_{10}$ strip also satisfies
\begin{equation}\label{eq:retB-S10-strip}
 \frac{10987/10000}{4\alpha^2}
 \left(\frac{5\lambda}{2}+\beta\right)
 \frac{\lambda-\mu_*}{\mu_*}
 =\frac{21137254516197687}{4120000000000000}
 <5.137<7.160.
\end{equation}
The lifted $S_{10}$ argument satisfies
$q+h/\alpha\le11519839/3200000+h/\alpha<4$.
Thus $f\le f(4)=2e^\gamma\log3/4$ throughout the strip,
as used in this density bound.
These are finite bounds, including their stated domains.
The retained cells also leave a reserve. Put
\begin{equation}\label{eq:retB-S10-surplus}
\begin{gathered}
 R_{10}=\Bigl[\frac15,\frac{21}{100}\Bigr]\times\Bigl[\frac9{100},\frac1{10}\Bigr],
 \qquad \delta_{10}=\frac1{20000\alpha},\\
 c_{10}=\frac1\alpha\cdot\frac{7\delta_{10}}{480}\cdot\frac1{210}
 =\frac{146689}{320000000000},\qquad
 c_{10}-\frac{21137254516197687}{4120000000000000}\,\tau>10^{-7}.
\end{gathered}
\end{equation}
Lemma~\ref{lem:retB-lifted-direct} uses \eqref{eq:retB-S10-strip} and
\eqref{eq:retB-S10-surplus}, so its comparison over the whole of
$\mathcal D_{10}$ needs no deduction. Paper~II's $S_{10}$ selector entry
$7.160\tau$ is replaced in Paper~II by $g_\tau=2\cdot10^{-24}$, a summand of
\eqref{eq:retB-fixed-cost}. The fixed $g_\tau$ amount is deducted once.
The retained reserve above eliminates the $S_{10}$ selector deduction in
this comparison. The inherited $G_1$ peel and its label replacement are
jointly covered by that same summand.

We next define the finite profile arrays.
Let $X_0=[u,1]$. On this compact interval, use the zero extension of the
profile outside its displayed domain, and put
\begin{equation}\label{eq:retB-measure}
 m_N=\sum_{\substack{p\in\mathbb P,\ p>2,\ p\nmid N\\N^u\le p\le N}}
       \frac{\delta_{\log_Np}}{p-1},\qquad m(dt)=\frac{dt}{t}.
\end{equation}
The symbol $\delta_x$ denotes unit mass at $x$.
Define
\begin{equation}\label{eq:retB-Euler}
 V_N(z)=\prod_{\substack{p<z,\ p\in\mathbb P\\p\nmid N}}
           \left(1-\frac1{p-1}\right),
\end{equation}
and
\begin{equation}\label{eq:retB-direct-normalization}
\begin{aligned}
 X_D&=\sum_{v\in\mathbb P}
        \mathbf1_{N/L^{30}<v<(1-\varepsilon)N},\\
 r_{D,N}&=\frac{2e^\gamma\alpha X_DV_N(N^\alpha)}{W}.
\end{aligned}
\end{equation}
For $j=4,5,7$, set
\begin{equation}\label{eq:retB-direct-base-kernels}
\begin{aligned}
 k_j^0(x)&=-\frac{\mathbf1_{I_j}(x)}{\alpha}\widehat F(s(x)),\\
 k_j^h(x)&=-\frac{\mathbf1_{I_j}(x)}{\alpha}
                       \widehat F(s(x)+h/\alpha).
\end{aligned}
\end{equation}
For $j=10$, set
\begin{equation}\label{eq:retB-pair-base-kernels}
\begin{aligned}
 k_{10}^0(x,y)&=\frac{\mathbf1_{\mathcal D_{10}}}{\alpha}\widehat f(q(x,y)),\\
 k_{10}^h(x,y)&=\frac{\mathbf1_{\mathcal D_{10}}}{\alpha}
                            \widehat f(q(x,y)+h/\alpha).
\end{aligned}
\end{equation}
Put $e_j=1$ for $j=4,5,7$, and $e_{10}=2$.
The original and lifted arrays are
\begin{equation}\label{eq:retB-direct-arrays}
\begin{aligned}
 P_j&=\int k_j^0\,d(m^{e_j}),&
 S_{j,N}^0&=r_{D,N}\int k_j^0\,d(m_N^{e_j}),\\
 S_{j,N}^h&=r_{D,N}\int k_j^h\,d(m_N^{e_j}),&
 B_{j,N}&=S_{j,N}^h-S_{j,N}^0.
\end{aligned}
\end{equation}
The powers denote product measures.
These formulas use $P_j$ from \eqref{eq:retB-six-comparison}.
Write $d_j=k_j^h-k_j^0$. For $j=4,5,7$, and then $j=10$,
\begin{equation}\label{eq:retB-direct-kernels}
\begin{aligned}
 d_j(x)&=\frac{\mathbf1_{I_j}(x)}{\alpha}
  [\widehat F(s(x))-\widehat F(s(x)+h/\alpha)],\\
 d_{10}(x,y)&=\frac{\mathbf1_{\mathcal D_{10}}}{\alpha}
  [\widehat f(q(x,y)+h/\alpha)-\widehat f(q(x,y))].
\end{aligned}
\end{equation}
Thus $\Delta_j=\int d_j\,d(m^{e_j})$.
Define the two conversion errors:
\begin{equation}\label{eq:retB-direct-Lambda}
 \Lambda^0_{j,N}=S^0_{j,N}-P_j,\qquad
 \Lambda_{j,N}=B_{j,N}-\Delta_j\qquad(j\in\mathcal H).
\end{equation}

\begin{lemma}[Finite profile shift]\label{lem:retB-finite-profile}
For the arrays above,
\begin{equation}\label{eq:retB-H6-use}
 S_{j,N}^h=P_j+\Lambda^0_{j,N}+\Delta_j+\Lambda_{j,N}
 \ge P_j+\Delta_j-|\Lambda^0_{j,N}|-|\Lambda_{j,N}|.
\end{equation}
\end{lemma}
\begin{proof}
Substitute the two definitions in \eqref{eq:retB-direct-Lambda}.
\end{proof}
Both absolute conversion terms appear as separate summands of $E'$.
We do not assume that Paper~II's original envelope contains the
saved-array error. Its stage 8 converts the actual-level array before
stage 9 domination \cite[\S6.1, stage table]{PaperII}.
The arithmetic application uses the preceding level domination.

\paragraph{Restoration.}
The four rows of $\mathcal H$ are Paper~II's rows $S_4$, the lower part
of $S_5$, $S_7$ and the retained part of $S_{10}$. Their stage-two
source-restoration shares are assigned once to the source parents $S_4$,
$S_5$, $S_7$ and $S_{10}$. The literal $S_5$ split at $7109/28000$
is made before the structured sieve or Heath--Brown expansion, as above.
Both parts retain the parent $S_5$ for this restoration accounting
\cite[(2.28), \S6.1, stage table and the list of S10, S5 and S4/S7 endpoints]{PaperII}.
There is no separate restoration for the lower part of $S_5$. The upper part
of $S_5$ is the outer chain $\mathrm{T5}$ of Lemma~\ref{lem:retB-S5top},
which the lift does not change. Each tuple of an actual source is assigned
to one entry by Paper~II's first-failure rule \cite[(6.4d) and the two
sentences after it]{PaperII}. This section changes no source and no stage
from one to eight, so the assignment is Paper~II's. Appendix~\ref{subsec:app-retB-rows} lists the finite
correspondence between these rows and Paper~II's error entries. It is an index of that
assignment, and each entry occurs once.

\begin{lemma}[Lifted direct comparison]\label{lem:retB-lifted-direct}
Let $j\in\mathcal H$, and let $Q_{j,8}$ be Paper~II's stage-eight quantity
for this row. It is the continuum mass of the actual-level main term after the
prime-measure and Euler conversion. There each actual cell carries its
actual relative level $\ell_c$. Put $Q^h_{j,9}=P_j+\Delta_j=\int
k_j^h\,d(m^{e_j})$. Then $Q_{j,8}\ge Q^h_{j,9}$.
\end{lemma}

\begin{proof}
Stages one to eight are Paper~II's. Paper~II passes from $Q_{j,8}$ to
$Q_{j,9}=P_j$ by pointwise domination of the actual-level integrand by the
specified profile. The cells and the continuum measure are the same
\cite[\S6.1, stage table; \S2.4, after (2.31); \S2.5, after (2.40)]{PaperII}.
We repeat that domination with the lifted profile.
For each B-labelled parent contributing to $S_4$, $S_7$ or $S_5$-low,
we first rederive the buffer used below. Write
\[
 m(t)=\min\left\{\frac{1+t}{2},\frac{57}{64}-\frac{3\rho t}{2}\right\},
 \qquad t_{\rm dd}=\frac{25}{32(1+3\rho)}.
\]
Paper~II gives $\Theta_c\ge m(t_1)$. The branch comparison is the exact
identity
\[
 \frac{1+t_1}{2}
 -\left(\frac{57}{64}-\frac{3\rho t_1}{2}\right)
 =\frac{1+3\rho}{2}(t_1-t_{\rm dd}).
\]
Thus the decreasing branch is active exactly when $t_1>t_{\rm dd}$.
Equality is assigned to A, where both bounds agree. This active branch test
is applied to every B-labelled parent. The branch is never inferred from the
largest completed coordinate alone. Paper~II keeps the inherited label through
clipping. In S4 and S7 the B parents are separate intersections of their
respective row domains with $t_1>t_{\rm dd}$. No common parent interval is
asserted. S5-low has $1/4\le r<t_*=7109/28000$ and
$t_1\le r<\rho t_1$, so $t_1>1/(4\rho)>t_{\rm dd}$ because $\rho<8$.
It is therefore entirely B. Let a parent in S4 or S7 be A-labelled.
Then $t_1\le t_{\rm dd}$, and Paper~II's first branch in (2.17f)
\cite[(2.17f)]{PaperII} gives
$\Theta_c\ge(1+t_1)/2$. Take an actual exponent $r\in[t_1,\rho t_1)$.
The S4 and S7 row domains are clipped by $r\ge\mu_*$
\cite[\S2.1, fixed-cutoff rows]{PaperII}, and the retained bin gives
$r\ge t_1$.
The residual is therefore
\[
 2(t_1-\mu_*)+\frac52(r-t_1)
 =2(r-\mu_*)+\frac12(r-t_1)\ge0.
\]
\[
 \ell_c-\left({57\over64}-u-{5r\over2}\right)
 \ge u-\eta_L-{25\over64}+2t_1+{5\over2}(r-t_1)
 = {19\over10000}+2(t_1-\mu_*)+{5\over2}(r-t_1)
 \ge {19\over10000}.
\]
Here $\mu_*=25/128$ is the retained cutoff parameter; the symbol $\mu(n)$
elsewhere denotes the Möbius function.
Thus the A pieces also have the required $19/10000-3\tau$ buffer. This is a row-local calculation.
It does not merge the distinct S4 and S7 parent domains. For any B parent,
Paper~II's floor gives
\begin{equation}\label{eq:retB-B-branch-buffer}
 \Theta_c\ge {57\over64}-{3\over2}\rho t_1,\qquad
 \ell_c-\left({57\over64}-u-{5r\over2}\right)
 \ge (u-\eta_L)+{5\over2}(r-t_1)-{3\over2}\tau t_1
 \ge {19\over10000}-3\tau .
\end{equation}
The first bound is Paper~II's linear branch of the full-anatomy floor
\cite[(2.17d)--(2.17f)]{PaperII}. Using
$\ell_c=\Theta_c-\eta_L-t_1$, $r\ge t_1$, $u-\eta_L=19/10000$, and
$t_1<t_*<1$ proves the second line by direct subtraction. The B pieces of
S4 and S7 use the same row-local calculation. Their parent intervals
stay distinct.

For $j=4,5,7$ the row coefficient is negative. At a point $x$ of an actual
cell, the actual sieve argument is $\ell_c/\alpha$ and, by
\eqref{eq:retB-H6-margins},
\[
 \frac{\ell_c}\alpha\ge s(x)+\frac h\alpha+\frac{10^{-4}-3\tau}\alpha
 >s(x)+\frac h\alpha .
\]
These are Paper~II's margins $19/10000-3\tau$ for $S_4$, $S_7$ and the
$B$ branch of the lower part of $S_5$. Here $h$ has moved into the profile
\cite[\S2.4, display before ``The S4/S7 A/B levels give the same buffer'']{PaperII}.
Since $F$ is decreasing, $-\widehat F(\ell_c/\alpha)\ge-\widehat F(s(x)+h/\alpha)$.
For $j=10$ the coefficient is positive, and on the retained cells
\[
 \frac{\ell_c}\alpha\ge q(x,y)+\frac h\alpha
 +\frac{10^{-4}-(5\lambda/2+\beta)\tau}\alpha>q(x,y)+\frac h\alpha,
\]
which is Paper~II's
$\alpha(s_{\mathrm{live}}-s_{\mathrm{saved}})\ge(u-\eta_L)-\tau((5/2)\lambda+\beta)$
with $h$ moved into the profile \cite[\S2.5, proof after (2.40)]{PaperII}.
Since $f$ is nondecreasing,
$\widehat f(\ell_c/\alpha)\ge\widehat f(q(x,y)+h/\alpha)$.
Paper~II uses the same monotonicity of $F$ and $f$ for these
comparisons \cite[\S2.4]{PaperII}.
For $j=4,5,7$ the cells cover $I_j$, and integrating the pointwise
inequality against $m$ gives $Q_{j,8}\ge\int k_j^h\,dm=P_j+\Delta_j$.

For $j=10$ the retained cells need not cover $\mathcal D_{10}$. Let
$R_\tau\subset\mathcal D_{10}$ be their union, put
$a_{10}=5\lambda/2+\beta$, and write $T_h(E)=\int_Ek_{10}^h\,d(m^2)$. Then
\[
 Q_{10,8}=\frac1\alpha\sum_{c\ \mathrm{retained}}\int_c
 \widehat f(\ell_c/\alpha)\,d(m^2).
\]
\emph{Strip.} Let $(x,y)\in\mathcal D_{10}\setminus R_\tau$ lie in a cell
with labels $t_1\le x$, $t_2\le y$. The cell is dropped, so
$y\ge t_2>\operatorname{sel}(t_1)\ge B_0(t_1)-a_{10}\tau\ge B_0(x)-a_{10}\tau$.
This uses the two deficit bounds of Paper~II, \S2.5, and the fact that $B_0$ is
nonincreasing \cite[\S2.5, proof of (2.40)]{PaperII}. Also $y<B_0(x)$. On
$\mathcal D_{10}$ one has $11/5\le q$ and $q+h/\alpha<4$, so
$\widehat f(q+h/\alpha)\le\widehat f(4)=\log3/4$. With $y\ge\alpha$ and
$\int_{\mu_*}^\lambda dx/x\le(\lambda-\mu_*)/\mu_*$, this gives
$T_h(\mathcal D_{10}\setminus R_\tau)\le D_*\tau$, where $D_*$ is the
constant in \eqref{eq:retB-S10-strip}.

\emph{Reserve.} On a retained cell the margin in
\eqref{eq:retB-H6-margins} satisfies $10^{-4}-a_{10}\tau>1/20000$, so
$\ell_c/\alpha\ge q+h/\alpha+\delta_{10}$. Here
$2<q+h/\alpha<q+h/\alpha+\delta_{10}<4$. For \(2<s<4\),
$\widehat f(s)=\log(s-1)/s$ and
\[
 \widehat f'(s)=\frac{s/(s-1)-\log(s-1)}{s^2}
 >\frac{4/3-11/10}{16}=\frac7{480},
\]
since $\log3<11/10$. As $f$ is nondecreasing,
$\widehat f(\ell_c/\alpha)-\widehat f(q+h/\alpha)\ge7\delta_{10}/480$ at
every point of $R_\tau$.

\emph{A fixed retained rectangle.} Let $(x,y)\in R_{10}$, in a cell with
labels $t_1\in(x/\rho,x]$, $t_2\in(y/\rho,y]$. Then $\mu_*<t_1<\lambda$ and
$\alpha<t_2<\beta$. On $[\mu_*,21/100]$ we have $B_0=\beta$, because
$57/64-u-11\alpha/5-(5/2)(21/100)>\beta$. Hence
$\operatorname{sel}(t_1)\ge\beta-a_{10}\tau>1/10\ge t_2$, and the cell is
retained: $R_{10}\subset R_\tau$. Also
$m^2(R_{10})=\log\frac{21}{20}\log\frac{10}9\ge\frac1{21}\cdot\frac1{10}
=\frac1{210}$.

Combining the three steps with the pointwise inequality above,
\[
 Q_{10,8}\ge T_h(R_\tau)+c_{10}
 \ge T_h(\mathcal D_{10})-D_*\tau+c_{10}>T_h(\mathcal D_{10})
 =P_{10}+\Delta_{10},
\]
by \eqref{eq:retB-S10-surplus}. No deduction is made. The reserve is the
part of Paper~II's level margin left after the lift $h$. It exists on cells
of $S_{10}$ only and is used only here.
\end{proof}

The finite arrays $S^0_{j,N},S^h_{j,N}$ are not used in this inequality.
The errors $|\Lambda^0_{j,N}|$ and $|\Lambda_{j,N}|$ are nonnegative and are
kept in $E'$ so that $E'$ agrees with \eqref{eq:intro-retained-error}.
They tend to zero by Lemma~\ref{lem:retB-conversion}.

\subsection{The internal joint shift}

The joint occurrence is $2S_{11}+S_{12}+S_{13}$.
Its six-branch function is Paper~II, (4.70)
\cite[(4.70)--(4.72)]{PaperII}:
\begin{equation}\label{eq:retB-theta}
\begin{aligned}
 v_1(t)&=\min\left\{\frac{1+t}{2},\frac{32-34t}{39},
                              \frac{57}{64}-\frac{3t}{2}\right\},\\
 v(t)&=v_1(t)-u,\qquad \psi(x,y)=x\mathbf1_{x\ge y},\\
 \Theta_*(a,b,c)&=\max\{v(a),v(b),v(a+b),v(a+b+c),\\
 &\hspace{34mm}\psi(v(a+c),a+2b+c),
                  \psi(v(b+c),2a+b+c)\}.
\end{aligned}
\end{equation}
Subtract $u$ before testing each guard.
Set
\begin{equation}\label{eq:retB-G-cone}
 \mathcal C=\{(t_1,t_2,t_3):u\le t_3<t_2<t_1<\mu_*\},
 \qquad
 a_3(\mathbf t)=\frac{\Theta_*(\mathbf t)-t_1-t_2-t_3}{t_3}.
\end{equation}
Shift the output of $\Theta_*$ by $h$.
Keep its guards and actual source unchanged:
\begin{equation}\label{eq:retB-G-shift}
\begin{aligned}
 a_3^s(\mathbf t)&=a_3(\mathbf t)+\frac h{t_3},\\
 J&=\iiint_{\mathcal C}\frac{1-f(a_3(\mathbf t))}
                      {t_1t_2t_3^2}\,d\mathbf t,\\
 J_s&=\iiint_{\mathcal C}\frac{1-f(a_3^s(\mathbf t))}
                      {t_1t_2t_3^2}\,d\mathbf t.
\end{aligned}
\end{equation}
Use $f(s)=0$ for $s<2$.
Keep its true values at every larger argument.
The new margins are
\begin{equation}\label{eq:retB-G-margins}
 g_s=\frac1{10000}-\frac{225}{256}\tau>2\mu_*\tau,
 \qquad
 g_{s,\mu_*}=\frac1{10000}-\frac{75}{256}\tau>0.
\end{equation}
The second margin concerns a parent interval crossing $\mu_*$.
Both follow by subtracting $h$ in Paper~II, (4.72); see
Lemma~\ref{lem:retB-shifted-retention}.
They hold throughout $0<\tau\le10^{-90}$.
If $f(a_3^s)>0$, the retained-cell condition follows:
\begin{equation}\label{eq:retB-G-retention}
 \ell_c\ge\Theta_*+h-t_1-t_2-t_3+g_s
       >2t_3+2\mu_*\tau\ge2\rho t_3.
\end{equation}
Thus the canonical restriction still loses zero.
At fixed numerator, use the increase of $sf(s)$.
This is the same comparison as Paper~II, (4.72).

\begin{lemma}[Shifted retention]\label{lem:retB-shifted-retention}
The margins \eqref{eq:retB-G-margins} hold for $0<\tau\le10^{-90}$.
On the original joint source, every positive shifted lower contribution
lies in an original retained whole cell of Paper~II, (4.19). On such a
cell the actual lower main term dominates the shifted canonical one.
\end{lemma}

\begin{proof}
Use an original G1-3 cell with inherited labels
$t_1\ge t_2\ge t_3>0$, and order its distinct actual prefix primes as
$a_1>a_2>a_3$. Put $x_i=\log_Na_i$. Unequal parent bins enforce the
actual order, while within an equal-label group we order the primes.
Hence $t_i\le x_i<\rho t_i$, including clipped bins.
No strict inequality between the labels is required.

Let $\Theta_c$ be the original full-anatomy infimum and
$\ell_c=\Theta_c-\eta_L-\sum_it_i$. Put
$B_s(x)=\Theta_*(x)+h-\sum_ix_i$.
The function $v_1$ is $3/2$-Lipschitz. On the closed ordered parent
a coordinate changes by at most $\mu_*\tau$. A sum of at most
three coordinates changes by at most $3\mu_*\tau$.
If a saved guarded branch is active at $x$, its buffered guard has
nonnegative gap, including equality. Removing its buffer adds $u$.
Throughout the parent its unbuffered gap is therefore at least
$u-(9/2)\mu_*\tau-4\mu_*\tau>0$.
Thus it stays available in the full anatomy. Take an unguarded or
active guarded saved branch. The corresponding full branch is at least
its value at $x$ plus $u-(9/2)\mu_*\tau$. Taking their maximum and then
the parent infimum preserves this lower bound. Ties between maximizing
branches cause no loss. Therefore
\[
 \ell_c\ge B_s(x)+g_s+\sum_i(x_i-t_i),\qquad
 g_s=u-\eta_L-h-(9/2)\mu_*\tau
     =10^{-4}-(225/256)\tau.
\]
For a parent crossing $\mu_*$, the largest exponent is within
$\mu_*\tau$ of $\mu_*$, and $v_1(\mu_*)=153/256$.
Every saved branch is at most $153/256-u$, so the crossing comparison
gives instead $10^{-4}-(75/256)\tau\ge g_s$.
Clipping only shrinks the infimum domain. These arguments also apply
when two or all three labels coincide.
For $0<\tau\le10^{-90}$ we have $g_s>2\mu_*\tau>0$.

If $f(B_s(x)/x_3)>0$, then $B_s(x)>2x_3$ and hence
\[
 \ell_c>2x_3+2\mu_*\tau
       \ge2t_3+2\tau t_3=2\rho t_3.
\]
Thus the whole original cell is retained by Paper~II, (4.19).
On positive support $\ell_c>B_s(x)>0$ and $0<t_3\le x_3$ imply
\[
 \frac{\ell_c}{t_3}\ge\frac{B_s(x)}{x_3},\qquad
 f(\ell_c/t_3)\ge f(B_s(x)/x_3),\qquad
 \frac{f(\ell_c/t_3)}{t_3}
      \ge\frac{f(B_s(x)/x_3)}{x_3}.
\]
These use only nonnegativity and monotonicity of $f$. Equal labels
invalidate none of the weak inequalities.
Also $E=N^{t_3}\le a_3$, so $V_N(E)\ge V_N(a_3)$.
At zero shifted contribution the required comparison follows from
nonnegativity without any assertion that $B_s(x)>0$.

Keep the original orbit normalization. If the equal-label groups are
$G_h$, put $M_c=\prod_h|G_h|!\in\{1,2,6\}$.
For a full labelled realization $\mathbf a$, write
$\mathbf a^\downarrow$ for its decreasing rearrangement within the
equal-label groups. The shifted comparison above is evaluated at
$\mathbf x^\downarrow$ and is invariant under these group permutations.
For every nonnegative, invariant main-mass weight $w_c$, therefore,
\[
 \frac1{M_c}\sum_{\rm full}K_P(\mathbf a)w_c(\mathbf a)
       V_N(E)f(\ell_c/t_3)
 \ge
 \frac1{M_c}\sum_{\rm full}K_P(\mathbf a)w_c(\mathbf a)
       V_N(a_3^\downarrow)
       f(B_s(\mathbf x^\downarrow)/x_3^\downarrow).
\]
Each distinct ordered triple has exactly $M_c$ labelled realizations.
The prefix mask $K_P$ and its correction are kept exactly as in
Paper~II, (4.22)--(4.23). No repeated-prefix realization is reclassified.
The tail mask and all remainder coefficients are likewise unchanged.
This proves the finite lower-main-term comparison with the cell mass
divided by $M_c$. The corresponding continuum comparison uses the
preceding inequality with the $1/t_3$ and $1/x_3$ factors.
Equal-label bins are kept. Their strictly ordered actual
triples lie in the original cone. The orbit identity converts
the full-block sum to that ordered sum exactly.
Only actual coordinate diagonals have zero continuum measure.
The proof imposes no $N\to\infty$ disappearance of equal-bin regions.
\end{proof}

Put $H_{\log}=\log(\mu_*/u)$.
Paper~II's remaining terms satisfy
\begin{equation}\label{eq:retB-G-residual}
\begin{aligned}
 0\le R_G\le R_{\mathrm{up}}
 ={}&500(F(30)-1)+500H_{\log}(1-f(30))\\
    &+250H_{\log}^2(F(30)-1).
\end{aligned}
\end{equation}
Its constant terms cancel exactly
\cite[(4.73)--(4.75)]{PaperII}:
\begin{equation}\label{eq:retB-G-cancellation}
 500\left(1-H_{\log}+\frac{H_{\log}^2}{2}\right)
 -\iiint_{\mathcal C}\frac{d\mathbf t}{t_1t_2t_3^2}
 =\frac1\mu_*.
\end{equation}
The shifted upper comparison is therefore
\begin{equation}\label{eq:retB-G-profile}
 G_1^s=\frac1\mu_*+R_G+J_s,
 \qquad
 \Phi_{G_1}^s=-K_GG_1^s,
 \qquad
 K_G=\frac{2I_{11}+I_{12}+I_{13}}{2e^\gamma}.
\end{equation}
The source-mass integrals are \cite[(4.77)]{PaperII}
\begin{equation}\label{eq:retB-G-masses}
\begin{aligned}
 I_{11}={}&\int_\lambda^{1/3}
          \frac{\log((1-2t)/t)}{t(1-t)}\,dt\\
 &+\int_\lambda^{1/4}\frac{dt}{t}
        \int_t^{(1-t)/3}
          \frac{\log((1-t-2v)/v)}{v(1-t-v)}\,dv,\\
 I_{12}={}&\int_\alpha^{1/3}
          \frac{\log(2-3t)}{t(1-t)}\,dt,\\
 I_{13}={}&\int_\beta^\lambda
          \frac{\log((1-t-\lambda)/\lambda)}{t(1-t)}\,dt.
\end{aligned}
\end{equation}
We have $I_{11},I_{12},I_{13}\ge0$ and $I_{12}>0$.
Thus $K_G>0$.
The finite shifted enclosure gives
\begin{equation}\label{eq:retB-G-finite}
 0\le G_1^s\le\frac{6073714261}{10^9}
 \quad\text{(shifted mesh $n_{\mathrm{mesh}}=512$)}.
\end{equation}
Paper~II's mass enclosure gives $0\le K_G\le K_{\mathrm{hi}}$
\cite[(A.10d)]{PaperII}. The exact constants are
\begin{equation}\label{eq:retB-exact-Khi}
 K_{\mathrm{hi}}=
 \frac{3989480791070578466390443009830933753748442947478110112223}{10^{58}}.
\end{equation}
Paper~II, (4.84), gives
\begin{equation}\label{eq:retB-exact-LG1}
 L_{G1}=-
 \frac{243950428265194656254206921939947762868871361874409552095097}
 {100433627766186892221372630771322662657637687111424552206336}.
\end{equation}
Define
\begin{equation}\label{eq:retB-M5}
 M_{5,\mathrm{internal}}
   =-K_{\mathrm{hi}}\frac{6073714261}{10^9}-L_{G1}.
\end{equation}
Their exact substitution gives
\begin{equation}\label{eq:retB-M5-value}
\begin{gathered}
 0.0058749499450263575220<M_{5,\mathrm{internal}},\\
 M_{5,\mathrm{internal}}<0.0058749499450263575221.
\end{gathered}
\end{equation}
The finite enclosure and identity are proved.
They concern the shifted profile integral.
The required actual-source inequality is
\begin{equation}\label{eq:retB-G-use}
 -\frac{2S_{11}+S_{12}+S_{13}}{W}
 \ge L_{G1}+M_{5,\mathrm{internal}}
       -E^D_{G1}-|\Lambda_{G,N}|.
\end{equation}

\begin{proof}[Proof of \eqref{eq:retB-G-use}]
Paper~II proves
\[
 \frac{2S_{11}+S_{12}+S_{13}}W\le K_GG_1+E^D_{G1},\qquad
 G_1=\frac1{\mu_*}+R_G+J,
\]
from (4.73)--(4.77) \cite[(4.75), (4.78)]{PaperII}. Run the same argument
with $J$ replaced by $J_s$. The first three residual rows and the bound
$R_G\le R_{\mathrm{up}}$ do not change, and neither does the cancellation
\eqref{eq:retB-G-cancellation}. In the fourth row,
Lemma~\ref{lem:retB-shifted-retention} replaces Paper~II's comparison after
(4.72). The canonical restriction loses zero, and the actual lower main term
dominates the shifted one. The source and its weights $(2,1,1)$ do not
change, and neither do the peel (4.45)--(4.57) and the mass conversion
(4.76)--(4.78). Hence
\[
 \frac{2S_{11}+S_{12}+S_{13}}W\le K_GG_1^s+E^D_{G1}
\]
with the same $E^D_{G1}$. By Proposition~\ref{prop:app-m5-cap},
$0\le G_1^s\le6073714261/10^9$, which is \eqref{eq:retB-G-finite}, with the
residual bound supplied by \eqref{eq:retB-G-residual}. By
Paper~II, (A.10d), $0\le K_G\le K_{\mathrm{hi}}$. Both factors are
nonnegative, so by \eqref{eq:retB-M5}
\[
 -K_GG_1^s\ge-K_{\mathrm{hi}}\frac{6073714261}{10^9}=L_{G1}+M_{5,\mathrm{internal}}.
\]
This gives \eqref{eq:retB-G-use} even without $|\Lambda_{G,N}|$. That term
is nonnegative and is kept so that the error agrees with
\eqref{eq:intro-retained-error}.
\end{proof}
The existing $E^D_{G1}$ is a disjoint component of the original
$F_{\mathrm{net}}+E_{\mathrm{van}}$ allowance.
Its fixed peel $145664566.4567\tau$ is covered by $g_\tau$,
and its fixed $C_E\sigma$ share by $g_{\mathrm{AA}}$.
Its vanishing part belongs to $E_{\mathrm{van}}$.
When replacing the joint chain, carry these same components once.
The reconciliation in \eqref{eq:retB-fixed-cost} remains paid.
These components are the $G_1$ comparator term $NC_{\mathrm{sf}}$ of
Paper~II, (4.64)--(4.65), and the $G_1$ share of the joint bound. That
bound treats $\Xi_U$ and $\Xi_A$ in Paper~II, (B.1)--(B.7). Both are
attached to the original source measure and not to the profile. Hence the
shift neither creates nor removes them.

\subsection{Conversion of the finite arrays}

We give the remaining definitions before stating the lemma.
All representation variables in the following table are prime.
Their product is $e$, with $\varepsilon N<e<N$.
Let $R_{11,3}(e)$ and $R_{11,4}(e)$ count the first two forms.
Let $R_{12}(e)$ and $R_{13}(e)$ count the last two forms.
Put $R_{11}=R_{11,3}+R_{11,4}$.
These are Paper~II's representation weights
\cite[(4.4), source table preceding (4.4)]{PaperII}.

\begin{table}[ht]
\caption{The internal representation counts}\label{tab:retB-representations}
\small
\renewcommand{\arraystretch}{1.15}
\begin{tabular}{@{}llp{.64\linewidth}@{}}
\hline
Count & Product & Restrictions\\
\hline
$R_{11,3}$ & $p_1p_2p_3$ &
$N^\lambda\le p_1<N^{1/3}$, $p_1<p_2$, $p_1p_2^2<N$, $p_3\ge p_2$.\\
$R_{11,4}$ & $p_1p_2p_3p_4$ &
The same conditions on $p_1,p_2$; $p_2\le p_3\le p_4$.\\
$R_{12}$ & $p_1p_2p_3$ &
$N^\alpha\le p_1<N^{1/3}\le p_2$, $p_1p_2^2<N$, $p_3\ge p_2$.\\
$R_{13}$ & $p_1p_2p_3$ &
$N^\beta\le p_1<N^\lambda\le p_2$, $p_1p_2^2<N$, $p_1p_2p_3^2\ge N$.\\
\hline
\end{tabular}
\end{table}
For the last form the endpoint orders are distinct.
Every permitted representation keeps its multiplicity.
Define the upper cut and masses by
\begin{equation}\label{eq:retB-source-normalizations}
\begin{aligned}
 N_{\mathrm{cut}}&=N(1-L^{-333}),\\
 X_j^{\mathrm{cut}}&=\sum_eR_j(e)
           \mathbf1_{\varepsilon N<e\le N_{\mathrm{cut}}}
                &&(j=11,12,13),\\
 X_G&=2X_{11}^{\mathrm{cut}}+X_{12}^{\mathrm{cut}}+X_{13}^{\mathrm{cut}},\\
 r_{G,N}(t)&=\frac{tX_GV_N(N^t)}{WK_G}.
\end{aligned}
\end{equation}
The cut precedes the sieve decomposition
\cite[(4.16)--(4.17)]{PaperII}.
The source-mass input used below is
\begin{equation}\label{eq:retB-source-mass-limits}
\begin{aligned}
 X_D&=\left(1-\varepsilon+o(1)\right)\frac NL,\\
 X_G&=\left((1-\varepsilon)(2I_{11}+I_{12}+I_{13})+o(1)\right)
                                                   \frac NL.
\end{aligned}
\end{equation}

\begin{lemma}[Source masses]\label{lem:retB-source-mass}
The formulas \eqref{eq:retB-source-mass-limits} hold.
\end{lemma}
\begin{proof}
The sum $X_D$ counts the primes in $(N/L^{30},(1-\varepsilon)N)$, so the
prime number theorem gives $X_D=(1-\varepsilon+o(1))N/L$.
For the representation weights of Table~\ref{tab:retB-representations},
Paper~II proves $X_j=((1-\varepsilon)I_j+o(1))N/L$ for the uncut masses,
$j=11,12,13$ \cite[(4.76)--(4.77)]{PaperII}. The cut strip
$N_{\mathrm{cut}}<e<N$ carries joint mass at most $100(NL^{-333}+1)$
\cite[(4.16)--(4.17)]{PaperII}, which is $o(N/L)$. Summing with the
weights $(2,1,1)$ gives the formula for $X_G$.
\end{proof}

On $X_0^3$, define the zero-extended kernel
\begin{equation}\label{eq:retB-G-gain-kernel}
 d_G(\mathbf t)=\mathbf1_{\mathcal C}(\mathbf t)
       K_G\frac{f(a_3^s(\mathbf t))-f(a_3(\mathbf t))}{t_3}.
\end{equation}
Its arrays and integral are
\begin{equation}\label{eq:retB-G-array}
\begin{aligned}
 B_{G,N}&=\iiint r_{G,N}(t_3)d_G(\mathbf t)\,d(m_N^3)(\mathbf t),\\
 \Delta_G&=\iiint d_G\,d(m^3)=K_G(J-J_s),\\
 \Lambda_{G,N}&=B_{G,N}-\Delta_G.
\end{aligned}
\end{equation}
Equivalently,
\begin{equation}\label{eq:retB-G-prime-sum}
\begin{aligned}
 B_{G,N}&=\frac{X_G}{W}
 \sum_{\substack{p_3<p_2<p_1,\ p_i\in\mathbb P,\ p_i\nmid N\\
           u\le\log_Np_3<\log_Np_2<\log_Np_1<\mu_*}}
 \frac{V_N(p_3)[f(a_3^s(\mathbf t))-f(a_3(\mathbf t))]}
      {(p_1-1)(p_2-1)(p_3-1)},\\
 &\hspace{38mm}t_i=\log_Np_i.
\end{aligned}
\end{equation}
These definitions specify all five additional conversion errors.

\begin{lemma}[Conversion of the five profile differences]
\label{lem:retB-conversion}
Fix $0<\sigma\le10^{-10}$. For $m_N,m$ in \eqref{eq:retB-measure}, the finite
arrays \eqref{eq:retB-direct-arrays} and \eqref{eq:retB-G-array} satisfy
\begin{equation}\label{eq:retB-conversion-limit}
 \Lambda_{j,N}\longrightarrow0\quad(j\in\mathcal H),\qquad
 \Lambda_{G,N}\longrightarrow0,
\end{equation}
and $\Lambda^0_{j,N}\to0$ for $j\in\mathcal H$.
Every limit ranges over even $N$ tending to infinity; the thresholds depend
on $\sigma$.
\end{lemma}

\begin{proof}
\emph{Step 1: convergence of the harmonic measures.}
Let $u\le a<b\le1$.
Mertens' theorem and partial summation give
\begin{equation}\label{eq:retB-weak}
 m_N([a,b))=\log\frac ba+o(1).
\end{equation}
Replacing $1/(p-1)$ by $1/p$ costs $O(N^{-u})$.
There are at most $L/\log2$ prime divisors of $N$.
Each omitted divisor above $N^u$ has mass $O(N^{-u})$.
Their total cost is $O(N^{-u}L)$.
The largest endpoint atom is $O(N^{-u})$.
The total masses therefore satisfy
\begin{equation}\label{eq:retB-total-mass}
 c_N:=m_N(X_0)\longrightarrow c_0:=m(X_0)=\log(1/u)>0.
\end{equation}
Let $g$ be continuous on $X_0$. Choose a finite interval partition and a
step function on it, uniformly within a prescribed error of $g$.
The integrals of the step function converge by \eqref{eq:retB-weak}. The
integral errors are at most the uniform error times $c_N+c_0$, which is
bounded by \eqref{eq:retB-total-mass}. Hence $m_N\Rightarrow m$
\cite[Theorem~2.2 and Example~2.3]{Billingsley1999}.
For $k=2,3$, rectangle masses are products of interval masses. The same
argument with rectangular step functions gives
\begin{equation}\label{eq:retB-product-weak}
 m_N^k\Rightarrow m^k\qquad(k=1,2,3).
\end{equation}
For a coordinate diagonal $D_{ij}=\{t_i=t_j\}$,
\begin{equation}\label{eq:retB-diagonal-mass}
 m_N^k(D_{ij})\le
  \left(\max_{p\ge N^u}\frac1{p-1}\right)c_N^{k-1}
 =O(N^{-u})\qquad(k\ge2).
\end{equation}
Thus strict prime order costs nothing in the limit, and no orbit factor
enters.

\emph{Step 2: the source and Euler normalizations.}
For $p>2$,
\begin{equation}\label{eq:retB-Euler-factorization}
 1-\frac1{p-1}=
 \left(1-\frac1p\right)\left(1-\frac1{(p-1)^2}\right).
\end{equation}
Mertens' product formula and the convergent second product yield
\begin{equation}\label{eq:retB-Euler-limit}
 V_N(N^t)=\frac{2e^{-\gamma}C(N)}{tL}(1+o(1)),
 \qquad u/\rho\le t\le1,
\end{equation}
uniformly in $t$.
The prime $2$ is omitted because $N$ is even.
The missing divisors $p\mid N$ above $N^t$ contribute
$\prod(1+1/(p-2))=1+O(N^{-u/\rho}L)$. The ordinary product error is
uniform for $N^t\ge N^{u/\rho}$.
By Lemma~\ref{lem:retB-source-mass} and \eqref{eq:not-normalization},
\begin{equation}\label{eq:retB-density-limits}
 r_{D,N}\longrightarrow1,\qquad
 \sup_{u\le t\le\mu_*}|r_{G,N}(t)-1|\longrightarrow0.
\end{equation}
For example, $WK_G=2e^{-\gamma}(1-\varepsilon)
(2I_{11}+I_{12}+I_{13})C(N)N/L^2$.
This is exactly the main factor in $tX_GV_N(N^t)$.
Thus both the truncation factor and the representation multiplicities
cancel once.

\emph{Step 3: restriction and bounded kernels.}
Write $\widehat m_N=m_N/c_N$ and $\widehat m=m/c_0$.
These are probability measures, and their products converge weakly.
Let $D$ have $m^k(\partial D)=0$, and let $g\mathbf1_D$ be bounded and
$m^k$-almost everywhere continuous. Apply the mapping form of the portmanteau
theorem \cite[Theorems~2.1 and~2.7]{Billingsley1999} to
$g\mathbf1_D$ and the probability measures. Then multiply by
$c_N^k\to c_0^k$. This gives
\begin{equation}\label{eq:retB-portmanteau}
 \int_Dg\,d(m_N^k)\longrightarrow\int_Dg\,d(m^k).
\end{equation}
The intervals $I_j$ have finitely many boundary points.
The boundary of $\mathcal D_{10}$ is piecewise affine, since $B_0$ is
continuous and piecewise affine. Both have zero $m^{e_j}$-measure.
On $I_j$ $(j=4,5,7)$ the argument $s(x)$ lies in a compact subinterval of
$(1,\infty)$, where $F$ is continuous. On $\mathcal D_{10}$,
$y<K-11\alpha/5-5x/2$ gives $q(x,y)>11/5$, and $q+h/\alpha<4$; there $f$ is
continuous. Hence $k_j^0,k_j^h$ and $d_j$ are bounded and continuous off a
null set, so \eqref{eq:retB-portmanteau} applies to them.
With \eqref{eq:retB-density-limits}, it gives
\begin{equation}\label{eq:retB-direct-convergence}
 S_{j,N}^0\longrightarrow P_j,\qquad
 B_{j,N}\longrightarrow\Delta_j\qquad(j\in\mathcal H).
\end{equation}

\emph{Step 4: the ordered three-dimensional kernel.}
The cone $\mathcal C$ has affine boundary faces, so its boundary is
$m^3$-null.
The function $\Theta_*$ is piecewise affine.
Its possible jumps lie on finitely many guard hyperplanes, and its branch
switches add finitely many affine faces.
On every remaining open piece, $a_3$ and $a_3^s$ are continuous.
The denominator satisfies $t_3\ge u>0$, so the arguments of $f$ remain in a
fixed compact interval. On it $f$, extended by zero below $2$, is
continuous. Since $0\le f\le1$, $|d_G|\le K_G/u$.
Thus $d_G$ is bounded and $m^3$-almost everywhere continuous.
Apply \eqref{eq:retB-portmanteau} to the zero extension on $X_0^3$:
\begin{equation}\label{eq:retB-G-portmanteau}
 \iiint d_G\,d(m_N^3)\longrightarrow
 \iiint d_G\,d(m^3)=\Delta_G.
\end{equation}
The normalization costs at most
\begin{equation}\label{eq:retB-G-normalization-error}
 \left|B_{G,N}-\iiint d_G\,d(m_N^3)\right|
 \le\sup_{[u,\mu_*]}|r_{G,N}-1|\,\|d_G\|_\infty c_N^3
 \longrightarrow0.
\end{equation}
The cone's coordinate-order inequalities are strict, so $B_{G,N}$ counts
strictly ordered triples, as in \eqref{eq:retB-G-prime-sum}.
Hence $B_{G,N}\to\Delta_G$, and with Step~3 all limits in
\eqref{eq:retB-conversion-limit} and $\Lambda^0_{j,N}\to0$ follow.
\end{proof}

\paragraph{Whole cells, retained cells and the label cutoff.}
For the four direct rows, the arithmetic identification is
Lemma~\ref{lem:retB-lifted-direct}. For the joint row it is the proof of
\eqref{eq:retB-G-use}. Both are made at stage eight on Paper~II's actual
cells, with their inherited signed coefficients. Hence no further whole-cell
kernel is needed. For the retained part of $S_{10}$, the actual whole cells
cover $\mathcal D_{10}$ up to a strip. Lemma~\ref{lem:retB-lifted-direct}
absorbs this strip by the retained reserve \eqref{eq:retB-S10-surplus}. For
$G_1$, the retained cells are those of Paper~II, (4.19). By
Lemma~\ref{lem:retB-shifted-retention}, the shifted profile has
no positive support outside them. The arrays \eqref{eq:retB-direct-arrays}
and \eqref{eq:retB-G-array} enter only through the nonnegative errors in
$E'$, which vanish by Lemma~\ref{lem:retB-conversion}.

The array \eqref{eq:retB-G-prime-sum} carries $V_N(p_3)$ at the actual
coordinate. The arithmetic lower main term carries $V_N(E)$,
$E=N^{t_3}$, at the inherited label. Since $E\le a_3$, $V_N(E)\ge V_N(a_3)$.
In the lower part of the joint upper bound this replacement is
favourable. We nevertheless bound it absolutely and include it once in
Paper~II's peel allowance (4.57). No fixed cost is added.
By \eqref{eq:retB-Euler-limit}, uniformly for actual label/coordinate pairs in the same cell with $u/\rho\le t_3\le x_3\le1$ and $x_3<\rho t_3$,
\[
 \frac{V_N(N^{t_3})}{V_N(N^{x_3})}=\frac{x_3}{t_3}(1+o(1))\le\rho(1+o(1)).
\]
For a weight $0\le w\le1$ on the strictly ordered cone, the replacement
therefore costs at most
\begin{equation}\label{eq:retB-label-cost}
\begin{gathered}
 C_{\mathrm{cut},N}\le(\tau+o_\sigma(1))T_N,\qquad
 T_N=\frac{X_G}W\sum_{\mathcal C}
       \frac{V_N(p_3)}{(p_1-1)(p_2-1)(p_3-1)}\longrightarrow K_GI_0,\\
 I_0=\iiint_{\mathcal C}\frac{d\mathbf t}{t_1t_2t_3^2}
    =500\Bigl(1-H_{\log}+\frac{H_{\log}^2}2\Bigr)-\frac1{\mu_*},
\end{gathered}
\end{equation}
where the sum runs over the primes of \eqref{eq:retB-G-prime-sum}. The limit
follows from Steps~2--4 of the proof of Lemma~\ref{lem:retB-conversion}
applied to the bounded kernel $\mathbf 1_{\mathcal C}K_G/t_3$. We evaluate $I_0$
by \eqref{eq:retB-G-cancellation}. Since $1<H_{\log}<459/100$ and
the quadratic increases for $H_{\log}>1$,
\[
 I_0<\frac{693381}{200},\qquad
 K_GI_0\le K_{\mathrm{hi}}\frac{693381}{200}<1400.
\]
At fixed $\tau$, eventually $o_\sigma(1)\le\tau$ and $T_N\le1400$, so
$C_{\mathrm{cut},N}\le2800\tau<4000\tau$. We include $4000\tau$ once.
Paper~II's peel bound before its final rounding is
\[
 P_{\mathrm{peel}}\tau=\frac{3\cdot502002\,(459/100)^3}{1-\varepsilon}\tau
 =\frac{8090814677193}{55550}\tau
\]
\cite[(4.47), (4.56)]{PaperII}, and its printed allowance is
$P_{\mathrm{cap}}\tau=(485500000000/3333)\tau$ \cite[(4.57)]{PaperII}.
The per-layer coefficient 502002 in Paper~II, (4.53), includes the
Selberg remainder. The printed allowance (4.57) uses that coefficient;
we therefore use its slack with the Selberg remainder already included.
Exact subtraction gives
\begin{equation}\label{eq:retB-peel-room}
 P_{\mathrm{cap}}-P_{\mathrm{peel}}=\frac{2555968421}{166650},\qquad
 P_{\mathrm{cap}}-P_{\mathrm{peel}}-4000=\frac{1889368421}{166650}>0.
\end{equation}
Hence the peel and the label replacement together stay below
$P_{\mathrm{cap}}\tau<145664566.4567\tau$, the entry already contained in
$g_\tau$. The fixed deduction $F_{\mathrm{net}}$ is unchanged.

\subsection{The resulting lower comparison}

\begin{remark}[Inputs]\label{rem:retB-inputs}
The proof of Proposition~\ref{prop:retB-comparison} uses the cited
statements of Paper~I and Paper~II. They are applied to Paper~II's unchanged sources.
It uses Paper~I, Theorem~2.8, through
Fact~\ref{fact:spectral-implication}, and the prime Siegel--Walfisz theorem,
Paper~I, (2.21). It uses Lemma~\ref{hyp:ra-ms-transfer}(ii) at
Paper~I, (3.48), and
the finite enclosures of Propositions~\ref{prop:app-h6-four-profiles}
and~\ref{prop:app-m5-cap}. The optional endpoint improvement also uses
Lemma~\ref{lem:retB-six-endpoints}. The table of
Appendix~\ref{subsec:app-retB-rows} is an index, not an input.
\end{remark}

\begin{proposition}[Retained comparison]\label{prop:retB-comparison}
Fix $0<\sigma\le10^{-10}$ satisfying Paper~II, (1.27) and assume
Lemma~\ref{lem:retB-six-endpoints}.
Then \eqref{eq:retB-goal} holds, and $E'(N;\sigma)\to0$ as $N\to\infty$.
Without that lemma, \eqref{eq:retB-goal} holds with $g_6$ replaced by $0$;
the same limit for $E'(N;\sigma)$ holds.
\end{proposition}

\begin{proof}
Write each chain of $\mathcal I$, including $\mathrm{T5}$, by
\eqref{eq:retB-telescope} and each joint group by \eqref{eq:retB-joint}. Lemma~\ref{lem:retB-base-errors} covers the
differences at stages one to eight and the remaining Paper~II differences at
stages nine and ten. Their sum is at least $-(F_{\mathrm{net}}-\delta_{G1})-E_{\mathrm{van}}$.
For $j\in\mathcal H$, replace the stage-nine object by $Q^h_{j,9}=P_j+\Delta_j$.
Then Lemma~\ref{lem:retB-lifted-direct} gives $Q_{j,8}-Q^h_{j,9}\ge0$.
For $j\in\mathcal J$, Lemma~\ref{lem:retB-six-endpoints} gives the
stage-ten difference $P_j-N_j\ge0$, so the endpoint rises by $N_j-O_j$.
Only the endpoint changes, so $g_6$ enters once.
By Proposition~\ref{prop:app-h6-four-profiles},
$\sum_{j\in\mathcal H}\Delta_j>572004/10^8>5719/10^6$, which is
\eqref{eq:retB-H6-finite}.
For the joint occurrence, \eqref{eq:retB-G-use} replaces Paper~II's $G_1$
bound. Its $E^D_{G1}$ is the $G_1$ part of $F_{\mathrm{net}}+E_{\mathrm{van}}$
and is counted once.
These modifications concern disjoint profile terms, so their gains enter once.
The endpoint sum \eqref{eq:retB-base-endpoints} and the floor
\eqref{eq:retB-paperII-floor} give $D_{\mathrm{II}}+\overline J_6+2\overline H$.
Hence
\begin{equation}\label{eq:retB-sum}
\begin{aligned}
 \frac{R_{\mathrm{fix}}+S_2}{W}
 \ge{}&D_{\mathrm{II}}+g_6+\overline J_6+2\overline H
      +\frac{5719}{10^6}+M_{5,\mathrm{internal}}\\
 &-E_{\mathrm{van}}(N;\sigma)
      -\sum_{j\in\mathcal H}
          \bigl(|\Lambda^0_{j,N}|+|\Lambda_{j,N}|\bigr)-|\Lambda_{G,N}|.
\end{aligned}
\end{equation}
The nine conversion errors are nonnegative, so \eqref{eq:retB-sum} is weaker
than the same bound without them. This is \eqref{eq:retB-goal}.
By Lemmas~\ref{lem:retB-base-errors} and~\ref{lem:retB-conversion},
$E'(N;\sigma)\to0$; the threshold depends on the fixed $\sigma$.
No fixed increment is added. The only new fixed cost is the label
replacement, and by \eqref{eq:retB-peel-room} it lies inside Paper~II, (4.57).
Without Lemma~\ref{lem:retB-six-endpoints}, take $N_j=O_j$; then the
same argument gives \eqref{eq:retB-sum} with $g_6$ replaced by $0$.
\end{proof}

\section{\texorpdfstring{The sixth sum at \(N/4\)}{The sixth sum at N/4}}\label{sec:root-a}

Put \(L=\log N\) and fix \(0<\sigma\le10^{-10}\) for the common
error allocation. No estimate in the sixth-sum construction depends on \(\sigma\). Fix
\begin{equation}\label{eq:ra-parameters}
\begin{gathered}
 \alpha=\frac{100}{1149},\qquad t_0=\frac{57}{224},\qquad
 v_{\mathrm{RA}}=\frac{5227}{10000},\qquad \eta_{\mathrm m}=\varepsilon=\frac1{10000},\\
 J=[t_0,1/3),\quad J_t=[t_0,1/3-\eta_{\mathrm m}),\quad
 J_{\mathrm{tail}}=[1/3-\eta_{\mathrm m},1/3),\\
 I_L=[2,N/4],\qquad I_H=(N/4,(1-\varepsilon)N),\qquad z_s=N^\alpha.
\end{gathered}
\end{equation}
Write \(p\in N^K\) for \(p\in\mathbb P\) with \(\log p/L\in K\).
Paper II, (1.7), gives the original sum \cite[(1.7)]{PaperII}:
\begin{equation}\label{eq:ra-original-s6}
 S_6=\sum_{\substack{p\in N^J\\(p,N)=1}}
 \ \sum_{\substack{m\ge1,\ e=N-pm\in\mathbb P\\2\le e<(1-\varepsilon)N}}
 \mathbf1_{P^-(m)\ge z_s}\mathbf1_{(m,Np)=1}.
\end{equation}
Each prime mark remains a separate occurrence.

\subsection{Three upper weights}

For squarefree \(d=\ell_1\cdots\ell_k\), order its primes decreasingly.
The ordinary upper Rosser weight is
\begin{equation}\label{eq:ra-rosser-weight}
 \mu_d^+(D)=\mu(d)
 \prod_{\substack{1\le j\le k\\j\ \mathrm{odd}}}
 \mathbf1_{\ell_1\cdots\ell_{j-1}\ell_j^3<D},
 \qquad \mu_1^+(D)=1.
\end{equation}
The weight is zero outside its stated squarefree support.
Its divisor sum is a nonnegative upper sieve weight
\cite[Lemma 1, (3.2), (3.6)]{Iwaniec1980Rosser}
\cite[Lemmas 2--3, p.~313]{Iwaniec1980Remainder}.

For \(K\subseteq J\) and \(p\in N^K\), define
\begin{equation}\label{eq:ra-coefficient}
 c_{R,K}(pd)=\mu(d)\mathbf1_{pd<R}
 \prod_{\substack{1\le j\le k\\j\ \mathrm{odd}}}
 \mathbf1_{p\ell_1\cdots\ell_{j-1}\ell_j^3<R}.
\end{equation}
Require \(\ell_j<z_s\) and \((pd,2N)=1\).
Set the coefficient to zero otherwise.
Since \(t_0>\alpha\), the mark is unique.
Thus \(|c_{R,K}(q)|\le1\).
On supported moduli with \(p<R\), the gates give \(d<R/p\).
For \(d=1\) the gate product is empty, and \(p<R\) suffices.
Every marked prime has \(p<N^{1/3}\).
Each sieve level below is \(>N^{1/3}\) for large \(N\).
Hence
\begin{equation}\label{eq:ra-weight-equality}
 c_{R,K}(pd)=\mu_d^+(R/p),\qquad
 \varphi(pd)=(p-1)\varphi(d).
\end{equation}
The equality is claimed on supported moduli only.
Off that support the coefficient is zero.

Choose fixed logarithmic exponents. Put
\begin{equation}\label{eq:ra-factories}
\begin{aligned}
 R_L&=N^{1/2}L^{-B_L},&
 R_{\mathrm{tail}}&=N^{1/2}L^{-B_{\mathrm{tail}}},\\
 Q&=N^{v_{\mathrm{RA}}}L^{-B},&Q_0&=N^{1/2}L^{-B_0},\\
 b_{\mathrm{all}}&=c_{R_L,J},&
 b_{\mathrm{tail}}&=c_{R_{\mathrm{tail}},J_{\mathrm{tail}}},\\
 c_{\mathrm{trim}}&=c_{Q,J_t},&
 c_{\mathrm{lo}}(q)&=c_{\mathrm{trim}}(q)\mathbf1_{q\le Q_0},\\
 \Gamma(q)&=c_{\mathrm{trim}}(q)\mathbf1_{q>Q_0}.
\end{aligned}
\end{equation}
The last two arrays keep the gates at \(Q\).

First omit \((m,Np)=1\) in \eqref{eq:ra-original-s6}.
Sieve the larger sequence by primes \(\ell\nmid2N\), \(\ell<z_s\).
Use the three weights on
\begin{equation}\label{eq:ra-three-pieces}
 I_L\times J,\qquad I_H\times J_{\mathrm{tail}},\qquad I_H\times J_t.
\end{equation}
Their divisor sum is the finite weighted count
\begin{equation}\label{eq:ra-mix-definition}
\begin{aligned}
 A_{\mathrm{mix}}={}&
 \sum_q b_{\mathrm{all}}(q)
       \sum_{\substack{e\in I_L\cap\mathbb P\\e\equiv N\ (q)}}1\\
 &+\sum_q\{b_{\mathrm{tail}}(q)+c_{\mathrm{trim}}(q)\}
       \sum_{\substack{e\in I_H\cap\mathbb P\\e\equiv N\ (q)}}1.
\end{aligned}
\end{equation}
The upper-weight comparison
\begin{equation}\label{eq:ra-positive-comparison}
 S_6\le A_{\mathrm{mix}},\qquad
 Z_6:=A_{\mathrm{mix}}-S_6\ge0,
\end{equation}
is proved in Lemma~\ref{lem:ra-upper-sieve}(a) below.
Positivity is used before expanding the divisor sums.

\subsection{The exact discrepancy identity}

For a bounded interval \(I\subset[2,\infty)\) and a finitely supported
modulus array \(a(q)\), put \(\operatorname{Li}(I)=\int_I dt/\log t\)
and define
\begin{equation}\label{eq:ra-discrepancies}
\begin{aligned}
 \Pi_I(q)&=\sum_{e\in I\cap\mathbb P}\mathbf1_{e\equiv N\pmod q},\\
 B_I(q)&=\frac1{\varphi(q)}
             \sum_{e\in I\cap\mathbb P}\mathbf1_{(e,q)=1},\\
 D_I(a)&=\sum_q a(q)
       \left\{\Pi_I(q)-\frac{\operatorname{Li}(I)}{\varphi(q)}\right\},\\
 P_I(a)&=\sum_q a(q)\{\Pi_I(q)-B_I(q)\},\qquad
 C_I(a)=D_I(a)-P_I(a).
\end{aligned}
\end{equation}
Also put
\begin{equation}\label{eq:ra-main-mass}
\begin{aligned}
 M_{\mathrm{mix}}={}&\operatorname{Li}(I_L)
              \sum_q\frac{b_{\mathrm{all}}(q)}{\varphi(q)}\\
 &+\operatorname{Li}(I_H)
              \sum_q\frac{b_{\mathrm{tail}}(q)+c_{\mathrm{trim}}(q)}{\varphi(q)}.
\end{aligned}
\end{equation}
Finite expansion gives
\begin{equation}\label{eq:ra-b9}
\begin{aligned}
 A_{\mathrm{mix}}={}&M_{\mathrm{mix}}+D_{I_L}(b_{\mathrm{all}})
     +D_{I_H}(b_{\mathrm{tail}})\\
 &+C_{I_H}(c_{\mathrm{trim}})+P_{I_H}(c_{\mathrm{lo}})
     +P_{I_H}(\Gamma).
\end{aligned}
\end{equation}
Since \(c_{\mathrm{trim}}=c_{\mathrm{lo}}+\Gamma\),
\begin{equation}\label{eq:ra-cancellation}
 C_{I_H}(c_{\mathrm{trim}})+P_{I_H}(c_{\mathrm{lo}})
       =D_{I_H}(c_{\mathrm{lo}})+C_{I_H}(\Gamma).
\end{equation}
Therefore
\begin{equation}\label{eq:ra-b9-f1}
\begin{aligned}
 A_{\mathrm{mix}}={}&M_{\mathrm{mix}}+D_{I_L}(b_{\mathrm{all}})
     +D_{I_H}(b_{\mathrm{tail}})+D_{I_H}(c_{\mathrm{lo}})\\
 &+C_{I_H}(\Gamma)+P_{I_H}(\Gamma).
\end{aligned}
\end{equation}
No low natural-coprime term is estimated separately.

For large \(N\), each prime in \(I_H\) exceeds \(Q\).
Hence
\begin{equation}\label{eq:ra-high-conversion}
 C_{I_H}(\Gamma)=
 \{\pi(I_H)-\operatorname{Li}(I_H)\}
          \sum_q\frac{\Gamma(q)}{\varphi(q)}.
\end{equation}
The elementary identity
\begin{equation}\label{eq:ra-totient-mass}
 \frac q{\varphi(q)}=\sum_{d\mid q}\frac{\mu(d)^2}{\varphi(d)},
 \qquad
 \sum_{q\le Q}\frac1{\varphi(q)}\ll\log(2Q)
\end{equation}
follows by multiplicativity and a convergent Euler product.
Thus Paper I, (2.21), bounds \eqref{eq:ra-high-conversion}
by \(O_A(NL^{-A})\) for each fixed \(A\)
\cite[(2.21)]{PaperI}.

For the three \(D\)-terms we use the maximal-endpoint prime form of
Bombieri--Vinogradov.
Here it is conditional on the prime Siegel--Walfisz input
\cite[(2.21)]{PaperI}. Define the prime progression error by
\begin{equation}\label{eq:ra-prime-error}
 \Delta_\pi(y;q,a)=\pi(y;q,a)-\frac{\operatorname{Li}(y)}{\varphi(q)}.
\end{equation}
Fact~\ref{fact:prime-estimates-maximal} gives, for each fixed \(S>0\),
\begin{equation}\label{eq:ra-prime-bv}
 \sum_{q\le X^{1/2}(\log X)^{-(S+60)}}
 \max_{(a,q)=1}\sup_{0\le y\le X}
 |\Delta_\pi(y;q,a)|\ll_S X(\log X)^{-S}.
\end{equation}
Here \(\operatorname{Li}(y)=0\) for \(0\le y<2\).
The level sits at \(X\), not at the variable endpoint \(y\),
so intervals with endpoints chosen separately for each
modulus cost at most twice this bound. This is the maximal prime statement.
We do not substitute the fixed-endpoint von Mangoldt statement of
Paper I, (2.20). The source and the fixed \(60\) logarithmic buffer
are those of Fact~\ref{fact:prime-estimates-maximal}.

Set \(X=N\), \(S=A+10\), and choose the half-level buffers so that
\begin{equation}\label{eq:ra-bv-buffer}
 B_L,\ B_{\mathrm{tail}},\ B_0\ge A+72.
\end{equation}
\begin{lemma}[The three low discrepancies]\label{lem:ra-n4-bv}
Under \eqref{eq:ra-bv-buffer},
\begin{equation}\label{eq:ra-three-bv-errors}
 |D_{I_L}(b_{\mathrm{all}})|+
 |D_{I_H}(b_{\mathrm{tail}})|+
 |D_{I_H}(c_{\mathrm{lo}})|\ll_A NL^{-A-10}.
\end{equation}
\end{lemma}
\begin{proof}
The arrays \(b_{\mathrm{all}}\), \(b_{\mathrm{tail}}\) and \(c_{\mathrm{lo}}\)
are supported on \(q<R_L\), \(q<R_{\mathrm{tail}}\) and \(q\le Q_0\),
respectively. Each is one-bounded. Since
\(B_L,B_{\mathrm{tail}},B_0\ge A+72>S+60\), every such \(q\) has
\(q\le N^{1/2}L^{-(S+60)}\).
This is the level of \eqref{eq:ra-prime-bv} at
\(X=N\). Every supported \(q\) is coprime to \(N\), so \(N\bmod q\) is a
reduced class. Since \(\pi(y;q,N)=\operatorname{Li}(y)=0\) for \(y<2\),
\[
 \Pi_{I_L}(q)-\frac{\operatorname{Li}(I_L)}{\varphi(q)}
   =\Delta_\pi(N/4;q,N).
\]
Write \(y^-\) for the left limit at \(y\). Then
\[
 \Pi_{I_H}(q)-\frac{\operatorname{Li}(I_H)}{\varphi(q)}
   =\Delta_\pi\bigl(((1-\varepsilon)N)^-;q,N\bigr)
    -\Delta_\pi(N/4;q,N).
\]
A left limit is a limit of values at points of \([0,N]\).
Hence each value on the right is at most
\(\sup_{0\le y\le N}|\Delta_\pi(y;q,N)|\).
Therefore the left side of
\eqref{eq:ra-three-bv-errors} is at most five times the left side of
\eqref{eq:ra-prime-bv}, with \(X=N\) and \(S=A+10\). This is the maximal
prime statement of Fact~\ref{fact:prime-estimates-maximal}.
\end{proof}


\subsection{The high coefficient}

Put
\begin{equation}\label{eq:ra-rough-arrays}
\begin{gathered}
 U=V=N^{1/3},\qquad z_0=N^{1/(\log\log N)^3},\\
 \mathcal R(n)=\mathbf1_{P^-(n)>z_0},\qquad \mathcal R(1)=1,\qquad
 \mathcal A(a)=\mu(a)\mathcal R(a),\\
 \mathcal C(t)=\mathcal R(t)
       \sum_{\substack{ab=t\\a\le U,\ b\le V}}\mu(a)\Lambda(b),\qquad
 \mathcal B(s)=\mathcal R(s)\sum_{\substack{b\mid s\\b>V}}\Lambda(b),\\
 \mathcal K_q(e)=\mathbf1_{e\equiv N\pmod q}
                   -\frac{\mathbf1_{(e,q)=1}}{\varphi(q)}.
\end{gathered}
\end{equation}
The von Mangoldt identity gives three sums:
\begin{equation}\label{eq:ra-three-bilinear-sums}
\begin{aligned}
 T_1&=\sum_q\Gamma(q)
   \sum_{\substack{a\le U,\ c\ge1\\ac\in I_H}}
   \mathcal A(a)\mathcal R(c)\frac{\log c}{\log(ac)}\mathcal K_q(ac),\\
 T_2&=-\sum_q\Gamma(q)
   \sum_{\substack{t\le UV,\ c\ge1\\tc\in I_H}}
   \frac{\mathcal C(t)\mathcal R(c)}{\log(tc)}\mathcal K_q(tc),\\
 T_{\mathrm{II}}&=\sum_q\Gamma(q)
   \sum_{\substack{a>U,\ s>V\\as\in I_H}}
   \frac{\mathcal A(a)\mathcal B(s)}{\log(as)}\mathcal K_q(as).
\end{aligned}
\end{equation}
These come from
\begin{equation}\label{eq:ra-vaughan-identity}
 \Lambda=\mu_{\le U}*\log-\mu_{\le U}*\Lambda_{\le V}*1
             +\Lambda_{\le V}+\mu_{>U}*\Lambda_{>V}*1.
\end{equation}
Roughness is multiplicative, including shared prime factors.
The short \(\Lambda\)-term vanishes on \(I_H\).

The detector retains proper prime powers whose bases exceed \(z_0\).
For \(e\in I_H\), hence \(e>z_0\),
\begin{equation}\label{eq:ra-proper-power-detector}
 \frac{\mathcal R(e)\Lambda(e)}{\log e}
 =\mathbf1_{e\in\mathbb P}
  +\sum_{\substack{P\in\mathbb P,\ j\ge2\\e=P^j}}
                                      \frac{\mathbf1_{P>z_0}}j.
\end{equation}
Consequently
\begin{equation}\label{eq:ra-b10}
\begin{aligned}
 P_{I_H}(\Gamma)&=T_1+T_2+T_{\mathrm{II}}-E_{\mathrm{pp}},\\
 E_{\mathrm{pp}}&=\sum_q\Gamma(q)
   \sum_{\substack{P\in\mathbb P,\ j\ge2\\P^j\in I_H}}
           \frac{\mathbf1_{P>z_0}}j\mathcal K_q(P^j).
\end{aligned}
\end{equation}
Let \(e\in I_H\).
A supported \(q\) with \(e\equiv N\pmod q\) divides
\(N-e\), and \(|\Gamma(q)|\le1\) by unique marking.
The comparator part is at most \(\sum_{q\le Q}1/\varphi(q)\). Hence \eqref{eq:ra-totient-mass}
gives
\begin{equation}\label{eq:ra-proper-power-bound}
 \sum_q|\Gamma(q)||\mathcal K_q(e)|\le d(N-e)+O(L),
 \qquad E_{\mathrm{pp}}\ll_\nu N^{1/2+\nu}L^2.
\end{equation}
Here \(N-e\ge\varepsilon N\), and \(0<\nu<1/2\) is fixed. The second bound
follows from the first, \(d(N-e)\ll_\nu N^{\nu}\) and
\(\sum_{j\ge2}\#\{P\in\mathbb P:P^j\le N\}\ll N^{1/2}\).
The high-coefficient estimate uses the \(P>z_0\) mask of
\(E_{\mathrm{pp}}\) in \eqref{eq:ra-b10}. That mask excludes bases
\(P\le z_0\),
so the displayed decomposition needs no separate add-back.

Proposition~\ref{hyp:ra-applications}(3) proves the high-coefficient estimate
\begin{equation}\label{eq:ra-ra1}
 P_{I_H}(\Gamma)\ll_{A,B,B_0}NL^{-A}
 \qquad(A>0\ \hbox{fixed}).
\end{equation}
Its modulus coefficient is exactly \eqref{eq:ra-factories}.
The estimate permits no joint arithmetic mask.
Table~\ref{tab:ra-rows} gives its decomposition.
Each power boundary belongs to the interval on its right.

\begin{table}[ht]
\caption{The high-coefficient decomposition}\label{tab:ra-rows}
\small
\begin{tabular}{@{}lll@{}}
\hline
Sum and range & Tested array & Required estimate\\
\hline
\(T_1\), all \(a\le U\) & \(\mathcal R(c)\) & Rough-factor estimate\\
\(T_2\), \(c<N^{.46}\) & \(\mathcal R(c)\) & Two-factor convolution\\
\(T_2\), \(N^{.46}\le c<N^{.50}\) & \(\mathcal C(t)/(2L)\) & Central estimate\\
\(T_2\), \(N^{.50}\le c<N^{.54}\) & \(\mathcal R(c)\) & Central estimate\\
\(T_2\), \(c\ge N^{.54}\) & \(\mathcal R(c)\) & Rough-factor estimate\\
\(T_{\mathrm{II}}\), \(a<N^{.46}\) & \(\mathcal A(a)\) & Two-factor convolution\\
\(T_{\mathrm{II}}\), \(N^{.46}\le a<N^{.50}\) & \(\mathcal B(s)/(2L)\) & Central estimate\\
\(T_{\mathrm{II}}\), \(N^{.50}\le a<N^{.54}\) & \(\mathcal A(a)\) & Central estimate\\
\(T_{\mathrm{II}}\), \(a\ge N^{.54}\) & \(\mathcal B(s)/(2L)\) & Two-factor convolution\\
\hline
\end{tabular}
\end{table}


\subsection{The three distribution estimates}

\begin{lemma}[Rough-factor distribution under parent-scale and strict-roughness conditions]\label{hyp:ra-rough}
Let \(L=\log N\), \(z_0=N^{1/(\log\log N)^3}\), and
\(Q=N^{5227/10000}L^{-B}\), with fixed \(B\). Assume the two explicit
conditions used by the direct rows:
\begin{enumerate}
 \item[(W1)] The original parent satisfies
 \(cN^{.54}\le H\le N\), for a fixed \(c>0\). In the \(T_1\) row
 \(H=N/a\le N\); in the right \(T_2\) row \(H=N/t\le N\).
 \item[(W2)] Use the strict roughness convention
 \(\mathcal R(n)=\mathbf1_{P^-(n)>z_0}\), as in
 \eqref{eq:ra-rough-arrays}. It is not replaced by
 \(P^-(n)\ge z_0\).
\end{enumerate}
For every fixed \(T_{\mathrm{save}}>0\), uniformly in real intervals \(I\subset(0,2H]\)
(so \(|I|\le2H\); open, closed or one-point),
supported \(q\le Q\) with \((q,N)=1\), and reduced residues \((b,q)=1\),
\begin{equation}\label{eq:ra-theorem-r}
 \sum_{n\in I}\mathcal R(n)
 \left\{\mathbf1_{n\equiv b\pmod q}
          -\frac{\mathbf1_{(n,q)=1}}{\varphi(q)}\right\}
 \ll_{T_{\mathrm{save}}}\frac{H}{qL^{T_{\mathrm{save}}}}.
\end{equation}
The interval may be arbitrarily short.
The saving is measured at its original parent. The estimate is an integer-sieve statement and uses no prime
distribution input.
\end{lemma}
\begin{proof}
Put \(\ell=\log\log N\) and
\(\mathcal P_q=\{p\le z_0:p\nmid q\}\). The Euler product at
\(s=1+1/\log z_0\) gives
\(h_q=\sum_{p\in\mathcal P_q}p^{-1}\le3\ell\) for sufficiently large
\(N\). Choose the fixed constant \(K\) after \(T_{\mathrm{save}}\), with
\begin{equation}\label{eq:ra-bonferroni-constant}
 K>3e,\qquad K\log\!\left(\frac K{3e}\right)>T_{\mathrm{save}}+4.
\end{equation}
Let \(k=2\lceil K\ell/2\rceil\). The two adjacent Bonferroni truncations
bracket the strict exclusion of all primes in \(\mathcal P_q\).
Their normalized gap satisfies
\begin{equation}\label{eq:ra-theorem-r-lattice}
 k=2\lceil K\log\log N/2\rceil,\qquad
 E_k\le\frac{h_q^k}{k!}
       \le\left(\frac{3e}{K}\right)^{K\ell}
       \le L^{-T_{\mathrm{save}}-4},\qquad
 (k+1)z_0^k=N^{O(K/\ell^2)}.
\end{equation}
Every divisor in the truncation is at most \(z_0^k\).
By CRT, each divisor
contributes its expected interval mass with error \(O(1)\).
This is uniform for open, closed, or singleton intervals.
The same expansion applies to the comparator after using
\[
 \mathbf1_{(n,q)=1}=\sum_{\substack{d\mid\operatorname{rad}(q)\\d\mid n}}\mu(d).
\]
Thus the main terms cancel. The accumulated lattice error is
\(N^{O(K/\ell^2)}q/\varphi(q)\), and
\(q/\varphi(q)\le\omega(q)+1\ll L\). Under (W1),
\(H/q\gg N^{.0173}L^B\). The Bonferroni remainder is at
most \((|I|/q)E_kL^2\le2HL^{-T_{\mathrm{save}}-2}/q\) because \(|I|\le2H\). Hence both are
absorbed by \(H/(qL^{T_{\mathrm{save}}})\) once the fixed truncation precision is raised. This proves \eqref{eq:ra-theorem-r}.

The strict convention in (W2) has no threshold atom.
An application might instead use \(P^-(n)\ge z_0\).
Then the only possible difference is the atom
\(z_0=P\) when \(z_0\) is prime. That weak convention requires its separate
CRT boundary correction. It is not used in \eqref{eq:ra-theorem-r}.
\end{proof}
The positive gap \(.54-v_{\mathrm{RA}}=173/10000\) absorbs the lattice error
in \eqref{eq:ra-theorem-r-lattice}.
For \(T_1\), the free scale is at least \(N^{2/3}/4\).
For the right \(T_2\)-part, it exceeds a fixed multiple of \(N^{.54}\).
Partial summation then gives
\begin{equation}\label{eq:ra-direct-row-total}
 \sum_{m\le M_{\mathrm{out}}}\sum_{q\le Q}\frac{N}{mqL^{T_{\mathrm{save}}}}
 \ll NL^{2-T_{\mathrm{save}}}.
\end{equation}
Both kernel terms remain in this estimate.
Here \(M_{\mathrm{out}}\) is the outer endpoint of the current row.
It is \(U\) in the \(T_1\) row and the
corresponding low-row endpoint in the \(T_2\) decomposition.


The other tested arrays need uniformity under coprimality
with an auxiliary integer. Write \(\tau(h)=\sum_{d\mid h}1\) and
\(\omega(h)=\#\{p\in\mathbb P:p\mid h\}\).
For any sequence \(f(n)\) on its original parent \((H,2H]\), define
\begin{equation}\label{eq:ra-sw-difference}
\begin{aligned}
 \Delta_f(I;k,b;h)={}&
 \sum_{\substack{n\in I,\ n\equiv b\pmod k\\(n,h)=1}}f(n)\\
 &-\frac1{\varphi(k)}
     \sum_{\substack{n\in I\\(n,kh)=1}}f(n) .
\end{aligned}
\end{equation}
\begin{lemma}[Prime-variable transfer]\label{lem:ra-prime-transfer}
Let \(N^{1/3}/8\le H\le 2N\), \(L=\log N\), and
\(z_0=N^{1/(\log\log N)^3}\). Let \(y_0\ge z_0\) and
\[
 f(n)=\sum_{\substack{mp=n\\p\in J_m}}\alpha(m)w(p),
\]
where \(p\) is prime and \(|\alpha(m)|\le1\).
For each \(m\) the
allowed primes lie in a real interval \(J_m\subset[y_0,\infty)\).
Assume \(\|w\|_\infty+\operatorname{Var}(w)\le W\) on \([y_0,2H]\)
and \(|f(n)|\le1\) on \((H,2H]\). For every fixed \(P>0\), uniformly
in intervals \(I\subset(H,2H]\), \(k,h\ge1\), and \((b,k)=1\),
prime Siegel--Walfisz, Paper~I (2.21), implies
\begin{equation}\label{eq:ra-prime-transfer}
 |\Delta_f(I;k,b;h)|
 \ll_P WH\left(L^{-P}+\frac{\omega(h)L}{y_0}\right)
 \ll_P WH\tau(h)L^{-P}.
\end{equation}
\end{lemma}
\begin{proof}
Put \(K=2P+4\). If \(k>L^K\), the two terms in the discrepancy are
bounded by a constant times \(\|f\|_\infty(H/k+1+H/\varphi(k))\).
For each represented \(n\), a prime divisor \(p\ge y_0\) determines at
most one pair \((m,p)\), since \(m=n/p\). Hence
\[
 |f(n)|\le W\frac{\log(2H)}{\log y_0}
 \ll W(\log L)^3,
\]
using \(y_0\ge z_0\), \(H\le2N\), and
\(\log z_0=L/(\log L)^3\). Since \(\varphi(k)\gg\sqrt{k}\) and
\(H\ge N^{1/3}/8\), the resulting discrepancy is
\(\ll_P WHL^{-P}\). The two spare powers in \(K/2=P+2\) absorb
\((\log L)^3\). The additive \(1\) is absorbed because
\(H\ge N^{1/3}/8\gg L^P\).

Now let \(k\le L^K\). Terms with \((m,kh)>1\) vanish from both sides.
Otherwise the prime progression is \(p\equiv bm^{-1}\pmod{k}\) on
\(J_m\cap I/m\).
Primes dividing \(h\) are excluded on the progression side and in the
comparator.
For each \(m\) the two exclusions differ by at most
\(2W\omega(h)\).
There are at most \(2H/y_0\) possible
values of \(m\), so these terms contribute
\(\ll WH\omega(h)/y_0\).

For the remaining primes define
\[
 E(t)=\pi(t;k,c)-\frac1{\varphi(k)}
       \sum_{\substack{a\bmod k\\(a,k)=1}}\pi(t;k,a).
\]
The reduced-class sum equals \(\pi(t)-\#\{p\le t:p\mid k\}\).
Thus this is the centered prime discrepancy without the primes dividing \(k\).
Apply (2.21) to the class \(c\) and average its uniform error over
all reduced classes.
This gives \(E(t)\ll_{A'}t(\log t)^{-A'}\), with no factor
of \(k\). Since \(t\ge y_0\ge z_0\),
\(\log t\ge L/(\log L)^3\). Hence, for all sufficiently large \(N\),
\(k\le(\log t)^{2K}\).
Therefore (2.21), with modulus exponent \(2K\)
and fixed saving \(A'\), applies throughout each interval of
integration. Stieltjes partial summation bounds each \(m\)-contribution
by \(O(W(H/m)(\log y_0)^{-A'})\). Sum over
\(m\le2H/y_0\) and take \(A'=2(P+K+2)\).
This gives \(O_P(WHL^{-P})\). The estimate is measured at the original parent
\(H\), including for short intervals. Finally,
\(L/y_0\ll_P L^{-P}\) because \(y_0\ge z_0\), and
\(\omega(h)\le\tau(h)\). This proves \eqref{eq:ra-prime-transfer}.
\end{proof}

\begin{lemma}[Four-family auxiliary-coprime Siegel--Walfisz]
\label{lem:ra-sw}
By prime Siegel--Walfisz, Paper~I (2.21), for each of
\(\mathcal R,\mathcal A,\mathcal C/(2L),\mathcal B/(2L)\), the bound
\begin{equation}\label{eq:ra-full-sw}
 |\Delta_f(I;k,b;h)|\ll_P H\tau(h)L^{-P},
 \qquad (b,k)=1,
\end{equation}
holds uniformly for every \(I\subset(H,2H]\) and \(k,h\ge1\).
The parents are the
central ranges in Table~\ref{tab:ra-rows} and the convolution parents
\(N^{1/3}/8\le H\le2N^{.46+O(1/L)}\).
The proof uses the lower
bound only to absorb the large-modulus lattice term and to exclude \(n=1\).
It uses \(H\le2N^{.54}\) for the remaining weight bounds. The divisor
exponent is the single fixed value \(C_{\mathrm{SW}}=1\), independent of
\(P\). The implied constant may depend on \(P\) and the fixed family
parameters, but not on \(N,H,I,k,h,b\). The error is measured at \(H\),
even when \(I\) is short.
\end{lemma}
\begin{proof}
For sufficiently large \(N\), every parent has
\(H\ge N^{1/3}/8>1\), so \(n=1\) is excluded. The rough and signed-rough
arrays have exact prime-variable decompositions.
For \(\mathcal R(n)\), take \(\alpha(m)=\mathcal R(m)\),
\(J_m=[P^+(m),\infty)\cap(z_0,\infty)\), and \(w(p)=1\).
Use \(P^+(1)=1\). The largest prime factor gives the unique term. For
\(\mathcal A(n)=\mu(n)\mathcal R(n)\),
take \(\alpha(m)=-\mu(m)\mathcal R(m)\),
\(J_m=(\max(P^+(m),z_0),\infty)\), and \(w(p)=1\). The strict
inequality excludes repeated largest primes, where \(\mu(n)=0\).

For \(\mathcal B(s)/(2L)\), the \(j=1\) terms in
\(b=p^j\mid s\), \(b>V\), use \(\alpha(m)=\mathcal R(m)\),
\(J_m=(V,\infty)\), and \(w(p)=\log p/(2L)\). For
\(\mathcal C(t)/(2L)\), the \(j=1\) terms use
\(\alpha(a)=\mu(a)\mathcal R(a)\mathbf1_{a\le U}\),
\(J_a=(z_0,V]\), and the same weight. In each case
\(\|w\|_\infty+\operatorname{Var}(w)\le1\).
The four tested arrays
are one-bounded on their parent ranges. For the last two arrays this follows
from \(\sum_{d\mid n}\Lambda(d)=\log n\) and
\(\log(2H)<2L\). On the central ranges \(H\le2N^{.54}\).
There the weight bound follows from \(y_0=V=N^{1/3}\) for \(\mathcal B\) and \(y_0=z_0\)
for \(\mathcal C\), since \(\log(2H)\le .54L+O(1)\).

For \(\mathcal B\) and \(\mathcal C\), write the remaining \(j\ge2\)
terms as \(f_{\ge2}\). They need not be nonnegative.
For instance, the
\(\mu(a)\) coefficient in \(\mathcal C\) may change sign.
Take absolute values before summing:
\[
 \sum_{n\le2H}|f_{\ge2}(n)|
 \ll \frac{H}{L}\sum_{p>z_0}\sum_{j\ge2}\frac{\log p}{p^j}
 \ll \frac{H\log z_0}{Lz_0}
 \ll \frac{H}{z_0}.
\]
For any sequence \(g\),
\(|\Delta_g|\le2\sum_{n\le2H}|g(n)|\).
Hence this remainder is
\(O_P(HL^{-P})\) for every fixed \(P\). Applying
Lemma~\ref{lem:ra-prime-transfer} to each prime term proves
\eqref{eq:ra-full-sw} with the same exponent \(C_{\mathrm{SW}}=1\).
\end{proof}
The mean-square estimate below has five uses.
Section~\ref{sec:k2} uses it for the zero mode of the exterior sum.
Section~\ref{sec:retained-side} uses it for the Class II estimate of
Paper~II inside the retained side.
Section~\ref{sec:l3} uses it for the row-three dispersion estimate of the
switched sum. These three apply Paper~I, Theorem~3.2.
In the large-norm case of Paper~I, (3.48), part (ii) below replaces
\cite[Theorem~0(a)]{BFI86}. The other two uses are the central rectangular
estimate of Lemma~\ref{hyp:ra-central} and the small-modulus norm-relative
step of Lemma~\ref{hyp:ra-l52}. Part (i) derives the conclusion of that
published theorem (Paper~I, Theorem~2.4).
Its inputs are the norm-relative condition
\eqref{eq:ra-mean-square-a2} and
an elementary large sieve, so the argument is self-contained
at this step. We do not correct \cite{BFI86}.
Paper~I's use of it is unaffected. The norm-relative premise is not deduced from the
absolute parent-scale estimate \eqref{eq:ra-full-sw} alone.
Let \(\xi\) be complex and supported in \([1,Z]\), \(Z\ge3\). Fix
\(B_{\mathrm{div}}\ge0\). Consider the following condition.
For every
\(P>0\), uniformly in integers \(h,k\ge1\) and reduced classes
\(b\bmod k\),
\begin{equation}\label{eq:ra-mean-square-a2}
 \left|\sum_{\substack{n\equiv b\ (k)\\(n,h)=1}}\xi_n
 -\frac1{\varphi(k)}\sum_{(n,kh)=1}\xi_n\right|
 \ll_P\|\xi\|_2Z^{1/2}\tau(h)^{B_{\mathrm{div}}}(\log(2Z))^{-P},
\end{equation}
The constants are common to the family of pairs \((Z,\xi)\) under
consideration.
Write \(\Delta_\xi(k,b)\) for the centered discrepancy at \(h=1\). The norm conclusion is
\begin{equation}\label{eq:ra-transferred-mean-square}
 \sum_{k\le Z(\log Z)^{-B_1(A)}}\sum_{b\bmod k}^{*}
       |\Delta_\xi(k,b)|^2
 \ll_A\|\xi\|_2^2 Z(\log Z)^{-A}.
\end{equation}
In the selected-factor form of Paper~I, (3.48),
\(\xi_n=\alpha_n1_{(n,r)=1}\), \(Z=\lceil2T_{\mathrm{int}}\rceil\), and
\begin{equation}\label{eq:ra-mean-square-transfer}
 \sum_{q\le3Q}\sum_{b\bmod q}^{*}|\Delta_\xi(q,b)|^2
 \ll N_{\mathcal I}T_{\mathrm{int}} L^{-A_3},\qquad
 3Q\le Z(\log Z)^{-B_F},\quad B_F=A_3+15,
\end{equation}
where \(N_{\mathcal I}\) is the interval mass in Paper~I, (3.8).
\begin{lemma}[Transferred mean-square estimate]
\label{hyp:ra-ms-transfer}
\begin{enumerate}
 \item[(i)] If \eqref{eq:ra-mean-square-a2} holds, then
 \eqref{eq:ra-transferred-mean-square} holds for every fixed \(A>0\).
 Here \(B_1(A)\) and the implied constant depend only on \(A\),
 \(B_{\mathrm{div}}\) and the common constants.
 \item[(ii)] Assume the hypotheses (3.1), (3.2), (3.6), (3.8) with its
 length condition, (3.11) and (3.25)--(3.27) of Paper~I. Then the estimate
 \eqref{eq:ra-mean-square-transfer} holds uniformly over the parameter
 boxes and selected factors of Paper~I, Theorem~3.2. Hence that theorem
 holds with this input in place of \cite[Theorem~0(a)]{BFI86} at (3.48).
\end{enumerate}
\end{lemma}
\begin{proof}
(i) Expand \(\Delta_\xi(k,b)\) in Dirichlet characters modulo \(k\) and
pass to primitive characters. For conductors up to a fixed power of
\(\log Z\), apply \eqref{eq:ra-mean-square-a2}.
Take the auxiliary divisor
\(h\) to be the imprimitive part of the modulus. For larger conductors
use the multiplicative large sieve inequality.
Its elementary proof uses
finite Fourier identities and the spacing of reduced fractions.
Summing over the moduli gives (i).
(ii) If \(\|\xi\|_2^2\le N_{\mathcal I}L^{-A_3-1}\), the elementary bound
in Paper~I, (3.48), applies. Otherwise (3.25), with auxiliary divisor
\(dr\), and (3.26) give \eqref{eq:ra-mean-square-a2} for
\(\xi_n=\alpha_n1_{(n,r)=1}\) with divisor exponent one. Also (3.6) gives
the modulus range in \eqref{eq:ra-mean-square-transfer}. Part (i) then
applies. The complete argument is in Appendix~
\hyperref[subsec:supp-lem-65]{\ref*{subsec:supp-lem-65}}, with its
row index in Appendix~
\hyperref[subsec:app-ra-rows]{\ref*{subsec:app-ra-rows}}. We assert no implication from
\eqref{eq:ra-full-sw} alone to the norm bound.
\end{proof}

For the three convolution rows, let
\begin{equation}\label{eq:ra-l52-scales}
 D_1=HN^{-\sigma_{\mathrm p}},\qquad
 D_2=QN^{w_{\mathrm{RA}}+\sigma_{\mathrm p}}/H,\qquad
 \sigma_{\mathrm p}=10^{-6},\quad w_{\mathrm{RA}}=9/1000,
\end{equation}
where \(N^{1/3}/8\le H\le2N^{.46+O(1/L)}\).
A supported modulus has
the unique form \(q=pd\) of \eqref{eq:ra-coefficient}.
Here \(p\) is the marked prime and \(d\) the squarefree cofactor.
\begin{equation}\label{eq:ra-l52-caps}
 r<D_1,\qquad s<D_2,\qquad rs=q<Q.
\end{equation}

\begin{lemma}[Deterministic allocation]\label{lem:ra-selector}
For each supported \(q=pd\), the following rule defines a unique coprime
split \(q=rs\) satisfying \eqref{eq:ra-l52-caps}. Put \(p\) in \(r\).
Order the cofactor primes exceeding \(N^{w_{\mathrm{RA}}}\) as
\(b_1>\cdots>b_m\).
Here \(m\le(v_{\mathrm{RA}}-t_0)/w_{\mathrm{RA}}=37553/1260<30\). Process them in this order.
Put \(b_i\) in \(r\) exactly when the current \(r b_i<D_1\), and put it
in \(s\) otherwise. Let \(r_{\rm big}\) and \(s_{\rm big}\) be the
resulting factors.
Order the remaining cofactor primes as
\(l'_1>\cdots>l'_k\). Set
\(r_0=r_{\rm big}\prod_{j=1}^k l'_j\). If \(r_0<D_1\), take
\(r=r_0\), \(s=s_{\rm big}\). Otherwise let \(j_*\) be the least index
for which \(r_0/(l'_1\cdots l'_{j_*})<D_1\), and take
\[
 r=r_0/(l'_1\cdots l'_{j_*}),\qquad
 s=s_{\rm big}l'_1\cdots l'_{j_*}.
\]
Call the resulting pair \(\operatorname{Sel}_H(q)=(r,s)\).
\end{lemma}
\begin{proof}
The marked prime fits, since
\(p<N^{1/3-\eta_{\rm m}}<D_1\), even at the fixed lower endpoint
\(H\ge N^{1/3}/8\). The accepted big-prime partial products stay below
\(D_1\).
Thus \(r_{\rm big}<D_1\), and the least index \(j_*\) exists
when needed. Write all cofactor primes in decreasing order as
\(l_1>l_2>\cdots\). The odd cube gates in \eqref{eq:ra-coefficient}
imply
\[
 p l_1\cdots l_{i-1}l_i^2<Q\qquad\text{for every }i:
\]
For odd \(i\), drop one factor from the cube gate.
For even \(i\), use
the preceding odd gate and \(l_i\le l_{i-1}\).

Suppose \(r_0<D_1\), and let \(b_i\) be the last big prime rejected
into \(s\). The rejection test gives \(r^{(i)}b_i\ge D_1\). The
two-gate inequality then gives
\[
 s=s_{\rm big}
 =\frac{p b_1\cdots b_i}{r^{(i)}}
 \le\frac{p b_1\cdots b_{i-1}b_i^2}{D_1}<\frac Q{D_1}<D_2.
\]
If no big
prime is rejected, \(s=1\). In the prefix case, minimality gives
\(r l'_{j_*}\ge D_1\), hence \(r\ge D_1N^{-w_{\mathrm{RA}}}\) and
\[
 s=q/r<QN^{w_{\mathrm{RA}}}/D_1=D_2.
\]
The factors partition the squarefree prime factors of \(q\).
Thus
\((r,s)=1\), \(rs=q<Q\), and \(r<D_1\).
\end{proof}

Define the exact joint coefficient on each dyadic Cartesian modulus shell by
\begin{equation}\label{eq:ra-global-tensor}
\begin{gathered}
 F_H(r,s)=
 \begin{cases}
 \Gamma(rs),&rs\text{ is supported and }\operatorname{Sel}_H(rs)=(r,s),\\
 0,&\text{otherwise},
 \end{cases}\\
 \sum_q\Gamma(q)\Phi(q)=\sum_{r,s}F_H(r,s)\Phi(rs),\\
 F_H(r,s)=\sum_\nu c_\nu f_\nu(r)g_\nu(s),\qquad
 |f_\nu|,|g_\nu|\le1,\qquad
 \sum_\nu|c_\nu|\ll L^{C_{\mathrm{tensor}}},\quad C_{\mathrm{tensor}}=2^{30}+91.
\end{gathered}
\end{equation}
Here \(\operatorname{Sel}_H\) is the rule in Lemma~\ref{lem:ra-selector}.
The expansion is exact on the full Cartesian parent, including zero values
off the selected pairs.

\begin{lemma}[Global tensor for the selector]\label{lem:ra-tensor}
The coefficient \(F_H\) has a rank-one expansion satisfying
\eqref{eq:ra-global-tensor}, with the sharper bound
\begin{equation}\label{eq:ra-tensor-norm}
 \sum_\nu|c_\nu|
 \le4(2^{30}-1)(3+L)^{78}\ll L^{79}.
\end{equation}
\end{lemma}
\begin{proof}
Fix a pair of dyadic shells. A branch fixes the merge order
\(\pi\in\{R,S\}^m\) of the at most 29 big cofactor primes. It also fixes the final
large-prime/prefix case.
In the prefix case it fixes the parity of the number
of factors assigned to \(s\). There are at most
\(4\sum_{m=0}^{29}2^m=4(2^{30}-1)\) branches. Within a branch, unary
conditions impose squarefreeness, coprimality with \(2N\) and the unique
marked prime \(p\) in \(r\). They also impose the bounds \(P^+(r/p)<z_s\) and
\(P^+(s)<z_s\), assigned big-prime counts, \(r<D_1\), \(s<D_2\), and
\(r_{\rm big}<D_1\). The remaining unary data are the case indicator, the parity of the small-prime
count, and the sign \(\Gamma(rs)=\mu(rs/p)=-\mu(r)\mu(s)\).

Every cross-condition is a comparison \(\phi(r)<\psi(s)\) or its weak
counterpart. Let \(K_0\) be the number of distinct values of \(\phi\) on
the \(r\)-shell. For ranks \(1\le a\le K_0\), \(0\le b\le K_0\), put
\(M=2K_0+1\) and
\[
 c_t=M^{-1}\sum_{u=0}^{K_0}e(-tu/M).
\]
Finite Fourier orthogonality gives the exact separator
\begin{equation}\label{eq:ra-rank-separator}
 \mathbf1_{a\le b}=\sum_{t\bmod M}c_t e(tb/M)e(-ta/M),
 \qquad \sum_{t\bmod M}|c_t|\le3+\log M.
\end{equation}
Rank \(\phi(r)\) among its shell values, and count the values below
\(\psi(s)\).
Each shell has at most \(Q\) values, so
\(\log M\ll L\). The cross-comparisons are exactly:
\begin{enumerate}
 \item the adjacent strict comparisons in the fixed merge order of the
 big primes (at most 28);
 \item for each \(S\)-letter \(b_i\), the greedy rejection test
 \(p(\prod_{b_j>b_i,\ b_j\in r}b_j)b_i\ge D_1\) (at most 29);
 \item the odd cube gates
 \(p l_1\cdots l_{i-1}l_i^3<Q\) at odd indices among the big primes
 (at most 15);
 \item the remaining odd cube gates on the small-prime part, grouped by
 the fixed parity into the two maxima of the \(r\)- and \(s\)-side
 partial products (two comparisons);
 \item in the prefix case, the prefix-order test
 \(\min P>\max R_{\rm small}\) and the least-prefix test
 \(r\min P\ge D_1\) (two comparisons);
 \item the support range \(Q_0<rs<Q\) (two comparisons).
\end{enumerate}
There are at most 78 comparisons in total. Once a merge branch is fixed,
each displayed test compares an \(r\)-only quantity with an
\(s\)-only quantity. The small-prime maxima combine all gates of the same
parity and side. The branch conditions are exactly the greedy tests, the
cube gates and the least-prefix condition.
Hence their union is precisely
\(F_H\). Multiplying the separators in
\eqref{eq:ra-rank-separator} gives \eqref{eq:ra-tensor-norm}. The unary
factors remain one-bounded. The bound in \eqref{eq:ra-tensor-norm} is per
dyadic shell pair. The resulting \(O(L^2)\) shell pairs are absorbed by the
chosen logarithmic precision.
\end{proof}

For the convolution rows, split \(s=dt\) with
\(P^+(d)\le z< P^-(t)\), where
\(z=N^{1/(\log\log N)^3}\) and \(y=N^{1/\log\log N}\). The aggregate
with \(d>y\) is bounded by Rankin's method. The mean-square modulus in the
proof of Lemma~\ref{hyp:ra-l52} is \(k=rd\), not the full \(q\).
Its scale is
\begin{equation}\label{eq:ra-l52-small-modulus}
 Q'=D_1y=HN^{-\sigma_{\mathrm p}+1/\log\log N},\qquad
 k\le5Q',\qquad Z=\lceil2H\rceil.
\end{equation}
The inner factor \(t\) is \(z\)-rough and is removed in \(L^2\), with cost
\(H^2L^C/z\). No bound for \(\tau(r)\) and no roughness condition on \(r\)
is used. The required range is
\begin{equation}\label{eq:ra-l52-range}
 5Q'\le Z(\log Z)^{-B_F},\qquad B_F=P_{\mathrm{MS}}+15,
\end{equation}
where \(P_{\mathrm{MS}}\) is the fixed source-level precision in the
supplementary Lemma~6.5(i). This holds since
\(Q'/H=N^{-\sigma_{\mathrm p}+o(1)}\). The full-modulus
range \(3Q\le Z(\log Z)^{-B_F}\) is not asserted for these parents.

\begin{lemma}[Two-factor convolution estimate]\label{hyp:ra-l52}
Let \(N^{1/3}/8\le H\le2N^{.46+O(1/L)}\) and \(HM\asymp N\). Let \(a_n\),
\(n\sim H\), be one-bounded, supported on \(P^-(n)>z\), and satisfy
\eqref{eq:ra-full-sw} at the original parent \(H\), where
\(z=N^{1/(\log\log N)^3}\). Let \(b_m\), \(m\sim M\), be one-bounded.
Let \(U_1,V_1\) be dyadic modulus scales, allowing the singleton scale
\(1/2\) for the integer \(1\), such that
\[
 U_1\le D_1,\qquad V_1\le D_2,\qquad U_1V_1\le N^{v_{\mathrm{RA}}}.
\]
Let \(f(r)\), \(r\sim U_1\), and \(g(s)\), \(s\sim V_1\), be arbitrary
one-bounded arrays.
Keep the unary restrictions \(r<D_1\), \(s<D_2\),
\((r,N)=(s,N)=1\). Neither a pointwise condition \(rs<Q\) nor a joint condition
\((r,s)=1\) is imposed on the pair. Then, for every fixed \(A>0\), uniformly in these
occupied-parent scales, \(f\), \(g\), unary interval cuts of \(a_n\) and
\(b_m\), and intervals \(I\subset[N/8,2N]\),
\begin{equation}\label{eq:ra-l52-input}
 \sum_{r,s}f(r)g(s)
 \sum_{\substack{n,m\\nm\in I}}a_n b_m\mathcal K_{rs}(nm)\ll_A NL^{-A}.
\end{equation}
\end{lemma}
\begin{proof}
We use the dispersion method of Paper~I, Theorem~3.2.
It runs on the
whole Cartesian shell \(r\sim U_1\), \(s\sim V_1\), with residue class
\(N\bmod rs\). Its occupied-parent scale conditions are
\(U_1\le HN^{-\sigma_{\mathrm p}}\),
\(V_1\le N^{v_{\mathrm{RA}}+w_{\mathrm{RA}}}/(HN^{-\sigma_{\mathrm p}})\), and
\(U_1V_1\le N^{v_{\mathrm{RA}}}\).
The second follows from the displayed
bound by \(D_2\) after dropping \(L^{-B}\). The tested sequence is
strictly \(z\)-rough. The unary supports carry \((rs,N)=1\) and the
individual caps. No pointwise product cap is needed for the rank-one
arrays.
After Cauchy's inequality in \(m\), the split \(s=dt\) above reduces the
modulus of the mean square to \(k=rd\) in the range
\eqref{eq:ra-l52-range}. The mean square is split at a fixed power of
\(L\) below its target. Below it the trivial bound suffices. Above it
\eqref{eq:ra-full-sw} gives \eqref{eq:ra-mean-square-a2} with a loss of
that power, so Lemma~\ref{hyp:ra-ms-transfer}(i) applies. The
prime Siegel--Walfisz input is Paper~I, (2.21). No joint condition \((r,s)=1\) is used, because \(r,s\) are not the
moduli of Paper~I, Theorem~2.8. Write \(s=dt\) as above and \(k=rd\). A
prime \(p\mid(r,s)\) with \(p\le z\) divides \(d\) and only enlarges \(k\).
The Chinese remainder step needs only that \(k,t_1,t_2\) are pairwise
coprime.
The smooth--rough split gives \((d,t_1t_2)=1\). A prime
\(p\mid(r,s)\) with \(p>z\) divides \(k\) and a rough \(t_i\). Such tuples
are the bad type \((k,t_1t_2)>1\). If \(p\mid t_2\), the second progression
\(kt_2\mid mn_2-N\) places \(n_2\) in one class modulo \(kp\). If
\(p\mid t_1\) only, exchange the two copies first; the absolute square
majorant is symmetric, and the same statement then applies. We therefore
remove these tuples before the Chinese remainder step.
The cost in the square is
\(\ll MHL^C[H/(zK)+1]\). After the factor \(KM\) of the
first Cauchy step this is
\(\ll NL^C(z^{-1/2}+(K/H)^{1/2})\).
That is \(O_A(NL^{-A})\), since
\(z=N^{1/(\log\log N)^3}\) and \(K/H\le N^{-\sigma_{\mathrm p}/2}\).
The zero-mode identity holds for all tuples, including those with gcd
failures. After the second Cauchy step, Paper~I, Theorem~2.8 bounds the nonzero
frequencies. It is applied with \(r_{\mathrm{int}}=tt'\), the product of two
\(z\)-rough inner variables, and \(s_{\mathrm{int}}=p_0\), the least prime not dividing
\(N\). Since \(p_0\ll\log N<z\), \((r_{\mathrm{int}},s_{\mathrm{int}})=1\), exactly as in Paper~I,
(3.82)--(3.83). This is the form of Paper~I, Theorem~2.8
(Sections~5.3 and~7.4) used here. Its statement allows arbitrary one-bounded
\(f,g\) with only \((rs,N)=1\). The product restriction
\(nm\in I\) is kept as in Paper~I, Theorem~3.2.
The finite product
partition used there also handles the endpoints.
Neither \eqref{eq:ra-mean-square-transfer} at the full modulus \(3Q\) nor
a norm-relative consequence of \eqref{eq:ra-full-sw} for small norms is
used. The complete argument is in Appendix~
\hyperref[subsec:supp-lem-68]{\ref*{subsec:supp-lem-68}}, with its
row index in Appendix~
\hyperref[subsec:app-ra-rows]{\ref*{subsec:app-ra-rows}}.
\end{proof}


The central rows use \(Y=N^y\), \(y=341/5000\).
Choose the largest divisor \(s\mid d\) with \(s\le Y\).
Since \(p<N^{1/3-\eta_{\mathrm m}}\) and \(q=pd>Q_0\),
\(d>Q_0/p\gg N^{1/6+\eta_{\mathrm m}}L^{-B_0}>Y\).
Hence a next divisor exists.
The cube gates control adjacent divisors and give
\begin{equation}\label{eq:ra-central-selector}
 \frac{Yq}{Q}<s\le Y,\qquad r=q/s<Q/Y.
\end{equation}
This allocation keeps a one-bounded joint modulus coefficient.
The required central estimate has
\begin{equation}\label{eq:ra-m53-range}
\begin{gathered}
 H_{\mathrm c}M_{\mathrm c}\asymp N,\qquad
 N^{.50}\ll H_{\mathrm c}\ll N^{.54},\\
 Q_1\ll N^{.4545},\qquad Q_2\ll N^{.0682}.
\end{gathered}
\end{equation}
\begin{lemma}[Central rectangular distribution]\label{hyp:ra-central}
Let \(x\) be even, \(H_{\mathrm c}M_{\mathrm c}\asymp x\), and let the
bounds in \eqref{eq:ra-m53-range} hold with \(N\) replaced by \(x\), with
fixed comparison constants. Let \(a_n,b_m\) be one-bounded on the complete
dyadic rectangle \(n\sim H_{\mathrm c},m\sim M_{\mathrm c}\). Let the
tested \(a_n\) satisfy \eqref{eq:ra-full-sw} at its original parent.
The divisor exponent is fixed, and the constants are uniform at every requested
precision. Write
\[
 \mathcal D(q,a)=\sum_{n\sim H_{\mathrm c},m\sim M_{\mathrm c}}a_nb_m
 \left(1_{nm\equiv a\ (q)}-\frac{1_{(nm,q)=1}}{\varphi(q)}\right).
\]
Then, for every fixed \(A>0\),
\begin{equation}\label{eq:ra-m53-input}
 \sum_{r\sim Q_1}\max_{(b,r)=1}
 \sum_{\substack{s\sim Q_2\\\mu(rs)^2=1,\ (rs,x)=1}}
 \max_{\substack{a\bmod rs,\ (a,rs)=1\\a\equiv b\ (r)}}
       |\mathcal D(rs,a)|\ll_A x(\log x)^{-A}.
\end{equation}
The constant is uniform over these families, rectangles, squarefree
modulus supports and unary interval restrictions of \(a_n\) and \(b_m\).
All errors are measured at the original parents. One outer residue
\(b\) is chosen before \(s\) varies.
In the application it is
\(N\bmod r\).
\end{lemma}
\begin{proof}
This is the rectangular dispersion theorem proved in Appendix~
\hyperref[subsec:supp-lem-69]{\ref*{subsec:supp-lem-69}} and indexed
in Appendix~\hyperref[subsec:app-ra-rows]{\ref*{subsec:app-ra-rows}}.
The restriction
\((rs,x)=1\) only removes nonnegative terms. Its proof opens the inner
modulus \(s\) by dispersion with the outer residue fixed. It uses the mean
square over all moduli of Lemma~\ref{hyp:ra-ms-transfer}(i), which gives
Paper~I, (2.17), under its hypothesis (2.16). The supplementary proof applies the mean square twice.
The second application is repeated over two
correlations. At each application a constant sequence
on the same parent is added. Its sign is chosen by the parallelogram
identity to make the tested sequence norm-large. The resulting sequence
and the constant sequence satisfy \eqref{eq:ra-mean-square-a2} uniformly.
Their squared norms are at most the original parent times a fixed
power of \(L\). Fix the requested logarithmic modulus buffer \(B_1\)
first.
Then choose the source-level parameter so that the moduli lie in the
range of \eqref{eq:ra-transferred-mean-square}. Apply part (i) to both
sequences,
subtract by linearity, and absorb that fixed power by increasing the
requested precision. This uses neither part (ii) nor a mean-square
conclusion for an arbitrary small-norm sequence from absolute
Siegel--Walfisz alone. The other inputs are Paper~I, (2.21) (prime
Siegel--Walfisz), (2.35) (Weil--Estermann)
and Lemma~3.4 (completed hyperbola count). It does not use Paper~I, Theorem~2.8. For the four tested
arrays of Table~\ref{tab:ra-rows}, Lemma~\ref{lem:ra-sw} supplies
\eqref{eq:ra-full-sw} with divisor exponent one.
\end{proof}
Lemma~\ref{hyp:ra-central} concerns complete dyadic rectangles and
one-bounded sequences. Lemma~\ref{lem:ra-product-region} below reduces the
four central rows, with their product cuts, to it.
The raw exponent checks are
\begin{equation}\label{eq:ra-central-margins}
\begin{aligned}
 7v_{\mathrm{RA}}+5y&=3.9999<4,\\
 3v_{\mathrm{RA}}-y-1&=.4999<.50,\\
 1-v_{\mathrm{RA}}+y&=.5455>.54.
\end{aligned}
\end{equation}
The three margins verify the numerical range only. They do not establish
the rectangular-support or family hypotheses for a selected product region.


Product cuts require an additional finite reduction. For \(e\in I_H\), put
\begin{equation}\label{eq:ra-restoration}
\begin{aligned}
 w(e)&=\frac{2L}{\log e},\\
 \mathcal A(a)\mathcal R(c)\frac{\log c}{\log(ac)}
 &=\left(\mathcal A(a)\mathcal R(c)\frac{\log c}{2L}\right)w(ac),\\
 -\frac{\mathcal C(t)\mathcal R(c)}{\log(tc)}
 &=\left(-\frac{\mathcal C(t)}{2L}\mathcal R(c)\right)w(tc),\\
 \frac{\mathcal A(a)\mathcal B(s)}{\log(as)}
 &=\left(\mathcal A(a)\frac{\mathcal B(s)}{2L}\right)w(as).
\end{aligned}
\end{equation}
Here \(\log c/(2L)\) is unary, and
\(\|w\|_\infty\le2L/(L-\log4)=O(1)\), and
\[
 \operatorname{Var}_{I_H}(w)
 =\frac{2L}{L-\log4}-\frac{2L}{L+\log(1-\varepsilon)}=O(1).
\]
Thus uniform product-prefix estimates restore the exact weights by
Stieltjes summation, including the minus sign of \(T_2\).
Use multiplicative cells of relative width \(L^{-D_{\mathrm m}}\).
The original-source boundary strips satisfy
\begin{equation}\label{eq:ra-strip}
 E_{\mathrm{strip}}\ll
       \bigl(NL^{-D_{\mathrm m}/2}+N^{1/2}\bigr)L^{15/2}.
\end{equation}
Each unnormalized source multiplicity is at most \(3d(e)\).
Cauchy and \(d(n)^4\le d_{16}(n)\) give \eqref{eq:ra-strip}.
Sum the boundary strips before the modulus expansion.
For a final saving \(A\), fix
\begin{equation}\label{eq:ra-precision-order}
 D_{\mathrm m}=2A+40,\qquad
 T_{\mathrm{save}}>A+C_{\mathrm{tensor}}+2D_{\mathrm m}+40.
\end{equation}
Then fix every Siegel--Walfisz and half-level precision.
Choose the common threshold last.

\begin{lemma}[Product-region reduction for the central rows]
\label{lem:ra-product-region}
Fix \(A>0\) and choose \(D_{\mathrm m}\) and \(T_{\mathrm{save}}\) as in
\eqref{eq:ra-precision-order}. Take one of the four central rows of
Table~\ref{tab:ra-rows}.
Keep its unary cuts, its normalized arrays from
\eqref{eq:ra-restoration} and the coefficient \(\Gamma\). For
\(N/4<u\le(1-\varepsilon)N\) let \(\mathcal S(u)\) be its normalized
discrepancy sum over \(N/4<nm\le u\), without the factor \(w\). Then
\begin{equation}\label{eq:ra-central-rows}
 \sup_{N/4<u\le(1-\varepsilon)N}|\mathcal S(u)|\ll_A NL^{-A},
\end{equation}
and the same bound holds for the sum over \(nm\in I_H\) with the factor
\(w(nm)\).
\end{lemma}
\begin{proof}
\emph{Selector.} Let \(q=pd\) be supported, with the primes of \(d\) in
decreasing order \(\ell_1>\cdots>\ell_k\). As in the proof of
Lemma~\ref{lem:ra-selector}, the cube gates give
\(p\ell_1\cdots\ell_{i-1}\ell_i^2<Q\) for every \(i\), hence
\[
 \frac{\ell_i}{\ell_{i+1}\cdots\ell_k}<\frac Qq\qquad(1\le i\le k).
\]
Consecutive divisors of a squarefree \(d\) have ratio at most
\(\max_i\ell_i/(\ell_{i+1}\cdots\ell_k)\). Write \(d=\ell_1M\).
The divisors of \(d\) form the two chains \(\mathcal D(M)\) and
\(\ell_1\mathcal D(M)\). If \(\ell_1>M\), the only new gap is
\(\ell_1/M\). If \(\ell_1\le M\), the ranges \([1,M]\) and
\([\ell_1,\ell_1M]\) overlap. A gap of the union below \(M\) lies inside a
gap of \(\mathcal D(M)\).
One above \(M\) lies inside a gap of
\(\ell_1\mathcal D(M)\). Induction on \(k\) gives the claim. Since
\(d>Y\), the largest divisor \(s\le Y\) of \(d\) has a successor
\(s'>Y\) with \(s'/s<Q/q\).
This proves \eqref{eq:ra-central-selector}.
The pair \((r,s)\) is determined by \(q\), and \(q=rs\).
Put \(F(r,s)=\Gamma(rs)\) on selected pairs and \(F(r,s)=0\) otherwise.
This coefficient is
one-bounded and supported on squarefree \(rs\) with \((rs,2N)=1\). Its
dyadic shells satisfy \(Q_1\ll N^{v_{\mathrm{RA}}-y}\) and
\(Q_2\ll N^{y}\), where \(v_{\mathrm{RA}}-y=909/2000\) and
\(y=341/5000\).
There are \(O(L^2)\) shell pairs.

\emph{Parents.} In the first central \(T_2\) row the tested variable and
array are \(n=t\) and
\(\mathcal C(t)/(2L)\).
In the second they are \(n=c\) and
\(\mathcal R(c)\). In the central \(T_{\mathrm{II}}\) rows they are
\(n=s\) and \(\mathcal B(s)/(2L)\) in the first.
In the second they are \(n=a\) and \(\mathcal A(a)\). The other variable carries the remaining one-bounded array,
with the sign of \(T_2\). On \(I_H\), every relevant pair of dyadic parents
has \(H_{\mathrm c}M_{\mathrm c}\asymp N\) and
\(N^{.50}\ll H_{\mathrm c}\ll N^{.54}\) with fixed constants.
There are
\(O(L^2)\) such pairs. Lemma~\ref{lem:ra-sw} gives \eqref{eq:ra-full-sw}
for the tested arrays at these parents.

\emph{Restriction.} Let \(a_n\) be a tested array on its parent and \(J\)
an interval.
On the same parent, \(a_n\mathbf1_J(n)\) satisfies
\(\Delta_{a\mathbf1_J}(I;k,b;h)=\Delta_a(I\cap J;k,b;h)\) for every
interval \(I\). It therefore satisfies \eqref{eq:ra-full-sw} with the same
constants, measured at the original parent. The other array may be
restricted to any interval.

\emph{Mesh.} Put \(\delta=L^{-D_{\mathrm m}}\). Divide each parent into
\(O(\delta^{-1})\) intervals with endpoint ratio at most \(1+\delta\).
All unary cut points of the row are endpoints. For fixed \(u\), call
a cell complete if it lies inside \(\{N/4<nm\le u\}\). Any other cell
meeting this region has all its products within a factor \((1+\delta)^2\)
of \(N/4\) or of \(u\). The part of the region it covers lies in the set
\(\mathcal E_u\) of integers \(e\in(N/4,(1-\varepsilon)N]\) with \(|e-N/4|\le3\delta N\) or
\(|e-u|\le3\delta N\).
Here \(\#\mathcal E_u\ll\delta N+1\). For a complete
cell apply Lemma~\ref{hyp:ra-central} with \(x=N\), saving \(T_{\mathrm{save}}\), the
restricted arrays, \(b=N\bmod r\) and \(a=N\bmod rs\). Both classes are
reduced since \((rs,N)=1\), and \(a\equiv b\pmod r\). Thus
\[
 \Bigl|\sum_{r,s}F(r,s)\mathcal D_{\mathrm{cell}}(rs,N)\Bigr|
 \le\sum_{r\sim Q_1}\max_{(b,r)=1}
   \sum_{\substack{s\sim Q_2\\\mu(rs)^2=1,\ (rs,N)=1}}
   \max_{\substack{a\bmod rs\\a\equiv b\ (r)\\(a,rs)=1}}
   |\mathcal D_{\mathrm{cell}}(rs,a)|
 \ll_{T_{\mathrm{save}}} NL^{-T_{\mathrm{save}}}.
\]
There are \(O(L^{2D_{\mathrm m}+4})\) choices of cell and shell pair, so the
 complete cells contribute \(O(NL^{-T_{\mathrm{save}}+2D_{\mathrm m}+4})\).

\emph{Strips.} Sum the incomplete cells together before expanding
the modulus coefficient. For fixed \(e\), the factorizations
\(e=nm\) in the row carry total normalized weight at most \(d(e)\).
Also \eqref{eq:ra-proper-power-bound} bounds the modulus sum by
\(d(N-e)+O(L)\). Their contribution is therefore
\[
 \ll\sum_{e\in\mathcal E_u}d(e)\{d(N-e)+L\}
 \ll(\#\mathcal E_u)^{1/2}N^{1/2}L^{15/2}
 \ll\bigl(N\delta^{1/2}+N^{1/2}\bigr)L^{15/2},
\]
by Cauchy's inequality twice, \(d(n)^4\le d_{16}(n)\) and
\(\sum_{n\le N}d_{16}(n)\ll NL^{15}\). This is \eqref{eq:ra-strip}.

\emph{Conclusion.} With \(D_{\mathrm m}=2A+40\) and
\(T_{\mathrm{save}}>A+C_{\mathrm{tensor}}+2D_{\mathrm m}+40\), the two contributions are
\(O(NL^{-A-C_{\mathrm{tensor}}-36})\) and \(O(NL^{-A-25/2}+N^{1/2}L^{15/2})\), uniformly
in \(u\). This proves \eqref{eq:ra-central-rows}. Since
\(\|w\|_\infty+\operatorname{Var}_{I_H}(w)=O(1)\), Stieltjes summation
against \(\mathcal S\) gives the weighted bound, with the sign of \(T_2\)
kept.
\end{proof}


\subsection{The main coefficient}

For primes \(\ell\nmid2N\), set \(g_N(\ell)=1/(\ell-1)\), and put
\(g_N(\ell)=0\) for \(\ell\mid2N\). Extend \(g_N\) multiplicatively to
squarefree integers, with \(g_N(\ell^j)=0\) for \(j\ge2\). For primes
\(e\equiv N\pmod p\) in the fixed reduced class modulo \(p\), and primes
\(\ell\nmid2Np\), the condition \(\ell\mid(N-e)/p\) has
density \(1/\varphi(\ell)=g_N(\ell)\). Put
\begin{equation}\label{eq:ra-euler-density}
 V_N(z_s)=\prod_{\substack{2<\ell<z_s\\\ell\nmid N}}
       \left(1-\frac1{\ell-1}\right)
 =\frac{2e^{-\gamma}C(N)}{\alpha L}(1+o(1)),
\end{equation}
with \(C(N)\) as in Paper~II, (1.2). For a marked prime \(p\in N^K\) and a
level \(R\), put \(D=R/p\) and \(s_p=\log D/(\alpha L)\).
Let \(P=P_N(z_s)\) be the product of the primes \(\ell<z_s\) with
\(\ell\nmid2N\).

\begin{lemma}[Upper-sieve application]\label{lem:ra-upper-sieve}
\begin{enumerate}
 \item[(a)] The comparison \eqref{eq:ra-positive-comparison} holds.
 \item[(b)] There is an absolute constant \(K_{\mathrm{dim}}\) such that,
 for every even \(N\),
 \begin{equation}\label{eq:ra-product-condition}
  \prod_{w\le\ell<z}(1-g_N(\ell))^{-1}
  \le\frac{\log z}{\log w}\left(1+\frac{K_{\mathrm{dim}}}{\log w}\right)
  \quad(z>w\ge2).
 \end{equation}
 \item[(c)] For the three pairs
 \((R,K)\in\{(R_L,J),(R_{\mathrm{tail}},J_{\mathrm{tail}}),(Q,J_t)\}\),
 uniformly in the mark and in even \(N\),
 \begin{equation}\label{eq:ra-finite-density}
  \sum_{d\mid P}\frac{\mu_d^+(R/p)}{\varphi(d)}
  \le V_N(z_s)\{F(s_p)+O(L^{-1/3})\},
 \end{equation}
 and, with \(R=N^rL^{-B_r}\),
 \begin{equation}\label{eq:ra-mark-density}
  \sum_q\frac{c_{R,K}(q)}{\varphi(q)}
  \le\frac{2e^{-\gamma}C(N)}{\alpha L}
  \left\{\int_K F\left(\frac{r-t}{\alpha}\right)\frac{dt}{t}+o(1)\right\}.
 \end{equation}
\end{enumerate}
\end{lemma}
\begin{proof}
We use Iwaniec's upper bound for the linear sieve
\cite{Iwaniec1980Remainder}. Let \(\mathcal A\) be a finite sequence and
\(D>1\) a level.
Let \(P(z)\) be the product of the sifting primes below \(z\). The
upper weight \(\mu_d^+\) of \cite[p.~313]{Iwaniec1980Remainder} is
\((-1)^k\) when \(d=\ell_1\cdots\ell_k\mid P(z)\),
\(\ell_1>\cdots>\ell_k\), and
\(\ell_1\cdots\ell_{2l}\ell_{2l+1}^3<D\) for every
\(0\le l\le(k-1)/2\).
Otherwise it is \(0\). With \(j=2l+1\) these are the
gates of \eqref{eq:ra-rosser-weight}, so \eqref{eq:ra-rosser-weight} is
this weight on \(d\mid P\). By the definition of \(W^+\)
\cite[p.~310]{Iwaniec1980Remainder} and Lemma~2 there
\cite[p.~313]{Iwaniec1980Remainder},
\[
 \sum_{d\mid P(z)}\mu_d^+|\mathcal A_d|=W^+\ge S(\mathcal A,P,z)
\]
for every \(D>1\), since \(W^+\) is \(S(\mathcal A,P,z)\) plus nonnegative
sifted counts. If the density satisfies \eqref{eq:ra-product-condition}
with a constant \(K\) and \(z\le D\), Lemma~3, (12), of
\cite[p.~313]{Iwaniec1980Remainder} gives, with \(s=\log D/\log z\),
\[
 \sum_{d\mid P(z)}\mu_d^+g(d)
 \le V(z)\{F(s)+O(e^{\sqrt K-s}(\log D)^{-1/3})\}.
\]

(a) Let \(e\) be a prime in one of the pieces \eqref{eq:ra-three-pieces}
and \(p\) a mark with \(pm=N-e\). If \(P^-(m)\ge z_s\), then
\((m,P)=1\). The inequality \(W^+\ge S\) for the one-term sequence
\(\{m\}\) at level \(D=R/p\) gives
\begin{equation}\label{eq:ra-upper-weight-positivity}
 \sum_{\substack{d\mid P\\d\mid m}}\mu_d^+(D)
   \ge\mathbf1_{(m,P)=1}\ge\mathbf1_{P^-(m)\ge z_s}.
\end{equation}
This holds for every \(D>1\), in particular when \(s_p=383/200<2\).
Omitting \((m,Np)=1\) in \eqref{eq:ra-original-s6} only enlarges the
count. The condition \(d\mid m\) is \(e\equiv N\pmod{pd}\), and
\(c_{R,K}(pd)=\mu_d^+(R/p)\) by \eqref{eq:ra-weight-equality}. The pieces
\eqref{eq:ra-three-pieces} partition \([2,(1-\varepsilon)N)\times J\). The
point \(e=N/4\) lies in \(I_L\) and \(t=1/3-\eta_{\mathrm m}\) lies in
\(J_{\mathrm{tail}}\). Summing \eqref{eq:ra-upper-weight-positivity} over
\(e\) and \(p\), each mark with its multiplicity, gives
\(S_6\le A_{\mathrm{mix}}\) by \eqref{eq:ra-mix-definition}.

(b) For odd \(\ell\),
\[
 \Bigl(1-\frac1{\ell-1}\Bigr)^{-1}
 =\Bigl(1-\frac1\ell\Bigr)^{-1}\Bigl(1-\frac1{(\ell-1)^2}\Bigr)^{-1}.
\]
The factors with \(\ell\mid2N\) equal one.
Omitting them only lowers the product. Mertens' theorem, with its uniform error term, bounds the
product of the first factors over \(w\le\ell<z\) by
\((\log z/\log w)(1+K_1/\log w)\). The product of the second factors
over \(\ell\ge\max(w,3)\) is at most \(1+K_2/\log w\). The constants
\(K_1,K_2\) are absolute. This gives \eqref{eq:ra-product-condition} with one
\(K_{\mathrm{dim}}\), independent of \(N\) and of \(\omega(N)\).

(c) The limiting ranges of \(s_p\) for the three pairs are
\[
 \Bigl(\frac{383}{200},\frac{12639}{4480}\Bigr],\qquad
 \Bigl(\frac{383}{200},\frac{1916149}{1000000}\Bigr],\qquad
 \Bigl(\frac{544243}{250000},\frac{43148397}{14000000}\Bigr],
\]
all inside \([383/200,3.083)\subset(1,\infty)\). The factors \(L^{-B_r}\)
move \(s_p\) by \(O(\log L/L)\). Hence \(z_s\le D\) and \(\log D\asymp L\)
for \(N\ge N_0(B_r)\).
Then (12) gives \eqref{eq:ra-finite-density}, with
an error \(O(L^{-1/3})\) uniform in the mark, the pair and \(N\). For
\eqref{eq:ra-euler-density}, use the factorization in (b) for odd
\(\ell<z_s\) and restore the factors with \(\ell\mid N\). The primes
\(\ell\mid N\) with \(\ell\ge z_s\) number at most \(1/\alpha\) and change
the product by a factor \(1+O(1/z_s)\). By \eqref{eq:ra-weight-equality},
\[
 \sum_q\frac{c_{R,K}(q)}{\varphi(q)}
 =\sum_{\substack{p\in N^K\\p\nmid N}}\frac1{p-1}
   \sum_{d\mid P}\frac{\mu_d^+(R/p)}{\varphi(d)}.
\]
The outer weights are positive, so \eqref{eq:ra-finite-density} may be
summed termwise. Since \(F\) is Lipschitz on the compact range above, the
prime number theorem and partial summation give \eqref{eq:ra-mark-density}.
Marks dividing \(N\) number at most \(1/t_0\), and endpoint atoms and the
logarithmic buffers contribute \(o(1)\).
\end{proof}
The delayed branch \(s_p>3\) of \(F\) occurs only for \(c_{\mathrm{trim}}\),
when \(t<t_{\mathrm{RA}}\), where \(t_{\mathrm{RA}}=v_{\mathrm{RA}}-3\alpha\).


Define
\begin{equation}\label{eq:ra-j-functional}
 J_K(r)=\frac1{2e^\gamma\alpha}
          \int_K F\left(\frac{r-t}{\alpha}\right)\frac{dt}{t}.
\end{equation}
The rational endpoint \(\overline J_6\) is as in Paper II
\cite[(A.13)]{PaperII}:
\begin{equation}\label{eq:ra-j6bar-comparison}
\begin{gathered}
 \overline J_6=
 \frac{13148581959157321533513247276782285310758770468886237619585}
      {10^{58}},\\
 J_J(1/2)=2\log(110/57)\le\overline J_6.
\end{gathered}
\end{equation}
The inequality is certified by outward interval evaluation.
No central approximation is rounded.
Since \(W=4(1-\varepsilon)C(N)N/L^2\), the interval masses give
\begin{equation}\label{eq:ra-main-normalized}
\begin{aligned}
 \frac{M_{\mathrm{mix}}}{W}
 &\le w_LJ_J(1/2)
       +w_H\{J_{J_{\mathrm{tail}}}(1/2)+J_{J_t}(v_{\mathrm{RA}})\}+o(1),\\
 w_L&=\frac{2500}{9999},\qquad w_H=\frac{7499}{9999}.
\end{aligned}
\end{equation}
Here \(w_L+w_H=1\).
No mark multiplicity or singular-series factor is suppressed.

On \(J_{\mathrm{tail}}\), the explicit upper-sieve formula gives
\begin{equation}\label{eq:ra-trim-debit}
\begin{aligned}
 0&\le w_H\{J_{J_{\mathrm{tail}}}(1/2)-J_{J_{\mathrm{tail}}}(v_{\mathrm{RA}})\}\\
 &=w_H\int_{1/3-\eta_{\mathrm m}}^{1/3}
       \frac{v_{\mathrm{RA}}-1/2}{t(1/2-t)(v_{\mathrm{RA}}-t)}\,dt\\
 &\le\frac{w_H\eta_{\mathrm m}(v_{\mathrm{RA}}-1/2)}
              {(1/3)(1/2-1/3)(v_{\mathrm{RA}}-1/3)}\\
 &=:d_{\mathrm m}=\frac{5106819}{31557955000}.
\end{aligned}
\end{equation}
The last integrand is increasing on this interval.
Put
\begin{equation}\label{eq:ra-exact-gain}
 \mathfrak g_{\mathrm{RA}}
   =w_H\{J_J(1/2)-J_J(v_{\mathrm{RA}})\}-d_{\mathrm m}.
\end{equation}
Combining \eqref{eq:ra-main-normalized} with \eqref{eq:ra-j6bar-comparison} gives
\begin{equation}\label{eq:ra-b11-integral}
 \frac{M_{\mathrm{mix}}}{W}
       \le\overline J_6-\mathfrak g_{\mathrm{RA}}+o(1).
\end{equation}

The remaining integral is elementary.
The linear-sieve functions $F$ and $f$ are as in
\cite[(1.8)--(1.9), pp.~172--173]{Iwaniec1980Rosser}.
Set \(t_{\mathrm{RA}}=v_{\mathrm{RA}}-3\alpha\). On the ranges here,
\begin{equation}\label{eq:ra-delay-formula}
 \frac{F(s)}{2e^\gamma}
  =\begin{cases}
      1/s,&1\le s\le3,\\[2pt]
      \displaystyle\frac1s\left\{1+
            \int_2^{s-1}\frac{\log(\xi-1)}\xi\,d\xi\right\},&3<s<5.
    \end{cases}
\end{equation}
This is the linear-sieve normalization
\cite[p.~202]{Iwaniec1980Rosser}\cite[(2.23)]{PaperI}.
Write the delayed contribution as
\begin{equation}\label{eq:ra-delayed-integral}
 \mathcal I_{\mathrm{RA}}=
 \int_{t_0}^{t_{\mathrm{RA}}}\frac{dt}{t(v_{\mathrm{RA}}-t)}
 \int_2^{(v_{\mathrm{RA}}-t)/\alpha-1}
                 \frac{\log(\xi-1)}\xi\,d\xi.
\end{equation}
Substitution gives
\begin{equation}\label{eq:ra-gain-elementary}
\begin{aligned}
 \mathfrak g_{\mathrm{RA}}+d_{\mathrm m}
 =w_H\Bigg\{&\int_{t_0}^{1/3}
   \left(\frac1{t(1/2-t)}-\frac1{t(v_{\mathrm{RA}}-t)}\right)dt\\
 &-\mathcal I_{\mathrm{RA}}\Bigg\}.
\end{aligned}
\end{equation}
Swap the order in \eqref{eq:ra-delayed-integral} to obtain
\begin{equation}\label{eq:ra-swapped-integral}
\begin{aligned}
 \mathcal I_{\mathrm{RA}}&=\frac1{v_{\mathrm{RA}}}
 \int_3^{s_{\max}}h(\zeta)\,d\zeta,\qquad
 s_{\max}=\frac{v_{\mathrm{RA}}-t_0}{\alpha},\\
 h(\zeta)&=\frac{\log(\zeta-2)}{\zeta-1}
 \log\frac{(v_{\mathrm{RA}}-\alpha\zeta)(v_{\mathrm{RA}}-t_0)}
               {\alpha\zeta t_0}.
\end{aligned}
\end{equation}
For \(n=20000\), set \(d=(s_{\max}-3)/n\). The composite midpoint sum
has error at most
\[
 \left|\int_3^{s_{\max}}h(\zeta)\,d\zeta
      -d\sum_{j=0}^{n-1}h(3+(j+1/2)d)\right|
 \le\frac{(s_{\max}-3)d^2}{24}\sup|h''|.
\]
Directed interval differentiation on \(400\) covering boxes gives
\(\sup|h''|\le0.720711111\). An independent directed-interval replay
encloses
\(\mathfrak g_{\mathrm{RA}}\) in
\([0.09895983533531812,0.09895983533543704]\).
Outward rounding therefore gives
\begin{equation}\label{eq:ra-gain-enclosure}
 0.0989598353353\le\mathfrak g_{\mathrm{RA}}
                      \le0.0989598353355.
\end{equation}
Thus we may use the certified rational floor
\begin{equation}\label{eq:ra-rational-floor}
 g_{\mathrm{RA,lo}}=\frac{989598353}{10^{10}}=0.0989598353,
 \qquad \mathfrak g_{\mathrm{RA}}\ge g_{\mathrm{RA,lo}}.
\end{equation}
The exact integral \(\mathfrak g_{\mathrm{RA}}\) and its certified floor
\(g_{\mathrm{RA,lo}}\) are distinct.

\begin{proposition}[Sixth-sum source applications]\label{hyp:ra-applications}
On the exact arrays above, the following hold.
\begin{enumerate}
 \item The comparison \eqref{eq:ra-positive-comparison}, and the density
 bound \eqref{eq:ra-mark-density} for the three pairs \((R,K)\).
 \item The three low discrepancy bounds \eqref{eq:ra-three-bv-errors}.
 \item For every fixed \(A>0\), each of the nine rows of
 Table~\ref{tab:ra-rows}, with its exact weight, is \(O_A(NL^{-A})\).
 Consequently the high-coefficient estimate \eqref{eq:ra-ra1} holds.
\end{enumerate}
\end{proposition}
\begin{proof}
Items 1 and 2 are Lemmas~\ref{lem:ra-upper-sieve} and~\ref{lem:ra-n4-bv}.
For item 3 fix \(A\). Choose in this order: \(D_{\mathrm m}=2A+40\);
\(T_{\mathrm{save}}>A+C_{\mathrm{tensor}}+2D_{\mathrm m}+40\); the saving \(A+C_{\mathrm{tensor}}\) in
Lemma~\ref{hyp:ra-l52} and the saving \(T_{\mathrm{save}}\) in Lemma~\ref{hyp:ra-central}.
Then fix the Siegel--Walfisz precisions that these require, the buffers
\eqref{eq:ra-bv-buffer} and \(B\), and the constant \(K_{\mathrm{dim}}\).
Choose the common threshold last. The rows partition \(T_1+T_2+T_{\mathrm{II}}\):
\(T_1\) is all of \(a\le U\). The sums \(T_2\) in \(c\) and \(T_{\mathrm{II}}\)
in \(a>U\) are split at \(N^{.46}\), \(N^{.50}\) and \(N^{.54}\).
Each boundary goes to the interval on its right.

\emph{Direct rows.} These are \(T_1\) and the part of \(T_2\) with
\(c\ge N^{.54}\). Fix the outer variable \(a\), respectively \(t\). Terms
with \((a,q)>1\), respectively \((t,q)>1\), vanish in \(\mathcal K_q\).
Otherwise \(\mathcal K_q(ac)\) is the kernel of \eqref{eq:ra-theorem-r} for
\(c\) in the reduced class \(N\bar a\bmod q\), and similarly for \(t\). The
parent is \(H=N/a\), respectively \(H=N/t\). The free variable runs over
\(H/4<c<(1-\varepsilon)H\), intersected with \(c\ge N^{.54}\) in the
\(T_2\)-row and with any prefix cut.
This interval lies in \((0,2H]\).
Condition (W1) holds.
For \(a\le U\) we have \(N^{2/3}\le H\le N\).
If the right \(T_2\)-row is nonempty, then
\(N^{.54}<N/t=H\le N\), since
\(N^{.54}\le c<(1-\varepsilon)N/t\). Thus Lemma~\ref{hyp:ra-rough} applies
with this \(H\). (Equivalently, split at \(H/2\) into
\((H/4,H/2]\) and \((H/2,(1-\varepsilon)H)\).
Each lies inside a dyadic parent
\((H',2H']\) with \(H'\in\{H/4,H/2\}\ge N^{.54}/4\). This only doubles the
constant.) The weights
\(\log c/\log(ac)\) and \(\mathcal C(t)/\log(tc)\) are bounded and
monotone in \(c\). Partial summation and \eqref{eq:ra-direct-row-total}
with \(T_{\mathrm{save}}>A+2\) bound both rows by \(O(NL^{-A})\).

\emph{Convolution rows.} These are \(T_2\) with \(c<N^{.46}\) (tested
\(\mathcal R(c)\), \(c>N^{1/3}/4\)), \(T_{\mathrm{II}}\) with
\(a<N^{.46}\) (tested \(\mathcal A(a)\), \(a>N^{1/3}\)), and
\(T_{\mathrm{II}}\) with \(a\ge N^{.54}\) (tested \(\mathcal B(s)/(2L)\),
\(N^{1/3}<s<2N^{.46+O(1/L)}\)). Each tested array contains the
strict roughness factor \(\mathcal R(n)=1_{P^-(n)>z}\).
This is explicit
for \(\mathcal R(c)\).
It follows from
\(\mathcal A(a)=\mu(a)\mathcal R(a)\) and from
\(\mathcal B(s)=\mathcal R(s)\sum_{b\mid s,\ b>V}\Lambda(b)\).
Lemma~\ref{lem:ra-sw} gives \eqref{eq:ra-full-sw} on these original
parents. Write
\(\sum_q\Gamma(q)\Phi(q)=\sum_{r,s}F_H(r,s)\Phi(rs)\) as in
\eqref{eq:ra-global-tensor}. Lemma~\ref{lem:ra-selector} places the
joint support in \eqref{eq:ra-l52-caps}. Discard a dyadic shell pair
\(r\sim U_1\), \(s\sim V_1\) if \(F_H\) vanishes identically there.
Otherwise a supported pair in it satisfies
\(U_1V_1<rs<Q\), \(U_1<D_1\), and \(V_1<D_2\).
Hence this occupied
shell satisfies all scale conditions of Lemma~\ref{hyp:ra-l52}.
Lemma~\ref{lem:ra-tensor} expands \(F_H\) on the whole Cartesian shell
into rank-one terms with \(\sum_\nu|c_\nu|\ll L^{79}\).
The unary caps and coprimality are kept. Each rank-one term is estimated on that
whole shell. Its factors need not preserve the joint condition
\(rs<Q\) or the condition \((r,s)=1\), since
Lemma~\ref{hyp:ra-l52} imposes neither. Apply Lemma~\ref{hyp:ra-l52} with saving \(A+C_{\mathrm{tensor}}\) to each
term on each occupied shell pair and each product prefix
\(I=(N/4,u]\). Since
\(C_{\mathrm{tensor}}\ge83\), the total is \(O(NL^{-A})\), uniformly in \(u\). The factor
\(w\) is then restored by Stieltjes summation, as in
\eqref{eq:ra-restoration}.

\emph{Central rows.} The four central rows are bounded by
Lemma~\ref{lem:ra-product-region}.

\emph{Proper powers.} By \eqref{eq:ra-b10} and
\eqref{eq:ra-proper-power-bound},
\(P_{I_H}(\Gamma)=T_1+T_2+T_{\mathrm{II}}+O(N^{1/2+\nu}L^2)\). This
proves \eqref{eq:ra-ra1}.
\end{proof}

\begin{proposition}[Sixth-sum estimate]\label{prop:ra-b11}
Fix \(\sigma\) in the range stated above before taking \(N\to\infty\).
Then, along even integers,
\begin{equation}\label{eq:ra-b11}
 \frac{S_6}{W}\le\overline J_6-\mathfrak g_{\mathrm{RA}}+o(1).
\end{equation}
In particular, the certified floor in \eqref{eq:ra-rational-floor} gives
\begin{equation}\label{eq:ra-b11-floor}
 \frac{S_6}{W}\le\overline J_6-g_{\mathrm{RA,lo}}
                         +E_{\mathrm{RA}}(N;\sigma),
 \qquad E_{\mathrm{RA}}(N;\sigma)\ge0,\qquad
 E_{\mathrm{RA}}(N;\sigma)\longrightarrow0.
\end{equation}
\end{proposition}

\begin{proof}
By Proposition~\ref{hyp:ra-applications}(1), \(S_6\le A_{\mathrm{mix}}\).
Use \eqref{eq:ra-b9-f1}, obtained from \eqref{eq:ra-b9} and
\eqref{eq:ra-cancellation}. Fix \(A>2\), so that \(NL^{-A}=o(W)\). The
three low discrepancies are \(O(NL^{-A-10})\) by
Proposition~\ref{hyp:ra-applications}(2).
The term
\eqref{eq:ra-high-conversion} is \(O_A(NL^{-A})\).
Also \(P_{I_H}(\Gamma)=O_A(NL^{-A})\) by
Proposition~\ref{hyp:ra-applications}(3). Hence
\(S_6/W\le M_{\mathrm{mix}}/W+o(1)\), and \eqref{eq:ra-b11-integral} gives
\eqref{eq:ra-b11}. Enlarge the vanishing error to make it nonnegative. The
certified floor \eqref{eq:ra-rational-floor} gives \eqref{eq:ra-b11-floor}.
\end{proof}

The precisions are fixed in the order \eqref{eq:ra-precision-order}.
All constants precede the common threshold.
Their dependence through prime Siegel--Walfisz may be ineffective.
The final section uses only \eqref{eq:ra-b11-floor}.
It uses the certified rational floor in \eqref{eq:bill-components}.
Its gain is independent of the common coefficient used in the
balanced-semiprime and switched comparisons.

\section{The balanced semiprime term}\label{sec:k2}

We bound the complete occurrence sum \(B_b(N)\).
The two prime orientations keep their multiplicities.
In this section, fix
\begin{equation}\label{eq:k2-parameters}
\begin{gathered}
 \frac13<b<\frac12,\qquad \varepsilon=\eta_L=10^{-4},\qquad
 u=\frac1{500},\qquad \mu_* =\frac{25}{128},\qquad L=\log N,\\
 0<\sigma\le10^{-10},\qquad \tau=\sigma^9,\qquad \rho=1+\tau,\\
 N'=N(1-L^{-333}),\\
 \mathcal N_K=(1-\varepsilon)U_C=\frac{W(N)}4,\\
 I(b)=\int_b^{1/2}\frac{dt}{t(1-t)}\\
 =\log\frac{1-b}{b}.
\end{gathered}
\end{equation}
The final argument uses \eqref{eq:intro-exponents}.
Every limit below fixes \(b\) and \(\sigma\) first.
For row three, write
\(\ell_0=53/412-\eta_L/2=264897/2060000>0\). This is the fixed floor of
Paper~I, \S3.1.

\subsection{The convolution source}

Define
\begin{equation}\label{eq:k2-source}
\begin{aligned}
 A_b(e)&=\#\{(m,p')\in\mathbb P^2:
 N^b\le m\le\sqrt N,\quad \varepsilon N<mp'=e\le N'\},\\
 X_b&=\sum_e A_b(e),\qquad
 B_b' =\sum_e A_b(e)1_{N-e\in\mathbb P},\\
 E_{\mathrm{top}}&=B_b-B_b',\qquad
 0\le A_b(e)\le2,\qquad
 0\le E_{\mathrm{top}}\le2(NL^{-333}+1).
\end{aligned}
\end{equation}
The last bound counts integers in the removed interval.
A prime square contributes one occurrence.
No new ordering or coprimality restriction is imposed.

For a finite sequence \(F\), put
\begin{equation}\label{eq:k2-natural-kernel}
 \Delta_F(q,a)=\sum_eF(e)
 \left(1_{e\equiv a\ (q)}-\frac{1_{(e,q)=1}}{\varphi(q)}\right),
 \qquad (a,q)=1.
\end{equation}
Both terms use the full support.
For this source, PNT and partial summation give
\begin{equation}\label{eq:k2-source-mass}
 X_b=\frac{(1-\varepsilon)N}{L}I(b)+o_b(N/L).
\end{equation}
This is the calculation of Paper~II,
\cite[(5.44)--(5.51)]{PaperII}.
The moving top removes at most \(2(NL^{-333}+1)\) mass.

The required convolution estimate is
\begin{equation}\label{eq:k2-bv-source}
 \sum_{q\le Q}\max_{(a,q)=1}|\Delta_{A_b}(q,a)|
 \ll_{A,b}NL^{-A},\qquad
 Q\le N^{1/2}L^{-B_{\mathrm{BV}}(A,b)},
\end{equation}
for every fixed \(A>0\).
Its proof keeps the complete occurrence source and the moving top.
Fix a top \(T\in\{N,N'\}\). Before the discrepancy estimate, compare the
exact cell sources \eqref{eq:k2-cell-source} with their upper and lower
convolutions. Each signed child keeps the natural
coprimality comparator \(1_{(e,q)=1}/\varphi(q)\). It is not replaced
after summing the blocks. The step from
\(T=N\) to \(T=N'\) costs the endpoint bands
\eqref{eq:k2-boundary-cost}. It also costs the removed top mass in
\eqref{eq:k2-source}. The actual moving-top application is proved below.

We use the following convolution input.
Let \(a_m,b_n\) have supports \(m\sim M,n\sim Y\).
Put \(x=MY\ge3\), and fix \(\eta,B_{\mathrm{div}}>0\).
Assume
\begin{equation}\label{eq:k2-convolution-sw}
\begin{gathered}
 x^\eta\le M,Y\le x^{1-\eta},\\
 \left|\sum_{\substack{n\equiv l\ (k)\\(n,h)=1}}b_n
 -\frac1{\varphi(k)}\sum_{(n,hk)=1}b_n\right|
 \ll_{A_2}\|b\|_2Y^{1/2}d(h)^{B_{\mathrm{div}}}
                       (\log(2Y))^{-A_2}
\end{gathered}
\end{equation}
uniformly for every fixed \(A_2>0\), \(h,k\ge1\),
and \((l,k)=1\).
Here \(d(n)\) is the divisor function.
Then Fact~\ref{fact:bv-full-support} gives
\begin{equation}\label{eq:k2-convolution}
 \sum_{q\le Q}\max_{(a,q)=1}|\Delta_{a*b}(q,a)|
 \ll_A\|a\|_2\|b\|_2x^{1/2}(\log x)^{-A},
 \qquad Q\le x^{1/2}(\log x)^{-B_{\mathrm{BV}}}.
\end{equation}
The constants depend on all family constants in
\eqref{eq:k2-convolution-sw}.
This specializes Fact~\ref{fact:bv-full-support}. It is
Paper~II, Lemma~5.1 \cite[(5.23)--(5.28)]{PaperII}.
It applies to complete rectangular supports.

We make the interval reduction explicit.
Fix a target \(A>2\), and put \(H_g=A+4\).
Choose
\begin{equation}\label{eq:k2-interval-grid}
\begin{gathered}
 R_g=\left\lceil
 \frac{(1/2-b)L}{\log(1+L^{-H_g})}\right\rceil,\qquad
 \xi=N^{(1/2-b)/R_g},\qquad y_j=N^b\xi^j,\\
 J_j=[y_{j-1},y_j),\qquad Y_j=y_j.
\end{gathered}
\end{equation}
Close the final upper endpoint.
For the fixed top \(T\in\{N,N'\}\), put
\begin{equation}\label{eq:k2-convolution-envelopes}
\begin{aligned}
 b_j(m)&=1_{m\in\mathbb P\cap J_j},\\
 a_{j,T}^+(p')&=1_{p'\in\mathbb P}
       1_{\varepsilon N/Y_j\le p'\le T\xi/Y_j},\\
 a_{j,T}^-(p')&=1_{p'\in\mathbb P}
       1_{\varepsilon N\xi/Y_j<p'<T/Y_j}.
\end{aligned}
\end{equation}
For \(F,G\) on the positive integers write
\((F*G)(e)=\sum_{mn=e}F(m)G(n)\), and put
\(A_{j,T}^{\pm}=a_{j,T}^{\pm}*b_j\). The exact cell source is
\begin{equation}\label{eq:k2-cell-source}
 A_{j,T}(e)=1_{\varepsilon N<e\le T}
       \sum_{\substack{m\in\mathbb P\cap J_j,\ p'\in\mathbb P\\mp'=e}}1.
\end{equation}
The definitions give the pointwise envelopes
\(A_{j,T}^-\le A_{j,T}\le A_{j,T}^+\).
Their difference is supported in multiplicative boundary bands at
\(\varepsilon N\) and \(T\). Put
\(\mathcal E_Q(F)=\sum_{q\le Q}\max_{(a,q)=1}|\Delta_F(q,a)|\).
Let
\(B_T(e)=\sum_j(A_{j,T}^+(e)-A_{j,T}^-(e))\). Then \(B_T\ge0\), with
support in
\([\varepsilon N/\xi,\varepsilon N\xi]\cup[T/\xi,T\xi]\). Each integer
has at most two ordered factorizations into primes. Since
\(1<\xi\le1+L^{-H_g}\) and \(T\in\{N,N'\}\) satisfies \(T\le N\),
\[
 \sum_e B_T(e)\le
 2\{\varepsilon N(\xi-\xi^{-1})+T(\xi-\xi^{-1})+4\}
 \ll NL^{-H_g}+1.
\]
Since \(\xi-\xi^{-1}\le2(\xi-1)\le2L^{-H_g}\), counting integers in the
two complete boundary bands gives the displayed mass bound.
Paper~II, (5.47), gives the aggregate boundary discrepancy
\begin{equation}\label{eq:k2-boundary-cost}
 \mathcal E_Q\!\left(\sum_j(A_{j,T}^+-A_{j,T}^-)\right)
       \ll NL^{-H_g+1}+Q.
\end{equation}
The \(Q\) term includes the bounded-multiplicity changes at exact
endpoints. Split the second factor into dyadic blocks \((M,2M]\):
\[
 a_{j,T,M}^{\pm}(p')=a_{j,T}^{\pm}(p')1_{M<p'\le2M}.
\]
For each block meeting the support,
\(\varepsilon N/(2Y_j)<M\le T\xi/Y_j\), so
\(M\asymp_{\varepsilon}N/Y_j\), with only
\(O_{\varepsilon}(1)\) blocks. Here \(n\sim Z\) means
\(Z<n\le2Z\), as in Section~\ref{sec:classical-inputs}. For the short
factor, the theorem-scale parent is $Y=Y_j/2$; since $\xi<2$ eventually,
$J_j\subset(Y_j/2,Y_j]$ lies in that parent. The norms are
\begin{equation}\label{eq:k2-interval-norms}
 \|a_{j,T,M}^+\|_2\ll\sqrt M,\qquad
 \|a_{j,T,M}^-\|_2\le\|a_{j,T,M}^+\|_2,\qquad
 \|b_j\|_2\ll\sqrt{Y_j}L^{-(H_g+1)/2}.
\end{equation}
The prime number theorem on these intervals also gives a matching
lower bound. It is \(\|b_j\|_2\gg\sqrt{Y_j}L^{-(H_g+1)/2}\) for all sufficiently
large \(N\). This lower bound normalizes the conversion of the absolute
prime Siegel--Walfisz error to
\eqref{eq:k2-convolution-sw}.
The prime-interval form of \eqref{eq:k2-convolution-sw} holds uniformly
in \(j,T,M\). For a fixed target \(A>2\), write
\(K_{\mathrm{SW}}(A)\) for the
chosen logarithmic level. For auxiliary moduli \(k\le L^{K_{\mathrm{SW}}(A)}\), this is prime
Siegel--Walfisz at arbitrarily high fixed logarithmic precision.
The divisor loss from excluding primes dividing the auxiliary integers is
absorbed into \(d(h)^{B_{\mathrm{div}}}\). For \(k>L^{K_{\mathrm{SW}}(A)}\), use
\[
 |\Delta|\le |J|/k+1+\pi(J)/\varphi(k)
\]
and \(\sup_{k>L^{K_{\mathrm{SW}}(A)}}\varphi(k)^{-1}
\ll L^{-K_{\mathrm{SW}}(A)}\log\log(3L^{K_{\mathrm{SW}}(A)})\), taking
\(K_{\mathrm{SW}}(A)\) larger than the
requested saving. For every fixed target \(A>2\), choose
\(H_g=A+4\). The grid has \(O_b(L^{H_g+1})\) intervals. The prime
number theorem gives the norm bounds \eqref{eq:k2-interval-norms}.
With \(A_{\mathrm{block}}=A+H_g/2+2\), Fact~\ref{fact:bv-full-support}
and Paper~II, (5.23)--(5.28), bound the rectangular blocks by
\(O_{A,b}(NL^{-A-3/2})\). The boundary term in
\eqref{eq:k2-boundary-cost} is \(O(NL^{-A-3}+Q)\). Increasing the
single fixed logarithmic level buffer absorbs the \(x\asymp_b N\)
constants. Thus the complete moving-top source, including the
actual top \(T=N'\), satisfies \eqref{eq:k2-bv-source}. No fixed-width
truncation estimate is substituted.


\subsection{Four directed sieve rows}

Let \(\mathcal P_N=\{p\in\mathbb P:p\nmid2N\}\).
Put
\begin{equation}\label{eq:k2-sifted-source}
 S_b(d;z)=\sum_{d\mid N-e}A_b(e)\,1_{((N-e)/d,\,P_N(z))=1},\qquad
 P_N(z)=\prod_{p<z,\ p\in\mathcal P_N}p .
\end{equation}
There is no prefix--cofactor coprimality condition.
Both source primes exceed \(N^{\mu_*}\) eventually.
Thus, for every supported modulus \(d\),
\begin{equation}\label{eq:k2-local-density}
 \sum_{d\mid N-e}A_b(e)
 =\frac{X_b}{\varphi(d)}+\Delta_{A_b}(d,N),\qquad
 (d,2N)=1,\quad P^+(d)<N^{\mu_*}.
\end{equation}
Here \(P^+(d)\) is the largest prime factor, with \(P^+(1)=1\).
For the squarefree sieve divisors, use the multiplicative density
\begin{equation}\label{eq:k2-squarefree-density}
 g(p)=\frac1{p-1}\quad(p\in\mathcal P_N),\qquad
 g(p^j)=0\quad(j\ge2),\qquad
 g(p)=0\quad(p\notin\mathcal P_N).
\end{equation}
This is the squarefree-divisor density of the linear-sieve application.
That input requires its zero prime-power tail. The
full arithmetic density is \(g(p^j)=1/\varphi(p^j)\) for
\(p\in\mathcal P_N\) and \(j\ge1\). No estimate uses its prime-power
tail. The two densities agree at active primes. These are the only local
factors used by the main-mass calculation.
In Paper~I, Theorem~2.6 notation, this density is
\[
 \omega(p)=\frac p{p-1}\quad(p\in\mathcal P_N),\qquad
 \omega(p)=0\quad(p\notin\mathcal P_N),\qquad
 \omega(p^a)=0\quad(a\ge2).
\]
Thus \(\omega(p)/p=g(p)\), and condition (2) has zero left side. For
\(p\ge3\),
\[
 \left(1-\frac1{p-1}\right)^{-1}
 =\left(1-\frac1p\right)^{-1}
   \left(1+\frac1{p(p-2)}\right).
\]
The product of the second factors over \(p\ge w\) is
\(1+O(1/w)=1+O(1/\log w)\). Hence Mertens' product estimate gives
condition (1) with exponent one and one fixed \(K_{\mathrm{dim}}\ge2\).
We may take \(L_{\omega}=2\). These are the dimension and prime-power
inputs of Iwaniec's linear sieve \cite{Iwaniec1980Remainder}.

If \(N-e\) is prime, it is at least \(NL^{-333}\).
It therefore exceeds \(N^{\mu_*}\) eventually.
Hence
\begin{equation}\label{eq:k2-source-comparison}
 B_b'\le S_b(1;N^{\mu_*}).
\end{equation}

Use the ordered primes
\(N^u\le a_j<\cdots<a_1<N^{\mu_*}\)
in \(\mathcal P_N\).
The exact decomposition of
\eqref{eq:source-buchstab}--\eqref{eq:source-upper-slack} gives
\begin{equation}\label{eq:k2-buchstab-upper}
 B_b\le E_{\mathrm{top}}+T_0-T_1+T_2-R_3+P_3.
\end{equation}
Here
\begin{equation}\label{eq:k2-four-rows}
\begin{aligned}
 T_0&=S_b(1;N^u),\\
 T_j&=\sum_{N^u\le a_j<\cdots<a_1<N^{\mu_*}}
                 S_b(a_1\cdots a_j;N^u)\quad(j=1,2),\\
 R_3&=\sum_{\mathrm{ret}}S_b(a_1a_2a_3;N^{t_3}),\\
 P_3&=\sum_{\mathrm{ret}}\sum_{\substack{N^{t_3}\le r<a_3\\r\in\mathcal P_N}}
                       S_b(a_1a_2a_3r;N^{t_3}).
\end{aligned}
\end{equation}
Every prime in these sums belongs to \(\mathcal P_N\).
The last sum keeps the half-open cells it inherits.
The symbol \(\sum_{\mathrm{ret}}\) restricts by
\begin{equation}\label{eq:k2-retention}
 \Theta_{\mathfrak c}=\inf_{\overline{\mathfrak c}}H,\qquad
 \ell_{\mathfrak c}=\Theta_{\mathfrak c}-\eta_L-t_1-t_2-t_3,
 \qquad \ell_{\mathfrak c}\ge2\rho t_3.
\end{equation}
Here \(H\) is the full eight-branch packing function,
\eqref{eq:classical-packing-ceiling}.
The infimum uses the complete closed ordered cell.
It does not use a quadrature node.

The four directed applications are
\begin{equation}\label{eq:k2-aa-directions}
\begin{array}{c|c|c|c}
 \text{term}&\text{row sign}&\text{sieve direction}&\text{level exponent}\\ \hline
 T_0&+&\text{upper }F&1/2-\eta_L\\
 T_1&-&\text{lower }f&1/2-\eta_L-\rho t_1\\
 T_2&+&\text{upper }F&1/2-\eta_L-\rho(t_1+t_2)\\
 R_3&-&\text{lower }f&\ell_{\mathfrak c}
\end{array}
\end{equation}
The first three cutoffs are \(N^u\).
The fourth cutoff is \(N^{t_3}\).
The canonical functions are Iwaniec's linear-sieve functions
\cite{Iwaniec1980Remainder}.
Extend \(f\) by zero for \(s\le2\).
No shifted beta-sieve profile is substituted.

For each prefix product \(h\), let \(\mathcal A_h\) be the cofactor
sequence indexed by \(n=(N-e)/h\), with multiplicity \(A_b(e)\) when
\(h\mid N-e\). Put \(X_h=X_b/\varphi(h)\). Write \(z_h\) for the
row cutoff, \(D_h\) for its level, and \(s_h=\log D_h/\log z_h\).
Every prefix prime is at least \(z_h\), and each squarefree sieve
modulus \(t\mid P_N(z_h)\) has prime factors below \(z_h\). Thus
\((h,t)=1\), and the local identity \eqref{eq:k2-local-density} gives
the exact remainder identity
\begin{equation}\label{eq:k2-aa-prefix-remainder}
 r(\mathcal A_h,t)=\Delta_{A_b}(ht,N).
\end{equation}
Since \(\omega(t)/t=1/\varphi(t)\) and \((h,t)=1\), the local main
term is \(X_h\omega(t)/t=X_b/\varphi(ht)\), as required by
Theorem~2.6. For rows \(j=0,1,2\), put
\(S=t_1+\cdots+t_j\), with \(S=0\) for row zero, and
\(\theta=1/2-\eta_L\). The clipped exponent cells give
\[
 h\le N^{\rho S},\qquad D_h=N^{\theta-\rho S},\qquad
 S\le2\mu_*.
\]
Hence \(\log D_h/\log N\ge\theta-2\rho\mu_*>0.1092>2u\).
The upper rows satisfy \(N^u\le D_h\). The lower row satisfies
\(N^u\le D_h^{1/2}\). For sufficiently large
\(N\), with \(\sigma\) fixed, \(D_h\ge2\).
Apply Paper~I, Theorem~2.6, (2.30), with density
\(g(p)=1/(p-1)\) on \(\mathcal P_N\), \(g(p^a)=0\) for \(a\ge2\),
and sieve parameter \(\epsilon_s=\sigma\). With
\(\overline E=2C_{\mathrm{AA37}}\sigma\), the resulting directed
inequalities are
\begin{equation}\label{eq:k2-aa-directed-estimates}
\begin{aligned}
 S(\mathcal A_h,z_h)
 &\le X_hV_N(z_h)\{F(s_h)+E^+_{s_h}\}+R_1^+(h)+R^+(h),\\
 S(\mathcal A_h,z_h)
 &\ge X_hV_N(z_h)\{f(s_h)-E^-_{s_h}\}-R_1^-(h)-R^-(h),\\
 |E^{\pm}_{s_h}|&\le\overline E.
\end{aligned}
\end{equation}
Only the upper bound \(|E^{\pm}_{s_h}|\le\overline E\) is used. No sign of
\(E^{\pm}_{s_h}\) is assumed. Since \(X_hV_N(z_h)\ge0\), replacing
\(E^{\pm}_{s_h}\) by \(\overline E\) weakens each inequality.
The actual dimension and prime-power constants, the level lower bound,
and the fixed-\(\sigma\) threshold are in
\eqref{eq:k2-aa-decay}.

The superscript \(\mathrm{AA}\) below marks errors from this
linear-sieve theorem. It does not denote a separate family of profiles.
For row \(j\), sum over its actual prefix set \(\mathcal H_j\). Let
\(\mathcal F_j\) be \(F\) on the upper rows and \(f\) on the lower
rows. The definitions used in the signed equalities are
\begin{equation}\label{eq:k2-aa-row-definitions}
\begin{aligned}
 m_j&=\sum_{h\in\mathcal H_j}X_hV_N(z_h)\mathcal F_j(s_h),\\
 \mathfrak M_j&=\sum_{h\in\mathcal H_j}X_hV_N(z_h),\\
 \mathcal A_j^{\mathrm{AA}}&=\overline E\,\mathfrak M_j,\\
 r_j&=\sum_{h\in\mathcal H_j}\{R_1^{\nu_j}(h)+R^{\nu_j}(h)\},
 \qquad \nu_j=+\text{ on upper rows},\quad \nu_j=-\text{ on lower rows},\\
 z_j&=m_j+\mathcal A_j^{\mathrm{AA}}+r_j-T_j
       \quad\text{on upper rows},\\
 z_j&=T_j-m_j+\mathcal A_j^{\mathrm{AA}}+r_j
       \quad\text{on lower row }j=1,\\
 z_3&=R_3-m_3+\mathcal A_3^{\mathrm{AA}}+r_3.
\end{aligned}
\end{equation}
Summing \eqref{eq:k2-aa-directed-estimates} over prefixes proves
\(z_j\ge0\). Also \(\mathfrak M_j\ge0\), so
\(\mathcal A_j^{\mathrm{AA}}\ge0\). These nonnegativity
statements are proved at source level for the exact signed remainders in
\eqref{eq:k2-aa-prefix-remainder}. No sign of these remainders is assumed.
Write the directed equalities as
\begin{equation}\label{eq:k2-aa-equalities}
\begin{aligned}
 T_0&=m_0+\mathcal A_0^{\mathrm{AA}}+r_0-z_0,&
 T_1&=m_1-\mathcal A_1^{\mathrm{AA}}-r_1+z_1,\\
 T_2&=m_2+\mathcal A_2^{\mathrm{AA}}+r_2-z_2,&
 R_3&=m_3-\mathcal A_3^{\mathrm{AA}}-r_3+z_3,
 \qquad \mathcal A_j^{\mathrm{AA}},z_j\ge0.
\end{aligned}
\end{equation}
\begin{lemma}[Signed row identity]\label{lem:k2-row-identity}
With \(M_K=m_0-m_1+m_2-m_3\),
\[
 \Delta_{\mathrm{arith}}:=T_0-T_1+T_2-R_3-M_K
 =\sum_{j=0}^3\mathcal A_j^{\mathrm{AA}}+\sum_{j=0}^3r_j-\sum_{j=0}^3z_j,
\]
and every \(\mathcal A_j^{\mathrm{AA}}\) occurs with coefficient \(+1\).
\end{lemma}
\begin{proof}
Substitute \eqref{eq:k2-aa-equalities}. The lower rows carry the row
sign \(-\), which turns \(-\mathcal A_j^{\mathrm{AA}}-r_j+z_j\) into
\(+\mathcal A_j^{\mathrm{AA}}+r_j-z_j\).
\end{proof}
For rows \(j=0,1,2\), the total moduli in these remainders fit the
level of \eqref{eq:k2-bv-source}. A tail tuple satisfies
\(\prod_iD_i<D_h\). In the parity ending in a defining cubic condition,
remove two copies of the last label. In the other parity, use the
preceding cubic condition and the ordering of the last two labels. The
empty and one-label cases follow from the tuple definitions.
Each \(p_i\) lies in the interval
\(D_i\le p_i<\min(D_i^\rho,z_h)\), and
\(d<D_h^\sigma\). Thus, with the same \(S,\theta\) as above,
\[
 q=h d\prod_i p_i\le hD_h^{\sigma+\rho}
 \le N^{\rho S+(\sigma+\rho)(\theta-\rho S)}
 =N^{\theta+(\sigma+\tau)(\theta-\rho S)}
 \le N^{\theta+(\sigma+\tau)\theta}
 <N^{1/2-\eta_L/2}.
\]
This power buffer is below \(N^{1/2}L^{-B}\) for every fixed \(B\),
once \(N\) is large enough.

The factors of such a modulus that are at least \(z_h\) are exactly the
ordered prefix primes. They determine \(h\) and its clipped cell.
Factoring the remaining part at \(D_h^{\sigma^2}\) determines the
squarefree smooth divisor and the tail-prime multiset. The number of
tail labels is at most \(J_\sigma=\lceil\sigma^{-2}\rceil\).
Equal-label permutations contribute at most \(J_\sigma!\). Each
AA coefficient has modulus at most one. The number of label vectors
is less than \(\exp(8\sigma^{-3})\). Hence the total coefficient at each
modulus is at most
\(C_{\mathrm{K2},\sigma}=1+\exp(8\sigma^{-3})J_\sigma!\), independent of \(N\).
Apply \eqref{eq:k2-bv-source} with target \(A=4\). Divide by
\(\mathcal N_K=(1-\varepsilon)C(N)N/L^2\); here \(C(N)\ge C_2>0\)
by Paper~II, (1.2). This proves
\begin{equation}\label{eq:k2-low-row-remainders}
 \frac{|r_0|+|r_1|+|r_2|}{\mathcal N_K}\ll_{b,\sigma}L^{-2}.
\end{equation}
The low-row bound uses Fact~\ref{fact:bv-full-support}, the
prime Siegel--Walfisz theorem, the prime number theorem and the Paper~II
convolution inputs. It gives no estimate for the principal
masses or row three.

The \(m_j\) are the directed main terms.
The \(r_j\) are signed arithmetic remainders.
Thus, with \(M_K=m_0-m_1+m_2-m_3\),
\begin{equation}\label{eq:k2-signed-join}
\begin{aligned}
 B_b\le{}&E_{\mathrm{top}}+P_3+M_K
       +\sum_{j=0}^3\mathcal A_j^{\mathrm{AA}}+\sum_{j=0}^3r_j-\sum_{j=0}^3z_j\\
 \le{}&E_{\mathrm{top}}+P_3+M_K
       +\sum_{j=0}^3\mathcal A_j^{\mathrm{AA}}+\sum_{j=0}^3|r_j|.
\end{aligned}
\end{equation}
This step is finite algebra.
The directed source comparisons are proved separately.

The canonical row masses are explicit.
For a prefix cell \(\mathfrak c\), put
\[
 H_j(\mathfrak c)=
 \sum_{\substack{a_1>\cdots>a_j\\a_i\text{ in the cell}}}
       \frac1{\varphi(a_1\cdots a_j)}.
\]
Since the prefix primes exceed every sieve prime, the product over sieve
primes at each row is \(V_N(z_h)\). With the ordered cell sets
\(D_j\) from
\eqref{eq:k2-parent-functional},
\begin{equation}\label{eq:k2-row-main-masses}
\begin{aligned}
 m_0&=X_bV_N(N^u)F(s_0),\\
 m_1&=X_bV_N(N^u)\sum_{\mathfrak c\in D_1}
                 H_1(\mathfrak c)f(s_1(\mathfrak c)),\\
 m_2&=X_bV_N(N^u)\sum_{\mathfrak c\in D_2}
                 H_2(\mathfrak c)F(s_2(\mathfrak c)),\\
 m_3&=X_b\sum_{\mathfrak c\in D_3}\chi_{\mathfrak c}
       V_N(N^{t_3})H_3(\mathfrak c)f(\ell_{\mathfrak c}/t_3).
\end{aligned}
\end{equation}
The masses \(\mathfrak M_j\) are these sums with \(F\) and \(f\)
replaced by \(1\). In particular, the row-three retention indicator
\(\chi_{\mathfrak c}\) appears in both sums.


\subsection{The principal functional}

We display the numerical functional used for \(M_K\).
Set
\begin{equation}\label{eq:k2-canonical-profile}
\begin{aligned}
 v_1(s)&=\min\left\{\frac{1+s}{2},
             \frac{32-34s}{39},\frac{57}{64}-\frac32s\right\},\\
 v(s)&=v_1(s)-u,\qquad \psi(x,y)=x1_{x\ge y},\\
 \Theta_*(x,y,z)&=\max\{v(x),v(y),v(x+y),v(x+y+z),\\
 &\hspace{28mm}\psi(v(x+z),x+2y+z),
                   \psi(v(y+z),2x+y+z)\}.
\end{aligned}
\end{equation}
The subtraction of \(u\) precedes both guards.
This is Paper~II, (4.70)--(4.72)
\cite[(4.70)--(4.72)]{PaperII}.
Put
\begin{equation}\label{eq:k2-g-functional}
\begin{aligned}
 H_{\log}&=\log(\mu_*/u),\\
 R_{\mathrm{up}}&=500(F(30)-1)+500H_{\log}(1-f(30))
                    +250H_{\log}^2(F(30)-1),\\
 J_*&=\int_u^{\mu_*}\frac{dx}{x}
         \int_u^x\frac{dy}{y}
         \int_u^y\frac{1-f((\Theta_*(x,y,z)-x-y-z)/z)}{z^2}\,dz,\\
 G_*&=\frac1{\mu_*}+R_{\mathrm{up}}+J_*.
\end{aligned}
\end{equation}
We keep the true delay-function tail.
Paper~II, (4.73)--(4.75), gives the constant cancellation
\begin{equation}\label{eq:k2-constant-cancellation}
 \frac{1-H_{\log}+H_{\log}^2/2}{u}
 -\int_u^{\mu_*}\frac{dx}{x}\int_u^x\frac{dy}{y}
                         \int_u^y\frac{dz}{z^2}
 =\frac1{\mu_*}.
\end{equation}

The finite-parent comparison also needs its direction.
Fix a finite clipped partition \(\Pi\) of \([u,\mu_*)\).
Its cells have endpoint ratios at most \(\rho\).
Each clipped cell is a product of half-open exponent intervals \([t_i,t_i')\) with \(t_i'\le\rho t_i\). Here \(t_i\) is the lower endpoint, so \(N^{t_i}\le a_i<N^{\rho t_i}\), i.e. \(t_i\le x_i<\rho t_i\). The closure \(\overline{\mathfrak c}\) is used only for infima and lower bounds. The cutoffs \(N^{t_3}\) in \(R_3,P_3\) and the factors \(\rho t_1,\rho(t_1+t_2)\) in \(s_1,s_2\) use this convention.
For a clipped cell \(\mathfrak c\), let
\(w_{\mathfrak c}=\int_{\mathfrak c}\prod_i dx_i/x_i\).
Define
\begin{equation}\label{eq:k2-parent-functional}
\begin{aligned}
 s_0&=(1/2-\eta_L)/u,\\
 s_1(\mathfrak c)&=(1/2-\eta_L-\rho t_1)/u,\\
 s_2(\mathfrak c)&=(1/2-\eta_L-\rho(t_1+t_2))/u,\\
 E_\Pi&=\frac{F(s_0)-1}{u}
 +\sum_{\mathfrak c\in D_1}\frac{w_{\mathfrak c}}u(1-f(s_1(\mathfrak c)))\\
 &\qquad+\sum_{\mathfrak c\in D_2}\frac{w_{\mathfrak c}}u(F(s_2(\mathfrak c))-1),\\
 P_\Pi&=\frac{1-H_{\log}+H_{\log}^2/2}{u}+E_\Pi
 -\sum_{\mathfrak c\in D_3}
    w_{\mathfrak c}1_{\ell_{\mathfrak c}\ge2\rho t_3}
             \frac{f(\ell_{\mathfrak c}/t_3)}{t_3}.
\end{aligned}
\end{equation}
Here \(D_j=\{u\le x_j<\cdots<x_1<\mu_*\}\).
Each cell is intersected with this ordered domain.
Its weight includes the actual ordered multiplicity.
The required comparisons are
\begin{equation}\label{eq:k2-parent-comparison}
 0\le E_\Pi\le R_{\mathrm{up}},\qquad
 P_\Pi\le\frac1{\mu_*}+E_\Pi+J_*\le G_*.
\end{equation}
The last line is an upper comparison.
It asserts no equality between finite sums and integrals.

Compare the six canonical branches of \(\Theta_*\) with their matching
branches among the eight parent branches of the packing function \(H\) in
\eqref{eq:classical-packing-ceiling}. Adding \(u\) to an unguarded canonical
branch gives the matching parent branch. For an active canonical guarded
branch, the parent guard is also active, and adding \(u\) again gives the
matching parent branch. An inactive canonical guarded branch is zero and
cannot maximize the canonical ceiling, since its first unguarded branch is
positive on the ordered domain. The two extra parent branches can only
increase \(H\). Hence
\(H(x,y,z)\ge\Theta_*(x,y,z)+u\) pointwise on that domain. On a closed
anchored cell, each coordinate moves by at most \(\mu_*\tau\). Fix a point
and a maximizing canonical branch there, and use its matching parent branch.
If it is guarded, its guard has reserve \(u-7\mu_*\tau>0\) throughout the
cell. This selected branch changes by at most \((9/2)\mu_*\tau\) across the
cell. Therefore
\begin{equation}\label{eq:k2-parent-reserve}
\begin{aligned}
 \Theta_{\mathfrak c}&\ge\Theta_*(x,y,z)+u-(9/2)\mu_*\tau,\\
 \ell_{\mathfrak c}&\ge\Theta_*(x,y,z)-x-y-z+g,\\
 g&=u-\eta_L-(9/2)\mu_*\tau>2\mu_*\tau.
\end{aligned}
\end{equation}
Monotonicity of \(f(s)\), together with this reserve,
gives the directed kernel comparison
\begin{equation}\label{eq:k2-parent-kernel}
 \chi_{\mathfrak c}\frac{f(\ell_{\mathfrak c}/t_3)}{t_3}
 \ge
 \frac{f((\Theta_*(x,y,z)-x-y-z)/z)}{z}.
\end{equation}
Summing over cells gives the upper comparison in
\eqref{eq:k2-parent-comparison}. It does not identify a finite sum with
the integral.


Put \([x]_+=\max\{x,0\}\).
We need the arithmetic-to-integral estimate
\begin{equation}\label{eq:k2-main-mass}
 \frac{[M_K-\mathcal N_K\,2e^{-\gamma}I(b)G_*]_+}
      {\mathcal N_K}=o_{b,\sigma}(1).
\end{equation}
The quantity \(M_K\) contains only the canonical values of \(F\) and
\(f\). The linear-sieve errors \(E^{\pm}_{s_h}\) and the remainders
\(r_j\) lie outside it. They are paid in
\eqref{eq:k2-aa-mass} and \eqref{eq:k2-remainder-estimates}. Fix
\((b,\sigma)\). Three inputs enter. The first is the two-orientation
arithmetic-mass estimate for the source \eqref{eq:k2-source}. The second
is its four signed canonical AA
main masses \eqref{eq:k2-row-main-masses}. The third is the finite functional
\(P_\Pi\) in \eqref{eq:k2-parent-functional}. For all sufficiently
large even \(N\) they give
\[
 M_K/\mathcal N_K\le2e^{-\gamma}I(b)G_*+o_{b,\sigma}(1).
\]
Here the error tends to zero with \(b,\sigma\) and the cell count fixed,
as \(N\) tends to infinity. This estimate is for the finite functional in
\eqref{eq:k2-parent-functional}. It gives
\eqref{eq:k2-main-mass} directly, so
\(N_{\mathrm{mass}}(b,\sigma)\) can be the threshold of
this fixed-\((b,\sigma)\) statement. This bound controls only the
principal-mass term. It does not control the AA remainders, the optional
numerical certificate, the complete exterior comparison or
the endpoint terms. No sharper arithmetic-mass estimate is used here.


For the exterior comparison we use the common $G_0$ coefficient
\begin{equation}\label{eq:k2-certificate}
\begin{gathered}
 G_0=\frac{6088440363}{10^9},\qquad
 G_*\le G_0,\\
 2e^{-\gamma}G_0\le c_{K2,G0}
 :=\frac{6836825163875517756286846321650754372931094689}{10^{45}}.
\end{gathered}
\end{equation}
For \(G_*\le G_0\), use the tight endpoint
\(J_*\le0.967120\) of Paper~II, (A.10a). It holds for the exact profile
\eqref{eq:k2-g-functional} \cite[(4.73), (4.75), (A.10a)]{PaperII}.
Its weaker displayed endpoint is not used.
The Taylor-model delay certificate is used for this coefficient comparison
and is not an input to the shifted replay itself. It gives the explicit bounds
\[
 0\le F(30)-1<3\cdot10^{-49},\qquad
 0\le1-f(30)<3\cdot10^{-49}.
\]
Since \(0<H_{\log}=\log(3125/32)<5\),
\eqref{eq:k2-g-functional} then gives
\[
 R_{\mathrm{up}}< (500+500\cdot5+250\cdot25)3\cdot10^{-49}
 <10^{-44},\qquad
 G_*<\frac{128}{25}+\frac{967120}{10^6}+10^{-44}
 <\frac{6088440363}{10^9}=G_0.
\]
The finite delay calculation and this rational budget do not replace the
source-to-profile comparison or its fixed-parameter hypotheses.
For the scalar inequality, define
\begin{equation}\label{eq:k2-gamma-enclosure}
\begin{gathered}
 \gamma_-:=\frac{57721566490153286060651209008240243104215933593992}{10^{50}},\\
 \gamma_+:=\frac{57721566490153286060651209008240243104215933593993}{10^{50}},\\
 \gamma_-<\gamma<\gamma_+,\\
 c_{K2,G0}-2G_0\sum_{j=0}^{60}\frac{(-\gamma_-)^j}{j!}
       >\frac{75307}{10^{50}}>0.
\end{gathered}
\end{equation}
The rational enclosure follows from Euler--Maclaurin for \(H_n\) at
\(n=10^5\), through \(B_{16}\). The next-term remainder is bounded
by \(|B_{18}|/(18n^{18})\). Both \(\log n\) and \(H_n\) are enclosed
outwards by the certified exact-input record. For \(0<x<1\), the even Taylor
partial sum of \(e^{-x}\) is an upper bound. Thus the exact rational
comparison in \eqref{eq:k2-gamma-enclosure} proves
\(2e^{-\gamma}G_0<c_{K2,G0}\).


The switched comparison uses the same coefficient.
At \(b=2999993/7000000\), the common final-margin inequality
\eqref{eq:bill-margin-simple} has slack
\(0.000228954311267289\ldots>1/5000\) on the stated route \(g_6:=0\).
Both comparisons use the same $G_0$ bound. Scalar positivity does not
establish the required source estimates.

\subsection{The complete error bound}

The analytic hypotheses are stated below.
The sums are over the original applications.
The four principal linear-sieve errors satisfy
\begin{equation}\label{eq:k2-aa-mass}
\begin{gathered}
 \frac{\sum_{j=0}^3 \mathcal A_j^{\mathrm{AA}}}{\mathcal N_K}
 \le A_K(b)C_{\mathrm{AA}}\sigma,\\
 A_K(b)=2e^{-\gamma}I(b)W_{\mathrm{AA}},\qquad
 W_{\mathrm{AA}}=500\left(1+H_{\log}
                    +\frac{H_{\log}^2}2+\frac{H_{\log}^3}6\right).
\end{gathered}
\end{equation}
Let \(C_{\mathrm{AA37}}\) be the absolute constant in Paper~I,
Theorem~2.6, (2.31), and define
\begin{equation}\label{eq:k2-aa-constant}
 C_{\mathrm{AA}}=4C_{\mathrm{AA37}}.
\end{equation}
The coefficient four is the factor two in
\eqref{eq:k2-aa-density-mass} times the factor two in the common bound
\(|E_s^{\pm}|\le 2C_{\mathrm{AA37}}\sigma\). The sum over the four rows is
already included in \(W_{\mathrm{AA}}\). It is independent of \(\sigma,N\).
It is unrelated to the global tensor constant in
\eqref{eq:ra-global-tensor}. All four actual levels satisfy
\(D\ge N^{1/250}\). We use (2.31) with the actual fixed dimension constant
\(K_{\mathrm{dim}}\). For the squarefree density in
\eqref{eq:k2-squarefree-density}, the left side of (2.27) is zero. Hence
the actual prime-power constant is \(L_{\omega}=2\). This value is used
in the exponential below.
\begin{equation}\label{eq:k2-aa-decay}
 |E_s^{\pm}|\le C_{\mathrm{AA37}}
 \{\sigma+\sigma^{-8}e^{K_{\mathrm{dim}}+2}(\log D)^{-1/3}\}
 \le2C_{\mathrm{AA37}}\sigma
\end{equation}
once \(L\ge250\exp(3(K_{\mathrm{dim}}+2))\sigma^{-27}\).
Name the corresponding even threshold
\[
 N_{\mathrm{AA,dec}}(\sigma)
 :=2\left\lceil\frac12
 \exp\!\left(250\exp(3(K_{\mathrm{dim}}+2))\sigma^{-27}\right)
 \right\rceil.
\]
It is chosen at fixed \(\sigma\). Every
\(N\ge N_{\mathrm{AA,dec}}(\sigma)\) satisfies the displayed decay bound.
The four ordered harmonic row masses, before the common Euler and source
factor, are bounded by
\begin{equation}\label{eq:k2-aa-row-masses}
\begin{gathered}
 \frac1u,\qquad \frac{H_{\log}}u,\qquad
 \frac{H_{\log}^2}{2u},\\
 \text{row-3 mass}\le
 \rho\int_u^{\mu_*}\frac{dx}{x}\int_u^x\frac{dy}{y}
        \int_u^y\frac{dz}{z^2}
 \le\frac{\rho H_{\log}^2}{2u}
 <\frac{H_{\log}^3}{6u}.
\end{gathered}
\end{equation}
For the last inequality use \(3\rho<H_{\log}=\log(3125/32)\).
The factor \(\rho\) allows the anchored-bin endpoint enlargement. The first three
identities are the volumes of ordered logarithmic simplexes. For the
fourth, bound \(\int_u^y dz/z^2\le1/u\) and integrate over the ordered
\((x,y)\) simplex. The sum of the four displayed upper envelopes is \(W_{\mathrm{AA}}\),
the four-row envelope used here. It is distinct from Paper~II's
coarser global sieve-error bound.
Let \(H=H_{\log}=\log(3125/32)\). PNT and Mertens' theorems, uniformly
over the finitely many fixed-\(\sigma\) cells, give
\[
 V_N(N^a)=\frac{2e^{-\gamma}C(N)}{aL}(1+o(1)),\qquad
 H_j(\mathfrak c)=w_{\mathfrak c}+o(1).
\]
For the retained row, \(1/t_3\le\rho/x_3\) on each cell. Dropping
the retention restriction only increases its mass. Consequently
\begin{equation}\label{eq:k2-aa-density-derived}
\begin{aligned}
 \frac{\sum_{j=0}^3\mathfrak M_j}{\mathcal N_K}
 &\le2e^{-\gamma}I(b)
 \left\{\frac1u+\frac{H}{u}+\frac{H^2}{2u}+\rho J_4\right\}+o(1),\\
 J_4&=\frac{H^2/2-H+1}{u}-\frac1{\mu_*}
       \le\frac{H^2}{2u}.
\end{aligned}
\end{equation}
The bracket is smaller than
\(W_{\mathrm{AA}}=u^{-1}(1+H+H^2/2+H^3/6)\), with relative gap greater
than \(0.2842\) for \(0<\sigma\le10^{-10}\). The coarser bound
\(\rho J_4\le\rho H^2/(2u)\) still leaves relative gap greater than \(0.1723\).
Consequently, for each fixed \(b,\sigma\), there is a threshold
\(N_{\mathrm{AA,dens}}(b,\sigma)\) such that, for every even
\(N\ge N_{\mathrm{AA,dens}}(b,\sigma)\),
\(\sum_j\mathfrak M_j\le\mathcal N_K\,2e^{-\gamma}I(b)W_{\mathrm{AA}}\).
Define
\[
 N_{\mathrm{AA}}(b,\sigma)
 :=\max\{N_{\mathrm{AA,dens}}(b,\sigma),N_{\mathrm{AA,dec}}(\sigma)\}.
\]
Thus both the density estimate and \eqref{eq:k2-aa-decay} hold above
\(N_{\mathrm{AA}}(b,\sigma)\), as required for \eqref{eq:k2-aa-mass}.
The maxima defining \(T_K\) and \(T_r\) below include this combined threshold.
The displayed source-density bound below is weaker by a factor of two.
Let \(\mathfrak M_j\) be the nonnegative principal mass multiplying the
relative sieve error in row \(j\). The resulting bound is
\begin{equation}\label{eq:k2-aa-density-mass}
 \sum_{j=0}^3\mathfrak M_j\le
 2\mathcal N_K\,2e^{-\gamma}I(b)W_{\mathrm{AA}}.
\end{equation}
Multiplying it by
\(2C_{\mathrm{AA37}}\sigma\) gives \eqref{eq:k2-aa-mass} with
\eqref{eq:k2-aa-constant}. The Euler factors and the actual-source mass
identification are included above. The secondary decay is paid once.


The positive peel is bounded by the two-form argument of Paper~II,
\cite[(4.46)--(4.56)]{PaperII}, applied to the original source.
Put \(\mathcal B_m=\{e:\ m\mid N-e\}\) with multiplicity \(A_b(e)\).
For \(m=a_1a_2a_3r\) in \(P_3\), put \(E=N^{t_3}\) and \(w=N^u\le E\).
Write \(A_b=\sum_{h=1}^21_{\{A_b\ge h\}}\).
Every \(e\) in the support is a product of two primes at least \(N^b\).
Hence \(e\) is odd and \(w\)-rough. Write \(N-e=mk\). Then \(k\) is odd. A
prime \(p<w\) dividing \(k\) and \(N\) would divide \(e\), which is impossible.
Hence both forms \(k\) and \(N-mk\) are free of all primes below \(w\).
Therefore \(S_b(m;E)\) is at most the two-form Selberg majorant of
\cite[(4.48)--(4.53)]{PaperII} on \(1\le k\le N/m\), once for each layer.
All prime factors of \(m\) are at least \(E\ge w\) and lie in
\(\mathcal P_N\). Hence \((m,N)=1\), and \(\nu_N(p)\) is as in
\cite[(4.48)]{PaperII}. Retention gives, as in \cite[(4.55)]{PaperII},
\begin{equation}\label{eq:k2-peel-ranges}
 \log_Nm\le\rho(t_1+t_2+2t_3)\le\rho(\Theta_{\mathfrak c}-\eta_L)
 <\tfrac{153}{256}+2\tau,\qquad
 \log_N(w^2m)<1-0.0116,\qquad \log_Nw=u\ge\tfrac2{1001}.
\end{equation}
Consequently \cite[(4.47), (4.53), (4.54)]{PaperII} give, eventually and
uniformly in \(m\),
\begin{equation}\label{eq:k2-peel-inputs}
 S_b(m;E)\le2\cdot\frac{4.002}{u^2}\cdot\frac{U_C}{m}
 \le\frac{2008008\,U_C}{m},\qquad
 \max_hM_h\le2\tau,\qquad \sum_hM_h\le\frac{459}{100},
\end{equation}
since one layer is at most \((N/m)/\mathcal G_N(w)+O(w^2\log^5w)\) with
\(\mathcal G_N(w)=(1+o(1))u^2L^2/(4C(N))\), and
\(4.002/u^2=1000500<1004004=2\cdot502002\). The bin sums are
\cite[(4.56)]{PaperII}: \(\log(\mu_*/u)=\log(3125/32)<4.59\), and
\(a_3,r\) lie in one parent bin. Therefore
\begin{equation}\label{eq:k2-peel}
 \frac{P_3}{U_C}\le2008008\cdot2\left(\frac{459}{100}\right)^3\tau
 =\frac{24272444031579}{62500}\tau<388400000\,\sigma^9,
\end{equation}
and \(P_3/\mathcal N_K\le P_K\sigma^9\) with
\(P_K=388400000/(1-\varepsilon)=3884000000000/9999\), for
\(N\ge N_{\rm peel}(\sigma)\). The coefficient \(2008008\) includes the two
layers. No orientation factor is added.

The remaining estimates are
\begin{equation}\label{eq:k2-remainder-estimates}
\begin{aligned}
 \frac{|r_0|+|r_1|+|r_2|}{\mathcal N_K}&\ll_{b,\sigma}L^{-2},\\
 \frac{|r_3|}{\mathcal N_K}&\ll_{b,\sigma}
 N^{-\ell_0\sigma^2}L^3+N^{-b/4}L^7+L^{-260}+L^{-120}
 +L^{-2}+N^{-c_\sigma}L^{4024+C}
 =O_{b,\sigma}(L^{-2}).
\end{aligned}
\end{equation}
The first line is \eqref{eq:k2-low-row-remainders}. The second is
Proposition~\ref{prop:k2-row3} below. Its displayed terms come, in order,
from the squareful mask, proper powers, logarithm freezing and the
boundary. The last two are the combined logarithmic Type-II and Type-I-2
bound and the Type-II power-saving term.


\subsection{The row-three remainder}\label{subsec:k2-row3}

We bound the complete third-row remainder of the actual source
\eqref{eq:k2-source} by Paper~I, Theorem~3.2. Every other contribution is
bounded separately. Fix \(b,\sigma\). Write \(\Xi^{\sharp}_3\) for the
retained cells \(\mathfrak c\in D_3\) with \(\chi_{\mathfrak c}=1\). Its
size is \(O_\sigma(1)\). For \(\mathfrak c\in\Xi^{\sharp}_3\) put
\(D_{\mathfrak c}=N^{\ell_{\mathfrak c}}\), \(z_{\mathfrak c}=N^{t_3}\) and
\(u_{\mathfrak c}=D_{\mathfrak c}^{\sigma^2}\). Let \(I_i\) be the actual
half-open exponent interval of coordinate \(i\) of the clipped cell. Then
\(\mathfrak c\subset I_1\times I_2\times I_3\). Put
\(\mathcal P_i=\{a\in\mathcal P_N:\log_Na\in I_i\}\). For every
\(a_i\in\mathcal P_i\), \(N^{t_i}\le a_i<N^{\rho t_i}\) and \(a_i<N^{\mu_*}\),
also for a clipped interval. Write \(t_i'\le\min(\rho t_i,\mu_*)\) for the
upper endpoint of \(I_i\). On \(\overline{\mathfrak c}\) the actual
coordinates satisfy \(x_1\ge x_2\ge x_3\ge u\). Two of the intervals are
equal or disjoint.

\paragraph{Exact occurrence map.}
By \eqref{eq:k2-aa-prefix-remainder} and Paper~I, (2.30), (2.32)--(2.33),
applied to \(\mathcal A_h\) at level \(D_{\mathfrak c}\), cutoff
\(z_{\mathfrak c}\le D_{\mathfrak c}^{1/2}\) and \(\epsilon_s=\sigma\),
\begin{equation}\label{eq:k2-r3-occurrence}
 r_3=\sum_{\mathfrak c\in\Xi^{\sharp}_3}\ \sum_{k\in\mathcal K_{\mathfrak c}}
       \sum_qO_{\mathfrak c,k}(q)\,\Delta_{A_b}(q,N),
 \qquad \mathcal K_{\mathfrak c}=\{\mathrm{R1}\}\cup\mathcal D^-(D_{\mathfrak c}),
\end{equation}
\[
 O_{\mathfrak c,k}(q)=\sum_{\substack{a_1>a_2>a_3\\a_i\in\mathcal P_i}}
 \sum_{\substack{d<D_{\mathfrak c}^\sigma\\d\mid P_N(u_{\mathfrak c})}}
 \lambda^{(k)}_d
 \sum_{\substack{(p_j)_{j\le\rho_k}\\ p_j\in\mathcal Q_{k,j}\\
 p_i\ne p_j\ ((i,j)\in I(\mathbf D))}}
 1_{q=a_1a_2a_3d\,p_1\cdots p_{\rho_k}} .
\]
Here \(\lambda^{(\mathrm{R1})}_d=\phi_d^-\) with no tail
(\(\rho_{\mathrm{R1}}=0\)). For \(k=\mathbf D\),
\(\lambda^{(k)}_d=\lambda_d^-(\mathbf D)\) and
\(\mathcal Q_{k,j}=\mathcal P_N\cap[D_j,\min(D_j^{1+\tau},z_{\mathfrak c}))\).
All coefficients have modulus at most one and depend only on
\((\mathfrak c,k)\). The identity is exact at every supported modulus.
Let \(M_{\mathfrak c}\in\{1,2,6\}\) be the number of orderings of the
equal intervals among \(I_1,I_2,I_3\). Let \(\Gamma_{\mathfrak c,k}\)
be \(M_{\mathfrak c}^{-1}\) times the same sum, with \(a_i\in\mathcal P_i\)
independent and no tail distinctness. Then
\begin{equation}\label{eq:k2-native13-coefficient}
 O_{\mathfrak c,k}=\Gamma_{\mathfrak c,k}-\Xi_{T,\mathfrak c,k}
 -\Xi_{P,\mathfrak c,k},
\end{equation}
where \(\Xi_T\) collects unrestricted terms with
\(p_i=p_j\) for some forbidden pair \((i,j)\in I(\mathbf D)\).
\(\Xi_P\) collects the remaining unrestricted terms with
\(a_i=a_j\) for an equal-interval pair \(I_i=I_j\). A prefix prime is
at least \(z_{\mathfrak c}\), and every tail prime is below it. The primes of
\(d\) are below \(u_{\mathfrak c}\le D_j\). Unequal grid intervals are
disjoint. Hence every modulus in the support of \(\Xi_T\) or \(\Xi_P\) is
divisible by \(p^2\) for a prime \(p\ge u_{\mathfrak c}\ge u_0\). Here
\(u_0=N^{\ell_0\sigma^2}\).

\paragraph{Support, cutoffs and constants.}
Every supported modulus is coprime to \(2N\). The cells are clipped
to \([u,\mu_*)\), and \(v_{\rm pk}(x)=(1+x)/2\) on \([0,\mu_*]\). Thus
\(\Theta_{\mathfrak c}\ge(1+t_1)/2\). The actual coordinate bound
\(x_1\ge x_3\ge u\) on \(\overline{\mathfrak c}\) also gives
\(\Theta_{\mathfrak c}\ge(1+u)/2\). Hence
\begin{equation}\label{eq:k2-r3-ell}
 \ell_{\mathfrak c}\ge\max\{A_3-t_3,2\rho t_3\}\ge\tfrac23A_3
 =\tfrac{33109}{240000},\qquad A_3=\tfrac12-\eta_L-\tfrac32\mu_*,
\end{equation}
and \(\ell_{\mathfrak c}-\ell_0\ge231463/24720000\). By the cubic tail
conditions, \(\sum_js_j<\ell_{\mathfrak c}\) for
\(s_j=\log_ND_j\). Thus every supported modulus satisfies
\begin{equation}\label{eq:k2-r3-support}
 q<N^{\rho(t_1+t_2+t_3)+\sigma\ell_{\mathfrak c}+\rho\ell_{\mathfrak c}}
 \le N^{e_Q},\qquad
 e_Q=(\rho+\sigma)(H_{\max}-\eta_L)<\tfrac35,\qquad H_{\max}=\tfrac{153}{256}.
\end{equation}
Since the cell ceiling gives \(0<\ell_{\mathfrak c}\le
H_{\max}=153/256<1\) and \(0<\sigma\le10^{-10}\),
\(\ell_{\mathfrak c}\sigma^2<\sigma^2<u=1/500\). Thus
\(u_{\mathfrak c}=N^{\ell_{\mathfrak c}\sigma^2}<N^u\). The cutoff
comparison uses only this exponent inequality.

Every labelled prime is at least \(u_{\mathfrak c}\), so the number of
labels is below \(3/(5\ell_0\sigma^2)\). With
\(M_\sigma=\lceil3/(5\ell_0\sigma^2)\rceil\) and the vector count
\(\exp(8\sigma^{-3})\), one has
\(|\Gamma_{\mathfrak c,k}|,|\Xi_{T}|,|\Xi_P|\le M_\sigma!\), and
\(\sum_{\mathfrak c,k}
 \bigl(|\Gamma_{\mathfrak c,k}|+|\Xi_T|+|\Xi_P|\bigr)
 \le C'_{\mathrm{K2},\sigma}\).

\paragraph{The squareful terms.}
Since \(|\Delta_{A_b}(q,N)|\le2(N/q+1)+2N/\varphi(q)\), Paper~I, (B.5),
with \(Y=u_0\) and \(Q_{\rm supp}=N^{e_Q}\) gives
\begin{equation}\label{eq:k2-r3-squareful}
 \sum_{\mathfrak c,k}\sum_q|\Xi_{T,\mathfrak c,k}(q)
 +\Xi_{P,\mathfrak c,k}(q)|\,|\Delta_{A_b}(q,N)|
 \ll_\sigma N^{1-\ell_0\sigma^2}L .
\end{equation}

\paragraph{Prime detection and the two Heath--Brown expansions.}
Put \(z_1=N^{1/(\log L)^3}\), so \(z_1\ge z_0(X)\) for every occurring
box scale \(\varepsilon N/2\le X\le N\).
Both source primes exceed \(z_1\). Apply Paper~I, (B.7), to \(1_{\mathbb P}(m)\)
and to \(1_{\mathbb P}(p')\). Apply Paper~I, (B.9), with roughness \(z_1\), to
\(\Lambda_{z_1}(m)\) with cap \(U_1=2N^{1/6}\) and to \(\Lambda_{z_1}(p')\)
with cap \(U_2=2N^{(1-b)/3}\). The identities apply since
\(m\le N^{1/2}<U_1^3\) and \(p'\le N^{1-b}<U_2^3\). For patterns
\(j,j'\in\{1,2,3\}\) the tested index is a labelled tuple
\(\mathbf v=(\mathbf v^{(1)},\mathbf v^{(2)})\in\mathcal R_{z_1}^{2j+2j'}\),
at most twelve coordinates, with
\[
 m=\textstyle\prod\mathbf v^{(1)},\quad p'=\prod\mathbf v^{(2)},\quad
 e=mp',\quad N^b\le m\le N^{1/2},\quad \varepsilon N<e\le N',
\]
the Möbius coordinates of \(\mathbf v^{(i)}\) at most \(U_i\), and coefficient
\(c_j(\mathbf v^{(1)})c_{j'}(\mathbf v^{(2)})/(\log m\log p')\).
Here \(c_j\) is the coefficient after (B.10). For a fixed product \(e\), put
\begin{equation}\label{eq:k2-pointmass-discrepancy}
 \mathcal S_e(q;N):=\Delta_{\delta_e}(q,N)
 =1_{e\equiv N\ (q)}-\frac{1_{(e,q)=1}}{\varphi(q)}.
\end{equation}
Thus the Heath--Brown tuple map uses the point-mass discrepancy
\(\mathcal S_e(q;N)\). The exact weighted point-mass expansion is
\[
 \Delta_{A_b}(q,N)=\sum_e A_b(e)\mathcal S_e(q;N).
\]
The weight \(A_b(e)\) keeps both eligible ordered orientations, as in
\eqref{eq:k2-source}. A prime square is counted once. The proper-power remainder
\(\mathcal E_{\rm pp}\) consists of the pairs in which \(m\) or \(p'\) is a
proper prime power. For every fixed \((\mathfrak c,k)\),
\begin{equation}\label{eq:k2-r3-hb-identity}
 \sum_q\Gamma_{\mathfrak c,k}(q)\Delta_{A_b}(q,N)
 =\mathcal E_{\rm pp}
 +\sum_{j,j'}\sum_{\mathbf v}
   \frac{c_j(\mathbf v^{(1)})c_{j'}(\mathbf v^{(2)})}{\log m\log p'}
   \sum_q\Gamma_{\mathfrak c,k}(q)\mathcal S_{e}(q;N),
\end{equation}
exactly. Each tuple keeps the natural discrepancy of its own product.
The per-cell coefficient bound is the factorial bound above; the sum over
cells uses the separate aggregate bound defined there.
Since \(e<N\), \(|c_jc_{j'}|/(\log m\log p')\le9\) and
\(\sum_q|\Gamma(q)||\mathcal S_e(q;N)|\le C_{\mathrm{K2},\sigma} F(e)\) with
\(F(e)=\tau(N-e)+2(1+L)\). The tuple multiplicity is at most
\(\tau_{12}(e)\). The divisor bounds \(\tau_{12}^2\le\tau_{144}\),
\(\tau^4\le\tau_{16}\) and Cauchy--Schwarz give
\begin{equation}\label{eq:k2-r3-moments}
 \sum_{e<N}\tau_{12}(e)F(e)\ll NL^{73},\qquad
 \sum_{e<N}\tau_{12}(e)F(e)^2\ll NL^{79}.
\end{equation}
There are \(\ll N^{1-b/2}L\) proper-power pairs, and
\(\sum_{e<N}\tau(e)F(e)^2\ll NL^9\). Therefore
\begin{equation}\label{eq:k2-r3-pp}
 |\mathcal E_{\rm pp}|\ll_\sigma(N^{1-b/2}L)^{1/2}(NL^9)^{1/2}
 =N^{1-b/4}L^5 .
\end{equation}

\paragraph{Boxes, logarithm freezing and faces (H6).}
Put \(\delta=L^{-334}\). Use the equal-ratio mesh of Paper~I, (B.12),
with \(\delta_X\) replaced by \(\delta\), on every ambient coordinate axis.
Keep \(\{1\}\) as an atom. A box \(\mathcal B\) is a Cartesian product
of cells. Its lower endpoints are \(V_i\), and \(X=\prod_iV_i\). Let
\(\lambda_1,\lambda_2\) be the sums of \(\log V_i\) over the two sides.
Then \(\varepsilon N/2\le X\le N\) on every box meeting the source, so
\(K_{\rm scale}=1\). Replacing \(1/(\log m\log p')\) by
\(1/(\lambda_1\lambda_2)\) changes the weight by \(O(L^{-335})\), since
\(\lambda_1\ge bL/2\). By \eqref{eq:k2-r3-moments},
\begin{equation}\label{eq:k2-r3-log}
 |E_{\rm LOG}|\ll_\sigma NL^{72-334}.
\end{equation}
The faces are \(m=N^b\), \(m=N^{1/2}\), \(e=\varepsilon N\), \(e=N'\), and
the cap faces \(v_i=U_1\), \(v_i=U_2\) on Möbius coordinates. All are
singletons. On each, an active coordinate is at least \(N^{1/20}\).
Paper~I, Lemma~(B.14), with \(s\le12\), \(K_\gamma=1\) and
\(\eta_{\rm dist}\le12\delta\), gives at most \(\ll NL^{11-334}\) boundary
tuples. Cauchy--Schwarz with \eqref{eq:k2-r3-moments} gives
\begin{equation}\label{eq:k2-r3-bdry}
 |E_{\rm bdry}|\ll_\sigma NL^{45-334/2}.
\end{equation}
Thus (H6) holds with \(c_{\log}=72\), \(c_R=45\). On a nonboundary box every
source cut and cap is constant. The coefficient is the outer scalar
\(c_{\mathcal B}=(-1)^{j+j'}\binom3j\binom3{j'}(\log X)^2/(\lambda_1\lambda_2)\),
times
\(w(\mathbf v)=\prod_iw_i(v_i)\) with
\(w_i\in\{\mu_{z_1},1_{\mathcal R_{z_1}},(\log v/\log X)1_{\mathcal R_{z_1}}\}\).
Each side carries exactly one logarithmic coordinate, and
\((\log X)^2/(\lambda_1\lambda_2)
\le 1/(b(1-b))+o(1)\le4.5+o(1)\), hence the scalar norm bound
\(\lVert c_{\mathcal B}\rVert_\infty\le45\).
Each coordinate meets at most \(2L^{335}\) cells, so
\begin{equation}\label{eq:k2-total-coefficient-norm}
 \sum_{\mathfrak c,k}\sum_{j,j'}\sum_{\mathcal B}|c_{\mathcal B}|
 \ll_\sigma L^{4020}\le L^{4022},\qquad K_0=4022\le24C_{\rm box}+26=8042.
\end{equation}

\paragraph{Classification.}
Let \(e_i=\log_NV_i\), \(\mathcal W=[60\sigma,1/3+14\sigma]\), and
\(\kappa=(\log2+12\log(1+\delta))/L\le\sigma/202\). A box is Type~II if
some \(e_i\in\mathcal W\). Select the largest such \(e_i\), breaking ties
by the fixed coordinate order. Otherwise every Möbius coordinate is small, since
\(1/6+\kappa\) and \((1-b)/3+\kappa\) are below \(1/3+14\sigma\). On the
\(m\)-side, \(\sum e_i\ge b-\kappa>360\sigma\). On the \(p'\)-side,
\(\sum e_i\ge1/2-\kappa-\log(1/\varepsilon)/L>360\sigma\). These are the only side bounds
needed to rule out an all-small side. Together with the global Heath--Brown
sum bound \(\sum e_i<1+\kappa\), the threshold \(1/3+14\sigma\) permits
at most two large variables. Hence each side has exactly one large smooth
coordinate. Every other box has exactly two large smooth coordinates \(N_1\ge N_2\ge X^{1/3}\). It has at most ten small
coordinates, of product \(M\), with
\begin{equation}\label{eq:k2-r3-typeI2}
 \log_XM\le\frac{(600+1/202)\sigma}{1-\sigma/101}<6.02\cdot10^{-8},
 \qquad N_1N_2M\le2^{30}X,\qquad q<X^{0.5977}.
\end{equation}
There are no Type-I-1 boxes.

\paragraph{Type-II boxes: Paper~I, (H1)--(H5) and LEN.}
Fix \((\mathfrak c,k)\) and a Type-II box. Let its selected lower endpoint
be \(T\), put \(\xi=\log_NT\), and let the other lower endpoints have product \(M=X/T\).
\begin{description}
\item[H1.] \(C_{\mathcal B}=\Gamma_{\mathfrak c,k}\) has the form (3.4). Here
\(J\) is the three prefix labels and the tail labels. The prime sets are
\(\mathcal P_i\) and \(\mathcal Q_{k,j}\), \(u_L=u_{\mathfrak c}\ge u_0\),
and \(d\mid P_N(u_{\mathfrak c})\) is squarefree. Repeated primes are kept.
We have \(q\le X^{0.6}\) by \eqref{eq:k2-r3-support} once
\((3/5-e_Q)L\ge(3/5)\log(2/\varepsilon)\). Also \((q,2N)=1\). The
selector \(\chi_{\mathfrak c}\) is constant on the block. The removed
corrections are \(\Xi_T,\Xi_P\) and \(\mathcal E_{\rm pp}\). There is no
comparator correction, since the comparator is natural on every tuple.
\item[H2.] Take the upper exponents \(t_1'\ge t_2'\ge t_3'\) and
\(\min(\rho s_j,t_3)\). For \(d_*=\rho(\Theta_{\mathfrak c}-\eta_L)\), the
merged list satisfies \(\mathrm{WF}(d_*)\). The reserves are
\(\rho(1/2-\eta_L-3\mu_*/2)=33109\rho/160000\),
\(\rho(1/2-\eta_L-5\mu_*/2)=1859\rho/160000\), \(\rho(\ell_{\mathfrak c}-t_3)
\ge\rho(2\rho-1)t_3\). For a tail rank, the reserve is \(\rho\) times the strict
inequality \(\sum_{i<j}s_i+2s_j<\ell_{\mathfrak c}\) from the
\(\mathcal D^-\) conditions. Put \(h=d_*+60\sigma\), \(\delta_{\rm pk}=30\sigma\),
\(g=\xi-30\sigma\). Then \(h-1/2\ge9\cdot10^{-4}\),
\(\Theta_{\mathfrak c}-h>0.0000999939\), \(h\le H_{\max}\),
\[
30\sigma\le g\le1/3-16\sigma,
\qquad \delta_{\rm pk}=30\sigma,
\qquad g\le1/3-\delta_{\rm pk}/2.
\]
The upper requirement has reserve at least \(\sigma\), since
\((1/3-15\sigma)-(1/3-16\sigma)=\sigma\). Also
\(h\le\Theta_{\mathfrak c}\le\widetilde H(t_1',t_2',t_3')\) because
\((t_1',t_2',t_3')\in\overline{\mathfrak c}\) and \(t_1'\le\mu_*\).
Paper~I, Theorem~4.1, gives \(J=J_1\sqcup J_2\sqcup J_3\). Put \(d\) in
slot one, with \(\gamma_s\) as in Paper~I, (3.5). The tuple map gives
\(\Gamma_{\mathfrak c,k}=M_{\mathfrak c}^{-1}C_{\rm slot}(\gamma_1*\gamma_2*\gamma_3)\),
\(C_{\rm slot}\le M_\sigma!\), \(|\gamma_s|\le1\). Slots two and three are
\(u_0\)-rough, and empty slots carry the atom \(1_{v=1}\). The scalar
\(M_{\mathfrak c}^{-1}\) is an outer coefficient.
\item[H3.] Let \(Q,R,S\) be the scales of \(\gamma_1,\gamma_3,\gamma_2\). The
loads (4.3), the slot-one load \(\sigma\ell_{\mathfrak c}\le\sigma H_{\max}\)
of \(d\), and \(A_0=X^{1+\sigma/20}\) give reductions
\((7527/256,\,60,\,114421/1280,\,26247/256)\sigma\). They are, in turn, in
\(Q/T\), \(T^2R^2S\), \(A_0^{7/32}TQR^2S^5\), \((TR/S)^{7/16}TQR^2S^5\).
The required reductions are \(29,8,15,15\), with losses at most \(2\omega\),
\(\omega=\log_N(N/X)\le\sigma/202\). We also have \(\log_N(QRS)\le e_Q\), so
\(QRS<M_IX^{-89\sigma}\) with reserve \(>0.0691\). Next, \(N_I\ge QX^{\sigma/3}\)
with reserve \(29.069\sigma\), and \(M_I\ge X^{1/2}\). Thus (3.6)--(3.7)
hold. The selected coefficient is
\(\alpha_n=w_{\rm sel}(n)1_{I_T}(n)\in\{\mu(n),1,\log n/\log X\}1_{P^-(n)\ge z_1}1_{I_T}(n)\),
with \(\zeta=z_1\ge z_0(X)\) and \(I_T=(T,(1+\Delta)T]\). The complementary
coefficient \(\beta_{m^\sharp}\), the product of the other at most eleven
weights, satisfies \(|\beta_{m^\sharp}|\le\tau_{11}(m^\sharp)\le\tau_{23}(m^\sharp)\).
Here \(\tau_{23}=d_{23}\) is the 23-fold divisor function, with
\(d_{23}(n)=\sum_{n=n_1\cdots n_{23}}1\).
As in Paper~I, (B.22)--(B.23), with at most eleven complementary cells,
\(N_I\ge T/(6L^{334})\), \(M_I=M((1+\Delta)^{r}-1)\asymp ML^{-334}\),
\(m_0=M(1+(1+\Delta)^r)/2\le2X/T\), so \(C_{\rm ctr}=2\). On the doubled
support, \(mn\le X(1+\Delta)(1+\tfrac32((1+\Delta)^{11}-1))\le N'(1+20\delta)\),
because \(X\le e\le N'\) for every tested product of the box. Hence
\(N-mn\ge NL^{-334}(L-20)\ge1\) and \(|mn-N|\le N<X^3\), which is (3.10).
The integer convention of \(I_T\) is that of Paper~I, Lemma~B.22, since the
mesh and the roughness are the same. LEN follows from the condition
\begin{equation}\label{eq:k2-med-length}
 \sigma L\ge2^{130}(335)(\log L+1).
\end{equation}
It gives \(T^{2^{-126}}\ge12L^{334}\) and
\(N_I\ge(2T)^{1-2^{-126}}\).
\item[H4.] \(a=N\), \(0<N\le XL\le A_0\), and every supported modulus is
coprime to \(2N\).
\item[H5.] \(N^{60\sigma}\le T\le N^{1/3+14\sigma}\), and (3.99) gives
\(T\le X^{1/3+15\sigma}\) and \(T<X^{1/2}\).
\end{description}
The selected-factor Siegel--Walfisz bound is Paper~I, (3.25). It is
proved for these three coefficient classes with any \(\zeta\ge z_0\). In the
large-norm case of Paper~I, (3.48), the zero mode uses
Lemma~\ref{hyp:ra-ms-transfer}(ii) in place of \cite[Theorem~0(a)]{BFI86}.
Its premises (3.1), (3.2), (3.6), (3.8) with LEN, (3.11) and (3.25)--(3.27)
are those verified above. Theorem~2.8 of Paper~I enters only inside the
proof of Theorem~3.2, with the dictionary
\begin{equation}\label{eq:k2-theorem-28-dictionary}
 r_T=z,\quad s_T=p_0,\quad n_T=p_0m,\quad a_T=\varepsilon_yN,
 \quad N_T=8p_0H_*SF,\quad A_T=A_0,
\end{equation}
both signs, all nine smooth weights and the full bracket.
Paper~I, Theorem~3.2, (3.13), gives, uniformly in the box, the block and
the cell,
\begin{equation}\label{eq:k2-med-per-box}
 |S_{\mathcal B}^{\rm PSDT}|
 \ll_{\sigma,A'}XL^{-A'}+X^{1-c_\sigma}L^C,\qquad
 S_{\mathcal B}^{\rm PSDT}=\sum_q\Gamma_{\mathfrak c,k}(q)
 \sum_{\mathbf v\in\mathcal B}w(\mathbf v)\mathcal S_{e}(q;N).
\end{equation}
With \(A=4\), \(C_{\rm box}=334\), \(A'=A+24C_{\rm box}+31=8051\) and
\eqref{eq:k2-total-coefficient-norm}, Paper~I, (3.92), gives
\begin{equation}\label{eq:k2-med-aggregate}
 E_{\rm II}=\sum|c_{\mathcal B}|\,|S_{\mathcal B}^{\rm PSDT}|
 \ll_{b,\sigma}NL^{4022-8051}+N^{1-c_\sigma}L^{4022+C}.
\end{equation}
Here \(S_{\mathcal B}^{\rm PSDT}\) already contains
\(\Gamma_{\mathfrak c,k}\), hence its factor
\(M_{\mathfrak c}^{-1}\) from \eqref{eq:k2-native13-coefficient}. The
coefficient norm \eqref{eq:k2-total-coefficient-norm} sums only
\(|c_{\mathcal B}|\), so the orbit factor is paid exactly once.

\begin{lemma}[Relaxed Type-I-2 estimate]\label{lem:k2-I2-prime}
The source hypotheses are the Type-I-2
conditions (SH1)--(SH5) stated in Paper~I:
\begin{quote}
For a direct shell \([X,2X)\), assume the following source conditions
(SH1)--(SH5).
\begin{itemize}
\item[(SH1)] The coefficient \(\gamma(q)\) is independent of tested variables,
\((q,2N)=1\), \(|\gamma(q)|\le C_\sigma\), \(q\le N^{e_Q}\),
\(e_Q<153/256\), and \((a,q)=1\) on its support.
\item[(SH2)] The tested geometric source and cap restrictions are constant
on each nonboundary box; the roughness sieve is separate as in (SH3).
Smooth factor weights have \(\sup|w|+V(w)\ll1\).
\item[(SH3)] All HB factors are rough at the threshold \(z^*=z_0(X)\) or
\(z_0(N)\).
\item[(SH4)] The discrepancy is exactly (3.3).
\item[(SH5)] The assignment priority is proper prime powers, boundary boxes,
Type II, Type I-1, then Type I-2.
\end{itemize}
\end{quote}
Assume (SH1)--(SH5) for a nonboundary K2 box. Let the two large smooth
variables have nonnegative weights \(w_1,w_2\) and \(N_1,N_2\ge X^{1/3}\).
Let the remaining small variables, at most ten, have product \(M\). Assume
\[
 M\le X^{6.02\cdot10^{-8}},\qquad N_1N_2M\le2^{30}X,
 \qquad q<X^{0.5977},
\]
and \(|\alpha_m|\ll\tau_{10}(m)\le\tau(m)^9\). Then, for every fixed
\(A>0\), uniformly in the supported unit residues,
\[
\sum_{q\asymp Q}|\gamma(q)|\left|\sum_m\alpha_m
 \sum_{n_i\in I_i,\ P^-(n_i)\ge z^*}
 w_1(n_1)w_2(n_2)\mathcal S_{mn_1n_2}(q;a)\right|
\ll_{\sigma,A} XL^{-A}.
\]
The same statement holds for a fixed finite sum of signed weights. Split
each bounded real weight into its two nonnegative parts.
\end{lemma}

\begin{proof}
Restrict first to \((m,q)=1\), and write \(\Delta=C-M_q\) for the positive
congruence and comparator functionals. In the K2 application the inner
product is \(e=mn_1n_2\), the residue is \(N\), and \((q,N)=1\). If
\((m,q)>1\), then \((e,q)>1\), so both \(1_{e\equiv N\ (q)}\) and the
comparator \(1_{(e,q)=1}/\varphi(q)\) vanish. The omitted range contributes
zero. Apply the Brun--Bonferroni sandwich
of Paper~I, (3.102)--(3.104), at
\(D_f=X^{\epsilon_I/2}\), \(\epsilon_I=7.5\sigma\). For a requested
fixed logarithmic precision \(A_{\rm sv}\), put
\[
 s_f=\frac{\log D_f}{\log z^*},\quad
 E=\sum_{p<z^*}\frac1p,\quad J=2\lfloor s_f/2\rfloor,
\]
and increase the common threshold until
\(J\ge\max\{2,e^2E,A_{\rm sv}\log L\}\). Since
\(s_f\asymp\sigma(\log\log N)^3\) and \(E=O(\log\log N)\), this is
eventual at fixed \(\sigma\), and \(E^J/J!\le L^{-A_{\rm sv}}\).

For the two large factors put \(R_i=1_{P^-(n_i)\ge z^*}\) and
\(G_i=\Psi_i^+-\Psi_i^-\ge0\). Positivity gives
\[
 \Psi_1^+\Psi_2^+-G_1\Psi_2^+-\Psi_1^+G_2
 \le R_1R_2\le\Psi_1^+\Psi_2^+,
\]
so the three signed gaps and the two nonnegative comparator gaps
bound \(|\Delta(R_1R_2)|\), as in Paper~I, (3.113). Expand the
three gaps over \(d_i\le D_f\). The completed hyperbola estimate (3.109)
 gives \(O_{\epsilon_X}(D_f^2q^{1/2}X^{\epsilon_X})\). Here we fix
\(\epsilon_X=10^{-4}\). Squareful moduli are included
by the full-modulus Weil--Estermann bound.

The two one-dimensional sieve gaps satisfy
\[
 \sum_{n\in I_i}G_i(n)w_i(n)\ll_{A_{\rm sv}}N_iL^{-A_{\rm sv}}+D_f,
 \qquad \sum_{n\in I_i}\Psi_i^+(n)w_i(n)\ll N_i+D_f.
\]
Thus the total is bounded by
\[
 M D_f^2Q^{3/2}X^{\epsilon_X}L^C+XL^{C-A_{\rm sv}}
 +M D_f(N_1+N_2)L^C+M D_f^2L^C. \tag{I-2'}
\]
The ten-variable check uses the following upper bound. The expression is increasing
throughout \(0<\sigma\le10^{-10}\), so its maximum is evaluated at
\(\sigma=10^{-10}\):
\[
 \frac{600\sigma+\kappa}{1-\omega}
 \le\frac{(600+1/202)\sigma}{1-\sigma/101}
 \le\frac{(600+1/202)10^{-10}}{1-10^{-10}/101}
 =\frac{121201}{2019999999998}
 \approx6.0000495049564356926\cdot10^{-8}<6.02\cdot10^{-8},
\]
where \(\omega,\kappa\le\sigma/202\) absorb all fixed dyadic factors,
so \(1-\omega\ge1-\sigma/202\ge1-\sigma/101\), consistently with
\eqref{eq:k2-r3-typeI2}.
Eleven small variables would give
\(\frac{133321}{2019999999998}\approx
6.600049504957030\cdot10^{-8}\), which exceeds the allowed bound.
With \(q<X^{0.5977}\), the first exponent in (I-2') is at most
\[
 6.02\cdot10^{-8}+7.5\cdot10^{-10}+1.5(0.5977)+\epsilon_X
 =0.89665006095<0.897,
\]
leaving slack \(0.00034993905\). Also, using
\(N_1N_2M\le2^{30}X\), \(N_1\ge N_2\ge X^{1/3}\), and
\(D_f=X^{3.75\sigma}\),
\[
 M D_f(N_1+N_2)\le2^{31}D_fX/N_2
 \ll X^{2/3+3.75\cdot10^{-10}}<X^{0.668},
 \qquad
 M D_f^2\le X^{6.02\cdot10^{-8}+7.5\cdot10^{-10}}<X^{0.668}.
\]
The coefficient bound uses
\(\tau_{10}(m)\le\tau(m)^9\), so \(B_0=9\) is fixed. Finally, a family of
at most \(O_\sigma(L^K)\) boxes is absorbed by requesting
\(A_{\rm sv}\ge K+A+5\), after all fixed logarithmic losses. This proves
the lemma.
\end{proof}

\paragraph{Type-I-2 boxes: the complementary estimate.}
Lemma~\ref{lem:k2-I2-prime} is a relaxed form of Paper~I,
(3.111), for this source. Below we check each displayed hypothesis
against the actual K2 source.

For the present K2 Type-I-2 box, the nontrivial SH2 check is the
display
\begin{equation}\label{eq:k2-sh2-check}
\begin{gathered}
 1\le v\le X,
 \qquad w(v)\in\left\{1,\frac{\log v}{\log X}\right\},
 \qquad P^-(n_i)\ge z^*\ \text{is a separate tested condition},\\
 \sup|1|+V(1)=1,\qquad
 \sup\left|\frac{\log v}{\log X}\right|
   +V\!\left(\frac{\log v}{\log X}\right)
 \le 1+\frac{2\log(1+L^{-334})}{\log X}=O(1).
\end{gathered}
\end{equation}
Thus the two tested large factors have the required nonnegative weights and
bounded supremum plus variation. The roughness condition is not hidden inside
the weight.

For the Type-I-2 boxes above, (SH1) has
\(\gamma=\Gamma_{\mathfrak c,k}\). The support and bound are
\eqref{eq:k2-r3-support} and \eqref{eq:k2-native13-coefficient},
\(e_Q<153/256\), and the residue is the unit \(a=N\). The partition
construction proves the geometric constancy in (SH2) on each nonboundary
box. In this two-large-variable case the two tested large factors are the
selected smooth factors. The roughness restriction is kept as the separate
condition \(P^-(n_i)\ge z^*\) in the inner sum. Their smooth weights are then
individually \(1\) or \(\log v/\log X\), not \(\mu_{z_1}\). Hence they are
nonnegative with \(\sup|w|+V(w)\ll1\) as required by (SH2). The
\(\mu_{z_1}\) factor occurs only among the complementary small factors.
Condition (SH3) is \(z_1\ge z_0(X)\).
Condition (SH4) is the natural full-product discrepancy fixed above.
Condition (SH5) is the priority in the preceding boundary and Type-II
classification.

On a large-coordinate cell \((V,(1+\delta)V]\),
\(\delta=L^{-334}\), one has \(1\le V\le X\). The smooth weight
\(w(v)=1\) has supremum plus variation equal to 1. The alternative
\(w(v)=\log v/\log X\) is nonnegative and increasing, with
\[
 \sup |w|+V(w)
 \le 1+\frac{2\log(1+\delta)}{\log X}=O(1).
\]
The roughness restriction remains a separate tested condition under (SH3)
and inside the sieve, and it is not claimed to be constant on a box.

Thus the relaxed estimate of this chapter applies with
\(z^*=z_1\), the two large smooth weights nonnegative, and
\[
 |\alpha_{m^\sharp}|\le\tau_{10}(m^\sharp)\le\tau(m^\sharp)^9.
\]
Each of the at most ten small labels contributes one factor. Normalized
logarithmic weights \(\log v/\log X\) are at most one on the support.
Thus the coefficient hypothesis uses the fixed divisor exponent
\(B_0=9\), independent of \(N\). The first term of (I-2') uses the
following values. The slack is \(0.00034993905\)
at exponent \(0.89665006095\). The two final terms have the bound
\(2/3+3.75\cdot10^{-10}\). The absorption is
\(\omega,\kappa\le\sigma/202\), and the divisor bound is \(\tau_{10}\)
with \(B_0=9\). By \eqref{eq:k2-r3-typeI2}, Lemma~\ref{lem:k2-I2-prime}
applies to every K2 Type-I-2 box. Requesting \(A_I=K+A+5\) for the
fixed \(O_\sigma(L^K)\) box family (here \(K=4022\)) gives
\begin{equation}\label{eq:k2-r3-complementary}
 E_{\rm I2}\ll_{b,\sigma}NL^{-A}=NL^{-4}.
\end{equation}

\begin{proposition}[Row-three remainder]\label{prop:k2-row3}
Fix \(b\in(1/3,1/2)\) and \(0<\sigma\le10^{-10}\). Under Paper~I,
Theorems~2.8 and~3.2, with the relaxed Type-I-2
Lemma~\ref{lem:k2-I2-prime} and Lemma~\ref{hyp:ra-ms-transfer}(ii) at (3.48),
\[
 \frac{|r_3|}{\mathcal N_K}\ll_{b,\sigma}
 N^{-\ell_0\sigma^2}L^3+N^{-b/4}L^7+L^{-260}+L^{-120}+L^{-2}
 +N^{-c_\sigma}L^{4024+C}.
\]
\end{proposition}

\begin{proof}
By \eqref{eq:k2-r3-occurrence} and \eqref{eq:k2-native13-coefficient},
\(r_3\) is the \(\Gamma\)-sum minus the squareful sum
\eqref{eq:k2-r3-squareful}. By \eqref{eq:k2-r3-hb-identity}, the
\(\Gamma\)-sum is \(\mathcal E_{\rm pp}\) plus the tuple sum. Every tuple
lies in exactly one box. The boundary tuples give
\eqref{eq:k2-r3-bdry}. On the other boxes logarithm freezing costs
\eqref{eq:k2-r3-log}. Each box is of Type~II or Type~I-2. Sum
\eqref{eq:k2-r3-squareful}, \eqref{eq:k2-r3-pp}, \eqref{eq:k2-r3-log},
\eqref{eq:k2-r3-bdry}, \eqref{eq:k2-med-aggregate} and
\eqref{eq:k2-r3-complementary}, with \(O_\sigma(1)\) pairs
\((\mathfrak c,k)\). Dividing by
\(\mathcal N_K=(1-\varepsilon)C(N)N/L^2\ge(1-\varepsilon)C_2N/L^2\gg_b N/L^2\) gives the
claim.
\end{proof}

\begin{proposition}[Exterior source estimates]\label{prop:k2-applications}
Fix \(b,\sigma\) before \(N\). For the source \eqref{eq:k2-source}:
\begin{enumerate}
 \item \eqref{eq:k2-parent-comparison} and \eqref{eq:k2-main-mass} hold
 for the two-orientation source and four signed AA main terms in
 \eqref{eq:k2-row-main-masses} and \eqref{eq:k2-parent-functional};
 \item \eqref{eq:k2-aa-density-mass} and \eqref{eq:k2-aa-mass} hold for
 \(N\ge N_{\mathrm{AA}}(b,\sigma)\);
 \item \eqref{eq:k2-peel} holds for \(N\ge N_{\rm peel}(\sigma)\);
 \item \eqref{eq:k2-remainder-estimates} holds.
\end{enumerate}
\end{proposition}

\begin{proof}
Item 1 is \eqref{eq:k2-parent-reserve}--\eqref{eq:k2-parent-kernel} and the
derivation after \eqref{eq:k2-main-mass}. Item 2 is
\eqref{eq:k2-aa-decay}--\eqref{eq:k2-aa-density-derived}. Item 3 is
\eqref{eq:k2-peel-ranges}--\eqref{eq:k2-peel}. Item 4 is
\eqref{eq:k2-low-row-remainders} and Proposition~\ref{prop:k2-row3}.
\end{proof}


\begin{proposition}[Exterior comparison]\label{prop:k2-comparison}
For fixed \(b,\sigma\) there is a nonnegative \(r_K(N;b,\sigma)\to0\) such
that, for \(N\ge T_r(b,\sigma)\) with \(T_r\) as defined below,
\begin{equation}\label{eq:k2-final-fixed-sigma}
\begin{aligned}
 \frac{4B_b(N)}{W(N)}
 &\le c_{K2,G0}I(b)+A_K(b)C_{\mathrm{AA}}\sigma
          +P_K\sigma^9+r_K(N;b,\sigma),\\
 P_K&=\frac{388400000}{1-\varepsilon}
     =\frac{3884000000000}{9999}.
\end{aligned}
\end{equation}
Here
\(T_K=\max\{N_{\mathrm{mass}}(b,\sigma),N_{\mathrm{AA}}(b,\sigma),
N_{\rm peel}(\sigma),T_{\mathrm{len}}\}\). The threshold \(N_{\mathrm{mass}}\) is from
the fixed-parameter estimate at \eqref{eq:k2-main-mass}.
\(N_{\mathrm{AA}}\) is from
item~2 of Proposition~\ref{prop:k2-applications}. Both exist with
\(b,\sigma\) fixed. Finally, \(T_{\mathrm{len}}\) is the named length and
application threshold of Section~\ref{sec:notation}. In addition,
\begin{equation}\label{eq:k2-error-definition}
 r_K=\frac{E_{\rm top}+[M_K-\mathcal N_K2e^{-\gamma}I(b)G_*]_+
 +\sum_{j=0}^3|r_j|}{\mathcal N_K}
 =o_{b,\sigma}(1).
\end{equation}
In particular, let \(\widehat T_r\) be the least even threshold above
which \(r_K\le\delta_{\rm err}/10\), and set
\(T_r=\max\{N_{\mathrm{mass}}(b,\sigma),N_{\mathrm{AA}}(b,\sigma),
T_K,\widehat T_r\}\).
These thresholds exist at fixed \(b,\sigma\).
\end{proposition}

\begin{remark}[Single payment of each error]
In \cite[Theorem~2.6, (2.30)]{PaperI} the three summands
\(XV E_s\), \(R_1^{\nu}\) and \(R^{\nu}\) are distinct. The first is
paid only by \eqref{eq:k2-aa-mass}, the other two only by
\eqref{eq:k2-remainder-estimates}. The decaying part of
\eqref{eq:k2-aa-decay} and the \(o(1)\) in
\eqref{eq:k2-aa-density-derived} are absorbed by thresholds at fixed
\(\sigma\). The peel term \(P_3\) contains no linear-sieve
application. Thus
\[
 r_K=\frac{E_{\mathrm{top}}+\sum_j|r_j|
       +[M_K-\mathcal N_K2e^{-\gamma}I(b)G_*]_+}{\mathcal N_K},
\]
and no error is counted twice.
\end{remark}

\begin{proof}
Use \eqref{eq:k2-signed-join} and divide by \(\mathcal N_K=W/4\). The
principal term is bounded by \(\mathcal N_K2e^{-\gamma}I(b)G_*\) plus its
positive part. Also \(2e^{-\gamma}G_*\le c_{K2,G0}\) by
\eqref{eq:k2-certificate}--\eqref{eq:k2-gamma-enclosure}.
Proposition~\ref{prop:k2-applications}(2),(3) pay the sieve and peel terms
above \(T_r\). By definition, \(T_r\) includes
\(N_{\mathrm{mass}}(b,\sigma)\), \(N_{\mathrm{AA}}(b,\sigma)\) and
\(T_K\). The remaining terms form \(r_K\). First,
\(E_{\rm top}/\mathcal N_K\ll L^{-331}\). The positive part contributes
\(o_{b,\sigma}(1)\) by \eqref{eq:k2-main-mass}. The row remainders are
\(O_{b,\sigma}(L^{-2})\) by \eqref{eq:k2-remainder-estimates}. Hence
\(r_K=o_{b,\sigma}(1)\).
The convergence is ineffective through Siegel--Walfisz, so \(T_r\)
exists but is not computed.
\end{proof}

\section{The switched sum}\label{sec:l3}

We estimate the single term $-S_{\mathrm{sw}}/W$.
The normalization is \eqref{eq:not-normalization}.
Fix $\sigma$ as in \eqref{eq:not-sigma-choice}.
Put
\begin{equation}\label{eq:l3-constants}
\begin{gathered}
 \eta_L=10^{-4},\quad\varepsilon=10^{-4},\qquad
 H_{\max}=\frac{153}{256},\qquad
 \ell_0=\frac{264897}{2060000},\qquad \ell_1=\frac{109}{1000},\\
 N_{\mathrm{cut}}=N(1-L^{-333}),\qquad
 \mathcal P_N=\{p\in\mathbb P:p\nmid2N\}.
\end{gathered}
\end{equation}
All limits in this section fix $\sigma$ first.

\subsection{A common nonnegative source}

We use the prime-coordinate cells of Paper~II
\cite[(5.29a)--(5.31)]{PaperII}.
Their rules are
\begin{equation}\label{eq:l3-descriptors}
\begin{gathered}
 I_1=[\alpha,\beta),\quad I_2=[\beta,\lambda),\quad
 I_3=[\lambda,1/3),\quad I_4=[1/3,1],\\
 C_{\mathrm{L3}}=\frac{57}{64}-u-\frac{11\alpha}{5}
       =\frac{6408241}{9192000},\qquad
 v_\beta=\frac{2(C_{\mathrm{L3}}-\beta)}5,\\
 k_1+k_2+k_3+k_4=k,\quad 5\le k\le11,\quad
 k_1\ge m_1\ge\cdots\ge m_{k_2}\ge0,\\
 v_j\le v_\beta\ \Longrightarrow\ m_j=k_1,\\
 v_j>v_\beta\ \Longrightarrow\
 x_{m_j}<C_{\mathrm{L3}}-\frac52v_j\le x_{m_j+1},\\
 w_C=3+1_{k_1=0}-2k_1-2k_2-k_3+
            \binom{k_1}{2}+\sum_{j=1}^{k_2}m_j.
\end{gathered}
\end{equation}
The $x_i$ and $v_j$ increase in $I_1$ and $I_2$.
The outer sentinels are infinite.
Let $\mathcal C_+$ contain the cells with $w_C>0$.
For such a cell, put
\begin{equation}\label{eq:l3-source}
\begin{aligned}
 a_C(e)&=1_{\substack{e=p_1\cdots p_k,\ N^\alpha\le p_1<\cdots<p_k\\
             (\log_Np_1,\ldots,\log_Np_k)\in C}},\\
 A(e)&=\sum_{C\in\mathcal C_+}w_Ca_C(e)
                     1_{\varepsilon N<e<N_{\mathrm{cut}}},
                     &0\le A(e)&\le36,\\
 X_A&=\sum_eA(e),&Y_A(q)&=\sum_{(e,q)=1}A(e),\\
 C_A(q;N)&=\sum_eA(e)1_{q\mid N-e}.&&
\end{aligned}
\end{equation}
Since $\sum_jm_j\le k_1k_2$, we have
\[
 w_C\le3+1_{k_1=0}-2k_1-2k_2-k_3
             +\frac{k_1(k_1-1)}2+k_1k_2\le36.
\]
For $k_1\ge2$, maximize at $k_3=0$, $k_2=11-k_1$;
the resulting concave quadratic has maximum $36$ at $k_1=10,11$.
The cases $k_1=0,1$ give at most $4,1$.
The strict cells are disjoint, so this proves the bound for $A$.
Every prime is permitted in $A$.
In particular, $A$ has no factor $1_{(e,2N)=1}$.
Keep the selector $1_{\Omega(e)\ge3}$ in the switched source.
The required pointwise comparison is
\begin{equation}\label{eq:l3-domination}
 1_{\Omega(e)\ge3}(W_{13}(e))_+\le A(e)
 \quad\left(\begin{gathered}
 \mu_{\mathrm{Mob}}(e)^2=1,\quad P^-(e)\ge N^\alpha,\\
 \varepsilon N<e<N_{\mathrm{cut}}
 \end{gathered}\right).
\end{equation}
Here $\mu_{\mathrm{Mob}}$ denotes the M\"obius function.
Paper~II gives zero positive part when $\Omega(e)=3,4$
\cite[(1.7), (4.9), (5.29a)--(5.31)]{PaperII}.
\paragraph{Proof of \eqref{eq:l3-domination}.}
Let $e$ satisfy the side conditions and write $e=p_1\cdots p_k$,
$p_1<\cdots<p_k$, $x_i=\log_Np_i$.
If $(e,N)>1$, every occurrence in \eqref{eq:source-row-occurrences}
vanishes, because it requires $(d,N)=1$ and $(n/d,Nd)=1$
\cite[(1.7)]{PaperII}. Then $W_{13}(e)=0$. Assume $(e,N)=1$.
Since $N^{k\alpha}\le e<N$ and $12\alpha=1200/1149>1$, we have $k\le11$.
If $k\le2$, the selector vanishes. If $k\in\{3,4\}$, Paper~II, (4.9),
gives $W_{13}(e)\le0$. Its hypotheses are exactly squarefree $e$,
$(e,N)=1$, $\varepsilon N<e<N$ and $P^-(e)\ge N^\alpha$
\cite[(4.9)]{PaperII}.
Let $5\le k\le11$. The table after Paper~II, (5.29a), shows that the
three negative pair counts vanish, so
\[
 W_{13}(e)=3+1_{k_1=0}-2k_1-2k_2-k_3+\binom{k_1}2+\sum_{j=1}^{k_2}m_B(v_j),
\]
where $m_B(v)$ counts first-band coordinates below $B_{\mathrm{O1}}(v)$
\cite[(5.29a), (5.30)]{PaperII}. The bins determine $k_1,\ldots,k_4$.
If $v_j\le v_\beta$, then $B_{\mathrm{O1}}(v_j)=\beta$ exceeds every
first-band coordinate, so $m_B(v_j)=k_1$. If $v_j>v_\beta$, then
$B_{\mathrm{O1}}(v_j)=C_{\mathrm{L3}}-\frac52v_j$, and $m_B(v_j)$ is the unique $m_j$
with $x_{m_j}<C_{\mathrm{L3}}-\frac52v_j\le x_{m_j+1}$ (infinite sentinels).
It is nonincreasing in $j$. Thus $x$ lies in exactly one cell $C$ of
\eqref{eq:l3-descriptors}, and $W_{13}(e)=w_C$
\cite[(5.31)]{PaperII}. If $w_C\le0$, the left side of
\eqref{eq:l3-domination} is zero. Otherwise $C\in\mathcal C_+$ and
$a_C(e)=1$, so the left side equals $w_C=A(e)$. This proves
\eqref{eq:l3-domination}. The same argument, without the upper cut, gives
\begin{equation}\label{eq:l3-domination-nocut}
 1_{\Omega(e)\ge3}(W_{13}(e))_+\le A^0(e):=\sum_{C\in\mathcal C_+}w_Ca_C(e)\le36
\end{equation}
for squarefree $e$ with $P^-(e)\ge N^\alpha$ and $\varepsilon N<e<N$.

Put
\begin{equation}\label{eq:l3-sieved-source}
\begin{gathered}
 b_A(n)=A(N-n),\qquad
 P_N(z)=\prod_{\substack{p<z\\p\in\mathcal P_N}}p,\\
 S_v(z)=\sum_{\substack{n\ge1\\v\mid n}}b_A(n)1_{(n,P_N(z))=1},
 \qquad S_A=S_1(N^{\mu_*}),\\
 E_{\mathrm{src}}=(S_{\mathrm{sw}}-S_A)_+,\qquad
 Z_{\mathrm{src}}=(S_A-S_{\mathrm{sw}})_+.
\end{gathered}
\end{equation}
The definitions give the finite identity
\begin{equation}\label{eq:l3-source-id}
 S_{\mathrm{sw}}=S_A+E_{\mathrm{src}}-Z_{\mathrm{src}},
 \qquad E_{\mathrm{src}},Z_{\mathrm{src}}\ge0.
\end{equation}
Under \eqref{eq:l3-domination}, the required bound is
\begin{equation}\label{eq:l3-source-error}
 E_{\mathrm{src}}\le910N^{1-\alpha}
            +36(NL^{-333}+1)+36N^{\mu_*}.
\end{equation}
There are at most $2N^{1-\alpha}$ rough squareful integers.
Each costs at most $455$, by the pointwise weight bound of
Paper~II, proof of Lemma~1.1, immediately after (1.11)
\cite[Lemma~1.1, proof after (1.11)]{PaperII}.
The last two terms cover the top strip and small primes.
\paragraph{Proof of \eqref{eq:l3-source-error}.}
Write $e=N-p$ for $p\in\mathbb P$, $p<(1-\varepsilon)N$.
Every prime $p\ge N^{\mu_*}$ has $(p,P_N(N^{\mu_*}))=1$, and $A\ge0$.
Hence $S_A\ge\sum_{p\ge N^{\mu_*}}A(N-p)$, and
\[
 S_{\mathrm{sw}}-S_A\le\sum_{p\ge N^{\mu_*}}\bigl(s_{\mathrm{sw}}(e)-A(e)\bigr)_+
                +\sum_{p<N^{\mu_*}}s_{\mathrm{sw}}(e).
\]
Here $s_{\mathrm{sw}}(e)=0$ unless $(e,N)=1$ and $P^-(e)\ge N^\alpha$,
and always $0\le s_{\mathrm{sw}}(e)\le|W_{13}(e)|\le455$
\cite[Lemma~1.1, proof after (1.11)]{PaperII}.
Split the summands into three classes; overlaps only enlarge the bound.
(i) Non-squarefree $e<N$ with $P^-(e)\ge N^\alpha$ number at most
$\sum_{\ell\ge N^\alpha}N/\ell^2\le2N^{1-\alpha}$; they cost at most
$455\cdot2N^{1-\alpha}=910N^{1-\alpha}$.
(ii) Squarefree $e$ with $N_{\mathrm{cut}}\le e<N$ number at most
$NL^{-333}+1$; by \eqref{eq:l3-domination-nocut} they cost at most
$36(NL^{-333}+1)$.
(iii) The primes $p<N^{\mu_*}$ with squarefree $e$ cost at most $36N^{\mu_*}$,
again by \eqref{eq:l3-domination-nocut}.
For the remaining squarefree $e$ with $\varepsilon N<e<N_{\mathrm{cut}}$
and $p\ge N^{\mu_*}$, \eqref{eq:l3-domination} gives a zero summand.
This proves \eqref{eq:l3-source-error}. Since $W\gg N/L^2$,
$E_{\mathrm{src}}/W\ll L^{-331}$.

\subsection{Four rows and the lower-sieve cutoff}

Use the clipped intervals
$[N^t,N^{\rho t})\cap[N^u,N^{\mu_*})$.
Let $\mathcal A_j$ comprise the ordered prefixes
\[
 N^u\le a_j<\cdots<a_1<N^{\mu_*},\qquad a_i\in\mathcal P_N.
\]
Write $v=a_1\cdots a_j$.
The empty prefix has $v=1$.
Define
\begin{equation}\label{eq:l3-buchstab}
\begin{gathered}
 T_0=S_1(N^u),\quad
 T_1=\sum_{\mathcal A_1}S_{a_1}(N^u),\quad
 T_2=\sum_{\mathcal A_2}S_{a_1a_2}(N^u),\\
 T_3^{\mathrm{dyn}}=\sum_{\mathcal A_3}S_{a_1a_2a_3}(a_3),
 \qquad S_A=T_0-T_1+T_2-T_3^{\mathrm{dyn}}.
\end{gathered}
\end{equation}
This is Buchstab's finite identity
\cite[(4.18)]{PaperII}.
Its coefficient is
$1-k+\binom{k}{2}-\binom{k-1}{2}=0$ for $k\ge1$.
At $k=0$ it equals $1$.

At depth three, let $\Theta_c$ be the infimum
of $\widetilde H$ on the closed clipped parent. The underlying prime
cells stay half-open. The closure is used only for this infimum.
Use the whole anatomy function before choosing a branch
\cite[(2.17d)--(2.17e), (4.19)]{PaperII}.
Set
\begin{equation}\label{eq:l3-row3-level}
 \ell_c=\Theta_c-\eta_L-t_1-t_2-t_3,\quad
 D_c=N^{\ell_c},\quad z_c=N^{t_3},\quad
 \chi_c=1_{\ell_c\ge2\rho t_3}.
\end{equation}
Retained prefixes have $\chi_c=1$.
For these prefixes define
\begin{equation}\label{eq:l3-peel}
\begin{aligned}
 B_3&=\sum_{\mathcal A_3:\chi_c=1}S_v(z_c),&
 D_-&=\sum_{\mathcal A_3:\chi_c=0}S_v(a_3),\\
 D_+&=\sum_{\mathcal A_3:\chi_c=1}\{S_v(z_c)-S_v(a_3)\},&
 P_3&=\sum_{\mathcal A_3:\chi_c=1}
       \sum_{\substack{r\in\mathcal P_N\\z_c\le r<a_3}}S_{vr}(z_c),\\
 Z_{\mathrm{peel}}&=D_-+P_3-D_+\ge0.&&
\end{aligned}
\end{equation}
The union bound gives $D_+\le P_3$.
Therefore
\begin{equation}\label{eq:l3-buchstab-peel}
 S_A=T_0-T_1+T_2-B_3+P_3-Z_{\mathrm{peel}}.
\end{equation}
The favorable term enters exactly once
\cite[(4.45)]{PaperII}.

\subsection{Capping the first three levels}

At depths $j=0,1,2$, replace the relative level by
\begin{equation}\label{eq:l3-cap}
 T_c=\sum_{i=1}^jt_i,\quad
 \ell_c=\frac12-\eta_L-\rho T_c,\quad
 D_c=N^{\ell_c},\quad z_{j,c}=N^u.
\end{equation}
Here $T_c=0$ at depth zero.
Depth three keeps \eqref{eq:l3-row3-level}, with $z_{3,c}=z_c$.
The elementary bounds are
\begin{equation}\label{eq:l3-cap-margins}
\begin{gathered}
 \ell_c\ge\frac12-\eta_L-j\rho\mu_*\ge\ell_1,
 \qquad s_c=\ell_c/u\ge54.6375-10^{-87}>30,\\
 \text{continuum reference values for depths }0,1,2\text{ (the first three rows)}
 \text{ at }(\eta_L,\rho)=(0,1):
       250,\ \frac{4875}{32},\ \frac{875}{16},\\
 q<N^{\rho T_c}D_c^{\rho+\sigma}
   \le N^{1/2-\eta_L+(\sigma+\tau)/2}
   <N^{1/2-9\cdot10^{-5}}.
\end{gathered}
\end{equation}
The modulus includes the smooth divisor and every tail prime.
The exponent $\rho+\sigma$ is kept.
The cap changes no row-three level.

The lower floor changes the repeated-prime argument.
Put
\begin{equation}\label{eq:l3-rough-floor}
 u_1=N^{\ell_1\sigma^2},\qquad
 M_\sigma=\left\lceil\frac3{5\ell_0\sigma^2}\right\rceil,
 \qquad C_{\mathrm{L3},\sigma}=M_\sigma!.
\end{equation}
\paragraph{(P2.8b) Rough-coordinate and repeated-prime bounds.}
Assume the support and rough-coordinate hypotheses stated below.
These are the capped version of Paper~II's Lemmas B.1--B.4.
The coefficients and masks are defined in \eqref{eq:l3-masks}.
Every rough coordinate is at least $u_c=D_c^{\sigma^2}\ge u_1$.
On an active block, count these coordinates by $m_c$.
The capped modulus bound gives
\begin{equation}\label{eq:l3-rough-count}
 u_1^{m_c}\le q<N^{1/2},\qquad
 m_c<\frac{500}{109\sigma^2}
       \le\frac3{5\ell_0\sigma^2},\qquad
 \frac{m_c!}{M_c}\le C_{\mathrm{L3},\sigma}.
\end{equation}
The prefix orbit size is $M_c\in\{1,2,6\}$.
The small-prime part of $q$ fixes the smooth divisor.
The rough multiset fills at most $m_c!$ labelled positions.
Thus
\begin{equation}\label{eq:l3-cap-coefficient-bound}
 \frac{|C_{\mathfrak b}(q)|}{M_c}\le C_{\mathrm{L3},\sigma}.
\end{equation}
Define $\mathcal H_1=\{q<N^{1/2}:p^2\mid q\text{ for some }p\ge u_1\}$.
The forbidden and allowed-repeat masks satisfy
\begin{equation}\label{eq:l3-u1-masks}
 |\Xi_{U,\mathfrak b}(q)|+|\Xi_{A,\mathfrak b}(q)|
                \le C_{\mathrm{L3},\sigma}1_{\mathcal H_1}(q).
\end{equation}
Writing $q=vp^2$ and overcounting gives
\begin{equation}\label{eq:l3-u1-sums}
 \sum_{q\in\mathcal H_1}\frac1q\le\frac{2(1+L)}{u_1},
 \qquad \#\mathcal H_1\le\frac{2N^{1/2}}{u_1}.
\end{equation}
The tested measure has $0\le A\le36$ and support below $N$.
For $N\ge N_{\mathrm{dom}}$, $q/\varphi(q)\le1+L$ throughout $q<N^{1/2}$.
Consequently
\begin{equation}\label{eq:l3-u1-test}
 |\Delta_P(q)|,\ |d_A(q)|
 \le36\left(\frac Nq+1+\frac N{\varphi(q)}\right)
 \le36\left(\frac{N(2+L)}q+1\right).
\end{equation}
Per block, these bounds give
\begin{equation}\label{eq:l3-u1-block}
 \sum_q\bigl(|\Xi_U(q)\Delta_P(q)|+|\Xi_A(q)d_A(q)|\bigr)
       \le\frac{216C_{\mathrm{L3},\sigma} N(1+L)^2}{u_1}.
\end{equation}
There are finitely many blocks at fixed $\sigma$.
Division by $W\gg N/L^2$ proves the conditional estimate
\begin{equation}\label{eq:l3-repeated-cap}
 |E_{\mathrm{sq}}|/W=O_\sigma(N^{-\ell_1\sigma^2}L^4).
\end{equation}
These are Paper~II, Lemmas B.1--B.4 and (B.3)--(B.7),
with $u_0$ replaced by $u_1$ on rows $0,1,2$
\cite[Appendix~B.1, Lemmas B.1--B.4, (B.3)--(B.7)]{PaperII}.
Row three keeps $u_0=N^{\ell_0\sigma^2}$, since its level floor
in \eqref{eq:l3-row3-global-margins} exceeds $\ell_0$.
\paragraph{Verification for the capped rows.}
Fix a block in row $j\le2$. Prefix primes lie in
$[N^{t_i},N^{\rho t_i})\cap[N^u,N^{\mu_*})$, so their product is below
$N^{\rho T_c}$. The tail tuple satisfies $\prod_iD_i<D_c$ by its last
cubic constraint, and the tail primes are below $D_i^\rho$; also
$d<D_c^\sigma$. Thus
\[
 \log_Nq<\rho T_c+(\rho+\sigma)\ell_c
 =(\rho+\sigma)(\tfrac12-\eta_L)-\rho(\rho+\sigma-1)T_c
 \le(\rho+\sigma)(\tfrac12-\eta_L)<\tfrac12-9\cdot10^{-5}.
\]
Every prime of $q$ lies in $\mathcal P_N$ and $d\mid P_N(u_c)$, so
$(q,2N)=1$. Tail primes are at least $D_i\ge D_c^{\sigma^2}=u_c$.
Prefix primes are at least $N^u>u_c$, because $\ell_c\sigma^2<u$.
Also $u_c=N^{\ell_c\sigma^2}\ge N^{\ell_1\sigma^2}=u_1$.
For an active block choose a nonzero realization. Then
$u_1^{m_c}\le q<N^{1/2}$, so $m_c<1/(2\ell_1\sigma^2)=500/(109\sigma^2)$,
and $1/(2\ell_1)\le3/(5\ell_0)$. This is \eqref{eq:l3-rough-count}.
The primes of $q$ below $u_c$ determine $d$. The rough multiset fills
the $m_c$ labelled positions in at most $m_c!$ ways, and $|\lambda_d|\le1$.
This gives \eqref{eq:l3-cap-coefficient-bound}. It also gives
\eqref{eq:l3-u1-masks}. A forbidden tuple repeats a prime $\ge u_c$.
An allowed non-squarefree tuple has $p^2\mid q$ with $p\ge u_c$,
since $d$ is squarefree and coprime to the rough part. The grid count of
Paper~II, Lemma~B.1(4), gives finitely many blocks at fixed $\sigma$.
These are Paper~II, Lemmas~B.1--B.2, with $(u_0,N^{3/5})$ replaced by
$(u_1,N^{1/2})$ \cite[Appendix~B.1, Lemmas~B.1--B.2]{PaperII}.
Lemma~B.4 with $u_1$ and $Q=N^{1/2}$ gives \eqref{eq:l3-u1-sums}.
Lemma~B.3, with $A\le36<A_{\max}$, support in $[1,N)$ and
$q/\varphi(q)\le1+\log_2q\le1+L$, gives \eqref{eq:l3-u1-test}.
Then $\sum_{q\in\mathcal H_1}(N(2+L)/q+1)\le6N(1+L)^2/u_1$ gives
\eqref{eq:l3-u1-block}. Row three is Paper~II, Lemmas~B.1--B.4 as
printed: its modulus bound $q<N^{e_Q}<N^{3/5}$ and floor
$\ell_c>\ell_0$ are \eqref{eq:l3-row3-global-margins}. Hence
\eqref{eq:l3-repeated-cap} holds for all four rows.

\subsection{The signed sieve identity}

Put $\nu_j=(-1)^j$.
For all four rows, define
\begin{equation}\label{eq:l3-common-sifting-ratio}
 s_c=\frac{\ell_c}{\log_N z_{j,c}}
 =\begin{cases}\ell_c/u,&j=0,1,2,\\ \ell_c/t_3,&j=3.\end{cases}
\end{equation}
Let $K_j=F$ for even $j$, and $K_j=f$ for odd $j$.
For a prefix $v$, set
\begin{equation}\label{eq:l3-density}
\begin{gathered}
 \mathcal X_v=X_A/\varphi(v),\qquad
 \omega_v(p)=\frac p{p-1}1_{p\in\mathcal P_N},
 \qquad \omega_v(p^h)=0\quad(h\ge2),\\
 V_N(z)=\prod_{\substack{p<z\\p\in\mathcal P_N}}
                     \left(1-\frac1{p-1}\right),\qquad
 \Delta_0(q)=C_A(q;N)-\frac{X_A}{\varphi(q)}.
\end{gathered}
\end{equation}
All cutoff primes lie below the prefix primes.
Thus the density product on the sieve support is $V_N$.
The squarefree remainder is $\Delta_0(vd)$.
Apply Fact~\ref{fact:structured-sieve} with
$\epsilon_s=\sigma$, $\delta_s=\tau$, and $L_\omega=2$.
Choose a common nonnegative majorant
\begin{equation}\label{eq:l3-sieve-error}
 E_{s,c}:= C_E\{\sigma+
     \sigma^{-8}e^{K_{\mathrm M}+2}(\ell_cL)^{-1/3}\}.
\end{equation}
The signed sieve error satisfies \(|E_s|\le E_{s,c}\).
\paragraph{Verification of the sieve data.}
Fix a row, a cell and a prefix $v$, and sieve the finite multiset
$\{n:v\mid n\}$ with multiplicity $b_A(n)\ge0$ by $\mathcal P_N$ below
$z=z_{j,c}$. Every prefix prime is at least $z$: on rows $0,1,2$ it is at
least $N^u=z$, and on row three $a_3\ge N^{t_3}=z_c$. Hence $(d,v)=1$
for $d\mid P_N(z)$. In particular $\omega_v(d)=\prod_{p\mid d}p/(p-1)$
for these $d$; the primes $p\mid v$ are all $\ge z$ and are never sifted.
For squarefree such $d$,
\[
 |\mathcal A_d|=\sum_eA(e)1_{vd\mid N-e}=C_A(vd;N),\qquad
 \frac{\omega_v(d)}d\mathcal X_v=\frac{X_A}{\varphi(v)\varphi(d)}
                               =\frac{X_A}{\varphi(vd)},
\]
so $r(\mathcal A,d)=\Delta_0(vd)$. Each remainder modulus in
\eqref{eq:l3-remainders} is squarefree. Here $d$ uses primes below
$D_c^{\sigma^2}$. Tail primes are at least $D_c^{\sigma^2}$ and distinct
(equal labels by the mask, unequal grid intervals disjoint). Prefix
primes are at least $z$. Next, $0<\omega_v(p)=p/(p-1)<p$ for every
$p\in\mathcal P_N$ (all such $p$ are odd),
$\omega_v=0$ off $\mathcal P_N$, and $\omega_v(p^h)=0$ for $h\ge2$. Hence the
second condition in \eqref{eq:classical-density} holds with $L_\omega=2$.
For the first, $(1-\omega_v(p)/p)^{-1}=(1-1/(p-1))^{-1}$, and Paper~II,
(5.5), gives
\[
 \prod_{\substack{w\le p<z\\p\in\mathcal P_N}}\left(1-\frac1{p-1}\right)^{-1}
 \le\frac{\log z}{\log w}\left(1+\frac{K}{\log w}\right)\qquad(2\le w<z)
\]
with one absolute $K$, uniformly in $N$ \cite[(5.5)]{PaperII}.
Take $K_{\mathrm M}$ in \eqref{eq:l3-sieve-error} to be the maximum of
this $K$ and the constant of \eqref{eq:classical-mertens-density}. It
remains absolute and is fixed before $\sigma$. The product
$V(z)$ of Fact~\ref{fact:structured-sieve} is then $V_N(z)$.
The levels satisfy $z\le D_c$ on rows $0,2$ ($s_c>30$) and
$z\le D_c^{1/2}$ on rows $1,3$ ($s_c>30$, resp.\ $\chi_c=1$, i.e.\
$\ell_c\ge2\rho t_3$). Fact~\ref{fact:structured-sieve} with
$\epsilon_s=\sigma$, $\delta_s=\tau$, $D=D_c$ therefore applies to every
actual sequence, and \eqref{eq:l3-signs} follows with
$E_s\le E_{s,c}$, $\log D_c=\ell_cL$.

For each cell fix one coefficient construction.
Use it for every prefix realization in that cell.
Write its remainders as
\begin{equation}\label{eq:l3-remainders}
\begin{aligned}
 R_{1,v}^{\nu}&=\sum_{\substack{d<D_c^\sigma\\d\mid P_N(D_c^{\sigma^2})}}
                         \phi_d^\nu\Delta_0(vd),\\
 R_v^\nu&=\sum_{\mathbf D\in\mathcal D^\nu(D_c)}
   \sum_{\substack{d<D_c^\sigma\\d\mid P_N(D_c^{\sigma^2})}}
   \lambda_d^\nu(\mathbf D)
   \sum_{\substack{D_i\le p_i<\min(D_i^\rho,z_{j,c})\\
                p_i\in\mathcal P_N,\ \text{equal-label entries distinct}}}
                         \Delta_0(vd\textstyle\prod_i p_i).
\end{aligned}
\end{equation}
The coefficients satisfy $|\phi_d^\nu|,|\lambda_d^\nu(\mathbf D)|\le1$
\cite[Theorem~2.6, (2.29)--(2.33b)]{PaperI}.
Define the finite quantities
\begin{equation}\label{eq:l3-four-rows}
\begin{aligned}
 M_j&=\sum_c\sum_{v\in c}\frac{X_A}{\varphi(v)}
                   V_N(z_{j,c})K_j(s_c),\\
 E_j&=\sum_c\sum_{v\in c}\frac{X_A}{\varphi(v)}
                   V_N(z_{j,c})E_{s,c},\qquad
 R_j=\sum_c\sum_{v\in c}(R_{1,v}^{\nu_j}+R_v^{\nu_j}),\\
 Q_j&=M_j+\nu_j E_j+\nu_j R_j,\qquad
 (S_0,S_1,S_2,S_3)=(T_0,T_1,T_2,B_3),\\
 Z_{\mathrm{AA}}&=\sum_{j=0}^3\nu_j(Q_j-S_j),\qquad
 M_{\mathrm{prin}}=\sum_{j=0}^3\nu_jM_j,\\
 M_{\mathrm{err}}&=\sum_{j=0}^3E_j,\qquad
 R_{\mathrm{adv}}=\sum_{j=0}^3R_j.
\end{aligned}
\end{equation}
The sums keep strict prefix order.
Under the stated sieve hypotheses,
\begin{equation}\label{eq:l3-signs}
 Q_0\ge S_0,\quad Q_1\le S_1,\quad
 Q_2\ge S_2,\quad Q_3\le S_3,\quad Z_{\mathrm{AA}}\ge0.
\end{equation}
Multiplying the signs gives
\begin{equation}\label{eq:l3-aa-id}
 \sum_{j=0}^3\nu_jS_j
    =M_{\mathrm{prin}}+M_{\mathrm{err}}+R_{\mathrm{adv}}-Z_{\mathrm{AA}}.
\end{equation}
Every remainder in $R_{\mathrm{adv}}$ has sign $+1$.

\subsection{The natural coprime comparator}

A block $\mathfrak b=(j,c,f,\mathbf D)$ fixes a row,
a cell, and either remainder type $f$.
The tail vector is empty for $R_1$.
Prefix realizations stay inside the coefficient.
Let $K_P,K_T$ impose prefix and tail distinctness.
Let $C_{\mathfrak b}$ be the unrestricted labelled coefficient.
Define
\begin{equation}\label{eq:l3-masks}
\begin{aligned}
 C_{\mathfrak b}(q)&=\sum_{\mathrm{full\ labelled\ tuples}}
          \widetilde\lambda_d1_{q=d\prod p_i},\\
 O_{\mathfrak b}(q)&=\frac1{M_c}
      \sum_{\mathrm{full\ labelled\ tuples}}
          \widetilde\lambda_dK_PK_T1_{q=d\prod p_i},\\
 \Xi_{U,\mathfrak b}(q)&=\frac1{M_c}
      \sum_{\mathrm{full\ labelled\ tuples}}
          \widetilde\lambda_d(1-K_PK_T)1_{q=d\prod p_i},\\
 O_{\mathfrak b}&=C_{\mathfrak b}/M_c-\Xi_{U,\mathfrak b},\qquad
 \Xi_{A,\mathfrak b}=O_{\mathfrak b}1_{q\ \mathrm{not\ squarefree}}.
\end{aligned}
\end{equation}
The two comparator discrepancies are
\begin{equation}\label{eq:l3-comparator}
 \Delta_P(q)=C_A(q;N)-\frac{Y_A(q)}{\varphi(q)},\qquad
 d_A(q)=\Delta_0(q)-\Delta_P(q)
             =-\frac{X_A-Y_A(q)}{\varphi(q)}.
\end{equation}
For $\langle f,g\rangle=\sum_qf(q)g(q)$, put
\begin{equation}\label{eq:l3-adv-split}
\begin{aligned}
 E_{\mathrm{disp},f}&=\sum_{\mathfrak b:\mathrm{type}=f}
                \frac{\langle C_{\mathfrak b},\Delta_P\rangle}{M_c}
                     \quad(f=1,2),\\
 E_{\mathrm{sq}}&=\sum_{\mathfrak b}
       \{-\langle\Xi_{U,\mathfrak b},\Delta_P\rangle
                         +\langle\Xi_{A,\mathfrak b},d_A\rangle\},\\
 E_{\mathrm{nc}}&=\sum_{\mathfrak b}\sum_{q\ \mathrm{squarefree}}
                              O_{\mathfrak b}(q)d_A(q),\\
 R_{\mathrm{adv}}&=E_{\mathrm{disp},1}+E_{\mathrm{disp},2}
                                      +E_{\mathrm{sq}}+E_{\mathrm{nc}}.
\end{aligned}
\end{equation}
This is a finite coefficient identity
\cite[(4.22)--(4.26), Appendix~B.1--B.4]{PaperII}.
All moduli satisfy $(q,2N)=1$.
That restriction does not puncture $A$.

\subsection{The full-support bilinear application}

Split each prime coordinate into relative-width $L^{-334}$ intervals.
An interior box has constant order and descriptor predicates.
For $B=\prod_i[A_i,B_i)$, put
\begin{equation}\label{eq:l3-direct-prime}
\begin{gathered}
 T=\lceil A_1\rceil-1,\quad U_J=\lceil B_1\rceil-1,\quad
 I_T=(T,U_J],\quad X=\prod_iA_i,\\
 a_n=1_{\mathbb P}(n)1_{I_T}(n),\qquad
 b_m=\sum_{\substack{p_2<\cdots<p_k\\p_i\in[A_i,B_i)}}1_{m=p_2\cdots p_k},\\
 b_m\in\{0,1\},\qquad
 N^\alpha/2\le T<N^{1/5},\qquad \varepsilon N\le X\le N.
\end{gathered}
\end{equation}
The box contributes exactly $w_Ca_nb_m$.
For $T>2^{50}$ and $U_J\le2T$,
\begin{equation}\label{eq:l3-prime-option}
 a_n=1_{P^-(n)\ge T^{51/100}}1_{I_T}(n).
\end{equation}
A composite $n\le2T$ has a smaller prime divisor.
This is Paper~I, (3.8), option 1
\cite[(3.8)--(3.9)]{PaperI}.
Its logarithmic replacement error is zero.

Define the full-support discrepancy
\begin{equation}\label{eq:l3-box-discrepancy}
\begin{aligned}
 \Delta_B(q,N)&=\sum_{nm\equiv N\ (q)}a_nb_m
           -\frac1{\varphi(q)}\sum_{(nm,q)=1}a_nb_m,\\
 \mathcal S_B(\mathfrak b)&=\sum_qC_{\mathfrak b}(q)\Delta_B(q,N).
\end{aligned}
\end{equation}
There is no condition $nm\le X$ in these sums.
The exact interior decomposition is
\begin{equation}\label{eq:l3-interior-split}
\begin{aligned}
 D_{\mathrm{int}}^{(j)}
 &=\sum_{\mathfrak b:\mathrm{row}=j}\frac1{M_c}
           \sum_C\sum_{B\subset C\ \mathrm{interior}}w_C\mathcal S_B(\mathfrak b),\\
 E_{\mathrm{disp},1}+E_{\mathrm{disp},2}
 &=\sum_{j=0}^3D_{\mathrm{int}}^{(j)}+D_{\mathrm{bdry}}.
\end{aligned}
\end{equation}

For rows $0,1,2$, apply Fact~\ref{fact:bv-full-support}.
Its Siegel--Walfisz sequence must be the selected prime sequence.
The dictionary with Fact~\ref{fact:bv-full-support} is
$b_n^{\mathrm{F2}}:=a_n$, $a_m^{\mathrm{F2}}:=b_m$, and $Y:=T$.
Thus the prime sequence carries the Siegel--Walfisz hypothesis.
We need Paper~I, (3.25), in the form
\begin{equation}\label{eq:l3-prime-sw}
 \sum_{\substack{n\in I\\n\equiv l\ (k)\\(n,d)=1}}1_{\mathbb P}(n)
 -\frac1{\varphi(k)}\sum_{\substack{n\in I\\(n,dk)=1}}1_{\mathbb P}(n)
 \ll_{A_{\mathrm{SW}}}T\tau(d)L^{-A_{\mathrm{SW}}},\qquad(l,k)=1,
 \quad I\subset I_T.
\end{equation}
The normalized hypothesis also needs
\begin{equation}\label{eq:l3-prime-norm}
 \|a\|_2T^{1/2}\gg T L^{-168}.
\end{equation}
Use Fact~\ref{fact:prime-estimates}, specifically \eqref{eq:classical-sw}.
Its endpoint estimate gives, for every fixed $A>335$,
\begin{equation}\label{eq:l3-endpoint-prime-count}
 \pi(U_J)-\pi(T)=\operatorname{Li}(U_J)-\operatorname{Li}(T)
                         +O_A(TL^{-A}).
\end{equation}
Full mesh intervals satisfy $U_J-T\gg TL^{-334}$.
Thus $\#(I_T\cap\mathbb P)\gg TL^{-335}$ uniformly.
This implies \eqref{eq:l3-prime-norm}.
The $k=1$ discrepancy in \eqref{eq:l3-prime-sw} is identically zero.
The counting input is the endpoint estimate just displayed.
\paragraph{Proof of \eqref{eq:l3-prime-sw} and \eqref{eq:l3-prime-norm}.}
Let $N^\alpha/2\le T<N^{1/5}$, $U_J\le2T$, $(l,k)=1$, and let
$I=(y_1,y_2]\subset I_T$ be an arbitrary subinterval, so
$T\le y_1\le y_2\le U_J$. (Paper~I, (3.25), is quantified over all
such $I$ \cite[(3.25)]{PaperI}.)
The condition $(n,d)=1$ removes at most $\omega(d)$ primes from each sum.
If $k\le L^{A+2}$, apply \eqref{eq:classical-sw} at $y=y_2$ and $y=y_1$,
with $C=A+3$. Since $T\le y\le2T$, $\log y\asymp L$ and the error is
$O_A(TL^{-A})$ at each endpoint. Hence
$\pi(y_2;k,l)-\pi(y_1;k,l)=(\operatorname{Li}(y_2)-\operatorname{Li}(y_1))/\varphi(k)+O_A(TL^{-A})$.
The same holds for the unrestricted count divided by $\varphi(k)$,
up to the $\omega(k)$ primes dividing $k$.
If $k>L^{A+2}$, since $y_2-y_1\le T$ both terms are at most
$T/k+1+T/\varphi(k)\ll TL^{-A}$,
using $\varphi(k)\gg k/\log\log k$. This proves \eqref{eq:l3-prime-sw}
with $\tau(d)$, ineffectively, uniformly in $I\subset I_T$. For the norm,
$U_J-T\ge A_1L^{-334}-1\ge TL^{-334}/2$ on a full mesh interval, and
\eqref{eq:l3-endpoint-prime-count} with $A=336$ gives
$\#(I_T\cap\mathbb P)\ge TL^{-335}/4$ eventually. Hence
$\|a\|_2T^{1/2}\ge TL^{-335/2}/2\ge TL^{-168}$. Both bounds are uniform
over all interior boxes, since only $T$ enters.

Assume \eqref{eq:l3-prime-sw} and \eqref{eq:l3-prime-norm} uniformly.
Taking $A_{\mathrm{SW}}=A_0+168$ verifies
\eqref{eq:bvfs-sw} with divisor exponent one.
Both factor lengths exceed $x_B^{1/12}$ eventually.
Split each of the $k-1$ complementary prime intervals at any dyadic
boundary. This gives at most $2^{k-1}$ pieces. Localize each resulting
product on $m\sim M_B$. Since $(1+L^{-334})^{k-1}<2$ eventually,
there are at most two product dyadic intervals per piece.
This bounded splitting changes no logarithmic exponent.
Set $x_B=TM_B\asymp N$.
There are at most $T$ selected integers in $(T,2T]$ and at most
$2M_B$ integers in a product dyadic interval. Unique factorization gives
$\|a\|_2\|b\|_2\le\sqrt{2x_B}$.
The required application is
\begin{equation}\label{eq:l3-bv-application}
\begin{gathered}
 \frac{|C_{\mathfrak b}(q)|}{M_c}\le C_{\mathrm{L3},\sigma},\qquad
 Q<N^{1/2-9\cdot10^{-5}}\le x_B^{1/2}(\log x_B)^{-B_{\mathrm{BV}}},\\
 \|a\|_2\|b\|_2\le\sqrt{2}\,x_B^{1/2},\qquad
 |\mathcal S_B(\mathfrak b)|/M_c\ll_\sigma NL^{-3692}.
\end{gathered}
\end{equation}
Use $A=3692$, $B=1$, $\epsilon=1/12$ and
$B_{\mathrm{BV}}=B_{\mathrm{BV}}(3692,1,1/12)$ in \eqref{eq:bvfs}.
The source has $O_\sigma(L^{3689})$ weighted boxes.
There are finitely many blocks at fixed $\sigma$.
Since $W\gg N/L^2$, summation gives
\begin{equation}\label{eq:l3-bv-total}
 \frac{|D_{\mathrm{int}}^{(0)}|+|D_{\mathrm{int}}^{(1)}|
          +|D_{\mathrm{int}}^{(2)}|}{W}\ll_\sigma L^{-1}.
\end{equation}
\paragraph{Verification of \eqref{eq:l3-bv-application}.}
The coefficient bound is \eqref{eq:l3-cap-coefficient-bound}, proved above,
and the modulus range is \eqref{eq:l3-cap-margins}. Every supported $q$ is
prime to $2N$, so $a=N$ is a reduced class and the maximum in
\eqref{eq:bvfs} covers it. The selected sequence $a_n$ is supported in
$(T,2T]$, and $b_m$, after at most two dyadic pieces, in $(M_B,2M_B]$.
For $x_B=TM_B$, $x_B\in[\varepsilon N/8,2N]$ and
$\log T/\log x_B\in[\alpha-o(1),1/5+o(1)]\subset[1/12,11/12]$, because
$\alpha-1/12=17/4596>0$. The hypothesis \eqref{eq:bvfs-sw} holds with
$B=1$ by \eqref{eq:l3-prime-sw}, \eqref{eq:l3-prime-norm} and
$A_{\mathrm{SW}}=A_0+168$. Interior boxes have constant order predicates,
so $nm=e$ has one labelled representation and $b_m\in\{0,1\}$. Thus
$\|a\|_2\|b\|_2\le(2x_B)^{1/2}$. On each interior box,
$\mathcal S_B(\mathfrak b)=\sum_qC_{\mathfrak b}(q)\Delta_B(q,N)$ with
$\Delta_B=\sum_{\mathrm{pieces}}\Delta_{\mathrm{fs}}$, and
Fact~\ref{fact:bv-full-support} gives
$|\mathcal S_B(\mathfrak b)|/M_c\ll_\sigma NL^{-3692}$.
Fact~\ref{fact:bv-full-support} is the full-support bilinear estimate
derived from Bombieri--Friedlander--Iwaniec, Theorem~0(b). We verified its hypotheses
above for this source.

\subsection{The source-level row-three estimate}\label{subsec:l3-source-row3}

Row three keeps the full anatomy level.
Its factorization keeps every labelled occurrence.
The smooth divisor stays in the first factor.
For $d_*=\rho(\Theta_c-\eta_L)$ and $e_i=\rho t_i$,
the prefix inequalities are
\begin{equation}\label{eq:l3-prefix-margins}
\begin{aligned}
 d_*-2e_1&\ge\frac{33109}{160000}\rho,\\
 d_*-e_1-2e_2&\ge\frac{1859}{160000}\rho,\\
 d_*-e_1-e_2-2e_3&\ge\frac{33109}{480000}\rho
                                      \quad(\chi_c=1).
\end{aligned}
\end{equation}
The anatomy branch gives $\Theta_c\ge(1+t_1)/2$.
The condition $\chi_c=1$ gives the last inequality.
In particular, with $A_3=33109/160000$,
\begin{equation}\label{eq:l3-row3-global-margins}
\begin{gathered}
 \ell_c\ge\max\{A_3-t_3,2\rho t_3\}\ge\frac{33109}{240000},\qquad
 \ell_c-\ell_0\ge\frac{231463}{24720000}=0.0093633899676\ldots,\\
 \Theta_c\le H_{\max}=\frac{153}{256},\qquad
 e_Q=(\rho+\sigma)(153/256-\eta_L)<\frac35,\\
 \frac35-e_Q\ge\frac{391}{160000}
       -\frac{95609}{160000}(\sigma+\tau)>0.0024437499.
\end{gathered}
\end{equation}
Here $e_Q$ bounds the full modulus exponent
$\rho(t_1+t_2+t_3)+(\rho+\sigma)\ell_c$.
The inherited anatomy and rough-factor estimates are Paper~II,
(4.20) and Appendix~B, Lemma~B.1(1)--(2)
\cite[(4.20), Appendix~B, Lemma~B.1(1)--(2)]{PaperII}.
For the tail exponents $s_i$, the cubic conditions give
\begin{equation}\label{eq:l3-tail-margins}
 \sum_{i<l}s_i+2s_l<\ell_c,\qquad
 h=d_*+60\sigma,\qquad
 \Theta_c-h\ge\eta_L-\tau H_{\max}-60\sigma>0.
\end{equation}
Theorem~4.1 of Paper~I then supplies the load surpluses
\begin{equation}\label{eq:l3-packing-loads}
 \left(\frac{7527}{256},\ 60,\ \frac{114421}{1280},\
                         \frac{26247}{256}\right)\sigma.
\end{equation}
The required length data are
\begin{equation}\label{eq:l3-lengths}
\begin{gathered}
 1\le|mn-N|\le N<X^3,\quad
 N_I\ge\frac{T}{6L^{334}}\ge(2T)^{1-2^{-126}},\\
 X^{60\sigma}\le T\le X^{1/3+15\sigma},\quad
 QX^{\sigma/3}\le N_I,\quad QRS<M_IX^{-89\sigma},\quad M_I\ge X^{1/2}.
\end{gathered}
\end{equation}
Here $N_I$ is the selected integer interval length.
The complement has length scale $M_I$.
The letters $Q,R,S$ are the three factor scales.
Each condition holds only above its fixed-$\sigma$ application threshold.
The floor $2\lceil e^{2^{150}}/2\rceil$ is imposed separately.
See Paper~I, (3.6)--(3.9) and (4.9)
\cite[(3.6)--(3.9), Theorem~4.1, (4.9)]{PaperI}.
We verify them for this source below.

The remaining analytic estimate is
\begin{equation}\label{eq:l3-masked-use}
 |\mathcal S_B(\mathfrak b)|
 \ll_{\sigma,A_{\mathrm{loc}}}XL^{-A_{\mathrm{loc}}}
          +X^{1-\ell_0\sigma^2/4}L^{C_{\mathrm{loc}}}
 \quad(\mathfrak b\text{ in row }3).
\end{equation}
It is uniform over the permitted boxes and blocks.
Here $A_{\mathrm{loc}}>0$ is arbitrary;
$C_{\mathrm{loc}}$ is fixed by the local theorem.
Paper~I, Theorem~3.2, (3.13), supplies the local dispersion theorem.
Its Theorem~2.8 enters only in the inner spectral stage.
There we use the form (P1.7), Fact~\ref{fact:spectral-implication}, with the full bracket
\eqref{eq:t28-main-bound}. The inner substitution is
\begin{equation}\label{eq:l3-spectral-dictionary}
 r_T=z,\quad s_T=p_0,\quad n_T=p_0m,\quad
 a_T=\epsilon_yN\quad(\epsilon_y=\pm1),\quad
 A_T=X^{1+\sigma/20}.
\end{equation}
The auxiliary $z$ and $p_0$ are those of Paper~I's dispersion proof.
Both signs and the mixed cusp pair are kept.
Let $\chi$ be smooth, zero below zero and one above one.
For the fixed bump $\psi_0$, define
\begin{equation}\label{eq:l3-nine-weights}
\begin{aligned}
 \psi_{0,1}(v)&=\psi_0(v)\{1-\chi((v-0.8)/0.1)\},\\
 \psi_{0,2}(v)&=\psi_0(v)\chi((v-0.8)/0.1)
                                  \{1-\chi((v-1.4)/0.2)\},\\
 \psi_{0,3}(v)&=\psi_0(v)\chi((v-1.4)/0.2).
\end{aligned}
\end{equation}
They sum to $\psi_0$.
Their strict dyadic windows are
$(0.45,0.90]$, $(0.80,1.60]$, and $(1.40,2.80]$.
The two smooth modulus coordinates use all nine products
$\psi_{0,i}(b/D)\psi_{0,j}(c/C)$, $1\le i,j\le3$.
Multiply each by the fixed cutoffs on the supported $n_T,z,p_0$.
These cutoffs equal one on the actual coefficient support.
The derivative hypothesis is \eqref{eq:t28-seminorms} with
$\epsilon_0=\sigma/2000$ and $\epsilon=\sigma/100$.

Write $\theta_{\mathrm{sp}}=7/32$ and
$V_{\mathrm{sp}}=T^4R^4S^6$.
The coefficient in (P1.7) is zero outside the actual indices
$n_T=p_0m$, $r_T=z$, $s_T=p_0$.
Use its supported norm
\begin{equation}\label{eq:l3-supported-norm}
\begin{gathered}
 \sum_{n''\mid N^\infty}(n'')^{\theta_{\mathrm{sp}}}
                \|\widetilde b(n'')\|_2^2,\\
 \|\widetilde b(n'')\|_2^2
 =\sum_{\substack{n',r,s\\(n',N)=1}}|b_{n'n'',r,s}|^2.
\end{gathered}
\end{equation}
Here $b_{n,r,s}$ is the coefficient of (P1.7).
It is distinct from the complementary sequence $b_m$.
The norm includes no coefficients outside that support.
After substitution, the four bracket contributions divided by
$V_{\mathrm{sp}}$ are bounded by
\begin{equation}\label{eq:l3-four-spectral-terms}
\begin{aligned}
 K_A&=A_T^{\theta_{\mathrm{sp}}}\frac{TQR^2S^5}{X^2},
 &K_2&=\frac{A_T}{X}K_A,\\
 K_R&=\left(\frac{TR}{S}\right)^{2\theta_{\mathrm{sp}}}
                              \frac{TQR^2S^5}{X^2},
 &K_4&=\frac{TQR^2S^5}{X^2}\frac{T^2R^2}{X}.
\end{aligned}
\end{equation}
Only uniform $X^{o(1)}$ factors and powers of $p_0$ are suppressed.
These are precisely Paper~I, (3.81), (3.85), and (3.87)
\cite[(3.69), (3.81)--(3.87)]{PaperI}.
Each term carries the actual-use hypotheses.
\subsection{The row-three local dispersion estimate}\label{subsec:l3-local-row3}

The source-wide estimate also has boundary and comparator contributions.
The local argument gives only its interior row-three block.

We treat one row-three block with $\chi_c=1$ and one full
interior prime-coordinate mesh box. We do not assert the source-wide
boundary, comparator, or row-zero-to-two estimates.
Fix $\sigma>0$, $\tau=\sigma^9$, $\rho=1+\tau$, and a retained
closed clipped parent $c$, used only for the displayed infimum and endpoint
checks. The actual prime cells are half-open. Keep its inherited labels
$t_1\ge t_2\ge t_3$,
$\Theta_c=\inf_c\widetilde H$, $\ell_c=\Theta_c-\eta_L-t_1-t_2-t_3$,
$D_c=N^{\ell_c}$, $z_c=N^{t_3}$ and $u_c=D_c^{\sigma^2}$.
On the whole block $\chi_c=1$. The inequalities
\eqref{eq:l3-prefix-margins}--\eqref{eq:l3-row3-global-margins}
apply with the inherited labels, including clipped endpoints.
Here an interior source box is a full relative-width
$x=(\log N)^{-334}$ mesh box whose closure lies on one side of every
source, order and descriptor face. A clipped terminal mesh fragment
belongs to the boundary contribution.
LEN is not applied to it.

The prefix constants can be checked without an unclipped replacement.
Put $A_3=33109/160000$. Since $t_1,t_2\le25/128$ and
$\Theta_c\ge(1+t_1)/2$, the first two gaps are at least
$\rho(1/2-3(25/128)/2-\eta_L)=\rho A_3$ and
$\rho(1/2-5(25/128)/2-\eta_L)=\rho\,1859/160000$.
For the third, $\ell_c-t_3\ge A_3-2t_3$ and retention gives
$\ell_c-t_3\ge(2\rho-1)t_3$; their maximum is at least $A_3/3$.
Likewise $\ell_c\ge\max\{A_3-t_3,2\rho t_3\}\ge2A_3/3>\ell_0$.
Thus all three prefix WF inequalities have strict room. The whole
anatomy infimum is taken before this branch bound is used.

\paragraph{Exact coefficient and labelled factorization (H1--H2).}
Write $\mathfrak b=(3,c,f,\mathbf D)$. Give the three prefix
occurrences their independent prime sets
$[N^{t_i},N^{\rho t_i})\cap[N^u,N^{\mu_*})\cap\mathcal P_N$.
For $f=2$, append each tail occurrence, independently, with prime set
$[D_j,\min(D_j^\rho,z_c))\cap\mathcal P_N$.
For $f=1$ the tail list is empty. In both cases put $J$ equal to this
complete labelled list, and set $\lambda_d=\lambda_d^-(\mathbf D)$
or $\phi_d^-$, respectively. Thus, exactly,
\[
 C_{\mathfrak b}(q)=\sum_{\substack{d<D_c^\sigma\\d\mid P_N(u_c)}}
 \lambda_d\sum_{(p_j)_{j\in J}}1_{q=d\prod_{j\in J}p_j},
 \qquad |\lambda_d|\le1 .
\]
Neither distinctness mask is inserted here. The orbit divisor $M_c$,
the source weight $w_C$, and the two distinctness corrections remain
outside this coefficient, as in \eqref{eq:l3-masks}--\eqref{eq:l3-adv-split}.
In particular its supports include repeated-prime moduli, and
$(q,2N)=1$. The prime sets are the actual clipped sets, not their
unclipped parents. All labelled primes are at least $u_c\ge u_0$
with $u_0=N^{\ell_0\sigma^2}$, since $\ell_c>\ell_0$ and
$\ell_c\sigma^2<u$.

For packing use the actual upper exponents of the clipped prime sets.
The first three are at most $\rho t_i$; all tail upper exponents are
at most $t_3$. The prefix masks in \eqref{eq:l3-prefix-margins}
and the tail masks $\sum_{i<j}s_i+2s_j<\ell_c$ show that their
merged decreasing list satisfies $\mathrm{WF}(d_*)$, where
$d_*=\rho(\Theta_c-\eta_L)$. For a tail rank, this follows after
adding at most $\rho(t_1+t_2+t_3)$ to its tail inequality.
The top three upper exponents are the prefix endpoints, belonging to
the closed clipped parent. Hence their anatomy value is at least
$\Theta_c$. Put $h=d_*+60\sigma$, $\delta=30\sigma$ and
$g=\log_N T-\delta$.
Every point $\mathbf x$ of the closed clipped parent satisfies
$u\le x_1\le\mu_*=25/128$. On this interval the $v_{\rm pk}(x_1)$ branch of
Paper~I gives $\widetilde H(\mathbf x)\ge v_{\rm pk}(x_1)=(1+x_1)/2$;
hence $\Theta_c\ge(1+u)/2$. Therefore
 $h\ge\Theta_c-\eta_L+60\sigma
 \ge\frac{1+u}{2}-\eta_L>\frac12$; also
 $h<\Theta_c\le153/256$, since
$\Theta_c-h\ge\eta_L-\tau H_{\max}-60\sigma>0$.
Theorem 4.1 of Paper I applies to this same labelled list, with
$\delta\le g\le1/3-\delta/2$. Let its partition be
$J=J_1\sqcup J_2\sqcup J_3$, with the smooth divisor attached to
slot one. Define, using precisely the preceding prime sets,
\[
 \gamma_1(v)=\frac1{|J_1|!}\sum_{d,(p_j)_{J_1}}
             \lambda_d1_{v=d\prod_{J_1}p_j},\qquad
 \gamma_s(v)=\frac1{|J_s|!}\sum_{(p_j)_{J_s}}
             1_{v=\prod_{J_s}p_j}\quad(s=2,3).
\]
Empty slots two and three mean the atom $1_{v=1}$. Slot one always
keeps the $d$ sum, even if $J_1$ is empty. The finite tuple map gives
\[
 C_{\mathfrak b}=C_{\rm slot}(\gamma_1*\gamma_2*\gamma_3),\qquad
 C_{\rm slot}=\prod_s|J_s|!,\qquad |\gamma_s|\le1.
\]
The primes below $u_c$ determine $d$ uniquely. At most $|J_s|!$
assignments have a fixed slot product. The count allows repeated primes.
The rough coordinate count is less than $3/(5\ell_0\sigma^2)$,
so $C_{\rm slot}=O_\sigma(1)$ and
$|C_{\mathfrak b}(q)|/M_c\le M_\sigma!$, with
$M_\sigma=\lceil3/(5\ell_0\sigma^2)\rceil$.
For the tested boxes below, take $C_{\rm ctr}=2$. Uniformly for
sufficiently large $N$,
\[
 \frac{N_I}{TL^{-334}}\in[1/2,2],\qquad
 \frac{M_I}{(X/T)L^{-334}}\in[2,20].
\]

\paragraph{The same tested box (H3--H5).}
Use exactly \eqref{eq:l3-direct-prime} and
\eqref{eq:l3-box-discrepancy}: $X=\prod_iA_i$, $T=\lceil A_1\rceil-1$,
$U_J=\lceil B_1\rceil-1$, and
$\alpha_n=1_{\mathbb P}(n)1_{T<n\le U_J}$.
No product truncation is added to $\Delta_B$. Take $C_{\rm box}=334$;
$\log X=\log N+O(1)$ throughout this family.
For sufficiently large $N$, $T>2^{50}$ and $U_J\le2T$, so
\[
 \alpha_n=1_{P^-(n)\ge T^{51/100}}1_{T<n\le U_J}.
\]
Thus (3.8), option $1$, holds with $\zeta=T^{51/100}\ge z_0(X)$.
The complementary sequence is the actual $b_m\in\{0,1\}$ of
\eqref{eq:l3-direct-prime}; in particular $|b_m|\le\tau_{23}(m)$.
There is no Heath--Brown expansion, logarithm freezing, or proper-power
replacement error for this selected direct-prime representation.
The selected-factor Siegel--Walfisz bound is Paper I (3.25), uniform
in this stronger cutoff and every endpoint and auxiliary coprimality.
Its relative-norm alternative is kept as in (3.48). We assume no lower
bound on any coprimality-truncated selected norm.
In the large-norm case of (3.48), use Lemma~\ref{hyp:ra-ms-transfer}(ii).
Its premises
(3.1), (3.2), (3.6), (3.8) with LEN, (3.11), (3.25), (3.26) and (3.27)
are the data verified in this subsection.

Let $M=\prod_{i=2}^kA_i$, $M_+=\prod_{i=2}^kB_i$,
$M_I=M_+-M$, $m_0=(M+M_+)/2$, and
$\psi_I(m)=\phi_0((m-m_0)/M_I)$, with the fixed bump of Paper I (3.9).
Then $\psi_I=1$ on the full complementary product interval. Since
$5\le k\le11$ and the mesh boxes are full,
$M_I\asymp ML^{-334}$, $m_0\le2X/T$, and
$N_I=U_J-T\ge T/(6L^{334})$, where $L=\log N$.
For fixed $\sigma$ enlarge the threshold so that
$\sigma L\ge2^{130}(\log L+1)$; then
\[
 N_I\ge(2T)^{1-2^{-126}},\quad
 X^{60\sigma}\le T\le X^{1/3+15\sigma},\quad M_I\ge X^{1/2}.
\]
These follow from $N^\alpha/2\le T<N^{1/5}$ and
$\varepsilon N\le X\le N$. The integer convention is
$(T,U_J]$; no continuous interval length replaces $N_I$.

The doubled support lies in
$[M-M_I/2,M_++M_I/2]$. Since $k-1\le10$,
$(1+x)^{k-1}-1\le11x$ eventually. On this entire support and for
every selected integer, $mn\le N_{\rm cut}+25xN$. Consequently
\[
 1\le N-mn,\qquad N-mn\ge(L-25)NL^{-334}>0,\qquad
 |mn-N|\le N<X^3.
\]
This verifies nonresonance on the full doubled rectangle, including
both Cauchy--Schwarz copies. Checking only where $b_m\ne0$
would not suffice.

After localization, $Q,R,S$ denote the scales of
$\gamma_1,\gamma_3,\gamma_2$, respectively. Theorem 4.1 gives, before
fixed dyadic rounding and base conversion, reductions
\[
 (7527/256,\ 60,\ 114421/1280,\ 26247/256)\sigma
\]
in $Q/T$, $T^2R^2S$, $A_0^{7/32}TQR^2S^5$ and
$(TR/S)^{7/16}TQR^2S^5$, respectively, relative to their bounds
$1,N,N^2,N^2$. Here $A_0=X^{1+\sigma/20}$.
The first slot includes the additional load $\sigma\ell_c\le\sigma H_{\max}$.
Each listed reduction strictly exceeds the required losses
$29,8,15,15$. Hence fixed dyadic constants and
$\log_N(N/X)=O(1/L)$ are absorbed for fixed $\sigma$.
In particular the exact requirements (3.6)--(3.7) hold:
\[
 Q<TX^{-29\sigma},\quad N_I\ge QX^{\sigma/3},\quad
 QRS<M_IX^{-89\sigma},\quad M_I\ge X^{1/2},
\]
\[
 T^2R^2S<X^{1-8\sigma},\quad
 A_0^{7/32}TQR^2S^5<X^{2-15\sigma},\quad
 (TR/S)^{7/16}TQR^2S^5<X^{2-15\sigma}.
\]
For $QRS<M_IX^{-89\sigma}$ use the positive fixed gap
$4/5-153/256$ and $M_I\gg N^{4/5}L^{-334}$.
The complete modulus, including $d$ and all labelled primes, obeys
\[
 q<N^{e_Q},\qquad e_Q=(\rho+\sigma)(153/256-\eta_L)<3/5.
\]
Because $X\ge\varepsilon N$, impose
$(3/5-e_Q)\log N\ge(3/5)\log(1/\varepsilon)$ to obtain
$q\le X^{3/5}$. Take the moving frequency $a=N$ and
$K_{\rm scale}=1$: eventually
$0<N\le X\log N\le A_0$ and every supported modulus is prime to $2a$.
This proves H1--H5 with the support relaxation of Paper I Theorem 3.2.

\paragraph{Application and restoration.}
Apply Paper I Theorem 3.2, (3.13), to this full unrestricted coefficient
and this full-support natural-coprime discrepancy. For every fixed
$A_{\rm loc}>0$ it gives exactly
\[
 |\mathcal S_B(\mathfrak b)|
 \ll_{\sigma,A_{\rm loc}}X(\log X)^{-A_{\rm loc}}
       +X^{1-\ell_0\sigma^2/4}(\log X)^{C_{\rm loc}}.
\]
Constants are uniform in these boxes and in blocks with $\chi_c=1$ after fixing
$\sigma$ and source constants. The threshold is ineffective.
Theorem 3.1's smaller abstract support class is not used.
The proof includes all local restorations (3.32)--(3.50):
with $\epsilon=\ell_0\sigma^2/8$ and $A_{\rm work}=A_{\rm loc}+4$,
\[
 XL^{-A_{\rm work}}+XL^Cz_0^{-1/2}+XL^C(Q/N_I)^{1/2}
 +XX^{\epsilon/2}L^C\{u_0^{-1/2}+(Q/N_I)^{1/2}\}.
\]
Keep both terms in braces. They come from restoring the five
cross-copy gcd failures of (3.30), the close-pair set, and the zero-mode
main-term cancellation. The spectral
estimate alone does not give them. The $z_0$ term joins the arbitrary logarithmic saving,
and the remaining terms have savings exceeding $\ell_0\sigma^2/4$.

Theorem 2.8 enters only at Paper I (3.69)--(3.90). In that inner stage,
with $p_0=\min\{p:p\nmid2N\}$, set
$r_T=z$, $s_T=p_0$, $n_T=p_0m$, $a_T=\varepsilon_yN$,
$N_T=8p_0H_*SF$ and $A_T=A_0$. Both signs are kept. The coefficient
is exactly $\beta_{z,m}^{\chi,\varepsilon_y}$ of (3.68) on these
indices and zero elsewhere. In particular
\[
 \sum_{n''\mid N^\infty}(n'')^{7/32}
       \|\widetilde b(n'')\|_2^2
 \ll X^{o(1)}(H_*S)^2F,\quad
 \|\widetilde b(n'')\|_2^2
 =\sum_{\substack{n',r,s\\(n',N)=1}}|b_{n'n'',r,s}|^2,
\]
with no coefficients added outside that support. The unweighted norm
has the same bound. These are the endpoint norm bounds (3.70)--(3.81).
When slot two is a unit, $S=1/2$ and every nonzero-frequency block is
empty. Otherwise (3.82)--(3.85) identify the exact phase, permit
$p_0\mid c$, and use all nine dyadic weights from
\eqref{eq:l3-nine-weights}, with $\epsilon_0=\sigma/2000$ and
$\epsilon=\sigma/100$. The full bracket of
\eqref{eq:t28-main-bound}, including its moving-frequency and mixed-cusp
terms, yields all four quantities in \eqref{eq:l3-four-spectral-terms}.
No bracket term is discarded. The full bounds (3.87)--(3.90) then
control the inner spectral term. After that, the restoration terms
above and $C_{\rm slot}$ are put back.

Finally, this estimate concerns $C_{\mathfrak b}\Delta_B$.
The outer identity is still
\[
 R_{\rm adv}=E_{\rm disp,1}+E_{\rm disp,2}
             +E_{\rm sq}+E_{\rm nc}.
\]
Thus $M_c^{-1}$ and $w_C$ enter once. The forbidden and allowed-repeat
masks remain in $E_{\rm sq}$. The natural-to-$X_A$ comparator remains
in $E_{\rm nc}$. The source cutoff, peel, boundary and terminal-mesh
restorations, capped rows zero to two, and complete-family coefficient
counts require their own estimates. This local estimate alone
gives neither the complete L3 bound nor the RSW bound.

The outer application of Theorem~3.2 above supplies \eqref{eq:l3-masked-use}.
The scoped spectral implications in Section~\ref{sec:classical-inputs}
do not by themselves give it.

\subsection{The complete box family}

The required complete-family bounds are
\begin{equation}\label{eq:l3-boundary}
 \#\mathcal E_{\mathrm{bdry}}\ll NL^{-334},\qquad
 D_{\mathrm{bdry}}\ll_\sigma NL^{-331/2},\qquad
 \sum_B|w_B|\ll_\sigma L^{3689}.
\end{equation}
There are at most $122$ face families per $k$.
Their integer normals have maximum coefficient five.
Their coefficient sums are at most eleven.
\paragraph{Proof of the complete-family bounds.}
Use the unclipped geometric mesh $[A^{(j)},A^{(j+1)})$,
$A^{(j)}=(N^\alpha/2)(1+L^{-334})^j$, $j\ge0$, in every coordinate.
Boxes with a coordinate $\ge N^{1-4\alpha}$ carry no source mass, since
$k\ge5$ and every prime is at least $N^\alpha$. So there are no terminal
fragments, and every interior box is a full mesh box.
\emph{Box count.} Each coordinate has at most
$((1-\alpha)L+1)(L^{334}+1)+1\le L^{335}$ mesh intervals for $L\ge24$.
With $k\le11$, finitely many cells and $w_C\le36$,
$\sum_B|w_B|\le36|\mathcal C_+|\sum_{k=5}^{11}L^{335k}\ll L^{3685}\le L^{3689}$.
\emph{Boundary count.} Let $h'=\log(1+L^{-334})/L\le L^{-335}$.
A source integer $e$ in a boundary box has its vector within $h'$
in each coordinate of a face. A radial face forces $e$ into
$[y(1+L^{-334})^{-11},y(1+L^{-334})^{11}]$ with $y\in\{\varepsilon N,N_{\mathrm{cut}}\}$.
There are at most $2(23NL^{-334}+1)$ such $e$.
Every other face is $c_ix_i+c_jx_j=\theta$ with integers $|c_i|,|c_j|\le5$,
$c_i\ne0$ \cite[(5.29a)--(5.31) and the facet discussion after (5.40)]{PaperII}.
Write $e=p_im'$. For fixed $m'$, $p_i$ lies in an interval of relative
width at most $21L^{-334}$ below $N/m'$, which contains at most
$21L^{-334}N/m'+1$ integers. Here $m'$ is a product of at most ten primes
in $[N^\alpha,N)$, and $m'<N^{1-\alpha}$. By Mertens,
$\sum_{m'}1/m'\le\exp(\sum_{N^\alpha\le p<N}1/p)\le e/\alpha<32$ eventually.
Hence one face family contributes at most $700NL^{-334}+N^{1-\alpha}$.
The number of face families is finite ($\le122$ for each $k$). Thus
$\#\mathcal E_{\mathrm{bdry}}\le C_{\mathrm{bd}}NL^{-334}$ with an absolute
$C_{\mathrm{bd}}$. No numerical value is required below.
\emph{Boundary discrepancy.} Let $A_{\mathrm{bd}}$ be $A$ restricted to
boundary boxes. Then
$D_{\mathrm{bdry}}=\sum_{\mathfrak b}M_c^{-1}\langle C_{\mathfrak b},\Delta_{P,\mathrm{bd}}\rangle$.
By $|C_{\mathfrak b}|/M_c\le C_{\mathrm{L3},\sigma}$,
$\sum_q1_{q\mid N-e}\le\tau(N-e)$, Cauchy--Schwarz,
$\tau^2\le\tau_4$, $\sum_{n\le N}\tau_4(n)\le N(1+L)^3$ and
$\sum_{q\le N}1/\varphi(q)\le2(1+L)$,
\[
 |D_{\mathrm{bdry}}|\ll_\sigma36(\#\mathcal E_{\mathrm{bdry}})^{1/2}(N(1+L)^3)^{1/2}
             +36\,\#\mathcal E_{\mathrm{bdry}}(1+L)\ll_\sigma NL^{-331/2},
\]
with finitely many blocks at fixed $\sigma$. This proves
\eqref{eq:l3-boundary}, with the first bound in the form
$\#\mathcal E_{\mathrm{bdry}}\ll NL^{-334}$.
Under \eqref{eq:l3-masked-use}, take $A_{\mathrm{loc}}\ge3692$.
The $L^{3689}$ box count and $W\gg N/L^2$ then leave $L^{-1}$.
Equations \eqref{eq:l3-bv-total} and \eqref{eq:l3-boundary} then give
\begin{equation}\label{eq:l3-disp-total}
 (R_{\mathrm{adv}})_+/W=o_\sigma(1),
\end{equation}
provided \eqref{eq:l3-repeated-cap} and the coprime-comparator bound hold.

\subsection{The principal term and its price}

\paragraph{(P2.8d) The source-independent profile comparison.}
We use Paper~II, (4.70)--(4.75), in the following form.
Put
\begin{equation}\label{eq:l3-profile}
\begin{aligned}
 v_1(s)&=\min\left\{\frac{1+s}{2},\frac{32-34s}{39},
                                      \frac{57}{64}-\frac32s\right\},\\
 v_*(s)&=v_1(s)-u,\qquad
 \psi_*(x,y)=\begin{cases}x,&x\ge y,\\-\infty,&x<y,\end{cases}\\
 \Theta_*(x,y,z)&=\max\{v_*(x),v_*(y),v_*(x+y),v_*(x+y+z),\\
 &\hspace{14mm}\psi_*(v_*(x+z),x+2y+z),
                    \psi_*(v_*(y+z),2x+y+z)\}.
\end{aligned}
\end{equation}
The subtraction of $u$ precedes both guards.
This is the same guard convention as
\hyperref[eq:app-six-branch]{equation~\ref*{eq:app-six-branch}}.
It agrees with Paper~II's zero guard because an unguarded branch is
positive, as explained immediately after that display.
Define
\begin{equation}\label{eq:l3-Gstar}
\begin{gathered}
 H_{\log}=\log(\mu_*/u),\qquad
 a_3(\mathbf t)=\frac{\Theta_*(\mathbf t)-t_1-t_2-t_3}{t_3},\\
 J_*=\int_u^{\mu_*}\frac{dt_1}{t_1}\int_u^{t_1}\frac{dt_2}{t_2}
                 \int_u^{t_2}\frac{1-f(a_3(\mathbf t))}{t_3^2}\,dt_3,\\
 R_{\mathrm{up}}=500(F(30)-1)+500H_{\log}(1-f(30))
                         +250H_{\log}^2(F(30)-1),\\
 G_*=\frac1{\mu_*}+R_{\mathrm{up}}+J_*.
\end{gathered}
\end{equation}
Extend $f$ by zero below two.
These are Paper~II, (4.70)--(4.75)
\cite[(4.70)--(4.75)]{PaperII}.
Their source-independent comparison uses
\begin{equation}\label{eq:l3-profile-margin}
 g=(u-\eta_L)-\frac92\mu_*\tau>0.001899>2\mu_*\tau.
\end{equation}
Hence a positive canonical lower term lies in a cell with $\chi_c=1$.
The capped rows still satisfy $s_c>30$.
Their residuals are bounded by the same $R_{\mathrm{up}}$.
Thus the cap changes neither $G_*$ nor its price.

Let $\mathcal P_\Pi$ be the signed prefix profile:
\begin{equation}\label{eq:l3-principal-profile}
\begin{aligned}
 \mathcal P_\Pi={}&\frac1uF(s_0)
 -\frac1u\sum_{c:j=1}\operatorname{vol}_c f(s_c)
 +\frac1u\sum_{c:j=2}\operatorname{vol}_c F(s_c)\\
 &-\sum_{c:j=3,\chi_c=1}\frac{\operatorname{vol}_c}{t_3}f(\ell_c/t_3),
 \qquad \operatorname{vol}_c=\int_c\prod_i\frac{dt_i}{t_i}.
\end{aligned}
\end{equation}
The required mass and comparison statements are
\begin{equation}\label{eq:l3-principal-estimates}
\begin{gathered}
 X_A\le\frac NL\{(1-\varepsilon)I_{\mathrm{HL}}+O(L^{-1})\},\\
 M_{\mathrm{prin}}=X_A\frac{2e^{-\gamma}C(N)}L
                            (\mathcal P_\Pi+o_\sigma(1)),\qquad
 \mathcal P_\Pi\le G_*+o_\sigma(1).
\end{gathered}
\end{equation}
The first bound is Paper~II, (5.34)--(5.42).
The last comparison uses (4.72)--(4.75)
\cite[(4.72)--(4.75), (5.34)--(5.42)]{PaperII}.
\paragraph{Proof of \eqref{eq:l3-principal-estimates}.}
\emph{Mass.} The cells $\mathcal C_+$ are Paper~II's cells with $w_C\ge1$,
and the radial range $\varepsilon N<e<N$ is Paper~II's. The outer-box
majorant $f^+$ of Paper~II, (5.34), counts labelled factorizations into
all primes in $[N^\alpha,N)$, and $c(\mathbf j)\ge w_C$ for the box
$\mathbf j$ containing the ordered factor vector of $e$, which lies in
$C_N$. Hence $A\le f^+$ pointwise directly (no comparison with the
intermediate $g$ is used). Paper~II, (5.34)--(5.42), then gives
$X_A\le X^+=\frac NL\{(1-\varepsilon)I_{\mathrm{HL}}+O(L^{-1})\}$
\cite[(5.32), (5.34), (5.42)]{PaperII}. In particular $X_A\le N/L$
eventually, since $I_{\mathrm{HL}}\le\overline H<1$.
\emph{Principal term.} As in Paper~II, Lemma~B.9,
$V_N(z)=(1+o(1))2e^{-\gamma}C(N)/\log z$ uniformly for $z\ge N^h$,
$h=1/1000$ \cite[(B.9) and its proof]{PaperII}. The proof of (B.9) gives
an exact factorization. Mertens' theorem gives the asymptotic. The
primes $p\mid N$ with $p\ge z$ change the product by $1+O(\omega(N)/z)$.
Every actual cutoff is at least $N^{u/\rho^3}>N^h$ \cite[Lemma~B.8]{PaperII}. At fixed $\sigma$ the cells are
finitely many, with fixed exponent boxes. Mertens gives
$\sum_{v\in c}1/\varphi(v)=\operatorname{vol}_c+o(1)$, the primes
dividing $2N$ costing $O(\omega(N)N^{-u})$. Also $K_j(s_c)$ is bounded.
Summing $\nu_jM_j$ gives
$M_{\mathrm{prin}}=X_A\frac{2e^{-\gamma}C(N)}L(\mathcal P_\Pi+o_\sigma(1))$,
with $\log_Nz_{3,c}=t_3$ supplying the factor $1/t_3$.
\emph{Profile comparison.} Write $F=1+(F-1)$, $f=1-(1-f)$.
On rows $0,1,2$, $s_c>30$, and $F$ decreases and $f$ increases. With
$\sum_{j=1}\operatorname{vol}_c=H_{\log}$ and
$\sum_{j=2}\operatorname{vol}_c=H_{\log}^2/2$, the residuals are at most
$R_{\mathrm{up}}$ \cite[(4.73)]{PaperII}. For row three, if
$f(a_3(\mathbf x))>0$ at an actual point $\mathbf x$ of a cell, then
\eqref{eq:l3-profile-margin} gives $\chi_c=1$ and
$\ell_c\ge\Theta_*(\mathbf x)-x_1-x_2-x_3$ \cite[(4.72) and the two
paragraphs after it]{PaperII}. For $x_3\in[t_3,\rho t_3]$, monotonicity of
$sf(s)$ at fixed numerator $\ell_c$ gives
$f(\ell_c/t_3)/t_3\ge f(\ell_c/x_3)/x_3\ge f(a_3(\mathbf x))/x_3$.
Integrating over each retained cell,
$\sum_{j=3,\chi_c=1}\operatorname{vol}_cf(\ell_c/t_3)/t_3
 \ge\int\!\!\int\!\!\int f(a_3)\,dt_1dt_2dt_3/(t_1t_2t_3^2)$.
The constant parts give $1/\mu_*$ by \cite[(4.74)]{PaperII}. Hence
$\mathcal P_\Pi\le1/\mu_*+R_{\mathrm{up}}+J_*+o_\sigma(1)=G_*+o_\sigma(1)$,
which is \cite[(4.75)]{PaperII} for this source.

Use the same coefficient bound as the exterior comparison:
\begin{equation}\label{eq:l3-certificate}
\begin{gathered}
 G_*\le G_0=\frac{6088440363}{10^9},\qquad
 2e^{-\gamma}G_0\le c_{K2,G0}
 =\frac{6836825163875517756286846321650754372931094689}{10^{45}},\\
 I_{\mathrm{HL}}\le\overline H=\frac{389253}{5000000}.
\end{gathered}
\end{equation}
The functional bound \(G_*\le G_0\) is Paper~II,
\cite[(4.73), (4.75), (4.79), (A.10a)]{PaperII}, for the identical
canonical functional \eqref{eq:l3-Gstar}. Its scalar rounding is the
same as \eqref{eq:k2-certificate}. The mass bound is Paper~II, (A.11)
\cite[(A.11)]{PaperII}.
The optional unshifted certificate in
\hyperref[subsec:app-g24]{the unshifted-comparison subsection of
Appendix~B} is outside this proof.

The constant $c_{\mathrm{O1}}$ is optional, non-certified appendix
material outside the L3 proof.
In the optional scalar comparisons below, $c$ denotes the rational
from \eqref{eq:intro-exterior-constant}.
No functional comparison with $c_{\mathrm{O1}}$ is needed here.
Its rational definition is
\begin{equation}\label{eq:l3-exact-certificate-constant}
 c_{\mathrm{O1}}=\frac{\sum_{i=0}^7n_iB^i}{\sum_{i=0}^7d_iB^i},
 \qquad B=10^{25},
\end{equation}
with the integers in Table~\ref{tab:l3-certificate-integers}.
\begin{table}[ht]
\caption{Exact certificate constant}\label{tab:l3-certificate-integers}
\small
\begin{tabular}{@{}rrr@{}}
\hline
$i$ & $n_i$ & $d_i$\\
\hline
 0 & 4790995369589327356095267 & 5033118760359274030825472\\
 1 & 6963256824825772577191322 & 8266496527677353149933429\\
 2 & 8996423827217290029713216 & 9138360880901719297497815\\
 3 & 4916164251930472277375479 & 4671424972330288696699300\\
 4 & 329572634628519073712385 & 939590108467331323910141\\
 5 & 9317769528669888650546075 & 5634060699913223990647374\\
 6 & 8280830054071768648850166 & 2330146450712536415611653\\
 7 & 389762593301325772 & 57030507720277440\\
\hline
\end{tabular}
\end{table}
The exact rational comparison of $c$ and $c_{\mathrm{O1}}$ is
in Section~\ref{sec:final-bill}, \eqref{eq:bill-coefficient-cost}. It uses no exponential
comparison.
It is independent of the coarse enclosure for $U_0$.
The mass-integral bound is Paper~II, (A.11)
\cite[(A.11)]{PaperII}.
Put
\begin{equation}\label{eq:l3-price}
\begin{gathered}
 m_{\mathrm{sw}}=\frac{e^{-\gamma}}2I_{\mathrm{HL}}G_*,\qquad
 \Delta_{\mathrm{cert}}=\frac{c_{K2,G0}}4\overline H-m_{\mathrm{sw}}\ge0,\\
 C_{\mathrm{sw}}=\frac{650000000}{1111}C_E,\qquad
 P_{\mathrm{sw}}=\frac{48544888063158}{27775},\\
 a_{\mathrm{sw}}=C_{\mathrm{sw}}\sigma,\qquad
 p_{\mathrm{sw}}=P_{\mathrm{sw}}\sigma^9.
\end{gathered}
\end{equation}
The remaining required estimates are
\begin{equation}\label{eq:l3-small-estimates}
\begin{gathered}
 (M_{\mathrm{prin}}/W-m_{\mathrm{sw}})_+=o_\sigma(1),\\
 M_{\mathrm{err}}/W\le C_{\mathrm{sw}}\sigma+o_\sigma(1),\qquad
 P_3/W\le P_{\mathrm{sw}}\sigma^9+o_\sigma(1),\\
 |E_{\mathrm{nc}}|/W=O_\sigma(N^{-\alpha}L^4).
\end{gathered}
\end{equation}
The exact numerical multipliers are
\begin{equation}\label{eq:l3-multipliers}
 \frac{4000\cdot585}{4(1-10^{-4})}=\frac{650000000}{1111},\qquad
 \frac{36\cdot502002(459/100)^3}{1-10^{-4}}
                       =\frac{48544888063158}{27775}.
\end{equation}
\paragraph{(P2.8c) Density product and prefix sums.}
Put $h=1/1000$ and $H_N=\sum_{N^h\le p\le N}(p-1)^{-1}$.
The two bounds used in the first multiplier are
\begin{equation}\label{eq:l3-mass-prefix-sums}
\begin{gathered}
 V_N(z)\le\frac{4C(N)}{\log z}\le\frac{4000C(N)}L
                                  \quad(z\ge N^h),\\
 \sum_{\text{actual prefixes }v}\frac1{\varphi(v)}
       \le\sum_{j=0}^3H_N^j\le1+8+64+512=585.
\end{gathered}
\end{equation}
These are Paper~II, Appendix~B, Lemmas B.9 and B.11,
(B.9) and (B.11) \cite[Appendix~B, Lemmas B.9, B.11, (B.9), (B.11)]{PaperII}.
Each half-open prefix cell is counted once.
Using $X_A\le N/L$ from (P2.8d) now gives
\begin{equation}\label{eq:l3-secondary-mass-step}
 M_{\mathrm{err}}\le\frac{X_A4000C(N)}L\,585
               \{C_E\sigma+o_\sigma(1)\}.
\end{equation}
Division by $W=4(1-\varepsilon)C(N)N/L^2$ gives $C_{\mathrm{sw}}$.

\paragraph{(P2.8e) The upper mass and the peel.}
The source mass is \eqref{eq:l3-principal-estimates},
Paper~II, (5.34)--(5.42).
The native lower-sieve cutoff uses the finite peel
\eqref{eq:l3-peel} and the estimate
\begin{equation}\label{eq:l3-peel-restatement}
 \frac{P_3}{W}\le
 \frac{36\cdot502002(459/100)^3}{1-\varepsilon}\tau+o_\sigma(1).
\end{equation}
For the uniform two-form calculation, put $m=vr$ and $w=\min(z_c,N^\alpha)$.
The required ranges are
\begin{equation}\label{eq:l3-selberg-ranges}
\begin{gathered}
 \log_Nm\le153/256+2\tau,\qquad \log_Nw\ge2/1001,\\
 w^2m\le N^{253/256+(153/64)\tau}<N^{1-0.0116}.
\end{gathered}
\end{equation}
For each sieve prime $p<w$, every prime factor of $m=vr$ is at least
$z_c\ge w$, so $p\nmid m$. The local two-form root count is then
$\nu_N(p)=1$ for $p\mid N$, and $\nu_N(p)=2$ otherwise.
For squarefree $d$, define $g_N(d)$ multiplicatively by
$g_N(p)=\nu_N(p)/(p-\nu_N(p))$.
The denominator and one-layer bound are
\begin{equation}\label{eq:l3-selberg-layer}
\begin{gathered}
 \mathcal G_N(w)=\sum_{\substack{d\le w\\d\mid P(w)}}g_N(d)
             =\frac{\log^2w}{4C(N)}(1+o(1)),\\
 S^{(1)}_m(w)\le\frac{N/m}{\mathcal G_N(w)}+O(w^2\log^5w)
       \le502002\,\frac{2C(N)N}{mL^2}.
\end{gathered}
\end{equation}
Here $P(w)$ is the product of all primes below $w$.
The last step uses
$2.001/(2/1001)^2=501251.00025<502002$.
The exponent saving in \eqref{eq:l3-selberg-ranges} absorbs its remainder.
There are at most $36$ source layers.
For the parent-bin prime sums $M_h$, the needed estimates are
\begin{equation}\label{eq:l3-peel-harmonic-sums}
 M_h=\sum_{p\ \mathrm{in\ bin}\ h}\frac1p\le2\tau,\qquad
 \sum_hM_h\le\frac{459}{100},\qquad
 \sum_{a_1,a_2,a_3,r}\frac1{a_1a_2a_3r}
       \le2\left(\frac{459}{100}\right)^3\tau.
\end{equation}
The pair $a_3,r$ belongs to one parent bin.
Multiplication of \eqref{eq:l3-selberg-layer} by $36$ and this sum,
then division by $W$, gives \eqref{eq:l3-peel-restatement}.
These are the exact steps of Paper~II, (4.45)--(4.56),
with the layer ceiling changed from three to $36$
\cite[(4.45)--(4.56), (5.34)--(5.42)]{PaperII}.
The uniform two-form estimates are proved here under the stated ranges.
\paragraph{Proof of \eqref{eq:l3-small-estimates}.}
\emph{Principal.} Since $\mathcal P_\Pi\le G_*+o_\sigma(1)$ and $G_*>0$,
the mass bound gives
$M_{\mathrm{prin}}/W\le\frac{(1-\varepsilon)I_{\mathrm{HL}}2e^{-\gamma}(G_*+o(1))}{4(1-\varepsilon)}
=m_{\mathrm{sw}}+o_\sigma(1)$.
\emph{Secondary sieve.} $E_{s,c}\le C_E\sigma+o_\sigma(1)$ uniformly,
since $\ell_c\ge\ell_1$. Lemma~B.9, Lemma~B.11 and $X_A\le N/L$ give
\eqref{eq:l3-secondary-mass-step}. Every prefix lies in one row, so the
four depths together contribute at most $585$ \cite[(B.9), (B.11)]{PaperII}.
\emph{Peel.} For $m=vr$ in $P_3$, all primes of $m$ are at least $z_c$,
$(m,N)=1$ and $z_c\in[N^{u/\rho^3},N^{\rho\mu_*}]$. Since $\chi_c=1$,
$\log_Nm\le\rho(t_1+t_2+2t_3)\le\rho(153/256-\eta_L)<153/256+2\tau$
\cite[(4.55)]{PaperII}. Write $A=\sum_{h=1}^{36}1_{\{A\ge h\}}$.
Each layer consists of $N^\alpha$-rough $e$. With $N-e=mk$, both $k$ and
$N-mk$ are $w$-rough, $w=\min(z_c,N^\alpha)$. A prime $p<w$ dividing $k$
and $N$ would divide $e$, and $k$ is odd. The two-form Selberg bound
\cite[(4.48)--(4.54)]{PaperII} gives \eqref{eq:l3-selberg-layer} for each
layer, with remainder absorbed by \eqref{eq:l3-selberg-ranges}. The prime
sums \eqref{eq:l3-peel-harmonic-sums} are \cite[(4.56)]{PaperII}, since
$r$ and $a_3$ lie in one parent bin. Multiplying by $36$ gives
\eqref{eq:l3-peel-restatement}.
\emph{Comparator.} For squarefree $q$,
$|d_A(q)|=\varphi(q)^{-1}\sum_{(e,q)>1}A(e)$. A common prime $\ell$ of
$e$ and $q$ is at least $N^\alpha$, so with $q=\ell r$,
\[
 \sum_{q\ \mathrm{sf}}|O_{\mathfrak b}(q)d_A(q)|
 \le C_{\mathrm{L3},\sigma}\sum_{\ell\ge N^\alpha}\frac{36N}{\ell(\ell-1)}
     \sum_{r\le N}\frac1{\varphi(r)}\ll_\sigma N^{1-\alpha}L.
\]
Summing finitely many blocks and dividing by $W$ gives
$|E_{\mathrm{nc}}|/W=O_\sigma(N^{-\alpha}L^3)$, which implies the stated
$O_\sigma(N^{-\alpha}L^4)$. This is the coprime-comparator bound used in
\eqref{eq:l3-disp-total}.

\subsection{The finite signed identity and its consequence}

\paragraph{(P2.8a) Finite signed identities.}
For the finite source \eqref{eq:l3-source} and stated sieve data,
combine \eqref{eq:l3-source-id}, \eqref{eq:l3-buchstab-peel}
and \eqref{eq:l3-aa-id}.
Their conditional signed identity is
\begin{equation}\label{eq:l3-rsw-id}
 S_{\mathrm{sw}}=E_{\mathrm{src}}+P_3+M_{\mathrm{prin}}
   +M_{\mathrm{err}}+R_{\mathrm{adv}}
                   -(Z_{\mathrm{src}}+Z_{\mathrm{AA}}+Z_{\mathrm{peel}}).
\end{equation}
Define
\begin{equation}\label{eq:l3-remainder}
\begin{aligned}
 \delta_P&=P_3/W-p_{\mathrm{sw}},\qquad
 \delta_A=M_{\mathrm{err}}/W-a_{\mathrm{sw}},\qquad
 \delta_M=M_{\mathrm{prin}}/W-m_{\mathrm{sw}},\\
 R_{\mathrm{sw}}&=E_{\mathrm{src}}/W+\delta_P+\delta_A+\delta_M
                                                +R_{\mathrm{adv}}/W,\\
 V_{\mathrm{sw}}&=(R_{\mathrm{sw}})_+\le E_{\mathrm{src}}/W
     +(\delta_P)_++(\delta_A)_++(\delta_M)_++(R_{\mathrm{adv}})_+/W.
\end{aligned}
\end{equation}
Under the stated sieve hypotheses, all three $Z$ terms are nonnegative.
Consequently
\begin{equation}\label{eq:l3-final}
\begin{aligned}
 -\frac{S_{\mathrm{sw}}}{W}
 &=-m_{\mathrm{sw}}-a_{\mathrm{sw}}-p_{\mathrm{sw}}-R_{\mathrm{sw}}
       +\frac{Z_{\mathrm{src}}+Z_{\mathrm{AA}}+Z_{\mathrm{peel}}}{W}\\
 &\ge-\frac{c_{K2,G0}}4\overline H
          -C_{\mathrm{sw}}\sigma-P_{\mathrm{sw}}\sigma^9-V_{\mathrm{sw}}.
\end{aligned}
\end{equation}
The favorable $Z_{\mathrm{peel}}$ is used once.
Every adverse remainder remains in \eqref{eq:l3-remainder}.

Assume all estimates stated above for the same $A$.
Then \eqref{eq:l3-source-error}, \eqref{eq:l3-disp-total}
and \eqref{eq:l3-small-estimates} imply
\begin{equation}\label{eq:l3-remainder-limit}
 V_{\mathrm{sw}}=o_\sigma(1).
\end{equation}
\paragraph{Proof of \eqref{eq:l3-remainder-limit}.}
All estimates above are proved for the same $A$, the same cells and the
same blocks. By \eqref{eq:l3-source-error}, $E_{\mathrm{src}}/W\ll L^{-331}$.
By \eqref{eq:l3-small-estimates}, $(\delta_M)_+,(\delta_A)_+,(\delta_P)_+$
are $o_\sigma(1)$. For $R_{\mathrm{adv}}$, \eqref{eq:l3-adv-split} and
\eqref{eq:l3-interior-split} give
\[
 \frac{|R_{\mathrm{adv}}|}W\le
 \frac{|D^{(0)}_{\mathrm{int}}|+|D^{(1)}_{\mathrm{int}}|+|D^{(2)}_{\mathrm{int}}|
      +|D^{(3)}_{\mathrm{int}}|+|D_{\mathrm{bdry}}|+|E_{\mathrm{sq}}|+|E_{\mathrm{nc}}|}W .
\]
The terms are $O_\sigma(L^{-1})$ by \eqref{eq:l3-bv-total};
$O_\sigma(L^{-1})$ by \eqref{eq:l3-masked-use} with $A_{\mathrm{loc}}=3692$,
$|w_C|/M_c\le36$ and $\sum_B|w_B|\le L^{3689}$;
$O_\sigma(L^{-327/2})$ by \eqref{eq:l3-boundary};
and $O_\sigma(N^{-\ell_1\sigma^2}L^4)$, $O_\sigma(N^{-\alpha}L^4)$.
Hence $V_{\mathrm{sw}}=o_\sigma(1)$.
After fixing $\sigma_*$, choose the eventual threshold so that
\begin{equation}\label{eq:l3-remainder-share}
 V_{\mathrm{sw}}\le\frac1{200000}.
\end{equation}
That threshold may be ineffective.
It also covers
\begin{equation}\label{eq:l3-structural-thresholds}
\begin{gathered}
 N_{\mathrm{dom}}=2\left\lceil\frac12
  \exp\!\left(\max\{10^4,2\log2/(\ell_1\sigma_*^2)\}\right)\right\rceil,\\
 N_{z,0}=2\left\lceil\frac12
       \exp\!\left(2\exp((\ell_1\sigma_*^2)^{-1/3})\right)\right\rceil,
 \qquad N_{\mathrm{LEN,app}}(\sigma_*).
\end{gathered}
\end{equation}
These structural thresholds alone do not imply \eqref{eq:l3-remainder-share}.

Relative to the comparison $-S_{\mathrm{sw}}/W\ge-2\overline H$,
the constant difference is
\begin{equation}\label{eq:l3-gain}
\begin{aligned}
 \Delta D_{L3}&=\left(2-\frac{c_{K2,G0}}4\right)\overline H\\
 &=0.022638464724298154340603\ldots.
\end{aligned}
\end{equation}
The exact cost of using the common coefficient instead of the optional
coefficient is
\[
 \frac{\overline H}{4}(c_{K2,G0}-c)
 =0.000049489563023753677508253183146688036059673516909.
\]
This is scalar arithmetic. Its use in the final count requires \eqref{eq:l3-final}
and the common error bounds.
The capped-row application is \eqref{eq:l3-bv-application}--\eqref{eq:l3-bv-total}.
The row-three estimate is \eqref{eq:l3-masked-use}, proved in the
row-three subsection. All errors use the common source
\eqref{eq:l3-source}. The inputs are Paper~I, Theorem~2.8;
Fact~\ref{fact:bv-full-support}; Lemma~\ref{hyp:ra-ms-transfer}(ii) for
(3.48); and the prime estimates of Fact~\ref{fact:prime-estimates}.
Hence \eqref{eq:l3-final} holds at the
threshold of \eqref{eq:l3-remainder-share}. The constant \eqref{eq:l3-gain} is
a gross coefficient difference. It enters the final count only together
with the switched error $C_{\mathrm{sw}}\sigma_*+P_{\mathrm{sw}}\sigma_*^9
+V_{\mathrm{sw}}$ of \eqref{eq:l3-final}. In the final count
this error is bounded by the caps \eqref{eq:bill-sigma}
($C_{\mathrm{sw}}\sigma_*\le10^{-5}$, $P_{\mathrm{sw}}\sigma_*^9\le10^{-5}$)
and by \eqref{eq:l3-remainder-share}, in total at most $1/40000$.

\section{The final bound}\label{sec:final-bill}

We keep the notation of Sections~\ref{sec:introduction} and
\ref{sec:notation}.
All error terms below are nonnegative.
The analytic inputs are \textup{(A1)}--\textup{(A5)} of
Table~\ref{tab:intro-inputs}.

\subsection{Addition of the comparisons}

Put $\delta_{\mathrm{err}}=10^{-4}$.
The constant of \eqref{eq:intro-DB} is
\begin{equation}\label{eq:bill-constant}
\begin{aligned}
 D_B={}&D_{\mathrm{II}}+\frac{5719}{10^6}
       +M_{5,\mathrm{internal}}+g_{\mathrm{RA,lo}}\\
     &-\delta_{\mathrm{err}}+
       \left(2-\frac{c_{K2,G0}}{4}\right)\overline H,
 \qquad \overline H=\frac{389253}{5000000}.
\end{aligned}
\end{equation}
Here $D_{\mathrm{II}}$ is Paper~II's original lower constant
\cite[(6.3a)]{PaperII}.
The six endpoint gains $g_6$ of \eqref{eq:retB-six-gain} are not used.
The two profile gains enter separately.
In particular,
\begin{equation}\label{eq:bill-components}
\begin{aligned}
 D_{\mathrm{II}}&=\frac{45849916338009170593}
                         {25000000000000000000},\\
 M_{5,\mathrm{internal}}
     &=-K_{\mathrm{hi}}\frac{6073714261}{10^9}-L_{G1},\\
 g_{\mathrm{RA,lo}}&=\frac{989598353}{10^{10}}=0.0989598353.
\end{aligned}
\end{equation}
The last line is the conservative rational floor for the integral
gain $\mathfrak g_{\mathrm{RA}}$. It is not a central estimate.
Section~\ref{sec:root-a} separates it from the integral gain.
The certified integral enclosure is
$[0.0989598353353,0.0989598353355]$.
We take $K_{\mathrm{hi}}$ and $L_{G1}$ as defined in
\cite[(A.10d), (4.84)]{PaperII}.
Their exact values are
\begin{equation}\label{eq:bill-M5-inputs}
\begin{aligned}
 K_{\mathrm{hi}}&=
 \frac{3989480791070578466390443009830933753748442947478110112223}
      {10^{58}},\\[2pt]
 L_{G1}&=-
 \frac{243950428265194656254206921939947762868871361874409552095097}
      {100433627766186892221372630771322662657637687111424552206336}.
\end{aligned}
\end{equation}
Thus $M_{5,\mathrm{internal}}=M_n/M_d$ exactly, where
\begin{equation}\label{eq:bill-M5-integers}
 M_n=\sum_{i=0}^5m_i10^{20i},\qquad
 M_d=\sum_{i=0}^5d_i10^{20i}.
\end{equation}
The coefficients are printed in Table~\ref{tab:bill-M5-integers}.
\begingroup
\small
\setlength{\tabcolsep}{8pt}
\renewcommand{\arraystretch}{1.25}
\begin{longtable}{@{}rrr@{}}
\caption{Exact shifted-profile gain}\label{tab:bill-M5-integers}\\
\hline
$i$ & $m_i$ & $d_i$\\
\hline
\endfirsthead
\hline
$i$ & $m_i$ & $d_i$\\
\hline
\endhead
0 & 66922148532305907989 & 0\\
1 & 34893223604885136721 & 0\\
2 & 3973350097562772509 & 0\\
3 & 87266878234046404033 & 48635364229120000000\\
4 & 28374567121249055735 & 73384187692692674921\\
5 & 3998 & 680564\\
\hline
\end{longtable}
\endgroup
The exact rational $c_{\mathrm{O1}}$ is defined in
\eqref{eq:l3-exact-certificate-constant} and
Table~\ref{tab:l3-certificate-integers}.
It records the optional sharper certificate only.
The main count uses $c_{K2,G0}$ in both comparisons.
The common functional comparison is \eqref{eq:l3-certificate}.
For comparison with the optional sharper coefficient, exact rational
arithmetic gives
\begin{equation}\label{eq:bill-coefficient-cost}
 0<c-c_{\mathrm{O1}}<4.081\cdot10^{-46},\qquad
 0<\frac{c-c_{\mathrm{O1}}}{4}\overline H
                  <7.943\cdot10^{-48}.
\end{equation}
This is a rational comparison only. It gives no bound for $2e^{-\gamma}U_0$.
The further cost of the common coefficient relative to \(c\) is exactly
\[
 \frac{c_{K2,G0}-c}{4}\overline H
 =0.000049489563023753677508253183146688036059673516909.
\]
The resulting main switched gain is exactly
\[
 \Delta D_{L3}=0.02263846472429815434060381043792394540367262995113415.
\]
Neither optional certificate is required for \eqref{eq:bill-constant}.
Each gain is proved on the original arithmetic source of its comparison.
The gains $5719/10^6$ and $M_{5,\mathrm{internal}}$ come from
Proposition~\ref{prop:retB-comparison}, through
Propositions~\ref{prop:app-h6-four-profiles} and~\ref{prop:app-m5-cap}.
The gain $g_{\mathrm{RA,lo}}$ comes from Proposition~\ref{prop:ra-b11}, and
$\Delta D_{L3}$ from \eqref{eq:l3-final}. Table~\ref{tab:bill-binding} lists the bindings.

Table~\ref{tab:bill-components} itemizes the constant.
Its decimal column gives approximations only.
The defining formulas determine the exact sum.
\begingroup
\small
\renewcommand{\arraystretch}{1.3}
\begin{longtable}{@{}p{.48\linewidth}p{.43\linewidth}@{}}
\caption{The six terms in $D_B$}\label{tab:bill-components}\\
\hline
Term & Value\\
\hline
\endfirsthead
\hline
Term & Value\\
\hline
\endhead
$D_{\mathrm{II}}$ & $1.83399665352036682372$\\
$5719/10^6$ & $0.005719$\\
$M_{5,\mathrm{internal}}$ & $0.005874949945026357\ldots$\\
$g_{\mathrm{RA,lo}}$ & $0.0989598353$\\
$-\delta_{\mathrm{err}}$ & $-0.0001$\\
$\Delta D_{L3}=(2-c_{K2,G0}/4)\overline H$
 & $0.022638464724298154\ldots$\\
\hline
$D_B$ & $1.967088903489691335\ldots$\\
\hline
\end{longtable}
\endgroup

Adding \eqref{eq:intro-comparisons} and
\eqref{eq:intro-exterior} cancels the sixth endpoint:
\begin{equation}\label{eq:bill-cancellation}
 \overline J_6-\overline J_6=0,\qquad
 2\overline H-\frac{c_{K2,G0}}{4}\overline H
 =\left(2-\frac{c_{K2,G0}}{4}\right)\overline H.
\end{equation}
Thus their conditional consequence is
\begin{equation}\label{eq:bill-added}
 \frac{4G_b(N)}{W(N)}
 \ge D_B+\delta_{\mathrm{err}}-c_{K2,G0}I(b)-\mathcal E_B(N;b,\sigma).
\end{equation}
This uses the complete $S_2$ of Paper~II.
The finite identities keep their original signs and multiplicities.
The fixed errors in Paper~II are subtracted once
\cite[(6.4)--(6.6)]{PaperII},
and not again in \eqref{eq:bill-constant}.

\subsection{Ten error allocations}

Use the exterior error from \eqref{eq:bill-K-error}.
For the table below, a conservative majorant is convenient:
\begin{equation}\label{eq:bill-K-error}
\begin{aligned}
 E_K(N;b,\sigma)
 &=A_K(b)C_{\mathrm{AA}}\sigma
      +\frac{3884000000000}{9999}\sigma^9+r_K(N;b,\sigma)\\
 &\le 4A_K(b)C_{\mathrm{AA}}\sigma
      +\frac{3884000000000}{9999}\sigma^9
      +r_K(N;b,\sigma),\\
 A_K(b)&=2e^{-\gamma}I(b)W_{\mathrm{AA},K},\qquad
 P_K=\frac{3884000000000}{9999},\\
 r_K(N;b,\sigma)&=o_{b,\sigma}(1).
\end{aligned}
\end{equation}
Here $W_{\mathrm{AA},K}$ is the exterior remainder-weight bound.
It belongs to the source in Section~\ref{sec:k2}.
The extra factor four only enlarges this nonnegative error.
It leaves the principal constant unchanged.
Similarly, the switched comparison requires
\begin{equation}\label{eq:bill-sw-error}
\begin{aligned}
 E_{\mathrm{sw}}(N;\sigma)
 &\le \frac{650000000}{1111}C_E\sigma
       +\frac{48544888063158}{27775}\sigma^9
       +V_{\mathrm{sw}}(N;\sigma),\\
 V_{\mathrm{sw}}(N;\sigma)&=o_\sigma(1).
\end{aligned}
\end{equation}
Here $V_{\mathrm{sw}}$ is the positive part in
\eqref{eq:l3-remainder}.
The conditional limit is \eqref{eq:l3-remainder-limit}.
These two limits concern the actual sums stated there.
Both majorants are the error terms of \eqref{eq:k2-final-fixed-sigma} under
\textup{(A3)} and of \eqref{eq:l3-final} under \textup{(A4)}, with
$C_{\mathrm{sw}}$ from \eqref{eq:l3-price}. By \textup{(A5)},
$C_{\mathrm{AA}}$ and $C_E$ do not depend on $\sigma$. The exterior weight
$W_{\mathrm{AA},K}$ is $W_{\mathrm{AA}}$ of \eqref{eq:k2-aa-mass}:
\begin{equation}\label{eq:bill-AK-positive}
\begin{gathered}
 A_K(b)=2e^{-\gamma}I(b)W_{\mathrm{AA},K},\qquad
 W_{\mathrm{AA},K}=500\Bigl(1+H_{\log}+\frac{H_{\log}^2}2
                    +\frac{H_{\log}^3}6\Bigr),\\
 H_{\log}=\log\frac{\mu_*}{u}=\log\frac{3125}{32}>0.
\end{gathered}
\end{equation}
Since $I(b)>0$, this gives $0<A_K(b)<\infty$. In particular
$A_K(b)C_{\mathrm{AA}}\sigma\le4A_K(b)C_{\mathrm{AA}}\sigma$.

The retained and Root--A comparisons require
\begin{equation}\label{eq:bill-eventual-errors}
 E'(N;\sigma)=o_\sigma(1),\qquad
 E_{\mathrm{RA}}(N;\sigma)=o_\sigma(1).
\end{equation}
For Root--A, use $E_{\mathrm{RA}}(N;\sigma)$
from \eqref{eq:ra-b11-floor}.
The first includes every conversion and the saved-array contribution
$\sum_{j\in\mathcal H}|\Lambda^0_{j,N}|$ in
\eqref{eq:intro-retained-error}; that contribution is not charged again.
The pointwise error may use the coarser bound
\begin{equation}\label{eq:bill-pointwise-error}
 E_{\mathrm{exc}}(N)=\frac{4E_{\mathrm{pt}}(N)}{W(N)}
 \le\frac{2197N^{1-\alpha}}{W(N)}=o(1).
\end{equation}
Here $E_{\mathrm{pt}}$ is defined in
\eqref{eq:source-pointwise-error}. The explicit pointwise argument in
Section~\ref{sec:source-decomposition} gives $E_{\mathrm{pt}}=0$
for $N\ge N_{\mathrm{pt}}$. The coarser majorant retains the same error bound.
The inequality follows from \eqref{eq:source-first-lower}, using both
the pointwise bound and the exceptional-set cardinality.
Since $C(N)\ge C_2>0$ and $\alpha>0$, the coarser majorant tends to zero:
$N^{-\alpha}(\log N)^2\to0$.
No sharper exceptional coefficient is needed.
Proposition~\ref{prop:retB-comparison} gives $E'(N;\sigma)\to0$, and
Proposition~\ref{prop:ra-b11} gives $E_{\mathrm{RA}}(N;\sigma)\to0$ under
\textup{(A2)}, each at fixed $\sigma$.
Lemma~\ref{lem:source-pointwise} gives $E_{\mathrm{exc}}(N)=0$ for
$N\ge N_{\mathrm{pt}}$.

Fix $b$, all absolute constants, and the source geometry first.
Choose positive absolute constants $C_E^{\mathrm{L3}}$ and
$C_E^{\mathrm{II}}$ for the Section~\ref{sec:l3} and Paper~II sieve bounds. Define
\begin{equation}\label{eq:bill-CE-definition}
 C_E:=\max\{C_E^{\mathrm{L3}},C_E^{\mathrm{II}}\}.
\end{equation}
By \textup{(A5)}, these constants are independent of $\sigma$ and $N$.
Here $C_{\mathrm{AA}}=4C_{\mathrm{AA37}}$ by \eqref{eq:k2-aa-constant}. The
constant $C_E$ dominates the absolute constant of
\eqref{eq:classical-structured-sieve} used in \eqref{eq:l3-sieve-error}.
The underlying sieve bounds come from Paper~I, Theorem~2.6, (2.31)
\cite[Theorem~2.6, (2.31)]{PaperI}.
Choose one rational $\sigma_*>0$ satisfying
\begin{equation}\label{eq:bill-sigma}
\begin{aligned}
 0<\sigma_*&\le 10^{-10},\\
 \sigma_*&\le\frac{713332}{625\cdot10^{21}C_E},\\
 \sigma_*&\le\frac{1}{400000A_K(b)C_{\mathrm{AA}}},\\
 \sigma_*&\le\left(\frac{9999}{388400000000000000}\right)^{1/9},\\
 \sigma_*&\le\frac{1111}{65000000000000C_E},\\
 \sigma_*&\le\left(\frac{27775}{4854488806315800000}\right)^{1/9},\\
 \tau_*&=(\sigma_*)^9.
\end{aligned}
\end{equation}
By \eqref{eq:bill-AK-positive} and \textup{(A5)}, every right side in
\eqref{eq:bill-sigma} is positive and independent of $\sigma$.
Hence such a rational $\sigma_*$ exists.
The second restriction preserves Paper~II's fixed error bound.
It is exactly
\begin{equation}\label{eq:bill-paperII-AA}
 \sigma_*\le\frac{4g_{\mathrm{AA}}}{C_EW_*},\qquad
 g_{\mathrm{AA}}=\frac{178333}{6250000000},\qquad W_*=10^{14}
 \quad\text{\cite[(1.26), (1.27), (6.6)]{PaperII}}.
\end{equation}
The other four restrictions control new fixed errors.
Freeze every finite partition and logarithmic precision next.
Only then choose the common threshold.

The ten allocations are printed in Table~\ref{tab:bill}.
The two zero rows are unused capacity.
\begingroup
\small
\renewcommand{\arraystretch}{1.65}
\setlength{\extrarowheight}{2pt}
\begin{longtable}{@{}r p{.60\linewidth} p{.25\linewidth}@{}}
\caption{The error bound $\delta_{\mathrm{err}}=10^{-4}$}
\label{tab:bill}\\
\hline
 & Error term & Allocation\\
\hline
\endfirsthead
\hline
 & Error term & Allocation\\
\hline
\endhead
1 & $4A_K(b)C_{\mathrm{AA}}\sigma_*$ & $\delta_{\mathrm{err}}/10$\\
2 & $\tfrac{3884000000000}{9999}\tau_*$ & $\delta_{\mathrm{err}}/10$\\
3 & $\tfrac{650000000}{1111}C_E\sigma_*$ & $\delta_{\mathrm{err}}/10$\\
4 & $\tfrac{48544888063158}{27775}\tau_*$ & $\delta_{\mathrm{err}}/10$\\
5 & $E'(N;\sigma_*)+V_{\mathrm{sw}}(N;\sigma_*)$
    & $\tfrac{\delta_{\mathrm{err}}}{20}+\tfrac{\delta_{\mathrm{err}}}{20}$\\
6 & $0$ & $\delta_{\mathrm{err}}/10$\\
7 & $E_{\mathrm{RA}}(N;\sigma_*)$ & $\delta_{\mathrm{err}}/10$\\
8 & $r_K(N;b,\sigma_*)$ & $\delta_{\mathrm{err}}/10$\\
9 & $E_{\mathrm{exc}}(N)$ & $\delta_{\mathrm{err}}/10$\\
10 & $0$ & $\delta_{\mathrm{err}}/10$\\
\hline
\end{longtable}
\endgroup
Its capacities sum exactly:
\begin{equation}\label{eq:bill-sum}
 4\frac{\delta_{\mathrm{err}}}{10}
 +\left(\frac{\delta_{\mathrm{err}}}{20}
       +\frac{\delta_{\mathrm{err}}}{20}\right)
 +5\frac{\delta_{\mathrm{err}}}{10}
 =\delta_{\mathrm{err}}.
\end{equation}
The first four bounds follow from \eqref{eq:bill-sigma}.
The five remaining nonzero bounds follow from the stated limits.
Each limit is taken at fixed $\sigma_*$.

\begin{proposition}[Error budget]\label{prop:bill-budget}
Assume \textup{(A1)}--\textup{(A5)}, and let $\sigma_*$ be rational with
\eqref{eq:bill-sigma}. Then every even $N\ge N_0^*(\sigma_*)$ satisfies
\begin{equation}\label{eq:bill-budget}
 \mathcal E_B(N;b,\sigma_*)\le\frac8{10}\delta_{\mathrm{err}}
 <\delta_{\mathrm{err}}.
\end{equation}
\end{proposition}

\begin{proof}
Rows 1 and 3. The third and fifth lines of \eqref{eq:bill-sigma} give
\[
 4A_K(b)C_{\mathrm{AA}}\sigma_*\le\frac4{400000}=\frac{\delta_{\mathrm{err}}}{10},
 \qquad
 \frac{650000000}{1111}C_E\sigma_*\le\frac{650000000}{65000000000000}
 =\frac{\delta_{\mathrm{err}}}{10}.
\]
Rows 2 and 4. The first line gives $\tau_*\le10^{-90}$. This is below
$9999/388400000000000000$ and $27775/4854488806315800000$, so the fourth and
sixth lines hold, and
\[
 \frac{3884000000000}{9999}\tau_*\le\frac{\delta_{\mathrm{err}}}{10},\qquad
 \frac{48544888063158}{27775}\tau_*\le\frac{\delta_{\mathrm{err}}}{10}.
\]
Rows 5, 7 and 8. At the fixed $\sigma_*$, $E'\to0$ by \textup{(A1)},
$V_{\mathrm{sw}}\to0$ by \eqref{eq:l3-remainder-limit} under \textup{(A4)},
$E_{\mathrm{RA}}\to0$ by \textup{(A2)}, and $r_K\to0$ by \textup{(A3)}.
By \eqref{eq:not-little-o}, the thresholds
$T_{E'},T_V,T_{E_{\mathrm{RA}}},T_r,T_{E_{\mathrm{exc}}}$ of
\eqref{eq:not-common-threshold} exist,
and above them $E'\le\delta_{\mathrm{err}}/20$,
$V_{\mathrm{sw}}\le\delta_{\mathrm{err}}/20$,
$E_{\mathrm{RA}}\le\delta_{\mathrm{err}}/10$ and
$r_K\le\delta_{\mathrm{err}}/10$.
Row 9. Lemma~\ref{lem:source-pointwise} gives $E_{\mathrm{exc}}(N)=0$ for
$N\ge N_{\mathrm{pt}}$, so $T_{E_{\mathrm{exc}}}\le N_{\mathrm{pt}}$.
Rows 6 and 10 are zero.
Every threshold used is bounded by the common maximum in
\eqref{eq:not-common-threshold}. Also $N_{\mathrm{pt}}<N_{\mathrm{tr}}$
by Lemma~\ref{lem:bill-transfer}. Hence all of them hold for
$N\ge N_0^*(\sigma_*)$. By \eqref{eq:intro-error-identifications},
\eqref{eq:bill-K-error} and \eqref{eq:bill-sw-error},
\[
\begin{aligned}
 \mathcal E_B
 \le{}&\bigl(E'+V_{\mathrm{sw}}\bigr)+E_{\mathrm{RA}}
   +\Bigl(\frac{650000000}{1111}C_E\sigma_*+\frac{48544888063158}{27775}\tau_*\Bigr)\\
   &+\Bigl(4A_KC_{\mathrm{AA}}\sigma_*+\frac{3884000000000}{9999}\tau_*+r_K\Bigr)
   +E_{\mathrm{exc}}.
\end{aligned}
\]
The bounds above give at most
$\frac{\delta_{\mathrm{err}}}{10}+\frac{\delta_{\mathrm{err}}}{10}
+\frac{2\delta_{\mathrm{err}}}{10}+\frac{3\delta_{\mathrm{err}}}{10}+0$,
which is below the nonzero allocations
$\frac{8\delta_{\mathrm{err}}}{10}$ of Table~\ref{tab:bill}.
\end{proof}

Take the named source, application and error thresholds listed in
Section~\ref{sec:notation}. The Section~\ref{sec:k2} definitions of
\(T_K\) and \(T_r\) already include the main-mass threshold
\(N_{\mathrm{mass}}(b,\sigma_*)\) and the AA threshold
\(N_{\mathrm{AA}}(b,\sigma_*)\). Thus the common maximum in \(\mathcal T_*\)
below includes both requirements at the fixed \(b,\sigma_*\).
The common threshold is the single definition
\begin{equation}\label{eq:bill-threshold}
 N_0^*(\sigma_*)\quad\text{as defined in \eqref{eq:not-common-threshold}}.
\end{equation}
Its existence depends on all those estimates. The floor alone supplies
none of them. By its definition and Lemma~\ref{lem:bill-transfer},
$N_0^*(\sigma_*)\ge\max\{N_{\mathrm{tr}},N_6\}>N_{\mathrm{pt}}$.
The Section~\ref{sec:notation} threshold definitions include $N_{\mathrm{pt}}$ through
$T_{\mathrm{src}}$ and $N_6$ through $T_{\mathrm{RA}}$. The bound $N_{\mathrm{tr}}$
enters directly in the common maximum. Lemma~\ref{lem:source-pointwise}
holds for $N\ge N_{\mathrm{pt}}$ and Lemma~\ref{lem:source-upper-sources}
for $N\ge N_6$.
Thus the order is
\begin{equation}\label{eq:bill-order}
 \text{constants}\ \longrightarrow\ \sigma_*
 \ \longrightarrow\ \text{finite choices}
 \ \longrightarrow\ N_0^*(\sigma_*)\ \longrightarrow\ N.
\end{equation}
No parameter $\sigma_*(N)$ is used.

Proposition~\ref{prop:bill-budget} gives
\begin{equation}\label{eq:bill-error-bound}
 \mathcal E_B(N;b,\sigma_*)\le\delta_{\mathrm{err}}
 \qquad(N\ge N_0^*(\sigma_*),\ N\in2\mathbb Z).
\end{equation}
Combining \eqref{eq:bill-added} and \eqref{eq:bill-error-bound} gives
\begin{equation}\label{eq:bill-positive-count}
 \boxed{\displaystyle
 \frac{4G_b(N)}{W(N)}\ge D_B-c_{K2,G0}I(b).}
\end{equation}
It holds for every even $N\ge N_0^*(\sigma_*)$ under
\textup{(A1)}--\textup{(A5)}.

\subsection{The numerical margin}

All constants in \eqref{eq:bill-constant} have exact definitions.
The exterior price uses the G0 coefficient in
\eqref{eq:k2-certificate}.
The same coefficient prices L3. No $G_*,U_0$ sharper domination
is used in this arithmetic or in the main-count hypotheses.
The constants have the following outward enclosures:
\begin{equation}\label{eq:bill-numeric-enclosures}
\begin{gathered}
 1.967088903489691335582703173494<D_B,\\
 D_B<1.967088903489691335582703173495,\\[2pt]
 1.966859949178424046530804965343<c_{K2,G0}I(b),\\
 c_{K2,G0}I(b)<1.966859949178424046530804965344,\\[2pt]
 0.000228954311267289051898208151<D_B-c_{K2,G0}I(b),\\
 D_B-c_{K2,G0}I(b)<0.000228954311267289051898208152.
\end{gathered}
\end{equation}
Here
\begin{equation}\label{eq:bill-price-definition}
 c_{K2,G0}=\frac{6836825163875517756286846321650754372931094689}{10^{45}},
 \qquad I(b)=\log\frac{4000007}{2999993}.
\end{equation}
The coefficient is defined in \eqref{eq:k2-certificate}.
The final comparison is \eqref{eq:k2-final-fixed-sigma}.
Hence
\begin{equation}\label{eq:bill-margin-simple}
 D_B-c_{K2,G0}I(b)>0.0002289>\frac1{5000}>0.
\end{equation}
The table below identifies each exact constant by its defining display and
lists the comparisons that use it unchanged.
\begingroup
\small
\renewcommand{\arraystretch}{1.3}
\begin{longtable}{@{}p{.27\linewidth}p{.31\linewidth}p{.34\linewidth}@{}}
\caption{Constants and the comparisons that use them}
\label{tab:bill-binding}\\
\hline
Constant & Definition & Comparison\\
\hline
\endfirsthead
\hline
Constant & Definition & Comparison\\
\hline
\endhead
$D_{\mathrm{II}}$ & \eqref{eq:intro-components}, \cite[(6.3a)]{PaperII}
 & Proposition~\ref{prop:retB-comparison}, by \eqref{eq:retB-paperII-floor}\\
$5719/10^6$ & \eqref{eq:retB-H6-finite}
 & Proposition~\ref{prop:retB-comparison}, by
   Proposition~\ref{prop:app-h6-four-profiles}\\
$M_{5,\mathrm{internal}}$, $K_{\mathrm{hi}}$, $L_{G1}$
 & \eqref{eq:bill-M5-inputs}, Table~\ref{tab:bill-M5-integers}
 & Proposition~\ref{prop:retB-comparison}, by
   Proposition~\ref{prop:app-m5-cap}\\
$g_{\mathrm{RA,lo}}$ & \eqref{eq:ra-rational-floor}
 & Proposition~\ref{prop:ra-b11}, \eqref{eq:ra-b11-floor}\\
$\overline H$, $c_{K2,G0}$ & \eqref{eq:bill-constant}, \eqref{eq:k2-certificate}
 & \eqref{eq:l3-final}, with \eqref{eq:l3-price}\\
$c_{K2,G0}$, $I(b)$ & \eqref{eq:bill-price-definition}
 & Proposition~\ref{prop:k2-comparison}, \eqref{eq:k2-final-fixed-sigma}\\
$A_K(b)$, $C_{\mathrm{AA}}$, $P_K$
 & \eqref{eq:bill-AK-positive}, \eqref{eq:k2-aa-constant},
   \eqref{eq:k2-final-fixed-sigma}
 & Table~\ref{tab:bill}, rows 1 and 2\\
$C_{\mathrm{sw}}$, $P_{\mathrm{sw}}$ & \eqref{eq:l3-price}
 & Table~\ref{tab:bill}, rows 3 and 4\\
$\delta_{\mathrm{err}}$ & \eqref{eq:bill-constant}
 & Proposition~\ref{prop:bill-budget}\\
\hline
\end{longtable}
\endgroup
No other constant enters $D_B$ or the error budget.

The logarithm enclosure uses rational arithmetic.
Put $t_{\log}=500007/3500000$ and
\begin{equation}\label{eq:bill-log-enclosure}
 S_m=2\sum_{k=0}^{m}\frac{t_{\log}^{2k+1}}{2k+1},
 \qquad
 S_m< I(b)<S_m+
       \frac{2t_{\log}^{2m+3}}{(2m+3)(1-t_{\log}^2)}.
\end{equation}
Here $I(b)=2\operatorname{arctanh}t_{\log}$ and $0<t_{\log}<1$.
Bounding each remaining denominator by $2m+3$ proves the upper bound.
Taking $m=24$ gives \eqref{eq:bill-numeric-enclosures}.
Thus no rounded logarithm enters the margin test.

For reference, put $\zeta=10^{-6}$ and $\mathfrak E_*=\exp(D_B/c_{K2,G0})$.
The buffered critical exponent is
\begin{equation}\label{eq:bill-critical-exponent}
 a_{\mathrm{buf}}^*
 =\frac{\mathfrak E_*+1}{(1-\zeta)\mathfrak E_*-\zeta}.
\end{equation}
To derive it, write $b_*=1/(1+\mathfrak E_*)+\zeta$.
The equation $a(1-b_*)=1$ gives
\eqref{eq:bill-critical-exponent}.
For the present constants,
\begin{equation}\label{eq:bill-critical-enclosure}
 1.749974884137892433433207484622
 <a_{\mathrm{buf}}^*
 <1.749974884137892433433207484623<\frac74.
\end{equation}
This enclosure also uses rational bounds.
For $x=D_B/c_{K2,G0}$ and $m=36$, put
\begin{equation}\label{eq:bill-exp-enclosure}
 T_m=\sum_{k=0}^{m}\frac{x^k}{k!},\qquad
 T_m<e^x<T_m+
 \frac{x^{m+1}}{(m+1)!}\frac{1}{1-x/(m+2)}.
\end{equation}
Here $0<x<1$.
The tail ratio is at most $x/(m+2)$.
The function in \eqref{eq:bill-critical-exponent} decreases with $\mathfrak E_*$.
Substitution of the two rational bounds proves
\eqref{eq:bill-critical-enclosure}.

The endpoint in \eqref{eq:bill-critical-enclosure} is equivalent to the
positive scalar margin at this fixed $b$. It supplies no theorem at a
smaller $a$ without new comparisons at the corresponding cutoff.
If the optional Lemma~\ref{lem:retB-six-endpoints} is used, $D_B$ increases by
$g_6=0.0020549454896034927680\ldots$ of \eqref{eq:retB-six-gain}. The margin
then becomes $0.0022838998008707818199\ldots$, and $a_{\mathrm{buf}}^*$ becomes
$1.7497494980881623276859\ldots$. Theorem~\ref{thm:main} does not use this.

For an optional arithmetic comparison, define
$D_{B,c}=D_B+(c_{K2,G0}-c)\overline H/4$.
Using $c$ in both K2 and L3 would require two further inputs.
This paragraph concerns the optional scalar comparison at the same fixed cutoff;
it supplies no new cutoff-dependent estimate. The first is
the separate sharper functional domination $G_*\le U_0$ from
\hyperref[eq:app-g24-inclusion]{equation~\ref*{eq:app-g24-inclusion}}. The second is a certified
$2e^{-\gamma}U_0\le c$. It does not follow from the coarse
interval in \hyperref[eq:app-g24-numbers]{equation~\ref*{eq:app-g24-numbers}}. The main proof uses
neither.
Its scalar endpoint is
\begin{equation}\label{eq:bill-active-critical-exponent}
 E_c=\exp(D_{B,c}/c),\qquad
 a^*_{\mathrm{buf},c}
 =\frac{E_c+1}{(1-\zeta)E_c-\zeta}.
\end{equation}
The same rational Taylor enclosure gives
\begin{equation}\label{eq:bill-active-critical-enclosure}
 1.749889173139093190128834678
 <a^*_{\mathrm{buf},c}
 <1.749889173139093190128834679<\frac74.
\end{equation}
Thus, on the optional sharper branch, $a^*_{\mathrm{buf},c}$ is
$1.7498891731390\ldots$. The main theorem uses the explicitly certified
G0 coefficient defined and bounded in \eqref{eq:k2-certificate}; the
display above is only the optional exponent endpoint.

\subsection{The transfer}

\begin{lemma}[Transfer]\label{lem:bill-transfer}
The parameters \eqref{eq:not-transfer-parameters} satisfy
\begin{equation}\label{eq:bill-transfer-parameters}
 \kappa_{\mathrm{tr}}=\frac34(1-b)-b=\frac7{4000000},\qquad
 N_{\mathrm{tr}}=10^{12000000/7}>10^{12}+1=N_{\mathrm{pt}}.
\end{equation}
For even $N\ge N_{\mathrm{tr}}$, every $n=rq\in\mathcal A_N$ with primes
$r\le q$ and $r<N^b$ satisfies
\begin{equation}\label{eq:bill-transfer}
 q>\varepsilon N^{1-b},\qquad
 q^{3/4}>\varepsilon^{3/4}N^{b+\kappa_{\mathrm{tr}}}\ge N^b>r.
\end{equation}
\end{lemma}

\begin{proof}
Since $a=7/4$, $a(1-b)-1=\frac34(1-b)-b$, and $b=2999993/7000000$ gives
$7/4000000$. Also $\varepsilon^{-(a-1)/\kappa_{\mathrm{tr}}}
=10^{4\cdot3000000/7}$, and $12000000/7>13$.
For $n=rq\in\mathcal A_N$ we have $n>\varepsilon N$, so
$q=n/r>\varepsilon N^{1-b}$ and
$q^{3/4}>\varepsilon^{3/4}N^{3(1-b)/4}=\varepsilon^{3/4}N^{b+\kappa_{\mathrm{tr}}}$.
Finally $N^{\kappa_{\mathrm{tr}}}\ge\varepsilon^{-3/4}$ exactly when
$N\ge N_{\mathrm{tr}}$.
\end{proof}

\section*{Proof of Theorem 1}

Use the established inputs \textup{(A1)}--\textup{(A5)} and fix a rational
$\sigma_*$ satisfying \eqref{eq:bill-sigma}. Let
$N\ge N_0^*(\sigma_*)$ be even. Then
$N\ge\max\{N_{\mathrm{tr}},N_6\}>N_{\mathrm{pt}}$.
Lemma~\ref{lem:source-pointwise} gives the first line of
\eqref{eq:intro-comparisons}. At $\sigma=\sigma_*$, the established
comparisons \textup{(A1)}, \textup{(A2)} and \textup{(A4)} give its other
three lines, and \textup{(A3)} gives \eqref{eq:intro-exterior}.
The zero-gain combination
Lemma~\ref{lem:source-trackb-combination} adds these four comparisons to
\eqref{eq:bill-added}. Proposition~\ref{prop:bill-budget} gives
\eqref{eq:bill-error-bound}, hence \eqref{eq:bill-positive-count}.
Equation~\eqref{eq:bill-margin-simple} then gives $G_b(N)>0$, since
$W(N)>0$ by \eqref{eq:not-normalization}.
Hence the set in \eqref{eq:intro-count} is nonempty.

Choose one counted integer $n=N-p$.
If $n$ is prime, take $q=n$ and $r=1$.
Otherwise $n=rq$ with primes $r\le q$ and $r<N^b$, and
Lemma~\ref{lem:bill-transfer} gives $r<q^{3/4}$.
Therefore $N=p+rq$ and $r\le q^{3/4}$. $\square$


\section*{Appendices}
Appendix~A proves three auxiliary estimates. Appendix~B records finite
calculations, numerical certificates, and exact arithmetic. All cross-references
to these appendices are internal to this document.
\appendix
\providecommand{\Li}{\operatorname{Li}}
\newenvironment{mainlemma}[2]%
  {\par\medskip\noindent\textbf{Lemma~#1}\ (#2)\textbf{.}\ \itshape}%
  {\par\medskip}

\section{Proofs of three estimates of Section~\ref{sec:root-a}}\label{app:ra-proofs}

This section proves Lemmas~6.5, 6.8 and~6.9.
Numbers (6.$k$) and~6.$k$ refer to Section~\ref{sec:root-a}. Numbers (\ref{sp:first})
and later refer to this section. The only inputs are Papers~I and~II and
the classical results cited at the point of use. Throughout,
$\tau_j$ is the $j$-fold divisor function and $\tau=\tau_2$.
Let $\omega(n)$ be the number of distinct prime divisors of $n$. Let $P^+(n)$ and $P^-(n)$
be its largest and least prime factors. The sum $\sum^*_{b\bmod k}$ runs over
reduced residues.

\subsection{Elementary inequalities}

For integers $j\ge1$, $k\ge1$ and $Y\ge1$,
\begin{equation}\label{sp:first}
 \tau(n)^j\le\tau_{2^j}(n),\qquad
 \sum_{n\le Y}\frac{\tau_k(n)}n\le(1+\log Y)^k,\qquad
 \sum_{n\le Y}\tau_k(n)\le Y(1+\log Y)^{k-1}.
\end{equation}
The first is multiplicative. At a prime power it is
$(a+1)^{j}\le\binom{a+2^j-1}{2^j-1}$, which follows by induction on $j$
from
\[
 (a+1)\binom{a+k-1}{k-1}\le\binom{a+2k-1}{2k-1}.
\]
The right side is the coefficient of $p^a$ in the convolution of two
$k$-fold divisor functions. Each of its $a+1$ terms is at least the left
factor $\binom{a+k-1}{k-1}$, since $(j'+i)(a-j'+i)\ge i(a+i)$ for
$0\le j'\le a$. The two sums in \eqref{sp:first} follow by writing
$n=n_1\cdots n_k$, enlarging the constraint to $n_i\le Y$, and summing
one factor last. Also
\begin{equation}\label{sp:phi-tau}
 \frac n{\varphi(n)}\le\tau(n),\qquad
 \varphi(mn)\ge\varphi(m)\varphi(n),\qquad
 \tau(mn)\le\tau(m)\tau(n),
\end{equation}
prime power by prime power, and $\tau(n)\ll_\epsilon n^\epsilon$.

For a finitely supported sequence $\xi$, a modulus $k$ and an integer
$h\ge1$, put
\begin{equation}\label{sp:disc}
 \Delta_\xi(k,b;h)=\sum_{\substack{n\equiv b\ (k)\\(n,h)=1}}\xi_n
 -\frac1{\varphi(k)}\sum_{(n,kh)=1}\xi_n,\qquad
 \Delta_\xi(k,b)=\Delta_\xi(k,b;1),
\end{equation}
\begin{equation}\label{sp:variance}
 V(k;\xi)=\sum_{b\bmod k}^{*}|\Delta_\xi(k,b)|^2 .
\end{equation}
If $\xi$ is supported in an interval of length at most $Y$, then
\begin{equation}\label{sp:projection}
 V(k;\xi)\le\sum_{b\bmod k}\Bigl|\sum_{n\equiv b\ (k)}\xi_n\Bigr|^2
 \le\Bigl(\frac Yk+1\Bigr)\|\xi\|_2^2 .
\end{equation}
Here $\Delta_\xi(k,b)$ is the deviation of the reduced class sums from
their mean. The variance is at most the second moment. Cauchy's
inequality in each class gives the second bound.

The sequence $e_n=\mathbf1_{H<n\le2H}$ satisfies, for all $k,h\ge1$ and
$(b,k)=1$,
\begin{equation}\label{sp:constant-sw}
 |\Delta_e(k,b;h)|\le3\tau(h).
\end{equation}
By inclusion--exclusion over $d\mid h$, $(d,k)=1$, the class count equals
$(H/k)\prod_{p\mid h,\,p\nmid k}(1-1/p)+O(\tau(h))$, with error at most
$\tau(h)$ in absolute value. The comparator has the same main term and error
at most $\tau(h)\tau(k)/\varphi(k)\le2\tau(h)$.

Let $\ell=\log\log N$, $z=N^{1/\ell^3}$ and $y=N^{1/\ell}$. Here $N$ is
the large parameter of the lemma at hand. For each fixed
$C\ge0$, Rankin's method gives
\begin{equation}\label{sp:rankin}
 \sum_{\substack{d>y\\P^+(d)\le z}}\frac{\tau(d)^C}d
 \ll_C(\log N)^{K_C}e^{-\ell^2}
\end{equation}
with a fixed $K_C$. With $\lambda=1/\log z$, the left side is at most
\[
 y^{-\lambda}\prod_{p\le z}\Bigl(1+\sum_{e\ge1}(e+1)^Cp^{-e(1-\lambda)}\Bigr),
 \qquad y^{-\lambda}=e^{-\ell^2}.
\]
For $p\le z$ one has $p^{\lambda}\le e$ and $1-\lambda\ge1/2$, so the factor
at $p$ is $1+O_C(1/p)$. Mertens' theorem gives
$\sum_{p\le z}1/p\le\ell+O(1)$. Hence the product is $\ll(\log N)^{K_C}$.
Since $e^{-\ell^2}=(\log N)^{-\ell}$, the right side of \eqref{sp:rankin}
is $O_{C,A}((\log N)^{-A})$ for every fixed $A$.

\subsection{Lemma 6.5}\label{subsec:supp-lem-65}

Recall the setting of Section~\ref{sec:root-a}. Let $\xi$ be complex, supported in
$[1,Z]$, $Z\ge3$, and fix $B_{\mathrm{div}}\ge0$. The hypothesis (6.35) is
that, for every $P>0$, uniformly in $h,k\ge1$ and reduced $b\bmod k$,
\[
 \left|\sum_{\substack{n\equiv b\ (k)\\(n,h)=1}}\xi_n
 -\frac1{\varphi(k)}\sum_{(n,kh)=1}\xi_n\right|
 \ll_P\|\xi\|_2Z^{1/2}\tau(h)^{B_{\mathrm{div}}}(\log(2Z))^{-P},
 \tag{6.35}
\]
with constants common to the family of pairs $(Z,\xi)$. The conclusion
(6.36) is
\[
 \sum_{k\le Z(\log Z)^{-B_1(A)}}\sum_{b\bmod k}^{*}
       |\Delta_\xi(k,b)|^2
 \ll_A\|\xi\|_2^2 Z(\log Z)^{-A}.
 \tag{6.36}
\]
In the selected-factor form of Paper~I, (3.48), $\xi_n=\alpha_n1_{(n,r)=1}$,
$Z=\lceil2T\rceil$, and (6.37) reads
\[
 \sum_{q\le3Q}\sum_{b\bmod q}^{*}|\Delta_\xi(q,b)|^2
 \ll N_{\mathcal I}T L^{-A_3},\qquad
 3Q\le Z(\log Z)^{-B_F},\quad B_F=A_3+15,
 \tag{6.37}
\]
where $N_{\mathcal I}$ is the interval mass in Paper~I, (3.8).

\begin{mainlemma}{6.5}{Transferred mean-square estimate}
\begin{enumerate}
 \item[(i)] If \textup{(6.35)} holds, then
 \textup{(6.36)} holds for every fixed \(A>0\). Here
 \(B_1(A)\) and the implied constant depend only on \(A\),
 \(B_{\mathrm{div}}\) and the common constants.
 \item[(ii)] Assume the hypotheses (3.1), (3.2), (3.6), (3.8) with its
 length condition, (3.11) and (3.25)--(3.27) of Paper~I. Then the estimate
 \textup{(6.37)} holds uniformly over the parameter
 boxes and selected factors of Paper~I, Theorem~3.2. Hence that theorem
 holds with this input in place of \cite[Theorem~0(a)]{BFI86} at (3.48).
\end{enumerate}
\end{mainlemma}

\begin{proof}[Proof of \textup{(i)}]
Put $E=\|\xi\|_2^2$, $\lambda=\log(2Z)$ and, for $1\le Q\le Z$,
\[
 S(Q)=\sum_{k\le Q}V(k;\xi).
\]
Since $2Z\le Z^2$, $\log Z\le\lambda\le2\log Z$.

\emph{Characters.} Orthogonality on the reduced classes gives
$\Delta_\xi(k,b)=\varphi(k)^{-1}\sum_{\chi\ne\chi_0}\overline{\chi}(b)
\sum_n\xi_n\chi(n)$. The centring removes exactly the principal character.
Hence
\begin{equation}\label{sp:parseval}
 S(Q)=\sum_{k\le Q}\frac1{\varphi(k)}\sum_{\substack{\chi\bmod k\\\chi\ne\chi_0}}
       \Bigl|\sum_n\xi_n\chi(n)\Bigr|^2 .
\end{equation}
Each $\chi\ne\chi_0$ modulo $k$ is induced by a unique primitive
$\psi\bmod f$ with $f\mid k$, $f>1$. Write $k=fe$. Then
$\chi(n)=\psi(n)\mathbf1_{(n,e)=1}$, since $\psi$ already vanishes on
primes of $f$. The pair $(e,\psi)$ determines $(k,\chi)$, and
$\varphi(fe)\ge\varphi(f)\varphi(e)$ by \eqref{sp:phi-tau}. Therefore
\begin{equation}\label{sp:primitive}
 S(Q)\le\sum_{e\le Q}\frac1{\varphi(e)}
 \sum_{1<f\le Q/e}\frac1{\varphi(f)}\sum_{\psi\bmod f}^{\rm prim}
 |C_e(\psi)|^2,\qquad
 C_e(\psi)=\sum_{(n,e)=1}\xi_n\psi(n).
\end{equation}

\emph{Small conductors.} Let $F\ge1$ and $1<f\le F$. Since
$\sum_{b}^{*}\psi(b)=0$,
\[
 C_e(\psi)=\sum_{b\bmod f}^{*}\psi(b)
 \Bigl\{\sum_{\substack{n\equiv b\ (f)\\(n,e)=1}}\xi_n
 -\frac1{\varphi(f)}\sum_{(n,ef)=1}\xi_n\Bigr\}.
\]
The bracket is the left side of (6.35) with $k=f$ and $h=e$. For a fixed
$a>0$ it is at most $C_aE^{1/2}Z^{1/2}\tau(e)^{B_{\rm div}}\lambda^{-a}$.
There are at most $\varphi(f)$ primitive characters modulo $f$, so
\begin{equation}\label{sp:small-cond}
 \sum_{1<f\le F}\frac1{\varphi(f)}\sum_\psi^{\rm prim}|C_e(\psi)|^2
 \le C_a^2EZ\tau(e)^{2B_{\rm div}}\lambda^{-2a}\sum_{f\le F}\varphi(f)^2
 \le C_a^2EZ\tau(e)^{2B_{\rm div}}F^3\lambda^{-2a}.
\end{equation}

\emph{Large conductors.} The multiplicative large sieve, Paper~I, (2.19),
applied to $c_n=\xi_n\mathbf1_{(n,e)=1}$ with support in $[1,Z]$ and
$\|c\|_2^2\le E$, gives
\[
 \sum_{f\le 2V}\sum_\psi^{\rm prim}|C_e(\psi)|^2\ll(V^2+Z)E .
\]
If $p_1<\cdots<p_s$ are the primes of $f$, then $p_i\ge i+1$ and
$f/\varphi(f)\le\prod_i(i+1)/i=s+1\le1+\log f/\log2$. Thus
$1/\varphi(f)\le3\lambda/V$ for $V<f\le2V$, and
\[
 \sum_{V<f\le2V}\frac1{\varphi(f)}\sum_\psi^{\rm prim}|C_e(\psi)|^2
 \ll\lambda E\Bigl(V+\frac ZV\Bigr).
\]
Summing over $V=2^jF<Q$ gives, uniformly in $e$,
\begin{equation}\label{sp:large-cond}
 \sum_{F<f\le Q/e}\frac1{\varphi(f)}\sum_\psi^{\rm prim}|C_e(\psi)|^2
 \ll\lambda E\Bigl(Q+\frac ZF\Bigr).
\end{equation}
If $F\ge Q$, this range is empty.

\emph{Summation over $e$.} Let $r=\lceil2B_{\rm div}+1\rceil$ and
$m=2^r$. By \eqref{sp:phi-tau} and \eqref{sp:first},
\[
 \sum_{e\le Q}\frac{\tau(e)^{2B_{\rm div}}}{\varphi(e)}
 \le\sum_{e\le Q}\frac{\tau(e)^r}e\le(1+\log Q)^m,\qquad
 \sum_{e\le Q}\frac1{\varphi(e)}\le(1+\log Q)^2 .
\]
Inserting \eqref{sp:small-cond} and \eqref{sp:large-cond} in
\eqref{sp:primitive} gives
\begin{equation}\label{sp:bdh-total}
 S(Q)\ll_aE\lambda^{m+3}\Bigl(ZF^3\lambda^{-2a}+\frac ZF+Q\Bigr).
\end{equation}
Put
\begin{equation}\label{sp:b1}
 B_1(A)=A+2^{\lceil2B_{\rm div}+1\rceil}+7,\qquad F=\lambda^{B_1},\qquad a=2B_1 .
\end{equation}
For $Q\le Z(\log Z)^{-B_1}\le2^{B_1}Z\lambda^{-B_1}$ each term in
\eqref{sp:bdh-total} is $\ll EZ\lambda^{m+3-B_1}\le EZ\lambda^{-A}$. This
is (6.36). The constants depend only on $A$, $B_{\rm div}$ and the constant
$C_{2B_1}$ of (6.35). For $B_{\rm div}=1$ one has $r=3$, $m=8$ and
$B_1(A)=A+15$.
\end{proof}

\begin{proof}[Proof of \textup{(ii)}]
Fix a box and a selected factor $r$ of Paper~I, (3.26). Put
$\xi_n=\alpha_n\mathbf1_{(n,r)=1}$, with $\alpha$ as in Paper~I, (3.8).
Then $\xi$ is supported in $I_T\subset(T,2T]\subset[1,Z]$, $Z=\lceil2T\rceil$,
$|\xi_n|\le1$, and $E=\|\xi\|_2^2\le N_{\mathcal I}+1$. Its discrepancy
\eqref{sp:disc} at $h=1$ is $\Delta(\alpha^{(r)};q,b)$ of Paper~I, (3.46).
The support of $\psi_0$ in Paper~I, (3.27), gives $q\le5Q/2<3Q$.

\emph{Low norm.} If $E\le N_{\mathcal I}L^{-A_3-1}$, then
\eqref{sp:projection}, $Q<T$ and $\log(3Q)\le L$ give
\[
 \sum_{q\le3Q}V(q;\xi)\le(N_{\mathcal I}(1+\log3Q)+3Q)E\ll TLE\le N_{\mathcal I}TL^{-A_3}.
\]
This is the first case of Paper~I, (3.48).

\emph{High norm.} Let $E>N_{\mathcal I}L^{-A_3-1}$. For $d,k\ge1$ and
$(l,k)=1$,
\[
 \Delta_\xi(k,l;d)=\sum_{\substack{n\in I_T,\ n\equiv l\ (k)\\(n,dr)=1}}\alpha_n
 -\frac1{\varphi(k)}\sum_{\substack{n\in I_T\\(n,dkr)=1}}\alpha_n ,
\]
which is the discrepancy of Paper~I, (3.25), with $I=I_T$ and auxiliary
divisor $dr$. Hence, for each $A_{\rm SW}>0$,
\begin{equation}\label{sp:pi03-sw}
 |\Delta_\xi(k,l;d)|\ll_{A_{\rm SW}}T\tau(d)\tau(r)L^{-A_{\rm SW}} .
\end{equation}
By Paper~I, (3.26), $r$ is $u_0$-rough and $r\le X^{0.6}\le N^{0.6}$. With
$\log u_0=\ell_0\sigma^2\log N$ from Paper~I, (3.2), the number $\Omega(r)$
of prime factors with multiplicity satisfies
$\Omega(r)\le0.6/(\ell_0\sigma^2)$, so
\[
 \tau(r)\le2^{\Omega(r)}=O_\sigma(1).
\]
By Paper~I, (3.8), $N_{\mathcal I}\asymp TL^{-C_{\rm box}}$, hence
\[
 E^{1/2}Z^{1/2}\gg TL^{-(C_{\rm box}+A_3+1)/2}.
\]
By Paper~I, (3.11), $T\ge X^{60\sigma}$, so $\log(2Z)\asymp_\sigma L$.
Take $A_{\rm SW}=P+(C_{\rm box}+A_3+1)/2$ in \eqref{sp:pi03-sw}. This gives
(6.35) for $\xi$ with $B_{\rm div}=1$, at every precision $P$. The
constants depend only on $\sigma$, $C_{\rm box}$, $A_3$, $P$ and the
constants of Paper~I, (3.25). These are common to all boxes and selected
factors.

By part (i) with $B_{\rm div}=1$ and $A=A_3$, the range parameter is
$B_1(A_3)=A_3+15=B_F$. Paper~I, (3.6), gives $Q<TX^{-29\sigma}$, hence
\[
 3Q<3TX^{-29\sigma}\le Z(\log Z)^{-B_F}
\]
for all sufficiently large $X$. Part (i) yields
\[
 \sum_{q\le3Q}V(q;\xi)\ll EZ(\log Z)^{-A_3}\ll N_{\mathcal I}TL^{-A_3}.
\]
This is (6.37), uniformly in the box and the selected factor.

In the proof of Paper~I, Theorem~3.2, \cite[Theorem~0(a)]{BFI86} is used
only to obtain the second case of (3.48). See the list of inputs after the
statement of that theorem in Paper~I. The two cases above give (3.48) for
every box and selected factor, with the same constants and quantifiers.
The rest of that proof is unchanged.
\end{proof}

\subsection{Lemma 6.8}\label{subsec:supp-lem-68}

Recall the notation of Section~\ref{sec:root-a}. Put $L=\log N$, $\ell=\log\log N$,
\[
 z=N^{1/\ell^3},\qquad y=N^{1/\ell},\qquad
 v=v_{\mathrm{RA}}=\frac{5227}{10000},\qquad w=\frac9{1000},\qquad
 \sigma_{\mathrm p}=10^{-6},\qquad \theta=\frac7{32},
\]
and, for an integer $q$ and a residue $e$,
\[
 \mathcal K_q(e)=\mathbf1_{e\equiv N\ (q)}-\frac{\mathbf1_{(e,q)=1}}{\varphi(q)}
 \tag{6.21}
\]
(the last line of (6.21)). For a parent $(H,2H]$ with
$N^{1/3}/8\le H\le2N^{.46+O(1/L)}$, (6.38) defines
\[
 D_1=HN^{-\sigma_{\mathrm p}},\qquad D_2=QN^{w+\sigma_{\mathrm p}}/H,
 \qquad Q=N^{v}L^{-B}.
 \tag{6.38}
\]
For a sequence $f$ on $(H,2H]$, (6.32) defines $\Delta_f(I;k,b;h)$ as
\eqref{sp:disc} with $n$ restricted to the interval $I$. Then (6.34) is the
bound
\[
 |\Delta_f(I;k,b;h)|\ll_PH\tau(h)L^{-P},\qquad(b,k)=1,
 \tag{6.34}
\]
uniformly for $I\subset(H,2H]$ and $k,h\ge1$.

\begin{mainlemma}{6.8}{Two-factor convolution estimate}
Let \(N^{1/3}/8\le H\le2N^{.46+O(1/L)}\) and \(HM\asymp N\). Let \(a_n\),
\(n\sim H\), be one-bounded, supported on \(P^-(n)>z\), and satisfy
\textup{(6.34)} at the original parent \(H\), where
\(z=N^{1/(\log\log N)^3}\). Let \(b_m\), \(m\sim M\), be one-bounded.
Let \(U_1,V_1\) be dyadic modulus scales, allowing the singleton scale
\(1/2\) for the integer \(1\), such that
\[
 U_1\le D_1,\qquad V_1\le D_2,\qquad U_1V_1\le N^{v_{\mathrm{RA}}}.
\]
Let \(f(r)\), \(r\sim U_1\), and \(g(s)\), \(s\sim V_1\), be arbitrary
one-bounded arrays, with the unary restrictions \(r<D_1\), \(s<D_2\),
\((r,N)=(s,N)=1\) in force. Neither a pointwise condition \(rs<Q\) nor a joint condition
\((r,s)=1\) is imposed on the pair. Then, for every fixed \(A>0\), uniformly in these
occupied-parent scales, \(f\), \(g\), unary interval cuts of \(a_n\) and
\(b_m\), and intervals \(I\subset[N/8,2N]\),
\[
 \sum_{r,s}f(r)g(s)
 \sum_{\substack{n,m\\nm\in I}}a_n b_m\mathcal K_{rs}(nm)\ll_A NL^{-A}.
 \tag{6.45}
\]
\end{mainlemma}

Write $H=N^h$. Fixed-factor enlargements of the parents change $h$ by
$O(1/L)$, so
\begin{equation}\label{sp:l52-h}
 \tfrac13-o(1)\le h\le\tfrac{23}{50}+o(1),\qquad
 M\gg N^{.54-o(1)} .
\end{equation}
All implied constants below may depend on $A$, on the fixed comparison
constants and on the constants in (6.34). They do not depend on $N$, $H$,
$f$, $g$, $a$ or $b$. All fixed precisions are chosen before the final
threshold for $N$.

\subsubsection*{Step 0: reductions}
A unary cut $a_n\mathbf1_J(n)$ satisfies (6.34) with the same constants,
since $\Delta_{a\mathbf1_J}(I;k,b;h)=\Delta_a(I\cap J;k,b;h)$. The error is
still measured at $H$. A unary cut of $b_m$ keeps it one-bounded. The
product cut $nm\in I$ is removed in Step~9. Until then we prove, for every
fixed $A>0$,
\begin{equation}\label{sp:l52-L}
 \mathcal S:=\sum_{r\sim U_1,\ s\sim V_1}f(r)g(s)
 \sum_{n\sim H,\ m\sim M}a_nb_m\mathcal K_{rs}(nm)\ll_ANL^{-A}.
\end{equation}

\subsubsection*{Step 1: a shifted divisor moment and the mean square}
For each fixed $j$ and each interval $J$ of length $O(M)$ in $(M/2,4M]$,
\begin{equation}\label{sp:shifted-moment}
 \sum_{\substack{n\sim H,\ m\in J\\mn\ne N}}\tau(|mn-N|)^j\ll_jMHL^{C_j}.
\end{equation}
Let $1\le u\le N^3$. Its prime factors above $N^{1/8}$ have total
multiplicity at most $24$. Group the others greedily into at most $25$
factors, each at most $N^{1/4}$. Submultiplicativity of $\tau$ gives
\[
 \tau(u)^j\le2^{24j}\sum_{\substack{d\mid u\\d\le N^{1/4}}}\tau(d)^{25j}.
\]
For fixed $n$ and $d$, put $h_0=(d,n)$. The congruence $d\mid mn-N$ has no
solution unless $h_0\mid N$, and otherwise $O(Mh_0/d+1)$ solutions $m\in J$.
Using $(d,n)\le\sum_{e\mid(d,n)}e$ and \eqref{sp:first},
\[
 \sum_{d\le N^{1/4}}\tau(d)^{25j}\Bigl(\frac{M(d,n)}d+1\Bigr)
 \ll L^{C}\bigl(M\tau(n)^{25j+1}+N^{1/4}\bigr).
\]
Summing over $n$ with \eqref{sp:first}, and using $N^{1/4}\le M$, proves
\eqref{sp:shifted-moment}.

Put $Z_H=\lceil2H\rceil$. For each fixed $P>0$ and every $K'\le HN^{-\sigma_{\mathrm p}/2}$,
\begin{equation}\label{sp:l52-ms}
 \sum_{k\le5K'}\tau(k)V(k;a)\ll_PH^2L^{-P}.
\end{equation}
First let $R>0$ be fixed. If $\|a\|_2^2\le HL^{-R}$, then
\eqref{sp:projection} gives
$\sum_{k\le5K'}V(k;a)\ll(H(1+\log5K')+K')\|a\|_2^2\ll H^2L^{1-R}$. If
$\|a\|_2^2>HL^{-R}$, then $\|a\|_2Z_H^{1/2}\ge HL^{-R/2}$. Taking
$I=(H,2H]$ in (6.34) at precision $P'+R/2$ gives (6.35) for $\xi=a$ with
$Z=Z_H$, $B_{\rm div}=1$ and common constants, since $\log(2Z_H)\asymp L$.
This argument is uniform at every requested precision $P$ (after choosing
$P'$ accordingly).
Lemma~6.5(i) with $A=P'$ and $B_1(P')=P'+15$ gives
\[
 \sum_{k\le Z_H(\log Z_H)^{-P'-15}}V(k;a)\ll\|a\|_2^2Z_H(\log Z_H)^{-P'}
 \ll H^2L^{-P'}.
\]
The range contains $k\le5K'$, because $5K'/H\le5N^{-\sigma_{\mathrm p}/2}$
beats every fixed power of $L$. Thus in both cases
$\sum_{k\le5K'}V(k;a)\ll H^2L^{-\min(R-1,P')}$. To insert $\tau(k)$, split at
$\tau(k)=L^{R'}$. Below the split the cost is $L^{R'}$. Above it,
$\tau(k)\le L^{-R'}\tau(k)^2$, and $V(k;a)\le(2H/k+1)(H+1)$ by
\eqref{sp:projection}; hence by \eqref{sp:first}
\[
 \sum_{\substack{k\le5K'\\\tau(k)>L^{R'}}}\tau(k)V(k;a)
 \ll L^{-R'}H^2\sum_{k\le5K'}\frac{\tau(k)^2}k\ll H^2L^{4-R'} .
\]
Choosing $R'=P+4$ and then $R,P'\ge P+R'+1$ proves \eqref{sp:l52-ms}.

For a $z$-rough integer $t\le N^{C}$ put $a^{(t)}_n=a_n\mathbf1_{(n,t)=1}$.
Since $\omega(t)\le\log t/\log z\ll\ell^3$ and each prime $p\mid t$ exceeds
$z$,
\[
 \|a-a^{(t)}\|_2^2\le\sum_{p\mid t}\Bigl(\frac{2H}p+1\Bigr)\ll\frac{HL}z ,
\]
and by \eqref{sp:projection} and \eqref{sp:first}, uniformly in $t$,
\begin{equation}\label{sp:l52-mask}
 \sum_{k\le5K'}\tau(k)V(k;a-a^{(t)})\ll\frac{H^2L^3}z .
\end{equation}

\subsubsection*{Step 2: large smooth parts}
Write $s=dt$ with $P^+(d)\le z<P^-(t)$. Consider the pairs $(r,s)$ with
$d>y$. Then the $z$-smooth part of $q=rs$ exceeds $y$. A modulus $q$ has at
most $\tau(q)$ factorizations $q=rs$, and $q\le4U_1V_1\le4N^v$. For
$(q,N)=1$, counting $m$ for each $n$ coprime to $q$ gives
\[
 \Bigl|\sum_{n,m}a_nb_m\mathcal K_q(nm)\Bigr|\ll\frac Nq+H+\frac N{\varphi(q)} .
\]
By \eqref{sp:phi-tau}, the contribution of these pairs to $\mathcal S$ is
\[
 \ll N\sum_{\substack{q\le4N^v\\ q=d'h',\ d'>y,\ P^+(d')\le z}}\frac{\tau(q)^2}q
 +H\sum_{q\le4N^v}\tau(q)
 \ll NL^{C}\sum_{\substack{d'>y\\P^+(d')\le z}}\frac{\tau(d')^2}{d'}
 +HN^vL .
\]
Here $\tau(q)^2\le\tau(d')^2\tau(h')^2$ and the sum over $h'$ is bounded
by \eqref{sp:first}. By \eqref{sp:rankin} the first term is
$O_A(NL^{-A})$. The second is at most $N^{.9827+o(1)}$, since
$h+v\le.46+.5227+o(1)$. It remains to treat $d\le y$.

\subsubsection*{Step 3: the first Cauchy square}
Split $d$ and $t$ into dyadic shells $d\sim D$, $t\sim S$, with $D\le y$,
$DS\asymp V_1$. There are $O(L)$ choices of $D$ and $O(1)$ of $S$. Put
\[
 \lambda_d(t)=g(dt)\mathbf1_{dt\sim V_1}\mathbf1_{P^-(t)>z}\mathbf1_{P^+(d)\le z},
 \qquad K=U_1D,\qquad k=rd .
\]
Then $|\lambda_d(t)|\le1$, $rs=kt$, and $(kt,N)=1$ on the supports. One
block of $\mathcal S$ is
\[
 \mathcal S_0=\sum_{r,d,m}f(r)b_m\sum_t\lambda_d(t)\sum_na_n\mathcal K_{kt}(mn).
\]
Fix a smooth $\psi\ge0$ supported in $[1/2,5/2]$, equal to $1$ on $[1,2]$.
Cauchy's inequality in $(r,d,m)$ gives
\begin{equation}\label{sp:l52-cs1}
 |\mathcal S_0|^2\le CKM\,\mathcal S',\qquad
 \mathcal S'=\sum_{r,d}\psi\Bigl(\frac r{U_1}\Bigr)\mathbf1_{(rd,N)=1}\mathbf1_{d\sim D}
 \sum_m\psi\Bigl(\frac mM\Bigr)
 \Bigl|\sum_t\lambda_d(t)\sum_na_n\mathcal K_{kt}(mn)\Bigr|^2 .
\end{equation}
The indicator $\mathbf1_{(rd,N)=1}$ is part of the original support and
is kept. The scale conditions and \eqref{sp:l52-h} give, for large $N$,
\begin{equation}\label{sp:l52-scales}
 \frac KH\le N^{-\sigma_{\mathrm p}/2},\qquad
 \frac{KS}M\le N^{-.0172},\qquad M\gg N^{.54-o(1)},
\end{equation}
because $K\le D_1y$, $KS\asymp U_1V_1\le N^v$ and $v-.54=-.0173$. In
particular $kt=o(M)$ and $K=o(H)$.

Expand $\mathcal S'=\mathcal S_3-2\operatorname{Re}\mathcal S_2+\mathcal S_1$
over tuples $\mathbf t=(r,d,t_1,t_2,n_1,n_2)$ with coefficient
\[
 c(\mathbf t)=\psi(r/U_1)\mathbf1_{(rd,N)=1}\mathbf1_{d\sim D}
 \lambda_d(t_1)\overline{\lambda_d(t_2)}a_{n_1}\overline{a_{n_2}},\qquad
 \omega(\mathbf t)=\mathbf1_{(n_1,kt_1)=(n_2,kt_2)=1}:
\]
\[
\begin{aligned}
 \mathcal S_3&=\sum_{\mathbf t}c\sum_m\psi(m/M)\mathbf1_{kt_1\mid mn_1-N}\mathbf1_{kt_2\mid mn_2-N},\\
 \mathcal S_2&=\sum_{\mathbf t}c\sum_m\psi(m/M)\mathbf1_{kt_1\mid mn_1-N}
   \frac{\mathbf1_{(mn_2,kt_2)=1}}{\varphi(kt_2)},\\
 \mathcal S_1&=\sum_{\mathbf t}c\sum_m\psi(m/M)
   \frac{\mathbf1_{(mn_1,kt_1)=1}\mathbf1_{(mn_2,kt_2)=1}}{\varphi(kt_1)\varphi(kt_2)} .
\end{aligned}
\]
In $\mathcal S_3$ the factor $\omega$ is automatic, since $(kt_i,N)=1$.

\subsubsection*{Step 4: restoration before the Chinese remainder theorem}
Call a tuple \emph{gcd-bad} if one of
\begin{equation}\label{sp:bad}
 (n_1,n_2)>1,\quad(t_1,t_2)>1,\quad(n_1,t_2)>1,\quad(n_2,t_1)>1,\quad(k,t_1t_2)>1
\end{equation}
 holds. Call a tuple \emph{near-diagonal} if
$|n_1-n_2|<HL^{-C_M}$; this condition is charged separately below. On the
other tuples $k,t_1,t_2$ are pairwise coprime. For a set $\mathcal B$ of
tuples define the formal terms
\begin{equation}\label{sp:formal}
 MT(\mathcal B)=\sum_{\mathcal B}\frac{c\,\omega\,M\hat\psi(0)}{kt_1t_2\varphi(k)},\qquad
 \mathcal Z(\mathcal B)=\sum_{\mathcal B}\frac{c\,\omega\,\mathbf1_{n_1\equiv n_2\ (k)}M\hat\psi(0)}{kt_1t_2}.
\end{equation}
We show
\begin{equation}\label{sp:restoration}
\begin{gathered}
 \mathcal S'=M\mathcal E+\mathcal Z(\mathrm{all})-MT(\mathrm{all})+O(R_0),\\
 R_0\ll_P\frac{MH^2}KL^{-P}+MHL^{C}\Bigl[\frac H{zK}+1\Bigr]
   +N^{o(1)}\Bigl[\frac{H^2}K+H^2S\Bigr]+N^{o(1)}\frac{(H+KS)^2}K ,
\end{gathered}
\end{equation}
where $\mathcal E$ is the nonzero-frequency sum on good tuples defined in
Step~6. The exponent $C$ is fixed after $P$ and the truncation parameters.

\emph{Resonant $m$.} If $mn_i=N$, then $m\mid N$. There are
$\tau(N)=N^{o(1)}$ such $m$. For fixed $(r,d)$ and such $m$, the
congruence terms over $(t,n)$ number at most $\sum_t(2H/(kt)+1)\ll H/K+S$.
The comparator terms are $\ll L^{C}H/K$. Hence these $m$ contribute
$\ll N^{o(1)}K(H/K+S)^2$ to $\mathcal S_1,\mathcal S_2,\mathcal S_3$ and to every subset of
tuples below. This is the last term of $R_0$.

\emph{Divisor truncation.} For the other $m$ put $u_i=|mn_i-N|\ge1$. Let
$\Sigma^*$ be the event $\tau(u_1),\tau(u_2)\le L^{C_\Sigma}$. On the part of
the complement with $\tau(u_1)\ge\tau(u_2)$, the triples $(r,d,t_1)$ with
$kt_1\mid u_1$ number at most $\tau(u_1)^3$. The $t_2\mid u_2$ number at
most $\tau(u_1)$, and $n_2$ lies in one class modulo $k$, with $O(H/K)$
choices. With $\mathbf1_{\tau(u_1)>L^{C_\Sigma}}\le\tau(u_1)L^{-C_\Sigma}$,
\eqref{sp:shifted-moment} gives a contribution to $\mathcal S_3$ of
\begin{equation}\label{sp:trunc}
 \ll L^{-C_\Sigma}\frac HK\sum_{m,n_1}\psi(m/M)\tau(u_1)^5\ll\frac{MH^2}KL^{C_5-C_\Sigma}.
\end{equation}
The other part is symmetric. Choose $C_\Sigma\ge C_5+P$.

\emph{Bad tuples in $\mathcal S_3$.} In this paragraph, ``bad'' means
$\Sigma^*$ intersected with the five gcd failures in \eqref{sp:bad}; the
near-diagonal set is not included in this term and is charged separately
below. On $\Sigma^*$ fix $(m,n_1)$. There are at
most $L^{3C_\Sigma}$ triples $(r,d,t_1)$ and at most $O(L)$ candidate primes
$p$ for a failure in \eqref{sp:bad}. Once $n_2$ is fixed, there are at most
$L^{C_\Sigma}$ choices of $t_2$. Every witnessing prime exceeds $z$: it divides a
$z$-rough $n_i$ or $t_i$. Since $(k,N)=1$ and $k\mid mn_1-N$, we have
$(m,k)=(n_1,k)=1$ and $n_2\equiv n_1\pmod k$. Each failure places $n_2$ in
one class modulo $kp$:
\begin{itemize}
 \item $p\mid(n_1,n_2)$: $p\nmid k$; combine $n_2\equiv0\ (p)$ with $n_2\equiv n_1\ (k)$.
 \item $p\mid(t_1,t_2)$, $p\nmid k$: $p\nmid m$, since $p\mid mn_1-N$ and $p\nmid N$;
 combine $mn_2\equiv N\ (p)$ with $n_2\equiv n_1\ (k)$. If $p\mid k$, then
 $kt_2\mid mn_2-N$ and $p\mid t_2$ give $kp\mid mn_2-N$.
 \item $p\mid(n_1,t_2)$: $p\nmid k$ and $p\nmid m$; combine $mn_2\equiv N\ (p)$ with $n_2\equiv n_1\ (k)$.
 \item $p\mid(n_2,t_1)$: $p\nmid k$; combine $n_2\equiv0\ (p)$ with $n_2\equiv n_1\ (k)$.
 \item $p\mid(k,t_1t_2)$: since $p>z\ge P^+(d)$, $p\mid r$. If $p\mid t_2$, then
 $kt_2\mid mn_2-N$ gives $kp\mid mn_2-N$. If $p\mid t_1$ only, exchange
 the two copies first.
\end{itemize}
The number of $n_2$ in one class modulo $kp$ is $O(H/(zK)+1)$. Hence
\begin{equation}\label{sp:bad-s3}
 \mathcal S_3(\text{bad})\ll MHL^{4C_\Sigma+1}\Bigl[\frac H{zK}+1\Bigr].
\end{equation}
The near-diagonal $|n_1-n_2|<HL^{-C_M}$ is counted in the same way, with
$H/z$ replaced by $HL^{-C_M}$. Choose $C_M\ge P+4C_\Sigma+1$. This gives the
first two terms of $R_0$. No joint condition $(r,s)=1$ has been used. A
prime $p\le z$ dividing $(r,s)$ divides $d$ and only enlarges $k$. A
prime $p>z$ dividing $(r,s)$ divides $k$ and some $t_i$. This is the last
case of \eqref{sp:bad}.

\emph{Bad tuples in the density terms.} For $\mathcal S_1$, $\mathcal S_2$ and
$MT$, dropping coprimality leaves the positive density
\[
 \ll L^{C_\varphi}\frac M{k^2t_1t_2},
\]
since $kt_1=o(M)$ and $q/\varphi(q)\le L^{C_\varphi}$ for $q\le N^{O(1)}$. (In
$\mathcal S_2$ the $m$-sum is nonzero only if $n_1$ is invertible modulo
$kt_1$, by $(k,N)=(t_1,N)=1$ and $kt_1\mid mn_1-N$. Then $m$ lies in one
class modulo $kt_1$.) For $\mathcal Z$ retain $\omega$ and
$n_2\equiv n_1\ (k)$, with weight $M/(kt_1t_2)$. In particular
$(n_1,k)=(n_2,k)=1$ there. Without the bad condition these weights sum to
$\ll L^{C_\varphi}MH^2/K$ and $\ll MH^2/K$, using $\sum_{t\sim S}1/t\le1$,
$k=rd\asymp K$ and, in $\mathcal Z$, $O(H/K)$ values of $n_2$ in one class
modulo $k$. As in $\mathcal S_3$ every witnessing prime $p$ exceeds $z$. In
each case below the prime $p$ divides $n_1$ or $t_1$, which is fixed first
together with $p$. Since $\omega(n_1),\omega(t_1)\ll L$, there are $O(L)$
choices of $p$. We use
\[
\begin{gathered}
 \#\{n_2\sim H:\ p\mid n_2\}\le\frac Hp+1,\qquad
 \sum_{\substack{t_2\sim S\\p\mid t_2}}\frac1{t_2}\le\frac1S\cdot\frac{2S}p=\frac2p,\\
 \#\{r\le\tfrac52U_1:\ p\mid r\}\le\frac{5U_1}{2p};
\end{gathered}
\]
the last two have no $+1$, because they count multiples of $p$ in
$[1,2S]$ and $[1,\frac52U_1]$.
\begin{itemize}
 \item $p\mid(n_1,n_2)$: fix $n_1$ and $p\mid n_1$. In $\mathcal S_1,\mathcal S_2,MT$
 there are $\le H/z+1$ values of $n_2$. In $\mathcal Z$, $p\nmid k$ since
 $(n_1,k)=1$; combining $n_2\equiv0\ (p)$ with $n_2\equiv n_1\ (k)$ leaves
 $O(H/(zK)+1)$ values of $n_2$.
 \item $p\mid(t_1,t_2)$: fix $t_1$ and $p\mid t_1$; the $t_2$-sum is at most
 $2/p\le2/z$. The $n_2$-count is unchanged; $p\mid k$ is allowed.
 \item $p\mid(n_1,t_2)$: fix $n_1$ and $p\mid n_1$; the $t_2$-sum is at most
 $2/z$, and the $n_2$-count is unchanged.
 \item $p\mid(n_2,t_1)$: fix $t_1$ and $p\mid t_1$. In $\mathcal S_1,\mathcal S_2,MT$
 there are $\le H/z+1$ values of $n_2$. In $\mathcal Z$, $p\mid k$ would
 contradict $(n_2,k)=1$, so $p\nmid k$; combining $n_2\equiv0\ (p)$ with
 $n_2\equiv n_1\ (k)$ leaves $O(H/(zK)+1)$ values of $n_2$.
 \item $p\mid(k,t_1t_2)$: the weights are symmetric in the two copies, so let
 $p\mid t_1$; fix $t_1$ and $p\mid t_1$. Since $p>z\ge P^+(d)$, $p\mid r$, and
 there are at most $5U_1/(2p)\le5U_1/(2z)$ values of $r$. The $d$-, $t_2$-
 and $n$-counts are unchanged.
\end{itemize}
In every case the bad part of each density term is
$\ll MHL^{C_\varphi+1}[H/(zK)+1]$. This is part of $R_0$. The near-diagonal is
counted in the same way, with $H/z$ replaced by $HL^{-C_M}$. There are
$O(HL^{-C_M}+1)$ values of $n_2$, and $O(HL^{-C_M}/K+1)$ in one class
modulo $k$. Hence the near-diagonal is $\ll MH^2L^{C_\varphi-C_M}/K+MHL^{C_\varphi}$. This is part of $R_0$ once also
$C_M\ge P+C_\varphi$. Thus all bad tuples may be removed from
$\mathcal S_1,\mathcal S_2,\mathcal S_3,\mathcal Z,MT$ at the cost $R_0$.

\emph{Good tuples.} On good tuples the two congruences in $\mathcal S_3$ are
compatible exactly on the support $\omega(\mathbf t)=1$ and
$n_1\equiv n_2\pmod k$. Then $m$ lies in one
class $\mu$ modulo $kt_1t_2$, with
\begin{equation}\label{sp:mu}
 \frac\mu{kt_1t_2}\equiv\frac{N\overline{n_1t_2}}{kt_1}+\frac{N\overline{n_2kt_1}}{t_2}\pmod1 .
\end{equation}
Poisson summation gives
\begin{equation}\label{sp:poisson-m}
 \sum_{m\equiv\mu\ (kt_1t_2)}\psi\Bigl(\frac mM\Bigr)
 =\frac M{kt_1t_2}\sum_{h\in\mathbb Z}\hat\psi\Bigl(\frac{hM}{kt_1t_2}\Bigr)
  e\Bigl(\frac{h\mu}{kt_1t_2}\Bigr).
\end{equation}
The term $h=0$ gives $\mathcal Z(\mathrm{good})$. All three expanded
summands vanish unless $\omega(\mathbf t)=1$. The following main-term
computations are on that support. For $\mathcal S_1$ the smooth
coprime $m$-sum is $M\hat\psi(0)\varphi(kt_1t_2)/(kt_1t_2)+O(\tau(kt_1t_2))$.
Division by $\varphi(kt_1)\varphi(kt_2)=\varphi(k)^2\varphi(t_1)\varphi(t_2)$
gives $MT(\mathrm{good})$ with total error $N^{o(1)}H^2/K$. For $\mathcal S_2$
expand $\mathbf1_{(m,t_2)=1}$ by M\"obius inversion. The main term is
$M\hat\psi(0)\varphi(t_2)/(kt_1t_2)$, with error $O(\tau(t_2))$.
Division by $\varphi(kt_2)=\varphi(k)\varphi(t_2)$ gives $MT(\mathrm{good})$
with total error $N^{o(1)}H^2S$. Completing $\mathcal Z$ and $MT$ back to all
tuples costs $R_0$. This proves \eqref{sp:restoration}.

\subsubsection*{Step 5: the zero mode}
Expanding the centred class counts gives the exact identity, valid for all
tuples including those with gcd failures,
\begin{equation}\label{sp:zero-mode}
 \mathcal Z(\mathrm{all})-MT(\mathrm{all})
 =M\hat\psi(0)\sum_{r,d,t_1,t_2}
 \frac{\psi(r/U_1)\mathbf1_{(rd,N)=1}\lambda_d(t_1)\overline{\lambda_d(t_2)}}{kt_1t_2}
 \sum_{b\bmod k}^{*}\Delta_{a^{(t_1)}}(k,b)\overline{\Delta_{a^{(t_2)}}(k,b)} .
\end{equation}
For fixed $(k,t_1,t_2)$ both sides equal
\[
 \sum_{b\bmod k}^{*}A_1(b)\overline{A_2(b)}
 -\frac1{\varphi(k)}\sum_{b\bmod k}^{*}A_1(b)\sum_{b\bmod k}^{*}\overline{A_2(b)},
 \qquad A_i(b)=\sum_{\substack{n\equiv b\ (k)\\(n,t_i)=1}}a_n .
\]
Use
$|u\bar v|\le(|u|^2+|v|^2)/2$ and
\[
 V(k;a^{(t)})\le2V(k;a)+2V(k;a-a^{(t)}).
\]
The sums $\sum_{t\sim S}1/t$ are $O(1)$, a given $k$ has at most $\tau(k)$
factorizations $k=rd$, and $k\le5K$. By \eqref{sp:l52-ms} with $K'=K$ and
\eqref{sp:l52-mask},
\begin{equation}\label{sp:zero-mode-bound}
 |\mathcal Z(\mathrm{all})-MT(\mathrm{all})|\ll_P\frac{MH^2}K\Bigl[L^{-P}+\frac{L^3}z\Bigr].
\end{equation}
Only $|\lambda_d|\le1$ is used, so $\lambda_d$ may vary with $d$. The factor
$1/K$ is cancelled by $KM$ in \eqref{sp:l52-cs1}.

By \eqref{sp:l52-cs1}, \eqref{sp:restoration}, \eqref{sp:zero-mode-bound} and
\eqref{sp:l52-scales}, the contributions to $|\mathcal S_0|$ of all terms
except $M\mathcal E$ are
\begin{equation}\label{sp:nonspectral}
 \ll NL^{-P/2}+NL^{C}\Bigl[z^{-1/2}+\Bigl(\frac KH\Bigr)^{1/2}\Bigr]
 +N^{1+o(1)}\Bigl[M^{-1/2}+\Bigl(\frac{KS}M\Bigr)^{1/2}\Bigr]+N^{.857}.
\end{equation}
For the last term, $M^{1/2}(H+KS)\le N^{(1+h)/2+o(1)}+N^{(1-h)/2+v+o(1)}$,
and $(1+h)/2\le.73+o(1)$, $(1-h)/2+v\le.856034$.

\subsubsection*{Step 6: nonzero frequencies}
Fix $d\sim D$ and let $\mathcal E_d$ be the part of $\mathcal E$ with this
$d$:
\[
 \mathcal E_d=\sum_{\substack{\text{good }\mathbf t,\ \omega(\mathbf t)=1\\n_1\equiv n_2\ (k)}}\frac{c(\mathbf t)}{kt_1t_2}
 \sum_{h\ne0}\hat\psi\Bigl(\frac{hM}{kt_1t_2}\Bigr)e\Bigl(\frac{h\mu}{kt_1t_2}\Bigr).
\]
Put $\eta=10^{-9}$ and
\begin{equation}\label{sp:hstar}
 H_*=C\frac{KS^2N^\eta}M ,
\end{equation}
with $C$ larger than every $kt_1t_2/(KS^2)$. Since
$|\hat\psi(\xi)|\ll_j|\xi|^{-j}$, the terms $|h|>H_*$ contribute $O(N^{-90})$
in total, taking $j>200/\eta$. If $H_*<1$, $\mathcal E_d$ is negligible.
This covers $t=1$, $S=1/2$, where $H_*\ll KN^\eta/M\le N^{-.07}$. Hence we may
assume $S\ge1$ and $H_*\ge1$.

\emph{Reciprocity.} For coprime $u,v$, $\bar u/v+\bar v/u\equiv1/(uv)\pmod1$.
Applied to the first term of \eqref{sp:mu}, with $u=n_1t_2$ and $v=kt_1$, it
gives
\[
 e\Bigl(\frac{h\mu}{kt_1t_2}\Bigr)
 =e\Bigl(\frac{Nh(n_1-n_2)\overline{n_2kt_1}}{n_1t_2}\Bigr)
  e\Bigl(\frac{Nh}{kt_1t_2n_1}\Bigr).
\]
Put $f_0=(n_1-n_2)/k$, a nonzero integer on good tuples. Cancelling $k$,
\begin{equation}\label{sp:phase}
 e\Bigl(\frac{Nhf_0\overline{n_2t_1}}{n_1t_2}\Bigr).
\end{equation}
The conditions on $t_1$ are
\begin{equation}\label{sp:inner-cond}
 (t_1,Nc)=1,\qquad(t_1,J)=1,\qquad b=n_2,\quad c=n_1t_2,\quad J=kn_2 .
\end{equation}
The first is kept. To separate the second, colour every prime in
$[z,N^3]$ blue or red independently with probability $1/2$. Since $t_1$ is
$z$-rough,
\begin{equation}\label{sp:colour}
 \mathbf1_{(t_1,J)=1}=\mathbb E_{\rm col}\Bigl[2^{\omega(t_1)+\omega_z(J)}
 \mathbf1_{\text{all primes of }t_1\text{ blue}}\,
 \mathbf1_{\text{all primes}\ \ge z\text{ of }J\text{ red}}\Bigr],
\end{equation}
where $\omega_z(J)$ counts the primes $\ge z$ of $J$. A shared prime makes
the event empty. Otherwise its probability is the reciprocal of the power of
$2$. Both counts are $O(\ell^3)$, so the multipliers are
$\exp(O(\ell^3))=N^{o(1)}$. Fix one colouring.

Split $f_0$ and $h$ into signed dyadic blocks $|f_0|\sim F$ and
$h=\varepsilon h_+,\ h_+\sim H_0>0$, with $F\ll H/K$ and
$H_0\le H_*$. Below, $h_+$ is renamed $h>0$. Fix $d$, this block, a
sign $\varepsilon$, and one prime colouring from \eqref{sp:colour}.
The \emph{outer} tuple is
$\mathbf o=(r,n_1,n_2,t_2)$, with $k=rd$,
$b=n_2$, $c=n_1t_2$ and $f_0=(n_1-n_2)/k$. Its support contains
$k\mid n_1-n_2$, the near-diagonal exclusion, and the conditions on
$(r,d,n_1,n_2,t_2)$ inherited from the good tuples and $\omega=1$.
It also requires every prime $\ge z$ of $J=kn_2$ to be red.
The remaining $t_1$-conditions are $(t_1,Nc)=1$, $t_1\sim S$ and
all primes of $t_1$ blue. The colouring enforces $(t_1,J)=1$.
Write
\[
 u=\frac{hS}{H_0t_1},\quad
 Y_2(\mathbf o)=\frac{\varepsilon H_0M}{kt_2S},\quad
 Y_1(\mathbf o)=\frac{\varepsilon NH_0}{kt_2n_1S},\quad
 W_{\mathbf o}(u)=w_0(u)\hat\psi(Y_2(\mathbf o)u)e(Y_1(\mathbf o)u),
\]
where $w_0$ is a fixed smooth cutoff supported in $(1/4,4)$ and equal
to $1$ on $[1/2,2]$. Since $H_0\le H_*$, $k\asymp K$,
$t_2\asymp S$, and $N/(n_1M)\asymp1$, uniformly in \emph{all}
outer tuples,
$|Y_1(\mathbf o)|+|Y_2(\mathbf o)|\ll N^\eta$.
On the original support, the remaining Fourier weight is exactly
$W_{\mathbf o}(hS/(H_0t_1))$.
Define the Mellin transform and its uniform majorant by
\[
 \widetilde W_{\mathbf o}(it)=\int_0^\infty
 W_{\mathbf o}(u)u^{it}\,\frac{du}{u},\qquad
 m(t)=C_{\psi,w_0}\min\{1,N^{2\eta}(1+|t|)^{-2}\}.
\]
The compact support, two integrations by parts, and the bounds on
$Y_1,Y_2$ give
\[
 \sup_{\mathbf o,t}\frac{|\widetilde W_{\mathbf o}(it)|}{m(t)}\le1,
 \qquad\int_{\mathbb R}m(t)\,dt\ll_{\psi,w_0}N^\eta,
\]
with the fixed constant $C_{\psi,w_0}$ chosen large enough.
Mellin inversion is the \emph{outer-dependent} identity
\[
 W_{\mathbf o}\Bigl(\frac{hS}{H_0t_1}\Bigr)
 =\frac1{2\pi}\int_{\mathbb R}m(t)
 \underbrace{\Bigl(\frac{H_0}{S}\Bigr)^{it}
 \frac{\widetilde W_{\mathbf o}(it)}{m(t)}}_{\text{outer coefficient}}
 \underbrace{h^{-it}t_1^{it}}_{\text{inner factors}}\,dt.
\]
In particular $\widetilde W_{\mathbf o}(it)/m(t)$ is not
an inner $t_1$-weight: it depends on $k,t_2,n_1$ and the fixed
sign and block.

\emph{The $t$-integral precedes Cauchy.}
Let $B_{\rm in}\ge\sup_{t_1}2^{\omega(t_1)}$ and
$B_{\rm out}\ge\sup_{\mathbf o}2^{\omega_z(kn_2)}$ on this block,
with $B_{\rm in}B_{\rm out}=\exp(O(\ell^3))=N^{o(1)}$.
For fixed real Mellin parameter $t$ set
\[
\begin{aligned}
 \lambda'_{v}(t)&=\frac{S}{v}\lambda_d(v)
 \frac{2^{\omega(v)}}{B_{\rm in}}
 \mathbf1_{\text{all primes of }v\text{ blue}}\,v^{it},\\
 \eta_h(t)&=h^{-it}\mathbf1_{h\sim H_0}\mathbf1_{h\le H_*},\\
 A_{\mathbf o}(t)&=\frac{KS}{kt_2}\,
 \psi(r/U_1)\overline{\lambda_d(t_2)}a_{n_1}\overline{a_{n_2}}
 \frac{2^{\omega_z(kn_2)}}{B_{\rm out}}
 \mathbf1_{\text{all primes }\ge z\text{ of }kn_2\text{ red}}\,
 \Bigl(\frac{H_0}{S}\Bigr)^{it}
 \frac{\widetilde W_{\mathbf o}(it)}{m(t)}.
\end{aligned}
\]
The factor $\mathbf1_{h\le H_*}$ is the one-bounded mask on a terminal
frequency block; it is independent of the outer tuple.
The sum over $\mathbf o$ is restricted to the outer support just
specified. Choose the fixed cutoff $\psi$ with $0\le\psi\le1$.
On the dyadic support $KS/(kt_2)\le2$. By the definitions of $m$, $B_{\rm in}$
and $B_{\rm out}$, every other factor above
has modulus at most one. Hence $|A_{\mathbf o}(t)|\le2$ and
$|\lambda'_v(t)|,|\eta_h(t)|\le1$, uniformly in the real $t$
and all original outer parameters.
Using $1/(kt_1t_2)=(KS^2)^{-1}(KS/(kt_2))(S/t_1)$,
the contribution of this block and colouring to $\mathcal E_d$ is
\[
 \frac{B_{\rm in}B_{\rm out}}{2\pi KS^2}
 \int_{\mathbb R}m(t)
 \sum_{\mathbf o}A_{\mathbf o}(t)
 \sum_{\substack{v\sim S\\(v,Nc)=1}}\ \sum_{h\sim H_0}
 \lambda'_v(t)\eta_h(t)
 e\Bigl(\frac{\varepsilon Nhf_0\overline{bv}}c\Bigr)\,dt.
\]
Here $(b,c)=1$ on the good support. The colour average and the
$O(L^2)$ signed dyadic blocks are taken \emph{after} the estimate.
Their cost is $N^{o(1)}$.

\emph{Second square, pointwise in $t$.} Group the outer tuples by
$(b,c,f_0)$. For fixed $(b,c,f_0)$, the factorization $n_1t_2=c$ has at most
$\tau(c)=N^{o(1)}$ choices. Then
$k=(n_1-b)/f_0$ and $r=k/d$ are determined. The number of outer
tuples is $\ll N^{o(1)}H^2S$. Let $\rho\ge0$ be a smooth majorant,
at least $1$ on their ranges $b\ll H$, $c\ll HS$.
Since $b=n_2$ is $z$-rough, $(b,p_0)=1$ for the least prime $p_0$
not dividing $N$, as $p_0\ll\log N<z$. Define, for each real $t$,
\begin{equation}\label{sp:cs2}
\begin{aligned}
 W_4(t)&=\sum_{\substack{(b,c)=1\\(b,p_0)=1}}
 \rho(b,c)\sum_{f_0\sim F}
 \Bigl|\sum_{\substack{v\sim S\\(v,Nc)=1}}
 \sum_{h\sim H_0}\lambda'_v(t)\eta_h(t)
 e\Bigl(\frac{\varepsilon Nhf_0\overline{bv}}c\Bigr)\Bigr|^2,\\
 |\mathcal E_d|&\ll N^{o(1)}\frac{H}{KS^{3/2}}
 \int_{\mathbb R}m(t)W_4(t)^{1/2}\,dt
 \le N^{\eta+o(1)}\frac{H}{KS^{3/2}}
 (W_4^*)^{1/2},\\
 W_4^*&=\sup_{\substack{\mathrm{signed\ blocks},\ \mathrm{colourings}\\t\in\mathbb R}}W_4(t).
\end{aligned}
\end{equation}
Here $W_4(t)$ is formed with a fixed block and colouring;
the supremum is taken afterwards. The displayed bound for the full
$\mathcal E_d$ follows from the colour expectation and the sum over
the $O(L^2)$ blocks. Their cost is a uniform $N^{o(1)}$. Thus no outer
coefficient is moved through Cauchy or the $t$-integral.
Cauchy is applied \emph{inside} the integral for each $t$.
No outer-dependent Mellin factor appears in $\lambda'_v(t)$.

For every fixed $t$, expand the square with $z_1=vv'$,
$k_0=h_1v'-h_2v$ and $y_1=f_0k_0$. Since
$|\lambda'_v(t)|,|\eta_h(t)|\le1$, the diagonal $k_0=0$ obeys,
uniformly in $t$,
\begin{equation}\label{sp:diag}
 W_{\rm diag}(t)\le N^{o(1)}H_*S\cdot H\cdot F\cdot HS
 \le N^{\eta+o(1)}\frac{H^4S^4}N .
\end{equation}
For $k_0\ne0$, write $m=|y_1|\ll FH_*S$ and separate the two
signs of $y_1$. The off-diagonal is a sum of expressions
\begin{equation}\label{sp:offdiag}
 \sum_{z_1,m}b^{\pm}_{z_1,m}(t)
 \sum_{\substack{(b,c)=(b,p_0)=(z_1,c)=1}}\rho(b,c)
 e\Bigl(\pm\frac{Nm\overline{z_1b}}c\Bigr),
\end{equation}
where $b^{\pm}_{z_1,m}(t)$ depends only on
$(d,S,F,H_0,\varepsilon,\text{colouring},t;z_1,m)$, and not on
$(b,c,f_0)$ or the original outer tuple. Its summands are products
of the one-bounded weights $\lambda'_v(t),\eta_h(t)$, and
$|b^{\pm}_{z_1,m}(t)|\le B_{z_1,\pm m}$ for the counting
majorant defined below. The support has $(z_1,Np_0)=1$.

\emph{Coefficient norm.} Fix $C>0$. For $1\le S\le N^C$, $1\le H'\le S$ and
$1/2\le F\le N^C$ let $B_{z_1,y_1}$ count $t,t'\sim S$, $(tt',N)=1$,
$1\le|h_1|,|h_2|\le H'$, $f_0\sim F$, with $tt'=z_1$ and
$f_0(h_1t'-h_2t)=y_1\ne0$. Write $y_N$ for the largest divisor of $y$ composed
of primes dividing $N$. Then
\begin{equation}\label{sp:coef-norm}
 \sum_{z_1,y_1\ne0}B_{z_1,y_1}^2(y_{1,N})^\theta\le N^{o(1)}FH'^2S^2 .
\end{equation}
At fixed $(z_1,y_1)$ there are at most $\tau(z_1)\tau(|y_1|)=N^{o(1)}$
choices of $(t,t',f_0)$. By Cauchy's inequality the left side is at most
\[
 N^{o(1)}\Bigl[\sum_{f\le2F}f_N^\theta\Bigr]\mathcal Q,\qquad
 \mathcal Q=\sum_{t,t'}\sum_{k_0\ne0}(k_{0,N})^\theta\mathcal N_{t,t'}(k_0)^2,
\]
where $\mathcal N_{t,t'}(k_0)$ counts $(h_1,h_2)$ with $h_1t'-h_2t=k_0$.
Since $f_N^\theta\le\sum_{d\mid f,\ d\mid N^\infty}d^\theta$, the $f$-sum is at
most $2FP_0$, with $P_0=\prod_{p\mid N}(1-p^{\theta-1})^{-1}$.

Write $g_0=(t,t')$, $t=g_0u$, $t'=g_0v$ and $S'=S/g_0$. A fibre of
$(h_1,h_2)\mapsto k_0$ is an arithmetic progression with step $(u,v)$. It
has at most $1+2H'/S'$ points, and at most one if $S'\ge2H'$. Also
\[
 (k_{0,N})^\theta=\sum_{d\mid k_0,\ d\mid N^\infty}\phi_\theta(d),\qquad
 0\le\phi_\theta(d)\le d^\theta,
\]
by telescoping each prime power. Let $A_d(t,t')$ count $(h_1,h_2)$ with
$k_0\ne0$ and $d\mid k_0$. Then
\begin{equation}\label{sp:Q-split}
 \mathcal Q\le\sum_{\substack{d\mid N^\infty\\d\le4H'S}}d^\theta
 \Bigl[\sum_{t,t'}A_d(t,t')+\sum_{t,t':\,S'<2H'}\frac{2H'}{S'}A_d(t,t')\Bigr].
\end{equation}
For the first sum write $(h_1,h_2)=e\,\mathbf w$ with $\mathbf w$ primitive,
$\|\mathbf w\|_\infty=j\le H'$. There are at most $8j$ such $\mathbf w$ and
$e\le H'/j$. Put $d_e=d/(d,e)$. Since $tt'$ is prime to $d$, a contributing
$\mathbf w$ has $(w_1,d_e)=1$ (and $w_1=0$ forces $d_e=1$). Counting $t$ and
then $t'$ in one class modulo $d_e$, and using
$\sum_{e\le E}(d,e)\le\tau(d)E$,
\begin{equation}\label{sp:Ad1}
 \sum_{t,t'}A_d(t,t')\le C\Bigl[\tau(d)\frac{H'^2S^2}d+H'^2S\Bigr].
\end{equation}
For the second sum $(g_0,d)=1$. There are $O(H'S'/d)$ nonzero levels of
$h_1v-h_2u$ divisible by $d$. Each has $O(H'/S')$ points, so
$A_d\le CH'^2/d$. Summing over $g_0$, with
$\sum_{g_0\le2S}(S/g_0+1)^2(g_0/S)\le S(8+\log2S)$,
\begin{equation}\label{sp:Ad2}
 \sum_{t,t':\,S'<2H'}\frac{2H'}{S'}A_d\le C(1+\log2S)\frac{H'^3S}d
 \le C(1+\log2S)\frac{H'^2S^2}d .
\end{equation}
Finally
\[
 \sum_{d\mid N^\infty}\tau(d)d^{\theta-1}=P_0^2,\qquad
 \sum_{\substack{d\mid N^\infty\\d\le4H'S}}d^\theta\le(4H'S)^{1/2}P_2,\qquad
 P_2=\prod_{p\mid N}(1-p^{\theta-1/2})^{-1}.
\]
Since $H'\le S$, $H'^2S(H'S)^{1/2}\le H'^2S^2$. Both $P_0$ and $P_2$ are
$N^{o(1)}$: with the primes of $N$ in increasing order,
$\log P_2\ll\sum_{i\le\omega(N)}p_i^{\theta-1/2}\ll(\log N)^{\theta+1/2}$.
This proves \eqref{sp:coef-norm}.

With $H'=H_*$ we have $H_*/S\ll KSN^\eta/M\le N^{-.0171}$ by
\eqref{sp:l52-scales}. For every real Mellin parameter $t$, fixed
$d,S,F,H_0,\varepsilon$ and colouring, the coefficient sequences
$b^{\pm}_{z_1,m}(t)$ in \eqref{sp:offdiag} satisfy the \emph{uniform}
weighted estimate
\[
 \sum_{z_1,m\ge1}(m_N)^\theta
 |b^{\pm}_{z_1,m}(t)|^2
 \le\sum_{z_1,y_1\ne0}B_{z_1,y_1}^2(y_{1,N})^\theta
 \ll N^{o(1)}F H_*^2S^2.
\]
This follows from \eqref{sp:coef-norm}, since the inner complex
weights have modulus at most one. In particular the implied constant
is independent of $t$ and of all original outer data
$(r,k,n_1,n_2,t_2,b,c,f_0)$. The dependence on these data was
placed in $A_{\mathbf o}(t)$ \emph{before} Cauchy.
Writing $\widetilde b_t^{\pm}(j)$ for the slices of Paper~I,
(2.43), with $n_T=p_0m$ and $a_T=\pm N$, gives
\begin{equation}\label{sp:slices}
 \sup_{t\in\mathbb R}\sum_{j\mid N^\infty}j^\theta
 \|\widetilde b_t^{\pm}(j)\|_2^2
 \ll N^{o(1)}F(H_*S)^2,
 \qquad \sup_{t\in\mathbb R}\|b_t^{\pm}\|_2^2
 \ll N^{o(1)}F(H_*S)^2.
\end{equation}
Since $(p_0,N)=1$, the $N$-part of $n_T=p_0m$ is precisely
$m_N$. This is the norm required by Paper~I, (2.42)--(2.43), for
every $t$, without any Mellin factor in the inner coefficients.

\emph{Application of Paper~I, Theorem~2.8.} In \eqref{sp:offdiag} put
\begin{equation}\label{sp:t28-map}
\begin{gathered}
 r_T=z_1\sim Z_1,\quad s_T=p_0,\quad R_T=Z_1,\quad S_T=p_0/2,\\
 d_T=b,\quad c_T=c,\quad n_T=p_0m,\quad
 Y:=FH_*S,\quad N_T=Cp_0Y\le N^{o(1)}Y,\quad p_0\ll\log N,\quad a_T=\pm N,
\end{gathered}
\end{equation}
with $A=N$. Since $t,t'\sim S$, $z_1=tt'\in(S^2,4S^2]$, which meets two
dyadic windows. Split the $z_1$-sum into $z_1\in(S^2,2S^2]$ and
$z_1\in(2S^2,4S^2]$. Apply the theorem to each, with $Z_1=S^2$ and
$Z_1=2S^2$ respectively. This costs a factor $2$ in the final bound, and
$Z_1\asymp S^2$ in both windows. Then $s_T\in(S_T,2S_T]$ and $S_T\ge3/2$,
since $N$ is even and $p_0\ge3$. The support has $(a_T,r_Ts_T)=1$. The
domain of the theorem is exactly
\[
 (z_1b,p_0c)=1\iff(b,c)=(z_1,c)=(z_1,p_0)=(b,p_0)=1,
\]
and $(z_1,p_0)=1$ holds by roughness. Exactly as in Paper~I,
(3.82)--(3.84), reducing the inverse modulo $c$ gives
\[
 e\Bigl(\frac{a_Tn_T\overline{z_1b}}{p_0c}\Bigr)=e\Bigl(\pm\frac{Nm\overline{z_1b}}c\Bigr),
\]
also when $p_0\mid c$. Split $\rho$ into finitely many smooth pieces on
strict dyadic windows in $(b,c)$. Multiply by fixed smooth cutoffs
equal to $1$ on the supported $n_T,z_1,p_0$. At $n_T\ge1$ use a cutoff
vanishing below $1/2$. All derivatives satisfy Paper~I, (2.39), with
uniform seminorms. Take $\epsilon=\epsilon_0=10^{-12}$. All parameters lie
in $[1,N^{10}]$, so $\mathcal X=2+ACDNRS\le3N^{60}$ in Paper~I, (2.41).
The prefactor there satisfies
$\mathcal X^{\epsilon+18\epsilon_0}=\mathcal X^{19\cdot10^{-12}}\le N^{1.2\cdot10^{-9}}$
for large $N$. Write $\delta_X=1.2\cdot10^{-9}$ for this fixed loss.
This loss and the factors $N^{\eta}=N^{10^{-9}}$ are
far below three power savings obtained below. These are $N^{-.00829}$, $N^{-.1317}$ and $N^{-.2914}$,
one for each term of the bound. It leaves the final
saving $N^{-.001}$ intact. Powers of $p_0$ are $N^{o(1)}$.

Put $B=H$ and $C_{\mathrm{aux}}=HS$. Then $Y\ll Z_1$, since
$Y/Z_1\ll N^\eta H^2S/N\le N^{-.008}$ (see \eqref{sp:reserves}). On a
nonempty $f_0$-block the near-diagonal removal gives $F\ge cHL^{-C_M}/K$,
hence $NY\ge cH^2S^3N^\eta L^{-C_M}\ge C_{\mathrm{aux}}BZ_1$. All exponents of $B,C_{\mathrm{aux}}$
in Paper~I, (2.42), are nonnegative. By (2.42) and \eqref{sp:slices},
\[
 J_*^2\le N^{o(1)}\bigl[N^\theta(C_{\mathrm{aux}}BZ_1+NY)+C_{\mathrm{aux}}BZ_1(C_{\mathrm{aux}}/Z_1)^{2\theta}+B^2Y\bigr]
 Z_1F(H_*S)^2;
\]
the regular spectral product is
$C_{\mathrm{aux}}(C_{\mathrm{aux}}B\sqrt{Z_1})^{2\theta}Z_1^{1-\theta}(BZ_1)^{1-2\theta}=C_{\mathrm{aux}}^{1+2\theta}BZ_1^{2-2\theta}$.
With $B^2<N$, $F\ll H/K$ and \eqref{sp:hstar},
\begin{equation}\label{sp:t28-out}
 \frac{J_*^2}{H^4S^6}\le N^{3\eta+o(1)}[L_A+L_R],\qquad
 L_A=\frac{N^\theta HKS^5}{N^2},\qquad
 L_R=\Bigl(\frac HS\Bigr)^{2\theta}\frac{HKS^5}{N^2}.
\end{equation}
Fix a real $t$. Apply the diagonal bound \eqref{sp:diag}, Paper~I,
Theorem~2.8 with the uniform coefficient norm \eqref{sp:slices}, and
\eqref{sp:t28-out}. The resulting estimate is the same for every real $t$. Thus their supremum $W_4^*$ satisfies
\begin{equation}\label{sp:Ed}
 |\mathcal E_d|\ll N^{\eta+o(1)}\frac H{KS^{3/2}}(W_4^*)^{1/2},\qquad
 \frac{W_4^*}{H^2S^3}\ll
 N^{\eta+o(1)}\frac{H^2S}N+
 N^{3\eta/2+\delta_X+o(1)}(L_A+L_R)^{1/2}.
\end{equation}
All Mellin and colouring costs remain $N^{\eta+o(1)}$ as recorded
in \eqref{sp:cs2}. No outer-dependent coefficient enters the norm.

\subsubsection*{Step 7: exponent reserves}
With $S\ll V_1/D$, $K=U_1D$ and the three scale conditions,
$H^2S\le H^2V_1$, $KS^5\le U_1V_1^5$ and
$H^{2\theta}KS^{5-2\theta}\le H^{2\theta}U_1V_1^{5-2\theta}$. The smooth part
$D$ only helps. Using $U_1\le HN^{-\sigma_{\mathrm p}}$,
$V_1\le N^{v+w+\sigma_{\mathrm p}}/H$ and $U_1V_1\le N^v$,
\begin{equation}\label{sp:reserves}
\begin{aligned}
 1-v-h-\sigma_{\mathrm p}-w&\ge\tfrac{8299}{1000000},\\
 2-\theta-5v+3h-4(\sigma_{\mathrm p}+w)&\ge\tfrac{65873}{500000},\\
 2-(5-2\theta)v+(3-4\theta)h-(4-2\theta)(\sigma_{\mathrm p}+w)&\ge\tfrac{13989529}{48000000},
\end{aligned}
\end{equation}
on $1/3\le h\le23/50$. The first minimum is at $h=.46$, the other two at
$h=1/3$. These are the savings of $H^2S/N$, $L_A$ and $L_R$. Hence
$H^2S/N\le N^{-.00829}$, $L_A\le N^{-.1317}$, $L_R\le N^{-.2914}$, and
\eqref{sp:Ed} gives
\begin{equation}\label{sp:Ed-final}
 |\mathcal E_d|\le C\frac{H^2}KN^{-.001}.
\end{equation}
Summing over the $O(D)$ values of $d$, $|\mathcal E|\le CDH^2N^{-.001}/K$.
Its contribution to $|\mathcal S_0|$ through \eqref{sp:l52-cs1} is
\[
 \bigl[KM\cdot M\cdot DH^2N^{-.001}/K\bigr]^{1/2}\le CND^{1/2}N^{-.0005}\le N^{1-.0004},
\]
since $D\le y=N^{o(1)}$.

\subsubsection*{Step 8: parameter order and conclusion of \eqref{sp:l52-L}}
Fix $A$. Choose $P\ge2A+20$; then $C_\Sigma$; then $C_M\ge\max(P+4C_\Sigma+1,\,P+C_\varphi)$. Here $C_\varphi$ is the absolute exponent with $q/\varphi(q)\le L^{C_\varphi}$. The smooth
weights, $\eta$, $\epsilon$ and $\epsilon_0$ are fixed numbers. Choose the
threshold for $N$ last. By \eqref{sp:nonspectral} and Step~7,
\[
 |\mathcal S_0|\ll NL^{-A-3}+N^{1-\sigma_{\mathrm p}/8}L^{C},
\]
since $(K/H)^{1/2}\le N^{-\sigma_{\mathrm p}/4}$ and $z^{-1/2}$ beats every
power of $L$. Sum over the $O(L)$ shells of $D$ and the $O(1)$ shells of
$S$, and add Step~2. This proves \eqref{sp:l52-L}, uniformly in $f,g,a,b$ and
the scales. Of $f$ and $g$ we used only one-boundedness and the
unary support $(rs,N)=1$.

\subsubsection*{Step 9: the product interval}
Let $I\subset[N/8,2N]$ and $C_{\mathrm{grid}}\ge1$. Divide $(H,2H]$ and
$(M,2M]$ into $O(L^{C_{\mathrm{grid}}})$ consecutive intervals of relative
length at most $L^{-C_{\mathrm{grid}}}$. Call a rectangle $J_1\times J_2$ interior if
$J_1J_2\subset I$, exterior if $J_1J_2\cap I=\emptyset$, and boundary
otherwise. On each interior rectangle apply \eqref{sp:l52-L} to
$a_n\mathbf1_{J_1}(n)$ and $b_m\mathbf1_{J_2}(m)$, with saving
$A+2C_{\mathrm{grid}}$. By Step~0 the first array satisfies (6.34) at the original
parent $H$. For fixed $n$, the $m$ in boundary rectangles lie in $O(1)$
intervals of total length $O(ML^{-C_{\mathrm{grid}}})$. A modulus $q\le4N^v$ has at
most $\tau(q)$ factorizations $q=rs$. Counting $m$ for each $n$,
\[
 \sum_{q\le4N^v}\tau(q)\sum_{n\sim H}
 \Bigl[\frac{ML^{-C_{\mathrm{grid}}}}q+1+\frac{ML^{-C_{\mathrm{grid}}}}{\varphi(q)}\Bigr]
 \ll NL^{C_1-C_{\mathrm{grid}}}+HN^vL ,
\]
with a fixed $C_1$. The last term saves a fixed power, since
$h+v\le.9827+o(1)$. Taking $C_{\mathrm{grid}}=A+C_1$ proves (6.45).
\qed

\subsection{Lemma 6.9}\label{subsec:supp-lem-69}

Recall (6.47):
\[
 H_{\mathrm c}M_{\mathrm c}\asymp N,\qquad
 N^{.50}\ll H_{\mathrm c}\ll N^{.54},\qquad
 Q_1\ll N^{.4545},\qquad Q_2\ll N^{.0682}.
 \tag{6.47}
\]

\begin{mainlemma}{6.9}{Central rectangular distribution}
Let \(x\) be even and \(H_{\mathrm c}M_{\mathrm c}\asymp x\). Let the
bounds in \textup{(6.47)} hold with \(N\) replaced by \(x\), with
fixed comparison constants. Let \(a_n,b_m\) be one-bounded on the complete
dyadic rectangle \(n\sim H_{\mathrm c},m\sim M_{\mathrm c}\). Let the
tested \(a_n\) satisfy \textup{(6.34)} at its original parent, with
one fixed divisor exponent and uniform constants at every requested
precision. Write
\[
 \mathcal D(q,a)=\sum_{n\sim H_{\mathrm c},m\sim M_{\mathrm c}}a_nb_m
 \left(1_{nm\equiv a\ (q)}-\frac{1_{(nm,q)=1}}{\varphi(q)}\right).
\]
Then, for every fixed \(A>0\),
\[
 \sum_{r\sim Q_1}\max_{(b,r)=1}
 \sum_{\substack{s\sim Q_2\\\mu(rs)^2=1,\ (rs,x)=1}}
 \max_{\substack{a\bmod rs,\ (a,rs)=1\\a\equiv b\ (r)}}
       |\mathcal D(rs,a)|\ll_A x(\log x)^{-A}.
 \tag{6.48}
\]
The constant is uniform over these families, rectangles, squarefree
modulus supports and unary interval restrictions of \(a_n\) and \(b_m\).
All errors are measured at the original parents. One outer residue
\(b\) is chosen before \(s\) varies. In the application it is
\(N\bmod r\).
\end{mainlemma}

Write $H=H_{\mathrm c}$, $M=M_{\mathrm c}$, $L=\log x$, $\ell=\log\log x$,
$z=x^{1/\ell^3}$, $y=x^{1/\ell}$. Let $c_0\ge0$ be the fixed divisor
exponent in the hypothesis, so that for every $D>0$, uniformly in intervals
$I\subset(H,2H]$, $k,h\ge1$ and $(u,k)=1$,
\begin{equation}\label{sp:c58-sw}
 |\Delta_a(I;k,u;h)|\le C_DH\tau(h)^{c_0}L^{-D}.
\end{equation}
Since $\Delta_{a\mathbf1_J}(I;k,u;h)=\Delta_a(I\cap J;k,u;h)$, a unary
restriction of $a$ satisfies \eqref{sp:c58-sw} with the same constants. A
unary restriction of $b$ keeps it one-bounded. We may therefore assume that
$a$ and $b$ are the restricted arrays. We keep $(rs,x)=1$ in (6.48)
and throughout the signed residue-difference argument. The
individual maxima in (6.48) are nonnegative, but the signed differences in
Step~2 are not. Enlarging their support does not preserve order. Every selected coefficient below is extended by zero outside
$\mu(rs)^2=1$ and $(rs,x)=1$. On this support $(r,s)=1$ and
$\varphi(rs)=\varphi(r)\varphi(s)$. Throughout,
$rs\le4Q_1Q_2\le Cx^{v}$ with $v=.5227$, and $M\gg x^{.46}$.

\subsubsection*{Step 1: the projected comparator}
For $t=rs$ and reduced $a\equiv b\pmod r$ put
$S_t(a)=\sum_{n,m}a_nb_m\mathbf1_{nm\equiv a\ (t)}$ and
\[
\begin{aligned}
 \Delta(a,b;r,s)&=S_t(a)-\frac{\varphi(r)}{\varphi(t)}
 \sum_{\substack{(nm,t)=1\\nm\equiv b\ (r)}}a_nb_m,\\
 \mathcal F(b;r,s)&=\frac{\varphi(r)}{\varphi(t)}\sum_{(nm,t)=1}a_nb_m
 \Bigl[\mathbf1_{nm\equiv b\ (r)}-\frac1{\varphi(r)}\Bigr].
\end{aligned}
\]
Then exactly
\begin{equation}\label{sp:proj-split}
 \mathcal D(t,a)=\Delta(a,b;r,s)+\mathcal F(b;r,s).
\end{equation}
With $a^{(s)}_n=a_n\mathbf1_{(n,s)=1}$,
\begin{equation}\label{sp:F-form}
 \mathcal F(b;r,s)=\frac1{\varphi(s)}\sum_{\substack{m\sim M\\(m,rs)=1}}b_m
 \Delta_{a^{(s)}}(r,b\overline m).
\end{equation}

\emph{Mean square for $a^{(s)}$.} The identity
$\Delta_{a^{(s)}}(I;k,u;h)=\Delta_a(I;k,u;sh)$ and \eqref{sp:phi-tau} give
\[
 |\Delta_{a^{(s)}}(I;k,u;h)|\le C_DH\tau(s)^{c_0}\tau(h)^{c_0}L^{-D}.
\]
Put
$h_s=a^{(s)}/\tau(s)^{c_0}$, which is one-bounded, and
$e_n=\mathbf1_{H<n\le2H}$. By the parallelogram identity
$\|h_s+e\|^2+\|h_s-e\|^2\ge2\|e\|^2$, there is $\varepsilon_s\in\{\pm1\}$
with
\[
 g_s=h_s+\varepsilon_se,\qquad \|g_s\|_2^2\ge\|e\|_2^2\ge H/2,\qquad
 \|g_s\|_2^2\le4\|e\|_2^2\le8H .
\]
Let $Z=\lceil2H\rceil$ and $c_1=\max(c_0,1)$. By \eqref{sp:c58-sw} for
$h_s$, by \eqref{sp:constant-sw} for $e$, and since
$H\le2^{1/2}\|g_s\|_2Z^{1/2}$ and $\|e\|_2Z^{1/2}\ge H/2$, both $g_s$ and $e$
satisfy (6.35) with $B_{\rm div}=c_1$. The constants are independent of $s$,
the sign and the restriction. Here $\log(2Z)\asymp L$, as
$H\gg x^{.50}$. Lemma~6.5(i) applies to both. Since
$\Delta_{a^{(s)}}=\tau(s)^{c_0}(\Delta_{g_s}-\varepsilon_s\Delta_e)$ and
$|u-w|^2\le2|u|^2+2|w|^2$, for every fixed $S_0>0$,
\begin{equation}\label{sp:B2}
 \sum_{k\le2Q_1}V(k;a^{(s)})\ll_{S_0}\tau(s)^{2c_0}H^2L^{-S_0}.
\end{equation}
The range condition $2Q_1\le Z(\log Z)^{-B_1(S_0)}$ holds for large $x$,
because $Q_1/H\ll x^{.4545-.50}$.

\emph{Bound for $\mathcal F$.} Each reduced class modulo $r$ contains at
most $M/r+1$ values $m\sim M$. The map $m\mapsto b\overline m$ permutes the
reduced classes. Cauchy's inequality in \eqref{sp:F-form} gives, uniformly
in $b$,
\[
 \Bigl|\sum_{m}b_m\Delta_{a^{(s)}}(r,b\overline m)\Bigr|^2
 \le2M\Bigl(\frac Mr+1\Bigr)V(r;a^{(s)}).
\]
Cauchy's inequality over $r\sim Q_1$ and \eqref{sp:B2} then give
\[
 \sum_{r\sim Q_1}\max_b\Bigl|\sum_{m}b_m\Delta_{a^{(s)}}(r,b\overline m)\Bigr|
 \ll\bigl[Q_1M(M/Q_1+1)\tau(s)^{2c_0}H^2L^{-S_0}\bigr]^{1/2}
 \ll\tau(s)^{c_0}HML^{-S_0/2},
\]
since $Q_1\le M$ by a fixed power. Summing over $s$ with
\eqref{sp:phi-tau} and \eqref{sp:first},
\begin{equation}\label{sp:E0}
 \sum_{r\sim Q_1}\max_{(b,r)=1}\sum_{s\sim Q_2}|\mathcal F(b;r,s)|
 \ll xL^{-S_0/2}\sum_{s\le2Q_2}\frac{\tau(s)^{c_0}}{\varphi(s)}
 \ll_A xL^{-A},
\end{equation}
choosing $S_0$ after $A$ (the full order is fixed in Step~8).

\subsubsection*{Step 2: residue differences}
For coprime $r,s$ there are exactly $\varphi(s)$ reduced lifts
$a'\bmod t$ of $b\bmod r$. Every $nm$ coprime to $t$ with
$nm\equiv b\ (r)$ lies in exactly one of them. Hence
\begin{equation}\label{sp:crt-average}
 \Delta(a,b;r,s)=\frac1{\varphi(s)}\sum_{\substack{a'\bmod t,\ (a',t)=1\\a'\equiv b\ (r)}}
 \bigl[S_t(a)-S_t(a')\bigr].
\end{equation}
For each $r$, choose one reduced class $b_r\pmod r$ maximizing the
masked sum of the projected differences, before $s$ varies. For each
$s$ with $\mu(rs)^2=1$ and $(rs,x)=1$, choose a reduced lift $a_{r,s}$
maximizing $|\Delta(a,b_r;r,s)|$. By \eqref{sp:crt-average} and the triangle
inequality, some reduced lift $a'_{r,s}\equiv b_r\pmod r$ satisfies
\[
 |S_{rs}(a_{r,s})-S_{rs}(a'_{r,s})|
 \ge |\Delta(a_{r,s},b_r;r,s)|.
\]
Choose $|\gamma_{r,s}|=1$ to make this difference nonnegative real, and
put $\gamma_{r,s}=0$ outside the masked squarefree support. By
\eqref{sp:proj-split}, \eqref{sp:E0} and \eqref{sp:crt-average}, it suffices
to prove, uniformly for all these residues and phases,
\begin{equation}\label{sp:Z-target}
 \mathcal Z_0^{(x)}=
 \sum_{r\sim Q_1}
 \sum_{\substack{s\sim Q_2\\\mu(rs)^2=1\\(rs,x)=1}}
 \gamma_{r,s}\bigl[S_{rs}(a_{r,s})-S_{rs}(a'_{r,s})\bigr]
 \ll_AxL^{-A}.
\end{equation}
This signed sum keeps the original mask.

\subsubsection*{Step 3: large smooth parts and large divisor functions}
We quote the completed-hyperbola lemma of Paper~I:
\begin{quote}
Let $c\bmod q$ be a unit. For integer
intervals $J_1,J_2$, put

\[
\mathcal D=\#\{(u,v)\in J_1\times J_2:uv\equiv c(q)\}
       -\frac{\#J_1^*\#J_2^*}{\varphi(q)}.
\tag{3.107}
\]

The Weil--Estermann bound \cite{Weil48,Estermann61}

\[
|S(h,k;q)|\le\tau(q)q^{1/2}(h,k,q)^{1/2}.
\tag{3.108}
\]

implies, uniformly in interval lengths and positions,

\[
|\mathcal D|\ll\tau(q)q^{1/2}\log^2(2q)
        +\tau(q)^2\log^2(2q)/\varphi(q).
\tag{3.109}
\]
\end{quote}
For \eqref{sp:P1}, set the lemma's $q=t=rs$,
$J_1=(H,2H]\cap\mathbb Z$, $J_2=(M,2M]\cap\mathbb Z$, and take
$c=a$ or $c=a'$. Both are units modulo $t$ by the reduced-lift condition.
Here $J_i^*=\{u\in J_i:(u,t)=1\}$. The source's $\mathcal D$ is an
unweighted counting discrepancy. It differs from our weighted
$\mathcal D(t,a)$. Define
\[
 C_t(c)=\#\{(n,m)\in J_1\times J_2:nm\equiv c\pmod t\}.
\]
The quoted lemma gives
$C_t(c)\le\#J_1^*\#J_2^*/\varphi(t)+O(E_t)$, where
$E_t=\tau(t)t^{1/2}\log^2(2t)+\tau(t)^2\log^2(2t)/\varphi(t)$.
Since $|a_n|,|b_m|\le1$, one has $|S_t(c)|\le C_t(c)$. This is the
only passage from unweighted counts to the arbitrary complex arrays.
Every pair counted has $(nm,t)=1$ automatically because $c$ is a unit.
No coprimality hypothesis on the arrays is inferred. Consequently,
\begin{equation}\label{sp:P1}
 |S_t(a)-S_t(a')|\le C_t(a)+C_t(a')
 \le\frac{2\#J_1^*\#J_2^*}{\varphi(t)}+O(E_t)
 \le\frac{8Cx}{\varphi(t)}+O(E_t),
\end{equation}
using $HM\le Cx$. The mask $(t,x)=1$ remains on the summation
support but is not a hypothesis of Lemma~3.4. In the smooth-part and
divisor-function eliminations below, we first take absolute values
termwise. Only then are the nonnegative majorants enlarged.
Write $s=k\varrho$ with $P^+(k)\le z<P^-(\varrho)$. If $k\ge y$, then the
$z$-smooth part of $t$ is at least $y$. A given $t$ has at most $\tau(t)$
factorizations $t=rs$. By \eqref{sp:P1}, \eqref{sp:phi-tau},
\eqref{sp:first} and \eqref{sp:rankin}, the sum of the absolute values
of these masked terms in $\mathcal Z_0^{(x)}$ is at most
\[
 \ll x\sum_{h'\le Cx^v}\frac{\tau(h')^2}{h'}
 \sum_{\substack{d'\ge y\\P^+(d')\le z}}\frac{\tau(d')^2}{d'}
 +\sum_{t\le Cx^v}\tau(t)E_t
 \ll_AxL^{-A}+x^{3v/2}L^{K},
\]
and $3v/2=.78405<1$. In the same way the pairs with $\tau(rs)>L^{D_\tau}$
contribute
\[
 \ll xL^{-D_\tau}\sum_{t\le Cx^v}\frac{\tau(t)^3}t+x^{3v/2}L^{K}\ll xL^{K'-D_\tau}+x^{3v/2}L^K .
\]
Choose $D_\tau=D_\tau(A)\ge K'+A$ after the requested precision $A$. From now on $k<y$, and every
remaining modulus and every divisor of it has $\tau\le L^{D_\tau}$.

\subsubsection*{Step 4: the first Cauchy square}
Rename the outer modulus $r$ as $q\sim Q:=Q_1$ and write the inner one as
$s=kr$, with $k\sim K\le y$ smooth and $r\sim R$ squarefree and $z$-rough.
Then $KR\asymp Q_2$. On the support of \eqref{sp:Z-target}, $q,k,r$ are
pairwise coprime and squarefree, $(qkr,x)=1$, and $QKR\ll x^v$.
There are $O(L^2)$ choices of $(K,R)$. Extend the selected coefficients
by zero and put
\[
 \Gamma_{qkr}=\gamma_{q,kr}\,\mathbf1_{\mu(qkr)^2=1}
             \mathbf1_{(qkr,x)=1},\qquad |\Gamma_{qkr}|\le1.
\]
The $\Gamma_{qkr}$ are identical in both choices $a,a'$ of residue.
For one shell the \emph{masked signed sum} is
\[
 \mathcal Z^{(x)}=
 \sum_{\substack{q\sim Q,\ k\sim K\\(qk,x)=1}}\sum_m b_m
 \sum_{\substack{r\sim R,\ n\\(r,x)=1}}\Gamma_{qkr}a_n
 \bigl[\mathbf1_{mn\equiv a_{qkr}\ (qkr)}
       -\mathbf1_{mn\equiv a'_{qkr}\ (qkr)}\bigr],
\]
where $a_{qkr}\equiv a'_{qkr}\equiv b_q\pmod q$ and $b_q$ does not
depend on $k,r$. Put
\[
 \mathcal F^{(x)}_{qk}(m)=
 \sum_{\substack{r\sim R,\ n\\(r,x)=1}}\Gamma_{qkr}a_n
 \bigl[\mathbf1_{mn\equiv a_{qkr}\ (qkr)}
       -\mathbf1_{mn\equiv a'_{qkr}\ (qkr)}\bigr].
\]
With $\psi$ as in Step~3 of the proof of Lemma~6.8,
\begin{equation}\label{sp:c58-cs1}
 |\mathcal Z^{(x)}|^2\ll QKM
 \sum_{\substack{q\sim Q,\ k\sim K\\(qk,x)=1}}
 \sum_m\psi(m/M)|\mathcal F^{(x)}_{qk}(m)|^2.
\end{equation}
For $A_i=A_i(q,k,r_i)$ chosen from $a_{qkr_i},a'_{qkr_i}$, the four
correlations, with signs $+,-,-,+$, are precisely
\[
 \begin{aligned}
 \mathcal C_x(A_1,A_2)={}&
 \sum_{\substack{q\sim Q,\ k\sim K\\(qk,x)=1}}
 \sum_m\psi(m/M)
 \sum_{\substack{r_1,r_2\sim R\\(r_1r_2,x)=1}}
 \Gamma_{qkr_1}\overline{\Gamma_{qkr_2}}\\[-2pt]
 &\qquad\times\sum_{n_1,n_2}a_{n_1}\overline{a_{n_2}}
 \mathbf1_{mn_1\equiv A_1\ (qkr_1)}
 \mathbf1_{mn_2\equiv A_2\ (qkr_2)}.
 \end{aligned}
\]
We prove at arbitrary fixed precision $P>0$ that
\begin{equation}\label{sp:C}
 \mathcal C_x(A_1,A_2)=MT_x+
 O_P\Bigl(\frac{MH^2}{QKL^{P}}\Bigr),
\end{equation}
where $MT_x$ contains the same mask and selected coefficients. It is
independent of the choices of $A_1,A_2$. Since the coefficients and
support are the same in all four correlations, their main terms cancel.
Equation \eqref{sp:c58-cs1} then gives
$|\mathcal Z^{(x)}|\ll xL^{-P/2}$ before the $O(L^2)$ shell sum.

Write $g=(r_1,r_2)$, $r_1=gu$, $r_2=gv$, $(u,v)=1$, with $g\sim G$ dyadic.
Put $d_0=qkg$ and $U=R/G$. Then $q,k,g,u,v$ are pairwise coprime.
The product of the two $\Gamma$-supports imposes
$(qkguv,x)=1$. This modulus-only restriction remains in the
correlation and its zero mode. It may be omitted only from a
termwise absolute-value majorant. It may not be omitted from a signed expression
or from an inner sum before squaring. The two congruences
$mn_1\equiv A_1\ (d_0u)$, $mn_2\equiv A_2\ (d_0v)$ are
soluble exactly when
\begin{equation}\label{sp:comp}
 n_2\equiv A_2\overline{A_1}\,n_1\pmod{d_0},\qquad(n_1,d_0u)=(n_2,d_0v)=1,
\end{equation}
and then $m$ lies in one reduced class modulo $l=d_0uv$. Modulo $q$ both
residues equal $b_q$, so $n_2\equiv n_1\pmod q$.

\subsubsection*{Step 5: large common factor}
Let $G\ge H/(QKL^{C_g})$, with $C_g$ fixed later. The target is
$\mathcal B=MH^2/(QK)\asymp xH/(QK)$. In the large-$G$ estimate
we first take absolute values of each weighted correlation term.
Hence its nonnegative counting majorant may be enlarged to all
$(q,k,g,u,v)$, while $\mathcal C_x$ itself stays masked.
For fixed $(q,k,g,u,v)$ the pairs
$(n_1,n_2)$ in \eqref{sp:comp} lie in a congruence lattice of determinant
$d_0$. Let
$\mathbf w$ be a shortest nonzero vector. If the line
$\mathbb R\mathbf w$ meets the positive rectangle, choose its sign so its
coordinates are positive; otherwise the $\beta=0$ contribution is empty.
Complete it to a reduced basis
$(\mathbf w,\mathbf w')$. The component of $\mathbf w'$ orthogonal to
$\mathbf w$ has length $h_\perp=d_0/|\mathbf w|\ge\sqrt3|\mathbf w|/2$. In a
disc of radius $CH$ the points $\alpha\mathbf w+\beta\mathbf w'$ with
$\beta\ne0$ have $|\beta|\ll H/h_\perp$ and, for fixed $\beta$,
$O(H/|\mathbf w|+1)$ values of $\alpha$. Such points exist only if
$h_\perp\ll H$. Then $|\mathbf w|\ll h_\perp\ll H$ and
$(H/h_\perp)^2\ll H^2/d_0$, so they number $O(H^2/d_0)$. Each such pair has $O(M/l+1)$ values of
$m$. Hence these points contribute
\[
 \ll QKGU^2\cdot\frac{H^2}{QKG}\Bigl(1+\frac M{QKGU^2}\Bigr)=H^2U^2+\frac{MH^2}{QKG}.
\]
For $\beta=0$, $(n_1,n_2)=\alpha\mathbf w$ with $\alpha\ge1$, and the two
congruences fix $\alpha m\le Cx$ in one unit class $c\pmod l$.
Here $l=d_0uv=qkguv\ge1$, and $c$ is a unit. Each original
congruence has a unit right-hand side $A_i$. Hence $m$ and $\alpha$
are units modulo every factor of $l$ on this line.
For $X\le Cx$, cover the positive integer hyperbola $\alpha m\le X$
by the $O(L)$ integer rectangles $J_{1,0}=\{1\}$,
$J_{2,0}=[1,\lfloor X\rfloor]$, and, for $j\ge1$ with $2^{j-1}<X$,
\[
 J_{1,j}=(2^{j-1},2^j]\cap\mathbb Z,\qquad
 J_{2,j}=[1,\lfloor X/2^{j-1}\rfloor]\cap\mathbb Z.
\]
These rectangles may overcount. This is harmless for the positive
quantity $\sum_{e\le X,e\equiv c\ (l)}\tau(e)$. Each has
$\#J_{1,j}\#J_{2,j}\ll X$, and their number is $O(L)$.
Apply the quoted Paper~I, Lemma~3.4 with its modulus
$q=l$, unit $c$, and these two integer intervals. Empty intervals
contribute zero. The positive counts give, uniformly in $c$ and $X$,
\begin{equation}\label{sp:TP}
 \sum_{\substack{e\le X\\e\equiv c\ (l)}}\tau(e)
 \ll\frac{XL}{\varphi(l)}+\Bigl[\tau(l)l^{1/2}+\frac{\tau(l)^2}{\varphi(l)}\Bigr]L^3
\end{equation}
uniformly in the unit $c$ and $X\le Cx$. A modulus $l$ has at most
$\tau_5(l)\le\tau(l)^4$ representations $l=qkguv$. Summing \eqref{sp:TP}
with \eqref{sp:first},
\[
 \ll xL^{C}+(QKGU^2)^{3/2}L^{C}.
\]
With $W=QK\le x^{.4545+o(1)}$, $R\ll x^{.0682}$ and $WR\ll x^{.5227}$, on
this range of $G$,
\[
 \frac{MH^2/(QKG)}{\mathcal B}=\frac1G\le x^{-.0455+o(1)},\quad
 \frac{H^2U^2}{\mathcal B}\ll\frac{W^3R^2L^{2C_g}}{xH}\le x^{-.0001+o(1)},\quad
 \frac x{\mathcal B}\le x^{-.0455+o(1)},
\]
and $(QKGU^2)^{3/2}\le(W^2R^2L^{C_g}/H)^{3/2}\le x^{.8181+o(1)}$, while
$\mathcal B\ge x^{1.0455-o(1)}$. So the large-$G$ part of
$\mathcal C_x$ is $O(\mathcal BL^{-P})$ at any requested fixed
precision $P$. The mask was removed only after taking absolute values.

\subsubsection*{Step 6: small common factor, Poisson and the zero mode}
Let $G<H/(QKL^{C_g})$. For one class $c\bmod l$,
\begin{equation}\label{sp:c58-poisson}
 \sum_{m\equiv c\ (l)}\psi\Bigl(\frac mM\Bigr)=\frac Ml\sum_{h\in\mathbb Z}
 \hat\psi\Bigl(\frac{hM}l\Bigr)e\Bigl(\frac{hc}l\Bigr),
\end{equation}
and the terms $|h|>x^{\eta}l/M$ contribute $O(x^{-100})$ per tuple. We take
$\eta=10^{-8}$.

The zero mode is $M\hat\psi(0)/l$ times the weighted number of compatible
$(n_1,n_2)$. It still carries the product of the two
$\Gamma$-supports, and thus $(qkguv,x)=1$. After taking absolute
values of the \emph{error}, its positive majorant can be enlarged.
This is the only use of the unmasked range in \eqref{sp:ZM}.
Put $f_v=a\mathbf1_{(n,v)=1}$. For fixed $n_1$, the
$n_2$-sum in \eqref{sp:comp} is a sum of $\overline{f_v}$ over one reduced
class $c(n_1)$ modulo $d_0$. Replace it by
$\varphi(d_0)^{-1}\sum_{(n_2,d_0v)=1}\overline{a_{n_2}}$. The error is
$\overline{\Delta_{f_v}(d_0,c(n_1))}$. The map $n_1\mapsto c(n_1)$ hits each
reduced class $O(H/d_0)$ times, so Cauchy's inequality gives
\[
 \sum_{n_1}|\Delta_{f_v}(d_0,c(n_1))|\ll Hd_0^{-1/2}V(d_0;f_v)^{1/2}.
\]
A given $d_0$ has at most $\tau_3(d_0)\le L^{2D_\tau}$ representations, and
$d_0\asymp D_0=QKG$. Multiplying by $M/(d_0uv)$ and summing over $u,v$,
the zero-mode error is
\begin{equation}\label{sp:ZM}
 \ll\frac{MH}{D_0}L^{2D_\tau}\sup_v\Bigl[\sum_{d_0\le C'D_0}V(d_0;f_v)\Bigr]^{1/2}.
\end{equation}
Since $v$ divides a remaining modulus, $\tau(v)\le L^{D_\tau}$. The identity
$\Delta_{f_v}(I;k,u;h)=\Delta_a(I;k,u;vh)$ and \eqref{sp:c58-sw} show that
$f_v$ satisfies \eqref{sp:c58-sw} at precision $D-c_0D_\tau$. As in Step~1, add $\varepsilon_ve$ to make the norm
at least $(H/2)^{1/2}$. Both sequences satisfy (6.35) with $B_{\rm div}=c_1$
and uniform constants. Lemma~6.5(i) and linearity give, uniformly in $v$,
\[
 \sum_{d_0\le C'D_0}V(d_0;f_v)\ll_{S_0}H^2L^{-S_0},
\]
because $C'D_0\le C'H/L^{C_g}\le Z(\log Z)^{-B_1(S_0)}$ once
$C_g>B_1(S_0)+1$. With $D_\tau=D_\tau(A)$ fixed in Step~3, choose first
$S_0>4D_\tau+2A$, then $C_g>B_1(S_0)+1$, then the threshold for $x$.
Since $D_0\ge cQK$, \eqref{sp:ZM} is
$O(\mathcal BL^{2D_\tau-S_0/2})=O(\mathcal BL^{-A})$.

What remains of the zero mode is the \emph{masked} main term
\begin{equation}\label{sp:c58-MT}
 MT_{x,G}=M\hat\psi(0)
 \sum_{\substack{q,k,g,u,v\\(qkguv,x)=1}}
 \Gamma_{qkgu}\overline{\Gamma_{qkgv}}
 \frac{\bigl[\sum_{(n_1,d_0u)=1}a_{n_1}\bigr]
       \bigl[\sum_{(n_2,d_0v)=1}\overline{a_{n_2}}\bigr]}
 {d_0uv\varphi(d_0)} ,
\end{equation}
with the previously imposed coprime/squarefree and dyadic conditions
understood. This expression depends on the modulus support and
$\Gamma$, but not on $A_1,A_2$. Put
$MT_x=\sum_{G\ \rm small}MT_{x,G}$.

\subsubsection*{Step 7: the second square}
It remains to bound the frequencies $h\ne0$ in the masked
correlation. Write
\[
 \frac Ml=\frac M{QKGU^2}\cdot\frac{QKGU^2}{qkguv};
\]
the second factor is a product of bounded weights in $(q,k,g)$, in $u$ and
in $v$. The mask factors exactly as
\[
 \mathbf1_{(qkguv,x)=1}=
 \mathbf1_{(qkg,x)=1}\mathbf1_{(u,x)=1}\mathbf1_{(v,x)=1}.
\]
Keep the first factor with the outer $(q,k,g,n_1,n_2)$ data and the
last two with the $(u,v,h)$-sum, together with
$\Gamma_{qkgu}\overline{\Gamma_{qkgv}}$. No masked term is reintroduced
into a signed sum.

Every $z$-rough $u,v\le x^2$ has $O(\ell^3)$ prime factors. Colour each prime
below $x^2$ red or blue independently. A coprime pair $(u,v)$ is covered if
all primes of $u$ are red and all of $v$ blue. This has probability at least
$2^{-O(\ell^3)}$. With $J_0=\lceil2^{O(\ell^3)}(10\log x+1)\rceil$ independent
colourings, the probability that one of the at most $x^4$ pairs is never
covered is $<1$. Fix such a list and assign each pair to its first
covering colouring. This partitions the coprime pairs into
$J\le J_0=x^{o(1)}$ classes, and $(uu',vv')=1$ for any two pairs in one
class.

On one class let $\mathcal Y$ be the sum over compatible $(n_1,n_2)$, $(u,v)$
and $0<|h|\le\mathsf F$, after removing $M/(QKGU^2)$, where
\begin{equation}\label{sp:Fcut}
 \mathsf F=Cx^{\eta}QKGU^2/M .
\end{equation}
If $\mathsf F<1$ there is nothing to bound. By CRT the class of $m$ in the
phase is
\[
 \frac c{l}\equiv\frac{A_1\overline{n_1}\,\overline v}{d_0u}
 +\frac{A_2\overline{n_2}\,\overline{d_0u}}v\pmod1 .
\]
Cauchy's inequality over the \emph{masked} outer data
$(q,k,g,n_1,n_2)$, $(qkg,x)=1$, with $n_2\equiv n_1\ (q)$ removes
$a$ and gives
\begin{equation}\label{sp:c58-cs2}
 |\mathcal Y|^2\ll H^2KG\,\mathcal D_3,
\end{equation}
since there are at most
$QKG\cdot O(H(H/Q+1))\ll H^2KG$ outer tuples ($Q<H$).
Here $\mathcal Y$ and $\mathcal D_3$ keep the $x$-mask.
The sum $\mathcal D_3$ runs over $(qkg,x)=1$ and $n_2\equiv n_1\ (q)$. Its terms are
a smooth majorant times the square of the $(u,v,h)$-sum. This inner sum
contains $\mathbf1_{(u,x)=1}\mathbf1_{(v,x)=1}$ and the selected
$\Gamma$ coefficients. The congruence modulo $kg$ stays inside the
square. In the expansion of this square each pair of copies keeps its
zero--one masks. The later diagonal, Weil and Ramanujan bounds are
termwise bounds after taking absolute values of the expanded
contributions. Hence they may discard zero--one factors.
We do not use the false inequality
$|\sum_{(uv,x)=1}z_{uv}|^2\le|\sum_{u,v}z_{uv}|^2$.

Expanding gives $(u,v,h)$ and $(u',v',h')$ with $0<|h|,|h'|\le\mathsf F$. Put
$A=A_1(u)$, $A'=A_1(u')$, $B=A_2(v)$, $B'=A_2(v')$ and
\[
 \Delta=hu'v'-h'uv .
\]
The two congruences modulo $kg$ must agree, or the term vanishes. Let
$t\equiv B\overline A\pmod{d_0}$. The $(n_1,n_2)$-sum is
\[
 \mathcal K=\sum_{\substack{n_2\equiv tn_1\ (d_0)\\(n_1,d_0uu')=1,\ (n_2,vv')=1}}
 \psi\Bigl(\frac{n_1}H\Bigr)\psi\Bigl(\frac{n_2}H\Bigr)
 e\Bigl(\frac{c_1\overline{n_1}}{d_0uu'}\Bigr)e\Bigl(\frac{c_2\overline{n_2}}{vv'}\Bigr),
\]
\[
 c_1\equiv(Ahu'v'-A'h'uv)\overline{vv'}\pmod{d_0uu'},\qquad
 c_2\equiv(Bhu'v'-B'h'uv)\overline{d_0uu'}\pmod{vv'}.
\]
All inverses exist, since $(uu',vv')=1$ and both are prime to $d_0$.
Completing first in $n_2$ and then in $n_1$,
\begin{equation}\label{sp:P2}
 \mathcal K=\frac{H^2}{d_0^2uu'vv'}\sum_{\ell_1,\ell_2\in\mathbb Z}
 \hat\psi\Bigl(\frac{\ell_2H}{d_0vv'}\Bigr)\hat\psi\Bigl(\frac{\ell_1H}{d_0uu'}\Bigr)
 S(c_2,\ell_2\overline{d_0};vv')\,S(c_1,\ell_1+\ell_2c_3;d_0uu'),
\end{equation}
with $c_3\equiv uu't\,\overline{vv'}\pmod{d_0uu'}$ and the Kloosterman sum
$S(a,b;m)$ of Paper~I, (2.35). The first completion produces the phase
$e(\ell_2tn_1\overline{vv'}/d_0)$, which combines with the $n_1$-phase to
give $c_3$. Let
\begin{equation}\label{sp:dual}
 \Lambda=C_1x^{\eta}QKGU^2/H .
\end{equation}
Frequencies with $|\ell_i|>\Lambda$ contribute $O(x^{-100})$ per tuple. If
$\Lambda<1$, only $(\ell_1,\ell_2)=(0,0)$ remains.

The target for $\mathcal D_3$ is $T_3=H^2U^4G/K$. By \eqref{sp:c58-cs2}, a
bound $\mathcal D_3\ll T_3x^{-\rho}$ gives $|\mathcal Y|\ll H^2U^2Gx^{-\rho/2}$.

\emph{Diagonal $\Delta=0$.} Then $uv\mid hu'v'$, so for fixed $(h,u',v')$
there are $x^{o(1)}$ choices of $(h',u,v)$. Bounding $\mathcal K$ trivially
by $O(H^2/d_0)$, as $d_0<H$,
\[
 \mathcal D_3^{\rm diag}\ll x^{o(1)}H^2\mathsf FU^2,\qquad
 \frac{\mathcal D_3^{\rm diag}}{T_3}\ll x^{\eta+o(1)}\frac{QK^2H}x\le x^{-.0055+\eta+o(1)}.
\]
The divisor-counting bound $x^{o(1)}$ for the tuple count immediately above is
loose. Only this $x^{o(1)}$ factor is used in the estimate for
$\mathcal D_3^{\rm diag}$.

\emph{Gcd bound.} Put $e=(u,u')$, $f=(v,v')$, $u_0=u/e$, $u_1=u'/e$,
$v_0=v/f$, $v_1=v'/f$. Then
\begin{equation}\label{sp:gcd-local}
 (c_1,uu')(c_2,vv')\le(h,u_0v_0)(h',u_1v_1)(c_1,e^2)(c_2,f^2).
\end{equation}
The four factors are supported on disjoint sets of primes. If $p\mid u$,
$p\nmid u'$, then $c_1\equiv Ahu'v'\overline{vv'}\pmod p$ with a unit
multiplier, so $(c_1,p)=(h,p)$. The cases $p\mid u'$ only, $p\mid v$ only and
$p\mid v'$ only are identical, with $(h',p)$, $(h,p)$ and $(h',p)$. If
$p\mid e$, the numerator of $c_1$ is $p(Aha'v'-A'h'a''v)$ with $u=pa''$,
$u'=pa'$. Thus the $p$-part of $(c_1,e^2)$ is at most $p^2$. The case
$p\mid f$ is similar. Modulo $q$ all four residues equal $b_q$, so
$c_1\equiv b_q\overline{vv'}\Delta\pmod q$ and $(c_1,q)=(\Delta,q)$. For
$p\mid kg$ we use only $(c_1,kg)\le kg$. Hence
\begin{equation}\label{sp:gcd-W}
 (c_1,d_0uu')(c_2,vv')\le kg\,(\Delta,q)\,(h,uv)(h',u'v')(u,u')^2(v,v')^2 .
\end{equation}

\emph{Long dual range, $\Lambda\ge1$, $\Delta\ne0$.} There are $O(\Lambda^2)$
pairs $(\ell_1,\ell_2)$. Paper~I, (2.35), gives
\[
 |S(a,b;m)|\le\tau(m)m^{1/2}(a,m)^{1/2}.
\]
Use the square root of \eqref{sp:gcd-W} and
$(kg)^{1/2}(QKG)^{1/2}\ll Q^{1/2}KG$. Then
\[
\begin{multlined}
 |\mathcal K|\ll x^{2\eta+o(1)}Q^{1/2}KGU^2\\
 \times(h,uv)^{1/2}(h',u'v')^{1/2}(u,u')(v,v')(\Delta,q)^{1/2}.
\end{multlined}
\]
The averages
\[
 \sum_{1\le|h|\le\mathsf F}(h,a)\le2\mathsf F\tau(a),\qquad
 \sum_{q\sim Q}(\Delta,q)\le CQ\tau(|\Delta|),\qquad
 \sum_{u,u'\sim U}(u,u')\ll U^2\log(2U)
\]
follow from $(h,a)=\sum_{d\mid(h,a)}\varphi(d)$ and counting multiples.
Summing over the $O(KG)$ choices of $(k,g)$,
\[
 \mathcal D_3^{\rm long}\ll x^{2\eta+o(1)}Q^{3/2}K^2G^2U^6\mathsf F^2,\qquad
 \frac{\mathcal D_3^{\rm long}}{T_3}\ll x^{4\eta+o(1)}\frac{Q^{7/2}K^5R^6}{G^3x^2}
 \le x^{-.00005+4\eta+o(1)},
\]
using $U=R/G$, $G\ge\tfrac12$, $K\le y$; the possible absolute factor
$G^{-3}\le8$ is absorbed in the implied constant and in $x^{o(1)}$, and
then
\[
 2-\tfrac72(.4545)-6(.0682)=.00005 .
\]

\emph{Short dual range, $\Lambda<1$, $\Delta\ne0$.} Now
\[
 \mathcal K=\frac{H^2\hat\psi(0)^2}{d_0^2uu'vv'}\,c_{vv'}(c_2)\,c_{d_0uu'}(c_1)+O(x^{-100}),
\]
with Ramanujan sums $c_m(a)=S(a,0;m)$, and $|c_m(a)|\le(a,m)$. By
\eqref{sp:gcd-W} each term is
\[
 \ll\frac{H^2}{Q^2KGU^4}(\Delta,q)(h,uv)(h',u'v')(u,u')^2(v,v')^2 .
\]
With $\sum_{u,u'\sim U}(u,u')^2\ll U^3$ and the averages above,
\[
 \mathcal D_3^{\rm short}\ll x^{o(1)}\frac{H^2\mathsf F^2U^2}Q,\qquad
 \frac{\mathcal D_3^{\rm short}}{T_3}\ll x^{2\eta+o(1)}\frac{H^2QK^3R^2}{Gx^2}
 \le x^{-.3291+2\eta+o(1)},
\]
since $2(.54)+.4545+2(.0682)-2=-.3291$.

\subsubsection*{Step 8: conclusion}
After the $x^{o(1)}$ factors the three ratios are at most
$x^{-.00005+5\eta}$. These factors include $J$, the divisor bounds,
$K\le y$ and fixed powers of $L$. Take $\rho=10^{-6}$. Then
$|\mathcal Y|\ll H^2U^2Gx^{-\rho/2}$ on each masked colouring class.
Restoring $M/(QKGU^2)$ and the $J=x^{o(1)}$ classes, the nonzero
frequencies contribute $\ll\mathcal Bx^{-\rho/3}$ to each small-$G$
block of $\mathcal C_x$. The modulus mask was preserved through both
Cauchy squares. Only termwise absolute-value majorants were enlarged.
There are $O(L)$ common-factor blocks. Together with Steps~5 and~6
this proves \eqref{sp:C} with the \emph{same} masked $MT_x$ in each
of its four choices of residues. Their signed combination is
$MT_x-MT_x-MT_x+MT_x=0$.
Fix a requested final exponent $A$. First choose the correlation
precision $P>2A+10$. Then choose
$D_\tau=D_\tau(P)\ge K'+P$ (Step~3). Then choose
$S_0>4D_\tau+2P$, also large enough to make \eqref{sp:E0} smaller
than $xL^{-A-10}$. Next take $C_g>B_1(S_0)+1$, and finally the
threshold for $x$. Equation \eqref{sp:c58-cs1} gives
$\mathcal Z^{(x)}\ll xL^{-P/2}$ on each shell. Sum over its
$O(L^2)$ choices of $(K,R)$ and add the masked Step~1 and Step~3
errors. This proves exactly (6.48), with $(rs,x)=1$ as stated.
Every constant depends only on $A$, the comparison constants in
(6.47) and the constants in \eqref{sp:c58-sw}.
\qed

\section{Finite calculations}\label{app:finite-calculations}

This appendix gives the finite quantities used above.
Each terminating decimal denotes an exact rational number.
All decimal interval endpoints are rounded outwards; decimal approximations
in explanatory prose are not additional certified bounds.
The comparisons with prime sums occur in the main text.

\subsection{Series bounds}

For \(|z|<1\),
\begin{equation}\label{eq:app-log-tail}
 \left|\operatorname{arctanh}z-
 \sum_{k=0}^{m-1}\frac{z^{2k+1}}{2k+1}\right|
 \le\frac{|z|^{2m+1}}{(2m+1)(1-z^2)}.
\end{equation}
Write any positive rational as \(t=2^a v\),
where \(2/3\le v<4/3\).
Then use
\begin{equation}\label{eq:app-log-reduction}
 \log t=a\log2+2\operatorname{arctanh}\frac{v-1}{v+1},
 \qquad \log2=2\operatorname{arctanh}(1/3).
\end{equation}
Also, for \(0\le z<1\),
\begin{equation}\label{eq:app-dilog-tail}
 0\le\operatorname{Li}_2(z)-\sum_{n=1}^m\frac{z^n}{n^2}
 \le\frac{z^{m+1}}{(m+1)^2(1-z)}.
\end{equation}
These inequalities enclose the elementary functions below.
Integer comparisons then give the printed rational bounds.

\subsection{Four profile differences}\label{subsec:app-h6-four-profiles}

Put
\begin{gather}
 \alpha=\frac{100}{1149},\quad \beta=\frac{50}{309},\quad
 \mu_* =\frac{25}{128},\quad \lambda=\frac{93}{400},\quad
 h_G=\frac9{5000},\quad K=\frac{57}{64}-\frac1{500},\notag\\
 \delta_h=\frac {h_G}\alpha,\qquad C=K-\frac{11}{5}\alpha,\qquad
 B(x)=\min\{\beta,\max\{\alpha,C-5x/2\}\}.
 \label{eq:app-h6-profile-parameters}
\end{gather}
Define the elementary functions
\begin{align}
 \widehat f(s)&=\frac{\log(s-1)}s &&(2\le s\le4),\notag\\
 \widehat F(s)&=\frac1s &&(1<s\le3),\label{eq:app-h6-profile-functions}\\
 s\widehat F(s)&=1+\log(s-2)\log(s-1)
 -\operatorname{Li}_2\!\left(\frac{s-2}{s-1}\right)\\
 &\quad-\frac{\log^2(s-1)}2+\operatorname{Li}_2(1/2)
 +\frac{\log^22}{2} &&(3<s<5),\notag
\end{align}
The two sieve normalizations satisfy
\[
 \widehat F=F/(2e^\gamma),\qquad \widehat f=f/(2e^\gamma).
\]
The raw multiplier $2/(e^\gamma\alpha)$, a signed row coefficient
$\pm1/4$, and $F=2e^\gamma\widehat F$ or
$f=2e^\gamma\widehat f$ combine to the displayed $1/\alpha$ factor.
The directed integrals below use these hatted functions. No further $1/4$ or
$2e^\gamma$ factor is applied.
Also $\operatorname{Li}_2(t)=\sum_{n\ge1}t^n/n^2$ for $0\le t<1$.
Write
\[
 D_F(s)=\widehat F(s)-\widehat F(s+\delta_h),\qquad
 D_f(s)=\widehat f(s+\delta_h)-\widehat f(s).
\]
For the three intervals in Table~\ref{tab:app-h6-profile-bounds}, put
\begin{align}
 \Delta_j&=\frac1\alpha\int_{I_j}
 D_F\!\left(\frac{K-5x/2}{\alpha}\right)\frac{dx}{x}
 \quad(j=4,5,7),\label{eq:app-h6-profile-negative}\\
 \Delta_{10}&=\frac1\alpha\int_{\mu_*}^\lambda\int_\alpha^{B(x)}
 D_f\!\left(\frac{K-5x/2-y}{\alpha}\right)
 \frac{dy\,dx}{xy}.
 \label{eq:app-h6-profile-positive}
\end{align}
These $\Delta_j$ are those of
equation~(5.22).
Closed or half-open interval endpoints have measure zero.
The curved domain is that of \cite[(2.38)--(2.40)]{PaperII}.

\begin{proposition}\label{prop:app-h6-four-profiles}
A directed interval computation certifies the following outward enclosure of the
four elementary integrals:
\begin{equation}\label{eq:app-h6-profile-total}
 \frac{572004}{10^8}<\Delta_4+\Delta_5+\Delta_7+\Delta_{10}
 <\frac{606486}{10^8}.
\end{equation}
In particular, their sum minus $10^{-20}$ exceeds $5719/10^6$.
Also,
\begin{equation}\label{eq:app-h6-profile-strip}
 \frac{\log3}{4\alpha^2}
 \left(\frac52\lambda+\beta\right)\frac{\lambda-\mu_*}{\mu_*}
 <\frac{21137254516197687}{4120000000000000}
 <\frac{5137}{1000}<\frac{21}{4}<\frac{179}{25}.
\end{equation}
\end{proposition}

\begin{table}[ht]
\centering
\caption{Outward decimal bounds from the directed computation. Every
endpoint is an exact terminating rational.}
\label{tab:app-h6-profile-bounds}
\begin{tabular}{cccc}
\hline
$j$ & $I_j$ & lower bound & upper bound\\
\hline
$4$ & $[25/128,1/4]$ & $0.00162068$ & $0.00169539$\\
$5$ & $[1/4,7109/28000]$ & $0.00040327$ & $0.00041325$\\
$7$ & $[25/128,93/400]$ & $0.00056463$ & $0.00059051$\\
$10$ & curved domain & $0.00313145$ & $0.00336572$\\
\hline
sum & & $0.00572003$ & $0.00606487$\\
\hline
\end{tabular}
\end{table}

\begin{proof}
The profiles are evaluated with directed interval arithmetic. An independent
verifier rechecks the directed output. Put
$h_G=9/5000$, $p(s)=-\widehat F'(s)$, and let $d$ be the right-branch
derivative profile for $\widehat f$, so $d(s)=\widehat f'(s)$ for $2<s<4$
and $d(2)=1/2$ by continuous extension.
The three one-dimensional rows each use 2048 equal cells:
\[
 x_k=a_j+\frac{k(b_j-a_j)}{2048}\qquad(0\le k\le2048).
\]
For a cell $[x_0,x_1]=[x_k,x_{k+1}]$, define
\[
 s_-=\frac{K-5x_1/2}{\alpha},\qquad
 s_+=\frac{K-5x_0/2+h_G}{\alpha}.
\]
Since $D_F(s)=\int_0^{h_G/\alpha}p(s+t)\,dt$, the code's lower and upper
cell contributions are exactly
\begin{equation}\label{eq:app-h6-profile-negative-cells}
 L=\frac{h_G}{\alpha^2}\frac{x_1-x_0}{x_1}p(s_+),\qquad
 U=\frac{h_G}{\alpha^2}\frac{x_1-x_0}{x_0}p(s_-).
\end{equation}
These hold because $s_-\le (K-5x/2)/\alpha+t\le s_+$ for every
$x\in[x_0,x_1]$ and $0\le t\le h_G/\alpha$; also $x_0\le x\le x_1$.

For the fourth integral, split at
\[
 x_B=\frac{506848823}{2366940000}.
\]
On each piece $[a,b]$ of $[\mu_*,x_B]\cup[x_B,\lambda]$, the S10 routine uses
$x_i=a+i(b-a)/96$ and $z_j=j/96$ for $0\le i,j\le96$.
For a cell $[x_i,x_{i+1}]\times[z_j,z_{j+1}]$, with $0\le i,j<96$, put
\[
 b_0=B(x_i),\quad b_1=B(x_{i+1}),\quad
 y_-=\alpha+(b_1-\alpha)z_j,\quad
 y_+=\alpha+(b_0-\alpha)z_{j+1},
\]
\[
 q_-=\frac{K-5x_{i+1}/2-y_+}{\alpha},\qquad
 q_+=\frac{K-5x_i/2-y_-+h_G}{\alpha}.
\]
The substitution is $y=\alpha+(B(x)-\alpha)z$. Its executed lower and
upper contributions are
\begin{equation}\label{eq:app-h6-profile-positive-cells}
 L=\frac{h_G}{\alpha^2}(x_{i+1}-x_i)(z_{j+1}-z_j)
       \frac{b_1-\alpha}{x_{i+1}y_+}d(q_+),\qquad
 U=\frac{h_G}{\alpha^2}(x_{i+1}-x_i)(z_{j+1}-z_j)
       \frac{b_0-\alpha}{x_iy_-}d(q_-).
\end{equation}
For every point in this cell and $0\le t\le h_G/\alpha$,
$q_-\le (K-5x/2-y)/\alpha+t\le q_+$, and
\[
 \frac{b_1-\alpha}{x_{i+1}y_+}
 \le\frac{B(x)-\alpha}{xy}
 \le\frac{b_0-\alpha}{x_iy_-}.
\]
Hence the two positive factors are bounded separately, as in the code.
The two rectangular grids contain $2\cdot96^2$ cells. The total
executed cell count is
\[
 3\cdot2048+2\cdot96\cdot96=24576.
\]

The endpoint bounds use the following branch formulas. For $2<s<4$,
\[
 d(s)=\frac{s/(s-1)-\log(s-1)}{s^2}>0,\qquad
 d'(s)=-\frac1{s(s-1)^2}
       -\frac{2\{s/(s-1)-\log(s-1)\}}{s^3}<0.
\]
They extend continuously to $s=2$ with $d(2)=1/2$ and right derivative
$d'_+(2)=-1$. Thus the branch profile is positive and decreasing on $[2,4]$.
For positivity: the numerator decreases, and at $s=4$ it is
$4/3-\log3>0$, since
$e^{4/3}>1+4/3+(4/3)^2/2+(4/3)^3/6>3$. For $s>3$, put
$A(s)=\widehat F(s)-\widehat f(s-1)$. Then
\[
 (sA(s))'=-s\,\widehat f'(s-1)<0,\qquad
 p(s)=A(s)/s,\qquad
 p'(s)=-2A(s)/s^2-\widehat f'(s-1)/s.
\]
At $s=3$ the two pieces give $p(3)=1/9$, so $p$ is continuous
there; the displayed differential identities are used on $s>3$.
The exact directed evaluation at
$s_{\max}=73930107/16000000$ gives
$A(s_{\max})>1597/100000>0$. Hence $A(s)>0$ throughout the required
range $3<s\le s_{\max}$, and there $p>0$ and $p'<0$. For $s\le3$,
$p(s)=1/s^2$ is also positive and decreasing. The S10 mixed-corner range is
\[
 \frac{17359415273}{7910400000}\le q_-\le q_+
 \le\frac{57930107}{16000000}<4,
 \qquad \frac{17359415273}{7910400000}>\frac{219}{100}.
\]
These inequalities justify using the decreasing positive $d$ in
\eqref{eq:app-h6-profile-positive-cells}. The lower bound refers only to the
executed mixed-corner interval.

The code evaluates the formulas by rational geometry and 256-bit Arb
ball arithmetic. In \texttt{profile\_direct.py}, logarithms use Arb's
\texttt{log()} operation. Put $w=s-2$. The dilogarithm uses Arb's
\texttt{polylog(2)} operation at the negative argument $-w$, together with
$\pi^2/12$. These are the code's evaluations of
\eqref{eq:app-h6-profile-functions}. By the Landen identity
\[
 \operatorname{Li}_2(-w)
 =-\operatorname{Li}_2\!\left(\frac{w}{1+w}\right)
  -\frac12\log^2(1+w),\qquad
 \operatorname{Li}_2(1/2)+\frac12\log^22=\frac{\pi^2}{12}
\]
the Arb expression equals the displayed formula. Arb supplies
outward ball enclosures. The code does not use the atanh and
dilogarithm tail inequalities above. The exact directed rational
endpoints are written to the result record. The terminating decimals in the
table are outward projections of each row's directed endpoint. Their displayed
column sums are $0.00572003$ and $0.00606487$. The sharper bound
\eqref{eq:app-h6-profile-total} adds the unrounded rational endpoints
in RESULT.json before decimal projection. Summing before rounding proves the stated
integer comparisons. Finally, 64 terms of
$\log3=2\sum_{n\ge0}2^{-(2n+1)}/(2n+1)$ and its geometric remainder give
$\log3<10987/10000$. This proves \eqref{eq:app-h6-profile-strip}.
\end{proof}

Proposition~\ref{prop:app-h6-four-profiles} concerns elementary integrals and
one strip coefficient. Its application to the prime sums requires the
source comparisons, uniform errors and a single charge of the strip.


\subsection{The shifted three-variable integral}\label{subsec:app-m5}

Here \(F,f\) are the linear-sieve functions.
Their normalization is
\begin{equation}\label{eq:app-delay-functions}
\begin{gathered}
 F(s)=2e^\gamma/s\quad(1\le s\le3),\qquad f(s)=0\quad(s\le2),\\
 (sF(s))'=f(s-1)\quad(s>3),\qquad(sf(s))'=F(s-1)\quad(s>2),\\
 0\le f\le1\le F,\qquad f\text{ is nondecreasing}.
\end{gathered}
\end{equation}
These classical properties are the inputs to this calculation.
Put \(u=1/500\), \(\mu_*=25/128\) and \(h_G=9/5000\).
Define
\begin{equation}\label{eq:app-six-branch}
\begin{aligned}
 v_1(s)&=\min\left(\frac{1+s}{2},\frac{32-34s}{39},
                         \frac{57}{64}-\frac{3s}{2}\right),&
 v(s)&=v_1(s)-u,\\
 \psi(x,y)&=\begin{cases}x,&x\ge y,\\-\infty,&x<y,\end{cases}\\
 \Theta_*(x,y,z)&=\max\{v(x),v(y),v(x+y),v(x+y+z),\\
 &\hspace{2.7em}\psi(v(x+z),x+2y+z),
               \psi(v(y+z),2x+y+z)\}.
\end{aligned}
\end{equation}
The unguarded branches include \(v(x)>0\) for \(u\le x\le\mu_*\).
Thus replacing the inactive value \(-\infty\) by \(0\)
leaves this maximum unchanged.
For \(d=0,h_G\), put
\begin{equation}\label{eq:app-G-functional}
\begin{aligned}
 J_d&=\int_{u\le z\le y\le x\le\mu_*}
 \frac{1-f((\Theta_*(x,y,z)+d-x-y-z)/z)}{xyz^2}
 \,dx\,dy\,dz,\\
 H_{\log}&=\log(\mu_*/u),\\
 R_{\mathrm{up}}&=500(F(30)-1)+500H_{\log}(1-f(30))
                   +250H_{\log}^2(F(30)-1),\\
 \mathcal G_d&=1/\mu_*+R_{\mathrm{up}}+J_d.
\end{aligned}
\end{equation}
The unshifted $d=0$ profile and functional come from
\cite[(4.70)--(4.75)]{PaperII}. The parameter $d$ and the shifted functional
$\mathcal G_d$ are defined in this paper. The retained-side argument uses the
instance $d=h_G$, constructed in
\eqref{eq:retB-G-use}. Paper~II's (5.4)--(5.5) are not the source of this
shifted functional.

\begin{proposition}\label{prop:app-m5-cap}
The shifted scalar satisfies
\begin{equation}\label{eq:app-m5-cap}
 0\le\mathcal G_{h_G}\le U_h<\frac{6073714261}{10^9}.
\end{equation}
Let \(F_h^+\ge F(30)\), \(f_h^-\le f(30)\),
and \(H_h^+\ge H_{\log}\) be the directed table endpoints.
Define
\begin{equation}\label{eq:app-directed-tail}
\begin{aligned}
 \mathcal R(X,Y,Z)&=500(X-1)+500Z(1-Y)
                             +250Z^2(X-1),\\
 R_h&=\mathcal R(F_h^+,f_h^-,H_h^+),\\
 U_h&=1/\mu_*+R_h+J_h^+.
\end{aligned}
\end{equation}
The construction has \(F_h^+\ge1\), \(f_h^-\le1\).
Thus \(R_{\mathrm{up}}\le R_h\). This directed endpoint comparison is the
monotonicity step of \cite[\S4.6, (4.73)]{PaperII}.
Table~\ref{tab:app-m5} encloses these finite upper bounds.
The constants $K_{\mathrm{hi}}$ and $L_{G1}$ are printed exactly in
\eqref{eq:retB-exact-Khi} and \eqref{eq:retB-exact-LG1}, with their
Paper~II locators.
The normalization identities are
\begin{equation}\label{eq:app-m5-normalization}
\begin{aligned}
 M_{5,\mathrm{internal}}
 &=-K_{\mathrm{hi}}\frac{6073714261}{10^9}-L_{G1},\\
 M_{5,\mathrm{internal}}+
 K_{\mathrm{hi}}\left(\frac{6073714261}{10^9}-U_h\right)
 &=-K_{\mathrm{hi}}U_h-L_{G1}.
\end{aligned}
\end{equation}
\end{proposition}

\begin{table}[ht]
\centering\small
\caption{Directed bounds for the shifted profile with $n_{\mathrm{mesh}}=512$.}
\label{tab:app-m5}
\begin{tabular}{@{}lll@{}}
\hline
Quantity & lower endpoint & upper endpoint\\
\hline
\(R_h\) & \(0.000004373236955\) & \(0.000004373236956\)\\
\(J_h^+\) & \(0.953709886776891\)
 & \(0.953709886776892\)\\
\(U_h\) & \(6.073714260013847\)
 & \(6.073714260013848\)\\
\(6073714261/10^9-U_h\)
 & \(0.000000000986152\)
 & \(0.000000000986153\)\\
\hline
\end{tabular}
\end{table}

\begin{proof}
For the outer partition put \(n=512\), \(x_0=u\), \(x_n=\mu_*\), and
\begin{equation}\label{eq:app-m5-mesh}
 x_i=2^{-48}\left\lfloor2^{48}
       \left(u+(\mu_*-u)\sqrt{i/n}\right)\right\rfloor
       \quad(1\le i<n).
\end{equation}
It is strictly increasing.
For the inner integral write
\begin{equation}\label{eq:app-inner}
 g_d(x,y)=\int_u^y
 \frac{1-f((\Theta_*(x,y,z)+d-x-y-z)/z)}{z^2}\,dz.
\end{equation}
Each branch residual decreases in \(x,y\).
Each guard's active set is a down-set.
Thus \(g_d\) is nondecreasing in \(x,y\),
guard jumps included.

Let \(w_i\) be an outward upper bound for
\(\log(x_i/x_{i-1})\), on the \(2^{-128}\) lattice.
Let \(v_{ij}\) enclose \(g_{h_G}(x_i,x_j)\) upwards
on that lattice.
The diagonal cell has half the logarithmic square's area.
Hence
\begin{equation}\label{eq:app-upper-corners}
 J_{h_G}\le J_h^+
 =\sum_{i=1}^n w_i\left(
 \sum_{j<i}w_jv_{ij}+\frac12w_iv_{ii}\right).
\end{equation}

Split the inner interval at every affine crossing.
Include every guard crossing and each branch kink.
On each resulting interval the argument is \(A/z+B\).
If \(\Phi(t)=\int_1^t(1-f(s))\,ds\), then
\begin{equation}\label{eq:app-inner-primitive}
 \int_p^q\frac{1-f(A/z+B)}{z^2}\,dz
 =\frac{\Phi(A/p+B)-\Phi(A/q+B)}{A}
 \quad(A\ne0).
\end{equation}
The trivial bound is \(1/p-1/q\).
The smaller of the two upper bounds may be used.
For \(A=0\), the argument is constant.

The delay equations give the primitive explicitly:
\begin{equation}\label{eq:app-Phi}
 P(t)=\frac{tF(t)}{2e^\gamma},\qquad
 \Phi(t)=t-1-2e^\gamma(P(t+1)-1)\quad(t\ge0).
\end{equation}
For \(t\ge30\), monotonicity gives
\begin{equation}\label{eq:app-Phi-tail}
 \Phi(30)\le\Phi(t)
 \le\Phi(30)+(t-30)(1-f(30)).
\end{equation}
We bound the trapezoid error. If $g$ is either $F$ or $f$ and
$[a,a+h_0]\subset[1,30]$ is one mesh panel, then
\begin{equation}\label{eq:app-trapezoid-remainder}
 \int_a^{a+h_0}g(s)\,ds-
 {h_0\over2}\bigl(g(a)+g(a+h_0)\bigr)
 =-{h_0^3\over12}g''(\xi),
 \qquad a<\xi<a+h_0,
\end{equation}
with one-sided interpretations at the integer breakpoints. The mesh has
$h_0=1/65536$ and every breakpoint is a mesh node. We prove the derivative
bound independently of the quadrature table. Write $C=2e^\gamma$. The
elementary inequality $\gamma\le H_6-\log6$ and the first three positive
terms of the atanh series give
\[
 \log2>\frac{842}{1215},\qquad \log3>\frac{263}{240},\qquad
 \gamma<\frac{12853}{19440}<\frac23<\log2.
\]
Thus $0<C<4$ and $\log2<1$. The defining functions satisfy
\[
 F(s)=C/s\ (1\le s\le3),\quad
 f(s)=0\ (1\le s\le2),\quad f(s)=C\log(s-1)/s\ (2<s\le3),
\]
\[
 sF(s)=C+\int_1^{s-1}f(t)\,dt,\qquad
 sf(s)=\int_1^{s-1}F(t)\,dt\quad(s>3).
\]
Induction on unit intervals gives $0\le F,f\le C$. If the delayed
functions are in this range, the integrals have length $s-2$, so
$F(s)\le C(s-1)/s\le C$ and $f(s)\le C(s-2)/s\le C$.
For either companion pair $(g,h)=(F,f)$ or $(f,F)$, differentiating on
open unit intervals gives
\[
 g'(s)=\frac{h(s-1)-g(s)}s,\qquad
 g''(s)=\frac{h'(s-1)}s+
        \frac{2(g(s)-h(s-1))}{s^2}\quad(s>3).
\]
For $f$, the first-derivative delay identity holds already for $t>2$.
For $F$ on $2<t<3$, use $F'(t)=-C/t^2$. Hence
$|F'(t)|,|f'(t)|\le C/t$ for $t>2$. Since $g(s),h(s-1)\in[0,C]$, for $s>3$
\[
 |g''(s)|\le\frac{C}{s(s-1)}+\frac{2C}{s^2}
 \le\frac{7C}{18}<\frac{14}{9}<12.
\]
On the initial intervals,
\[
 F''(s)=2C/s^3\ (1<s<3),\quad f''(s)=0\ (1<s<2),
\]
\[
 f''(s)=C\left(\frac{2\log(s-1)}{s^3}
       -\frac1{s(s-1)^2}-\frac2{s^2(s-1)}\right)\ (2<s<3).
\]
The formulas imply $0\le F''\le8$ on $[1,2]$,
$0\le F''\le1$ on $[2,3]$, and $-4\le f''\le\log2<1$ on $[2,3]$.
Inside the last parentheses each negative term has magnitude at most
$1/2$. The positive term is at most $\log2/4$. Therefore
$|g''|\le8<12$ on all smooth panel interiors of $[1,30]$, with one-sided
limits at the integer breakpoints.

The archived outward intervals and the independently computed exact rational
intervals are recorded with the reproducibility materials. The derivative
check starts from the proved invariant and then rebuilds the delay table.
It checks every computed unit-interval enclosure against the archived coarse
intervals. It confirms that the computed global bound is $\le12$, the
recorded value:
\[
\begin{array}{c|cc}
\text{range}&F''&f''\\ \hline
[1,2]&[0,8]&[0,0]\\
{}[2,3]&[-8,8]&[-4,4]\\
{}[3,30]&[-12,12]&[-12,12]
\end{array}
\]
The computed bound is 8. Hence the absolute local remainder is at most
$8h_0^3/12=(2/3)h_0^3\le h_0^3$. The verifier checks this radius
against the $h_0^3$ panel allowance before it accepts the regenerated
table.
The M5 computation builds the delay table with \texttt{DEN=65536},
\texttt{S=128}, 160-bit Arb
interval arithmetic, and 1,900,545 generated nodes. The shifted computation
then evaluates the exact \(n_{\mathrm{mesh}}=512\) triangular
grid, with \(512\cdot513/2=131,328\) inner-grid evaluations. It stores
one aggregated row per outer index. Hence \texttt{SHIFT-N512-fresh.json} contains 512
rows for the 131,328 evaluations.
The M5 computation uses no Taylor model and no width-\(1/32\) step.
The historical Taylor certificate provides background information and is not an input to this calculation. The table and
shifted rows are regenerated before the directed inclusion check. The
verifier compares recorded JSON fields and exact row values. It also checks
SHA-256 hashes of the canonical compact-JSON integer arrays used by the
real-inclusion verifier. These canonical-array hashes differ from hashes
of regenerated JSON files. File formatting is not used as a certificate.
Directed summation gives Table~\ref{tab:app-m5}.
The positive last row proves the strict cap. The 512 aggregated rows give the exact 131328-term sum
used for \(J_h^+\). The directed real-inclusion result records the directed
tail and the cap check. The reproducibility report documents the execution
boundary.
The two normalization identities follow by expansion.
\end{proof}

The proposition concerns the scalar functional only.
The prime-sum comparison is required separately in Section~\ref{sec:retained-side}.

\subsection{Unshifted comparison (historical and not used)}\label{subsec:app-g24}

\paragraph{Status: historical and non-certified.}
This historical unshifted comparison records an earlier \(n=8192\) calculation for
context only. No current computation reproduces it. None of its values,
table, constants or comparison is an input to the main proof.
The comparison concerns the unshifted functional
$G_*:=\mathcal G_0$. It is distinct from the main-line coefficient bound
$G_0=6088440363/10^9$. The main K2 and L3 comparisons use the latter
via \cite[(4.73), (4.75), (4.79), (A.10a)]{PaperII}. They do not use
$U_0$. The cited bound $G_*\le G_0$ is the coefficient used in the main
argument. It does not assert the separate, sharper $G_*\le U_0$ comparison.
This appendix proves that comparison only conditionally, as an optional one.
Its historical delay table was constructed separately, but no such table or
producer is part of the current proof package.
The panel width is \(h_0=1/65536\).
At 192-bit precision, each trapezoid receives radius \(h_0^3\).
On smooth panel interiors, \(|F''|,|f''|<12\).
The panels meet the integer breakpoints.
This bounds each local truncation remainder. The endpoint enclosures
pass through the delay recursion in Arb ball arithmetic.
Accumulated input radii are thus kept before outward rounding.
All table endpoints are then rounded outwards to \(2^{-128}\).
The primitive and true tail use \eqref{eq:app-Phi-tail}.
This table gives a historical directed bound for \(R_{\mathrm{up}}\).
It is kept only to reproduce the earlier non-certified calculation.
The separate earlier Taylor models are not part of the M5 computation.
Let \(F_0^+\ge F(30)\), \(f_0^-\le f(30)\),
and \(H_0^+\ge H_{\log}\) be these directed endpoints.
Define
\begin{equation}\label{eq:app-unshifted-tail}
 R_0=\mathcal R(F_0^+,f_0^-,H_0^+)
 \ge R_{\mathrm{up}}.
\end{equation}
The earlier finite calculation used \(n=8192\) in
\eqref{eq:app-m5-mesh}.
The outer sum has \(n(n+1)/2=33558528\) nodes.
Write \(J_0^+\) for that upper-corner sum.
Table~\ref{tab:app-g24-rows} groups its eight consecutive row ranges.
Each row lists a lower and upper bound for the sum.
The diagonal half-weights are already included.

\begin{table}[ht]
\centering\small
\caption{Historical non-certified contributions to the unshifted sum.}
\label{tab:app-g24-rows}
\begin{tabular}{@{}rrrr@{}}
\hline
Rows & nodes & lower endpoint & upper endpoint\\
\hline
\(1\)--\(1024\)&524800&0.000002727813799350&0.000002727813799351\\
\(1025\)--\(2048\)&1573376&0.001226034866061399&0.001226034866061400\\
\(2049\)--\(3072\)&2621952&0.019651743682986805&0.019651743682986806\\
\(3073\)--\(4096\)&3670528&0.080548405705743750&0.080548405705743751\\
\(4097\)--\(5120\)&4719104&0.146720281588383935&0.146720281588383936\\
\(5121\)--\(6144\)&5767680&0.200161682376550457&0.200161682376550458\\
\(6145\)--\(7168\)&6816256&0.242188572133550572&0.242188572133550573\\
\(7169\)--\(8192\)&7864832&0.275672089024516687&0.275672089024516688\\
\hline
\end{tabular}
\end{table}

The earlier, non-certified calculation reported
\begin{equation}\label{eq:app-g24-numbers}
\begin{aligned}
 J_0^+&<0.966171537191592960294580903427,\\
 R_0&<0.000004373236943357822206406687,\\
 U_0&=1/\mu_*+R_0+J_0^++\frac1{2\cdot10^{85}}\\
 &\in[6.086175910428536318116787310112545656706640625,\\
 &\hspace{2.4em}6.086175910428536318116787310112545656706640626].
\end{aligned}
\end{equation}
The term \(1/(2\cdot10^{85})\) bounds directed outer-sum rounding.
The shifted computation does not use this interval.
The comparison with the continuous functional is one-sided:
\begin{equation}\label{eq:app-g24-inclusion}
 \text{Historical comparison (not certified): }G_* = \mathcal G_0\le U_0.
\end{equation}
No claim $2e^{-\gamma}U_0\le c$ follows from the printed interval.
The finite row sums alone prove no prime-sum estimate.
The source comparison and its error terms remain required.

\subsection{The sixth-sum constant}\label{subsec:app-roota}

The certified finite enclosure is
\begin{equation}\label{eq:app-roota-enclosure}
 \mathfrak g_{\mathrm{RA}}
 \in[0.0989598353353,\ 0.0989598353355].
\end{equation}
Here \(\mathfrak g_{\mathrm{RA}}\) denotes the integral quantity
in Section~\ref{sec:root-a}.
The constant used in the final bound is
\begin{equation}\label{eq:app-roota-floor}
 g_{\mathrm{RA,lo}}=
 \frac{989598353}{10^{10}}=0.0989598353.
\end{equation}
Thus the certified enclosure gives
\begin{equation}\label{eq:app-roota-comparison}
 \frac{7499}{9999}
 \left(J_J(1/2)-J_J(5227/10000)\right)
 -\frac{5106819}{31557955000}
 =\mathfrak g_{\mathrm{RA}}>g_{\mathrm{RA,lo}}.
\end{equation}
Only this conservative lower constant enters the final bound.

\subsection{Computer checks}\label{subsec:app-checks}

Table~\ref{tab:app-checks} lists the supplementary scripts.
The second column states what each one checks.
The partitions and defining formulas are printed above.
The exact-arithmetic checks use Python's integer and fraction types.
Environment and timing details are recorded in the reproducibility
materials. They are not theorem inputs.
A full numerical regeneration is distinguished from a compact check.
The certificate bundle is available at the location given in the
Data availability statement; its SHA-256 manifest hash is
\texttt{81fc00441c3a3f0924995cd41263feddece827d9b40376dae531191c86f90ed3}.
The recorded M5 regeneration agrees with the exact 512-row certificate.
The reproducibility certificate set is independent. It consists of the H6
\texttt{RESULT.json}, the 512-row M5 output, the Root--A enclosure, and the
verifier scripts for those records. Each record is checked against the source
hashes and the manifest. None replaces an analytic estimate.

\begin{center}\footnotesize
\begin{minipage}{\textwidth}\centering
\captionof{table}{Scripts and finite checks.}\label{tab:app-checks}
\setlength{\tabcolsep}{4pt}
\begin{tabular}{>{\raggedright\arraybackslash}p{.25\textwidth}>{\raggedright\arraybackslash}p{.54\textwidth}>{\raggedleft\arraybackslash}p{.09\textwidth}}
Script & Statements checked & Time\tabularnewline\hline
\noalign{\vskip 2pt}
\shortstack[l]{Reproducibility\\replay verifier}
& manifests, source hashes, TeX structure, finite consistency checks,
 and the producer verification run
 & replay\tabularnewline
\shortstack[l]{H6/M5 producer\\verification}
 & executes the H6 computation and M5 $n_{\mathrm{mesh}}=512$ table, shifted probe
 and real-inclusion verifier, computing the derivative enclosure and comparing
 regenerated outputs
 & replayed\tabularnewline
\shortstack[l]{H6 directed\\verifier}
 & regenerates the directed H6 rows from the source inputs and Arb runtime
 & replayed\tabularnewline
\shortstack[l]{M5 derivative\\verifier}
 & exact rational unit-interval bounds for $F''$ and $f''$ and the panel radius
 & replayed\tabularnewline
\shortstack[l]{M5 table, shift and\\inclusion verifiers}
 & regenerates the $DEN=65536$ delay table, $n_{\mathrm{mesh}}=512$ rows, and real inclusion
 & replayed\tabularnewline
\shortstack[l]{Finite consistency\\verifier}
 & exact finite arithmetic, orbit, and sign checks for the patches
 & replayed\tabularnewline
\end{tabular}
\end{minipage}
\end{center}

The reproducibility verifier checks the recorded inputs and then runs the
producer verification.
That run regenerates the H6 output, the M5 delay table, the shifted rows
and the real inclusion on the recorded platform.
The earlier unshifted calculation used Python~3.12.3
and python-flint~0.9.0, with four workers.
Its slowest worker took \(7639\) seconds, rounded upwards.
The historical unshifted comparison is optional and not certified.
The main proof does not use it, and it does not affect the shifted
computation.
Neither these scripts nor their agreement establishes a uniform estimate.
The final proof requires the analytic comparisons in Sections~\ref{sec:retained-side}--\ref{sec:l3}.

\subsection{Rows of the retained side}\label{subsec:app-retB-rows}

The index set $\mathcal I$ of Section~\ref{sec:retained-side} consists of
Paper~II's nineteen original row chains and one outer chain. The outer
chain $\mathrm{T5}$ covers the upper part of $S_5$, Lemma~5.1. Sixteen of the row
chains are listed below, each
with its source parent and raw coefficient. The remaining three are
Paper~II's single-prefix comparisons for $S_3$, $S_8$ and $S_9$. They keep
Paper~II's own sources, restorations and error entries unchanged.
The four rows of $\mathcal H$ are restored to their parents
$S_4$, $S_5$, $S_7$ and $S_{10}$. The lower part of $S_5$ has no separate
restoration address. The four joint rows are restored once to
the joint source of $G_1$ with weights $(2,1,1)$.
\begin{center}
\small
\begin{tabular}{@{}llr|llr@{}}
\hline
Row & Parent & Coefficient & Row & Parent & Coefficient\\
\hline
S1-0 & S1 & $3$ & S4 & S4 & $-1$\\
S1-1 & S1 & $-3$ & S5-low & S5 & $-1$\\
S1-2 & S1 & $3$ & S7 & S7 & $-1$\\
S1-3 & S1 & $-3$ & S10-Ret & S10 & $1$\\
S2-0 & S2 & $1$ & G1-0 & G1 $(2,1,1)$ & $-1$\\
S2-1 & S2 & $-1$ & G1-1 & G1 $(2,1,1)$ & $1$\\
S2-2 & S2 & $1$ & G1-2 & G1 $(2,1,1)$ & $-1$\\
S2-3 & S2 & $-1$ & G1-3 & G1 $(2,1,1)$ & $1$\\
\hline
\end{tabular}
\end{center}
These sixteen rows, their seven source-parent restorations and four outer
terms carry $91$ error addresses. Each address occurs exactly once. One
outer term is the chain $\mathrm{T5}$. The others are the strict BV piece of
$S_2$, inside its joint group, and two $G_1$ comparator entries. The
same table has seven further rows with $38$ addresses. They treat
$S_3$, $S_8$ and $S_9$ another way, and this paper does not use them.
\subsection{Statements for the sixth sum}\label{subsec:app-ra-rows}

The sixth-sum argument uses the following statements.
\begin{center}
\small
\begin{tabular}{@{}>{\raggedright\arraybackslash}p{0.96\textwidth}@{}}
\hline
Statement\\
\hline
Lemma~6.2 (rough-factor distribution)\\
Lemma~6.4 (four-family Siegel--Walfisz); range \(N^{1/3}/8\le H\)\\
Lemmas~6.6 and~6.7 (selector and tensor, applied per shell pair)\\
Lemma~6.5 (transferred mean square)\\
Lemma~6.8 (two-factor convolution)\\
Lemma~6.9 (central rectangles)\\
Proper powers, (6.26)\\
Fact~3.2 (maximal prime estimates)\\
Paper~I, (2.21), as used in Paper~I, (3.25)\\
\hline
\end{tabular}
\end{center}


\section*{Use of AI tools}

The author devised the method of proof. Codex (OpenAI, Astra 6) derived and
checked the estimates, wrote the Python code for the certificates in
Appendix~B and ran blind reviews of the manuscript. Claude Code (Anthropic,
Claude Fable 5) checked the proofs, the numerical certificates and the
references. The author checked every step of the proof, every table and every
reference, and takes full responsibility for the content of the paper.

\section*{Statements and declarations}

\textbf{Competing interests.} The author has no competing interests.

\textbf{Funding.} No funding was received for this work.

\textbf{Data availability.} The supplementary scripts, exact inputs and
compact outputs are archived as
\path{paper-iii-public-certificates-v1.tar.gz} with release~v1 at
\url{https://github.com/ryzfsu/goldbach-1plus1.75-theorem/releases/download/v1/paper-iii-public-certificates-v1.tar.gz}.
\par\noindent The SHA-256 hash of the archive's internal
\texttt{SHA256SUMS} manifest is
\texttt{81fc0044}\allowbreak{}\texttt{1c3a3f09}\allowbreak{}\texttt{24995cd4}\allowbreak{}\texttt{1263fedd}\allowbreak{}\texttt{ece827d9}\allowbreak{}\texttt{b40376da}\allowbreak{}\texttt{e531191c}\allowbreak{}\texttt{86f90ed3}.

\begin{thebibliography}{99}

\bibitem{AtkinLehner1970}
A.~O.~L.~Atkin and J.~Lehner,
\emph{Hecke operators on $\Gamma_0(m)$},
Math. Ann. \textbf{185} (1970), 134--160.
\url{https://doi.org/10.1007/BF01359701}.

\bibitem{Billingsley1999}
P.~Billingsley, \emph{Convergence of Probability Measures}, 2nd ed.,
Wiley Series in Probability and Statistics, John Wiley \& Sons, New York,
1999.
\url{https://doi.org/10.1002/9780470316962}.

\bibitem{BFI86}
E.~Bombieri, J.~B.~Friedlander and H.~Iwaniec,
\emph{Primes in arithmetic progressions to large moduli},
Acta Math. \textbf{156} (1986), 203--251.
\url{https://doi.org/10.1007/BF02399204}.

\bibitem{Cai2002}
Y.~Cai, \emph{Chen's theorem with small primes},
Acta Math. Sinica (N.S.) \textbf{18} (2002), 597--604.
\url{https://doi.org/10.1007/s101140200168}.

\bibitem{Chen2002}
J.-R.~Chen, \emph{On the representation of a larger even integer as the
sum of a prime and the product of at most two primes}, in
\emph{The Goldbach Conjecture}, 2nd ed., World Scientific, 2002, 275--294.
\url{https://doi.org/10.1142/9789812776600_0021}.

\bibitem{Davenport1980}
H.~Davenport, \emph{Multiplicative Number Theory}, 2nd ed.,
revised by H.~L.~Montgomery, Graduate Texts in Mathematics, vol.~74,
Springer-Verlag, New York, 1980.
\url{https://doi.org/10.1007/978-1-4757-5927-3}.

\bibitem{DeligneBourbaki1971}
P.~Deligne, \emph{Formes modulaires et repr\'esentations $\ell$-adiques},
S\'eminaire Bourbaki 1968/69, expos\'e no.~355,
Lecture Notes in Mathematics, vol.~179, Springer-Verlag, 1971, 139--172.
\url{https://www.numdam.org/item/SB_1968-1969__11__139_0/}.

\bibitem{Deligne1974}
P.~Deligne, \emph{La conjecture de Weil. I},
Publ. Math. Inst. Hautes \'Etudes Sci. \textbf{43} (1974), 273--307.
\url{https://doi.org/10.1007/BF02684373}.

\bibitem{DI82}
J.-M.~Deshouillers and H.~Iwaniec, \emph{Kloosterman sums and Fourier
coefficients of cusp forms}, Invent. Math. \textbf{70} (1982), no.~2, 219--288.
\url{https://doi.org/10.1007/BF01390728}.

\bibitem{Drappeau2017}
S.~Drappeau, \emph{Sums of Kloosterman sums in arithmetic progressions, and
the error term in the dispersion method}, Proc. Lond. Math. Soc. (3)
\textbf{114} (2017), no.~4, 684--732.
\url{https://doi.org/10.1112/plms.12022}.

\bibitem{Weil48}
A.~Weil, \emph{On some exponential sums},
Proc. Natl. Acad. Sci. U.S.A. \textbf{34} (1948), 204--207.
\url{https://doi.org/10.1073/pnas.34.5.204}.

\bibitem{Estermann61}
T.~Estermann, \emph{On Kloosterman's sum},
Mathematika \textbf{8} (1961), 83--86.
\url{https://doi.org/10.1112/S0025579300002187}.

\bibitem{Iwaniec1980Remainder}
H.~Iwaniec, \emph{A new form of the error term in the linear sieve},
Acta Arith. \textbf{37} (1980), 307--320.
\url{https://doi.org/10.4064/aa-37-1-307-320}.

\bibitem{Iwaniec1980Rosser}
H.~Iwaniec, \emph{Rosser's sieve}, Acta Arith. \textbf{36} (1980),
no.~2, 171--202.
\url{https://doi.org/10.4064/aa-36-2-171-202}.

\bibitem{IwaniecSpectralMethods2002}
H.~Iwaniec, \emph{Spectral Methods of Automorphic Forms}, 2nd ed.,
Graduate Studies in Mathematics, vol.~53, American Mathematical Society,
Providence, RI, 2002.
\url{https://doi.org/10.1090/gsm/053}.

\bibitem{JacquetLanglands1970}
H.~Jacquet and R.~P.~Langlands,
\emph{Automorphic Forms on GL(2)}, Lecture Notes in Mathematics,
vol.~114, Springer-Verlag, Berlin--New York, 1970.
\url{https://doi.org/10.1007/BFb0058988}.

\bibitem{KimSarnak2003}
H.~Kim and P.~Sarnak, \emph{Refined estimates towards the Ramanujan and
Selberg conjectures}, Appendix~2 in H.~Kim, \emph{Functoriality for the
exterior square of $\mathrm{GL}_4$ and the symmetric fourth of
$\mathrm{GL}_2$}, J. Amer. Math. Soc. \textbf{16} (2003), no.~1, 139--183;
Appendix~2, 175--181.
\url{https://doi.org/10.1090/S0894-0347-02-00410-1}.

\bibitem{KiralYoung2017}
E.~M.~K{\i}ral and M.~P.~Young, \emph{Kloosterman sums and Fourier
coefficients of Eisenstein series}, Ramanujan J. \textbf{49} (2019),
no.~2, 391--409.
\url{https://doi.org/10.1007/s11139-018-0031-x}.
The cited proposition numbering refers to
\href{https://arxiv.org/abs/1710.00914v1}{arXiv:1710.00914v1} (2017).

\bibitem{LiLiu2026}
J.~Li and J.~Liu, \emph{Theorem $(1+1.9)$ on the Goldbach Conjecture},
\href{https://arxiv.org/abs/2606.05224v2}{arXiv:2606.05224v2}, 2026.

\bibitem{Lichtman21}
J.~D.~Lichtman, \emph{A modification of the linear sieve, and the count of
twin primes}, Algebra Number Theory \textbf{19} (2025), no.~1, 1--38.
\url{https://doi.org/10.2140/ant.2025.19.1};
\href{https://arxiv.org/abs/2109.02851v2}{arXiv:2109.02851v2}.

\bibitem{Lichtman23}
J.~D.~Lichtman, \emph{Primes in arithmetic progressions to large moduli,
and Goldbach beyond the square-root barrier},
\href{https://arxiv.org/abs/2309.08522v1}{arXiv:2309.08522v1}, 2023
(with an appendix by S.~Drappeau and J.~D.~Lichtman).

\bibitem{MaynardI}
J.~Maynard, \emph{Primes in arithmetic progressions to large moduli I:
Fixed residue classes}, Mem. Amer. Math. Soc. \textbf{306} (2025),
no.~1542.
\url{https://doi.org/10.1090/memo/1542}.

\bibitem{MaynardII}
J.~Maynard, \emph{Primes in arithmetic progressions to large moduli II:
Well-factorable estimates}, Mem. Amer. Math. Soc. \textbf{306} (2025),
no.~1543.
\url{https://doi.org/10.1090/memo/1543}.

\bibitem{MaynardIII}
J.~Maynard, \emph{Primes in arithmetic progressions to large moduli III:
Uniform residue classes}, Mem. Amer. Math. Soc. \textbf{306} (2025),
no.~1544.
\url{https://doi.org/10.1090/memo/1544}.

\bibitem{DLMF}
F.~W.~J.~Olver et al., eds., \emph{NIST Digital Library of Mathematical
Functions}, Release~1.2.8, September~15, 2026.
\url{https://dlmf.nist.gov/} (accessed November~5, 2026).

\bibitem{PaperI}
R.~Zhang, \emph{Shifted primes in arithmetic progressions to large moduli
with linear sieve remainder weights}, preprint, Zenodo, 2026.
\url{https://doi.org/10.5281/zenodo.23615525}.

\bibitem{PaperII}
R.~Zhang, \emph{Theorem (1 + 1.8) on the Goldbach conjecture}, preprint,
Zenodo, 2026.
\url{https://doi.org/10.5281/zenodo.23649264}.

\bibitem{Pascadi25}
A.~Pascadi, \emph{On the exponents of distribution of primes and smooth
numbers}, \href{https://arxiv.org/abs/2505.00653v2}{arXiv:2505.00653v2},
2025.

\bibitem{RunboLi2024}
R.~Li, \emph{On Chen's theorem, Goldbach's conjecture and almost prime
twins II}, Math. Rep. (Bucur.) \textbf{28(78)} (2026), nos.~1--2, 39--61.
\url{https://doi.org/10.59277/mrar.2026.28.78.1.2.39}.

\bibitem{Wu2004}
J.~Wu, \emph{Chen's double sieve, Goldbach's conjecture and the twin prime
problem}, Acta Arith. \textbf{114} (2004), no.~3, 215--273.
\url{https://doi.org/10.4064/aa114-3-2}.

\bibitem{Wu2008}
J.~Wu, \emph{Chen's double sieve, Goldbach's conjecture and the twin prime
problem, 2}, Acta Arith. \textbf{131} (2008), no.~4, 367--387.
\url{https://doi.org/10.4064/aa131-4-5}.

\end{thebibliography}


\end{document}
